% paper14/random_toeplitz_discrepancy.tex
%
\documentclass[11pt]{amsart}

\usepackage[T1]{fontenc}
\usepackage[utf8]{inputenc}
\usepackage{lmodern}
\usepackage{microtype}
\usepackage{amsmath,amssymb,amsthm,mathtools}
\usepackage{mathrsfs}
\usepackage{enumitem}
\usepackage[colorlinks=true,linkcolor=blue,citecolor=blue,urlcolor=blue]{hyperref}

\numberwithin{equation}{section}

\newtheorem{theorem}{Theorem}[section]
\newtheorem{proposition}[theorem]{Proposition}
\newtheorem{corollary}[theorem]{Corollary}
\newtheorem{lemma}[theorem]{Lemma}
\newtheorem{question}[theorem]{Question}
\theoremstyle{definition}
\newtheorem{definition}[theorem]{Definition}
\theoremstyle{remark}
\newtheorem{remark}[theorem]{Remark}

\newcommand{\C}{\mathbb C}
\newcommand{\R}{\mathbb R}
\newcommand{\Z}{\mathbb Z}
\newcommand{\T}{\mathbb T}
\newcommand{\M}{\mathrm M}
\newcommand{\Id}{I}
\newcommand{\op}{\mathrm{op}}
\newcommand{\HS}{\mathrm{HS}}
\newcommand{\Span}{\operatorname{span}}
\newcommand{\disc}{\operatorname{disc}}
\newcommand{\discop}{\operatorname{disc}_{\op}}
\newcommand{\Ent}{\mathsf E}
\newcommand{\Prob}{\mathbb P}
\newcommand{\Exp}{\mathbb E}
\newcommand{\supp}{\operatorname{supp}}
\newcommand{\diag}{\operatorname{diag}}
\newcommand{\wind}{\operatorname{wind}}
\newcommand{\wh}{\widehat}
\newcommand{\norm}[1]{\left\lVert #1\right\rVert}
\newcommand{\abs}[1]{\left\lvert #1\right\rvert}
\newcommand{\eps}{\varepsilon}
\newcommand{\mcA}{\mathcal A}
\newcommand{\mcB}{\mathcal B}
\newcommand{\mcF}{\mathcal F}

\DeclareMathOperator*{\argmin}{arg\,min}

\title[Operator-norm discrepancy of random Toeplitz contractions]{Operator-Norm Discrepancy of Random Banded Toeplitz Contractions}
\author{Yitzchak Shmalo}
\address{Einstein Institute of Mathematics, The Hebrew University of Jerusalem, Jerusalem, Israel}
\email{yitzchak.shmalo@gmail.com}
\date{June 2026}

\subjclass[2020]{Primary 15A60, 60B20; Secondary 15B05, 42A05, 46L53, 11K38}
\keywords{Matrix discrepancy, random Toeplitz matrices, operator norm, free probability,
intrinsic freeness, Kac polynomials, spectral edge}

\begin{document}

\begin{abstract}
For complex matrices $A_1,\ldots,A_n$ write
$\discop(A_1,\ldots,A_n)=\min_{\eps\in\{-1,1\}^n}\norm{\sum_i\eps_iA_i}_{\op}$.
We study this quantity when the $A_i$ are independent random banded Toeplitz or
Hankel contractions. The deterministic theory already settles two important
subregimes. When \(d\ge4w\) and the half-bandwidth \(w\)
is at most the number of matrices $n$, the banded theorem of the
companion paper \cite{Shmalo-paper12} gives
$\discop(A_1,\ldots,A_n)\le C\sqrt{n\log(e+w/n)}$ on \emph{every} realization, so
$\discop=O(\sqrt n)$ almost surely, with a constant free of the matrix dimension
$d$; no probabilistic input is used, and the multilevel version replaces $w$ by
$w^r$. The full-bandwidth theorem of the same companion gives the pathwise bound
\(C\sqrt{n\log(e+d/n)}\), so the random full-bandwidth question is already
true deterministically whenever \(d/n\) is bounded.  There is also a second
full-bandwidth endpoint: with the expected-norm normalization
\(A=T_d(\xi)/\Exp\|T_d(\xi)\|_{\op}\), every fixed sign pattern already gives
\(O(\sqrt n)\) with high probability, by a direct net bound for the random
Toeplitz symbol; every polynomial-sized deterministic list of sign patterns is
harmless by the same union bound, and the same holds for independent random
lists after conditioning on the list, even with only a polynomial expected
size bound.  More generally, a finite sign family of size \(M\) is uniformly
controlled at the entropy-sensitive scale
\(\sqrt{n(1+\log M/\log d)}\).  In the logarithmic sample window \(n=O(\log d)\),
all sign patterns satisfy the square-root bound simultaneously. Thus the
expected-normalized random full-bandwidth model is a theorem, not a conjecture.
Its sampled robust edge certificates also become actual full Toeplitz/Hankel
operator-norm certificates in the same unbounded-aspect range, after the
fixed-sign theorem supplies a uniform \(O(1)\) bound on the normalized summands.
What remains out of reach for the deterministic
contraction theory is the genuinely realization-normalized contraction model in
the unbounded-aspect regime \(d/n\to\infty\), where the available pathwise tools
still certify only the growing \(\sqrt{n\log(e+d/n)}\) scale.  We explain why the
intrinsic-freeness theory of Bandeira, Boedihardjo and van Handel does not by
itself certify \(O(\sqrt n)\) for this contraction problem: Toeplitz matrices are
highly delocalized in Hilbert--Schmidt norm, so the free-deviation correction is
not visibly lower order. The remaining sign-optimized contraction question is
reduced to a concrete edge-variance profile: if the normalized Fourier samples
have average square mass \(O(1/\log d)\) at every edge node, then independent
random signs already give \(O(\sqrt n)\), and this edge profile is verified in
the expected-normalized model when \(n\gtrsim\log d\).  A profile-normalized
variant also gives fixed-sign, sign-family entropy, coefficient-family
entropy, almost-sure triangular array, and explicit dimension-schedule bounds
for independent non-identical lags with arbitrary deterministic variance
envelopes.  A
normalizer-tail
criterion shows that, for any sub-Gaussian law, the same conclusion follows once
the one-matrix lower tail
\(\Prob\{\|T_d(\xi)\|_{\op}<a\,\Exp\|T_d(\xi)\|_{\op}\}\), multiplied by
\(\min\{n,\log^2(e+d)\}\), tends to zero.  On the same edge-profile event all scalar edge-grid
sign moments have the corresponding sub-Gaussian size, the sign-averaged
edge maximum is \(O(\sqrt n)\), short random search amplifies the signing
success probability with binomial success-count control, and good signings are Hamming-stable on balls of radius
\(\Theta(\log d)\), with the same binomial control for sampled robust centers;
in particular \(O_B(1)\) sampled signs give a robust certificate with failure
at most \(d^{-B}\).  A deterministic Borel transfer template upgrades every
eventual edge-grid robust-center count in the full-bandwidth range
\(d_m/n_m\to\infty\), \(\log d_m=o(n_m)\), to actual realization-normalized
Toeplitz/Hankel operator-norm robust certificates, with the same centers and
Hamming radii.  For the convex-concentrated coefficient classes this gives an
explicit high-probability operator-norm certificate theorem: \(O_B(1)\) sampled
signs contain a logarithmic-radius Hamming ball that is robust for both full
Toeplitz and full Hankel contractions; fixed-size sampled lists even contain a
\(1-\rho\) robust-center fraction, and every sign in a polynomial-size sampled
batch is robust with high probability.  The same convex-concentrated input has
an almost-sure schedule form: whenever
\(\sum_m\log^2(e+d_m)(e+d_m)^{-\beta}<\infty\) and the sampled-count tails are
summable, these robust operator certificates occur eventually almost surely;
stretched-exponential dimension schedules therefore support all-robust
polynomial sampled batches.  Strongly log-concave real or complex coefficient
laws are now recorded as another concrete unbounded class satisfying the same
convex-concentration input, and finite independent linear combinations of
convex-concentrated laws are also covered after variance normalization, even
when the resulting law is unbounded and not visibly log-concave.  Bounded
Gaussian scale mixtures with volatility bounded above and below are covered by a
separate conditional-Gaussian argument, giving non-log-concave variance-mixture
examples.  A
heterogeneous lower-normalizer library theorem further removes the
identically-distributed restriction: each matrix may be drawn from a different
coefficient law in any uniform convex, strongly log-concave, finite-convolution,
or bounded Gaussian scale-mixture library, and finite matrix-level mixtures are
included after conditioning on the selected laws.  Arbitrary positive
matrix-level radial amplitudes may also be added, even with heavy tails, because
they cancel pathwise under realization normalization.  The same heterogeneous
library input also supports the sampled
operator-certificate package: constant-size search, fixed-size robust fractions,
polynomial-size all-robust batches, and almost-sure sampled-batch schedules.
Finally, the uniform-library assumption can be weakened to an average
lower-tail budget: a sparse or budgeted population of bad coefficient laws is
allowed as long as
\(\sum_i p_{d,i}(a)\min\{1,\log^2(e+d)/n\}\to0\), and the same sampled
certificate package survives with this aggregate term.  A dyadic heterogeneous
version charges each low-normalizer law by its actual depth below \(aN_{d,i}\);
it requires only
\(\log(e+d)n^{-1}\sum_i\Exp Y_{d,i,a}^{\rm dy}\to0\), so top-scale
low-normalizer contamination may have density \(o(1/\log d)\).  A
heterogeneous inverse-moment corollary makes that dyadic budget checkable from
\(\log(e+d)n^{-1}\sum_i\nu_{d,i,a}^{(2)}\to0\), with the same
inverse-\(p\)-moment loss as in the iid case.  Its sampled-operator schedule
form gives almost-sure Toeplitz/Hankel robust centers under summable aggregate
inverse-moment matrix failure and sampled-count tails, and the corresponding
fraction form handles arbitrary sampled batch sizes, vanishing discarded
fractions, and all-robust polynomial batches.  A batch-maximal form also keeps
the entropy-sensitive whole-list maximum, both on the edge grid and in full
Toeplitz/Hankel operator norm; its concrete schedule form gives the
corresponding integer, polynomial-dimension, stretched-exponential, and
aggregate inverse-\(p\) thresholds, and the same thresholds persist under
random law-environment selection with summable empirical inverse-\(p\)-budget
exceptions; expectation and random-catalog criteria give directly checkable
ways to verify those environment tails, and bounded catalogs admit Bernstein or
Freedman checks without the extra Markov power loss, while a truncated
Freedman variant covers unbounded predictable catalogs with summable excess;
conditional moment envelopes make the truncation checkable, and an automatic
truncation choice turns this into one moment series, with an active-menu
version paying only the menu count in the Borel--Cantelli sums and concrete
active-prefix thresholds for polynomial or logarithmic menu growth; random
active counts are allowed after adding only a summable overflow tail, with
sub-Poisson proposal-count envelopes, predictable intensity budgets, and
moment, bounded-Freedman, truncated-Freedman, or automatic moment-truncation
random intensity envelopes giving direct checks, with deterministic prefixes
then chosen automatically from the same intensity envelopes.
The resulting end-to-end corollary couples that automatic prefix with the
active-menu catalog moment budget, so the sampled Toeplitz/Hankel certificates
follow from two explicit moment series; its schedule form gives all-integer,
polynomial-dimension, stretched-exponential polynomial-active, and
stretched-exponential logarithmic-active rate tests, with a sample-size form
that replaces the aggregate catalog moment series by direct lower bounds on
\(n_m\), and a compatibility form that identifies when those lower bounds still
fit the unbounded-aspect Toeplitz/Hankel window.  A per-stream variant replaces
the aggregate catalog moment exponent by the active-prefix exponent plus a
uniform row-catalog exponent.  A final first-moment-tuned variant replaces the
remaining mean-budget hypothesis by the actual per-stream first-moment decay
exponent and then chooses the inverse-\(p\) schedule parameter \(B\).
We also isolate the deterministic empirical
core: if \(M_a\) is the realized number of low-normalizer samples, then outside
the standard \(d^{-A}\) matrix events the profile is
\(O(n/\log(e+d)+M_a\log(e+d))\), giving sign scale
\(O(\sqrt{n+M_a\log^2(e+d)})\).  In particular, once
\(M_a=O(n/\log^2(e+d))\), the edge profile, random-sign tail, sign-averaged
maximum, randomized-search/success-count certificate, deterministic certificate, logarithmic
Hamming-stability with sampled robust-center tails, and all scalar edge-grid moments follow up to the standard \(d^{-A}\) exceptional
events.  A dyadic refinement charges a sample with
\(2^{-(r+1)}aN_d\le\|B_i\|_{\op}<2^{-r}aN_d\) only by a \(4^r\) weight and
keeps the Hilbert--Schmidt floor only for the final bottom bin; thus the
square-root profile follows from the weighted empirical condition
\(\sum_r4^rM_{a,r}+\ell M_{a,\infty}=O(n/\ell)\), \(\ell=\log(e+d)\), which is
strictly weaker than treating all low-normalizer samples as equally bad.  The
corresponding one-matrix dyadic-tail criterion uses only the expectation
\(\mu_{d,a}^{\rm dy}\) of this weighted depth: Markov gives matrix failure
\(O(d^{-A})+O(\ell\mu_{d,a}^{\rm dy})\), and in the large-sample window
\(\mu_{d,a}^{\rm dy}\le\theta/\ell\) gives the sharper excess cost
\(\exp\{-c n/\ell^2\}\).  An inverse-moment form bounds
\(\mu_{d,a}^{\rm dy}\) by the truncated mass
\[
  \mathbb E\!\left[
  \mathbf 1_{\{\|T_d(\xi)\|_{\op}<aN_d\}}
  \min\left\{\left(\frac{aN_d}{\|T_d(\xi)\|_{\op}}\right)^2,\ell\right\}
  \right],
\]
and inverse \(p\)-moments give the same test with an explicit logarithmic loss.
The finite inverse-moment form also carries the whole randomized certificate
package: sampled-list maxima, \(d^{-AM}\) success amplification, binomial counts,
Hamming-stable centers, and \(O_B(1)\)-sample robust certificates have the same
matrix exception \(d^{-A}+\ell\nu_{d,a}^{(2)}/R\), or the Bernstein replacement
when \(\nu_{d,a}^{(2)}=O(1/\ell)\).  The same inverse-moment matrix event
supports robust-fraction tests and polynomial-sized sampled batches: fixed
lists certify a \(1-\rho\) robust fraction with only an added \(d^{-S}\) term,
and \(M\le d^\beta\) batches are entirely robust with only an added
\(d^{-\sigma}\) term.  It also supports arbitrary-size sampled batches with a
discarded fraction \(\rho\ge e d^{-\tau}\), so the finite robust-center fraction
may be \(1-o(1)\) while the whole-list maximum keeps only its entropy factor.
The same inverse-moment input has a Borel form: summability of
\(\ell_m\nu_{d_m,a}^{(2)}\) gives eventual profiles, sampled robust
certificates, robust-fraction tests, polynomial-size all-robust batches, and
vanishing discarded-fraction tests for arbitrary sampled batches along
triangular arrays, with the Bernstein replacement whenever
\(\nu_{d_m,a}^{(2)}=O(1/\ell_m)\).  In the full-bandwidth asymptotic range
\(d_m/n_m\to\infty\), \(\log d_m=o(n_m)\), these edge-grid robust centers
transfer to actual realization-normalized Toeplitz and Hankel operator-norm
robust certificates.  Its concrete schedule form gives the
all-integer and polynomial-dimension consequences under polynomial
inverse-moment decay, the stretched-exponential schedule consequence under
logarithmic inverse-moment decay, and the critical \(\log^3d\) integer-window
consequence under \(O(1/\log d)\) truncated inverse-square mass.
Thus a shaped lower-tail curve
\(\Prob\{\|T_d(\xi)\|_{\op}<t aN_d\}\le L\ell^{-\gamma}t^\beta\),
\(\beta>2\), implies \(\mu_{d,a}^{\rm dy}\le C_{\beta}L\ell^{-\gamma}\);
at the critical density \(\gamma=1\), it already gives the same Chernoff-scale
transfer for every fixed budget \(R>C_\beta L\).  The same power-curve event also
carries the sampled random-search, success-count, and Hamming-stable
robust-center certificate package under this weaker critical density, and its
summable form gives almost-sure eventual randomized certificates along triangular
arrays.  If \(d_m^{-\alpha}\) is also summable, then \(O_\alpha(1)\) sampled
signs per dimension themselves eventually contain robust certificates, even in
the critical power-curve density.  Along stretched-exponential dimension
schedules \(d_m\ge\exp(c m^\kappa)\),
the Markov form is automatic from \(\gamma>1+1/\kappa\), while the critical
\(\gamma=1\) Bernstein form already follows from any superquadratic logarithmic
sample window \(n_m\gtrsim \ell_m^{2+\rho}\); on these sparse schedules the
sampled-list version needs no all-integer restriction \(\alpha>1\), since
\(\sum_m d_m^{-\alpha}<\infty\) for every \(\alpha>0\).  Along all integer
dimensions, the same critical \(\log^3 d\) window now also yields eventual
constant-size robust-fraction tests and all-robust polynomial batches, with the
natural \(S>1\) and \(\sigma>1\) summability exponents.  A quantitative version keeps the finite-size
failure probability \(O(d^{-A})+p_d(a)\min\{n,\log^2(e+d)\}\), and its summable
form gives an almost-sure eventual edge profile, including the same randomized
search and stable-center package, along triangular arrays.  If the sampled-list
failure term \(d_m^{-\alpha}\) is also summable, then \(O_\alpha(1)\) sampled
signs per dimension themselves eventually return logarithmically robust
certificates almost surely.  The same randomized certificate package inherits
the explicit logarithmic-window and polynomial lower-tail rates.  In particular,
polynomial tails with exponent \(B>1\) need no sparsification of the dimension
parameter: the profile holds eventually along the full sequence of integer
dimensions.  More generally, polynomial tails make this automatic along
polynomial dimension schedules \(d_m\gtrsim m^\kappa\), once \(B\kappa>1\) and
the auxiliary \(d^{-A}\) matrix-event exponent is summable; polynomial tails and
logarithmic tails of exponent greater than \(3\) make the same conclusion automatic along
exponentially sparse dimension schedules, while in
logarithmic sample windows \(n_m\lesssim\log^\gamma d_m\) along
stretched-exponential schedules \(d_m\ge \exp(c m^\kappa)\), the logarithmic
threshold improves to \(B>\min\{\gamma,2\}+1/\kappa\), with the endpoint
recovered on log-corrected schedules
\(d_m\ge\exp(c m^\kappa(\log m)^\lambda)\) when
\(\lambda(B-\min\{\gamma,2\})>1\).  At the remaining integral-test boundary,
an additional \(\log\log m\) correction with exponent above the reciprocal
threshold gives the same summability.  More generally, the sampled-count package
only needs the summability of
\((\log(e+d_m))^{\min\{\gamma,2\}-B}\) along the chosen schedule, and it is
enough to design any lower profile \(h_m\le\log(e+d_m)\) with
\(\sum_m h_m^{\min\{\gamma,2\}-B}<\infty\).  The same sparse schedules give
actual sampled robust certificates, constant-size
robust-fraction tests, vanishing discarded-fraction tests for arbitrary sampled
batches, and growing-batch count certificates for
polynomial-in-\(d_m\), subcritical log-corrected stretched-exponential, and
second-order log-corrected sign batches.  In the
large-sample window \(n/\log^2(e+d)\to\infty\), the bad-normalizer count can be
handled by Chernoff instead of Markov: if
\(p_d(a)\le \theta/\log^2(e+d)\) and one allows a fixed empirical budget
\(Rn/\log^2(e+d)\), \(R\ge1\) and \(R>\theta\), the excess count costs
\(\exp\{-I(R,\theta)n/\log^2(e+d)\}\).  In particular the weaker density bound
\(p_d(a)\le L/\log^2(e+d)\), with fixed \(L\), permits any relative excess
budget \((1+\eta)Ln/\log^2(e+d)\) at cost
\(\exp\{-Lh(\eta)n/\log^2(e+d)\}\), \(h(\eta)=(1+\eta)\log(1+\eta)-\eta\);
in particular it gives the same edge profile with failure
\(O(d^{-A})+\exp\{-c_L n/\log^2(e+d)\}\), and the same Chernoff-improved
exception controls the random-search and robust-center certificate package.
At finite size, the whole sampled-search package has joint matrix-and-sign
tails: sampled-list maxima cost \(e^{-u}\), and the counts of good signs and
robust Hamming centers have the same binomial upper tails as in the deterministic
edge-profile theorem.  Thus a fixed fraction of failures in any sampled batch
has an explicit binomial tail, simultaneously with the entropy-sensitive
sampled-list maximum.  In particular, \(O_B(1)\) sampled signs turn this into an
actual robust certificate with only the additional \(d^{-B}\) failure, while
\(O_{\rho,S}(1)\) sampled signs certify a \(1-\rho\) robust fraction at
additional failure \(d^{-S}\).  More
generally, polynomial-size sampled batches are entirely robust with polynomial
failure whenever the sign-count exponent beats the batch exponent.  The same
tails also have a Borel count form: on summable schedules, growing sign batches
eventually have the sampled-list maximum bound and all but the prescribed
number of sampled signs are good, robust centers.  With polynomial
lower-normalizer tails, polynomial-size batches are eventually entirely robust
along \(d_m\gtrsim m^\kappa\) once \((A-\beta)\kappa>1\).  In particular the fixed
robust-fraction test persists almost surely on every schedule for which the
corresponding \(d_m^{-S}\) confidence series is summable.  Under polynomial
lower-normalizer tails, this already gives robust fractions in every admissible
sample window: along \(d_m\gtrsim m^\kappa\) the confidence condition is
\(S\kappa>1\), and along all integer dimensions it is \(S>1\).  Adding the summable
sampled-list term \(\sum_m d_m^{-\alpha}<\infty\) turns the same Borel
transfer into actual \(O_\alpha(1)\)-sampled robust certificates.  The summable
Chernoff form also gives almost-sure eventual edge profiles, signing
certificates, scalar moments, random search, and robust centers.  In general it
is enough that \(\sum_m d_m^{-A}<\infty\) and
\(n_m\ge q\log^2(e+d_m)\log(e+m)\) with \(q\) above the Chernoff threshold.
In the full integer polynomial-tail branch the same count theorem allows
arbitrary sampled batches with robust-center fractions tending to one whenever
\(\rho_d\gtrsim d^{-\tau}\) for any \(\tau<A-1\).
Along stretched-exponential dimension schedules it already suffices to take
\(n_m\gtrsim\log^{2+\rho}d_m\) for any \(\rho>0\), and on those sparse
schedules the actual sampled-list version needs only \(\alpha>0\).  On stretched-exponential
schedules every \(S>0\) gives an eventual \(1-\rho\) robust fraction from
\(O_{\rho,S}(1)\) sampled signs.  The same sparse schedules support growing
sign batches: polynomial-in-\(d_m\) many
sampled signs are eventually all robust once the sign-tail exponent is larger
than the batch exponent, and both the power-curve and Chernoff sparse
count branches allow the discarded fraction \(\rho_m\) to tend to zero under
any power threshold \(\rho_m\gtrsim d_m^{-\tau}\) with \(\tau<B\).  Along all integer dimensions one may take
\(n(d)\gtrsim\log^3 d\), with \(\alpha>1\) for actual sampled lists, and in
the critical power-curve and Chernoff count branches arbitrary sampled batches
again have robust-center fractions tending to one whenever
\(\rho_d\gtrsim d^{-\tau}\) for any \(\tau<B-1\).
Convex concentration
verifies this lower tail, so the
realization-normalized full Gaussian, Rademacher, and bounded-coefficient
Toeplitz/Hankel contraction models themselves, real or complex, have
\(O(\sqrt n)\) discrepancy in the unbounded-aspect regime; strongly
log-concave laws and fixed finite independent linear combinations of these
convex-concentrated inputs are covered as well after variance normalization, as
are bounded Gaussian scale mixtures with volatility bounded above and below.
For sub-Gaussian laws without this concentration, convolution, or
scale-mixture input, the missing probabilistic input is the same
normalizer-tail/delocalization statement for realization-normalized Toeplitz
edges, not a scalar edge-moment obstruction; the near-circle phenomenon is the
one studied for a single matrix in the companion
paper \cite{Shmalo-paper5}. A final section records exact signed fourth- and
sixth-moment identities at a fixed Fourier edge, proves uniform phase-square
cancellation and scalar even-moment bounds for the expected-normalized model,
and states the higher-moment expansion that points toward a moment-method route.
\end{abstract}

\maketitle

\section{Introduction}

For self-adjoint matrices $A_1,\ldots,A_n$ the matrix discrepancy is
\[
\disc(A_1,\ldots,A_n)=\min_{\eps\in\{-1,1\}^n}\norm{\sum_{i=1}^n\eps_iA_i}_{\op},
\]
and for arbitrary complex matrices we use the same formula and write $\discop$.
The Matrix Spencer conjecture asks whether self-adjoint $d\times d$ contractions
$A_1,\ldots,A_n$ with $d\le n$ always satisfy $\disc=O(\sqrt n)$, with a constant
independent of $d$; the general case is open, and the best dimension-uniform
bounds carry a logarithmic loss from noncommutative Khintchine inequalities. We
refer to the companion paper \cite{Shmalo-paper12} for the deterministic theory
and to the surveys cited there for the background.

This paper concerns the \emph{random} version of the problem for one structured
family: independent random banded Toeplitz, or Hankel, contractions. The
question is whether the discrepancy of a typical realization is governed by the
favorable $\sqrt n$ scale or by the unfavorable $\sqrt{n\log d}$ scale that a
naive union bound would suggest. We separate the answer into two regimes by the
half-bandwidth.

For bounded bandwidth the answer is already settled, and it requires no
randomness. The deterministic banded theorem of \cite{Shmalo-paper12} bounds the
discrepancy of \emph{any} banded Toeplitz contractions, on a pathwise basis, by a
quantity that depends on the bandwidth but not on the dimension. Applying it on
each realization gives an almost-sure $O(\sqrt n)$ bound for half-bandwidth
$w\le n$ in dimensions \(d\ge4w\), dimension-free in $d$ (Theorem~\ref{thm:pathwise}). So for banded
inputs the existential ``do signs make the operator norm $O(\sqrt n)$'' is
already answered in the affirmative, deterministically. We state this cleanly
because it frames precisely what is left.

There is also a deterministic full-bandwidth endpoint that should be separated
from the genuinely random question.  The scalar Toeplitz theorem of
\cite{Shmalo-paper12} gives, for arbitrary full Toeplitz contractions,
\[
\discop(A_1,\ldots,A_n)\le C\sqrt{n\log(e+d/n)}.
\]
Consequently the full-bandwidth random question is already true on every
realization in any bounded-aspect-ratio regime \(d\le \Lambda n\), with a
constant depending only on \(\Lambda\).  What remains is the full-bandwidth
regime \(w=d-1\) with \(d/n\to\infty\).  In that regime the deterministic
pathwise method of \cite{Shmalo-paper12} reverts to the growing
\(\sqrt{n\log(e+d/n)}\) scale.  The expected-normalized model
\(A=T_d(\xi)/N_d\), \(N_d=\Exp\|T_d(\xi)\|_{\op}\), is easier than this pathwise
contraction problem: Theorem~\ref{thm:expected-full} below shows that any fixed
sign pattern already gives \(O(\sqrt n)\) with high probability, and
Corollary~\ref{cor:expected-sign-entropy} gives the entropy-sensitive uniform
version for finite sign families. Thus the expected-normalized full random
Toeplitz model is settled. What remains out of
reach is a high-probability, realization-wise signing theorem for honest
full-bandwidth random contractions in the unbounded-aspect regime. The generic
matrix Khintchine bound, using only \(\|A_i\|_{\op}\le1\), still gives the safe
upper bound \(\sqrt{n\log d}\).

One might hope that the sharp matrix-concentration theory of Bandeira,
Boedihardjo and van Handel \cite{BBvH}, which removes the Khintchine logarithm
whenever the summands behave as though freely independent, would certify
\(O(\sqrt n)\) for the contraction problem.  The direct route is not currently
justified, for a structural reason that we record in Remark~\ref{rem:bbvh}: an
expected-normalized Toeplitz matrix
is \emph{maximally delocalized}, with Hilbert--Schmidt norm of order
$\sqrt{d/\log d}$, so the free-deviation parameter $v(X)$ that controls the
correction term in \cite{BBvH} sits at the same order as the standard deviation
$\sigma(X)$, rather than below it. In that regime the $(\log d)^{3/4}$ error of
the intrinsic-freeness bound is at least the size of the main term, and the bound
collapses back to the Khintchine scale. We confirmed numerically that the ratio
$v(X)/\sigma(X)$ for signed Toeplitz superpositions stays bounded away from zero
across a range of dimensions, in the band $2.5$ to $5.3$, which is consistent
with the parameter count and rules out the free-probability shortcut.

Despite these obstacles we believe the realization-normalized full-bandwidth
contraction class is nonetheless well behaved. The worst-case factor
$\sqrt{\log(d/n)}$ is forced by adversarial Toeplitz matrices whose symbols
conspire to peak together; for independent random symbols there is no such
conspiracy. The remaining question is the following.

\begin{question}[Random full Toeplitz contractions]\label{q:main}
Let \(B_i=T_d(\xi^{(i)})\), \(1\le i\le n\), be independent full Toeplitz
matrices drawn from a fixed mean-zero sub-Gaussian unit-variance law, and set
\[
  A_i:=B_i/\norm{B_i}_{\op}.
\]
Suppose \(n\le d\le e^{o(n)}\) and \(d/n\to\infty\). Is
\[
\discop(A_1,\ldots,A_n)=O(\sqrt n)
\]
with probability tending to one as $n\to\infty$, with a constant independent of
$d$? The same question can be asked for the corresponding random Hankel model.
\end{question}

The evidence and the reduction are the content of
Section~\ref{sec:conjecture}. The deterministic theory of
\cite{Shmalo-paper12} already removes all $d$-dependence once $w\le n$, while
Theorem~\ref{thm:expected-full} removes the expected-normalized full-bandwidth
case, Corollary~\ref{cor:expected-sign-list} shows that every polynomial-sized
prescribed deterministic sign list, or independent random sign list with
polynomial expected size, is harmless, and
Corollary~\ref{cor:expected-sign-entropy} gives the sharper interpolation
according to the sign-family entropy \(\log|\mathcal S|\). In particular,
Corollary~\ref{cor:expected-all-signs} shows that in the logarithmic sample
window all sign patterns are harmless
simultaneously. Thus
Question~\ref{q:main} asks only whether realization-normalized typical
full-bandwidth contractions behave like moderately banded ones. After the
standard circulant approximation, signing Toeplitz matrices is signing the
$d$ Fourier samples of their symbols, and the discrepancy is the vector
discrepancy of those samples in $\ell^\infty$. Proposition~\ref{prop:variance-profile}
below shows that it is enough to prove a uniform empirical variance bound for
those samples, and Proposition~\ref{prop:expected-profile} proves that this
profile has exactly the expected scale in the easier expected-normalized model.
Corollary~\ref{cor:edge-to-operator-robust} records the deterministic robust
form of the reduction: any edge-grid Hamming ball at the \(O(\sqrt n)\) scale is
an actual \(O(\sqrt n)\) operator-norm robust certificate for full Toeplitz and
Hankel contractions in the unbounded-aspect asymptotic range.
Corollary~\ref{cor:borel-edge-to-operator-batches} packages the same observation
for all almost-sure sampled-count branches: eventual edge-grid robust batches
become actual Toeplitz/Hankel operator-norm robust batches with the same centers
and Hamming radii.
Corollary~\ref{cor:bounded-edge-to-operator-robust} extends the transfer to
arbitrary uniformly operator-bounded Toeplitz normalizations, and
Corollary~\ref{cor:expected-sampled-operator-certificates} applies it to the
expected-normalized sampled robust centers.
Theorem~\ref{thm:normalizer-tail} gives the precise remaining one-matrix input:
a lower tail whose product with \(\min\{n,\log^2(e+d)\}\) tends to zero
transfers the expected profile to realization normalization.
Corollary~\ref{cor:normalizer-tail-datadep} separates the deterministic
data-dependent input: with \(M_a\) low-normalizer samples, the profile is
\(O(n/\log(e+d)+M_a\log(e+d))\) and the sign scale is
\(O(\sqrt{n+M_a\log^2(e+d)})\).  The dyadic refinement
Corollary~\ref{cor:normalizer-tail-dyadic} replaces this single low-normalizer
count by a weighted budget \(W_a=\sum_r4^rM_{a,r}+\ell M_{a,\infty}\),
\(\ell=\log(e+d)\), so moderately low normalizers are charged by their actual
dyadic depth and the \(O(\sqrt n)\) scale follows from \(W_a=O(n/\ell)\).
Corollary~\ref{cor:normalizer-tail-dyadic-prob} converts this realized condition
into a one-matrix dyadic-tail input: if
\(\mu_{d,a}^{\rm dy}=\Exp W_a/n\) is \(o(1/\ell)\), Markov already gives the
square-root package with high probability, and if
\(\mu_{d,a}^{\rm dy}\le\theta/\ell\) then the excess budget has
Bernstein cost \(\exp\{-c n/\ell^2\}\).
Corollary~\ref{cor:normalizer-tail-invmoment} gives an equivalent practical
check through the truncated inverse-square normalizer mass, and through inverse
\(p\)-moments with the explicit factor \(\ell^{(2-p)_+/2}\).
Corollary~\ref{cor:normalizer-tail-invmoment-search} pushes the same
inverse-moment input through finite random search, success-count, Hamming
stability, and polynomial-confidence robust certificates.
Corollary~\ref{cor:normalizer-tail-invmoment-abundance} gives the corresponding
abundance, robust-fraction, and polynomial-batch versions under the same
inverse-moment matrix exception.
Corollary~\ref{cor:normalizer-tail-invmoment-vanishing-fractions} records the
finite arbitrary-batch version: the discarded fraction may be as small as a
power of \(d\) with no upper bound on the sampled batch size.
Corollary~\ref{cor:normalizer-tail-invmoment-borel-abundance} promotes these
count and robust-batch conclusions to almost-sure triangular-array statements
under the matching summable sign tails, with the same Chernoff-scale replacement
as the dyadic-tail criterion.
Corollary~\ref{cor:normalizer-tail-invmoment-borel-vanishing-fractions} then
lets the discarded robust-center fraction tend to zero for arbitrary sampled
batches, provided \(\rho_m\) stays above the natural binomial summability scale.
Corollary~\ref{cor:normalizer-tail-invmoment-borel-operator-batches} transfers
this inverse-moment sampled-batch package from edge-grid robust centers to
actual Toeplitz/Hankel operator-norm robust centers under the full-bandwidth
aspect assumptions.
The schedule form then extracts all-integer and polynomial-dimension
consequences from polynomial decay of this inverse mass, sparse
stretched-exponential consequences from logarithmic decay, and the critical
\(\log^3d\) integer window from an \(O(1/\ell)\) inverse-mass bound.
Corollary~\ref{cor:normalizer-tail-power-curve} gives a checkable sufficient
condition: a power lower-tail curve \(F_{d,a}(t)\lesssim \ell^{-\gamma}t^\beta\),
\(\beta>2\), forces \(\mu_{d,a}^{\rm dy}\lesssim \ell^{-\gamma}\).  In the
critical \(\gamma=1\) case the mass at the top cutoff \(aN_d\) may be only
order \(1/\ell\), rather than \(1/\ell^2\), provided the deeper tail has the
quadratic-plus decay needed to make the dyadic weights summable.
Corollary~\ref{cor:normalizer-tail-power-search} pushes this power-curve input
through the sampled random-search and Hamming-stable robust-center machinery, so
the same weaker critical-density hypothesis controls not only existence of a
good signing but also finite randomized certificate search.
Corollaries~\ref{cor:normalizer-tail-power-abundance},
\ref{cor:normalizer-tail-power-robust-fraction-confidence}, and
\ref{cor:normalizer-tail-power-poly-batch} record the corresponding finite
abundance, robust-fraction, and polynomial-batch all-robust forms under the
same power-curve matrix exception, with the Bernstein replacement in the
critical density.
Corollary~\ref{cor:normalizer-tail-power-borel-search} is the corresponding
almost-sure triangular-array form: summability of the power-curve matrix
exceptions gives eventual randomized certificates, and in the critical
\(\gamma=1\) density the Bernstein cost only requires
\(\sum_m\exp\{-c n_m/\ell_m^2\}<\infty\).
Corollary~\ref{cor:normalizer-tail-power-borel-sampled} adds the actual sampled
algorithmic conclusion: after the same power-curve matrix summability, adding
\(\sum_m d_m^{-\alpha}<\infty\) makes \(O_\alpha(1)\) sampled signs per
dimension eventually contain robust certificates almost surely.
Corollary~\ref{cor:normalizer-tail-power-borel-counts} gives the corresponding
growing-batch sampled-count theorem: summability of the sampled-list tail and
the binomial count tail makes the whole sampled list controlled and all but a
prescribed number of sampled signs good robust centers, eventually almost
surely.
Corollaries~\ref{cor:normalizer-tail-power-borel-fractions} and
\ref{cor:normalizer-tail-power-borel-poly-batch} promote the finite power-curve
abundance package to almost-sure robust-fraction and all-robust polynomial-batch
forms under the same matrix summability, with only the explicit sampled-sign
summability terms added.
Corollary~\ref{cor:normalizer-tail-power-borel-critical-integers} extracts the
full integer-dimension critical-density consequence: the \(\log^3 d\) sample
window gives not only eventual one-certificate searches, but also eventual
constant-size robust-fraction tests and all-robust polynomial batches.
Corollary~\ref{cor:normalizer-tail-power-borel-critical-vanishing-fractions}
uses the same all-integer critical window with the sampled-count theorem: for
arbitrary sampled batch sizes the discarded robust-center fraction may vary
with \(d\) and tend to zero, provided \(\rho_d\ge e d^{-\tau}\) eventually for
some \(\tau<B-1\).
Corollary~\ref{cor:normalizer-tail-power-borel-window} extracts the sparse
dimension schedules explicitly: if \(d_m\ge\exp(c m^\kappa)\), then
\(\gamma>1+1/\kappa\) is enough in the Markov form, and the critical
\(\gamma=1\) Bernstein form needs only \(n_m\ge C\ell_m^{2+\rho}\) for any
\(\rho>0\).
Corollary~\ref{cor:normalizer-tail-power-borel-window-sampled} is the actual
sampled-list sparse version: under the same sparse power-curve windows, any
fixed positive confidence exponent gives a constant number of sampled signs per
dimension which eventually contain logarithmically robust certificates almost
surely.
Corollary~\ref{cor:normalizer-tail-power-borel-window-counts} gives the
corresponding sparse growing-batch count theorem: after only the sampled-list
and binomial count summability terms, all but a prescribed number of sampled
signs are eventually good robust centers, and polynomial-in-\(d_m\) or
subcritical stretched-exponential batches are eventually entirely robust.
Corollary~\ref{cor:normalizer-tail-power-borel-window-growing-fractions} uses
the same count tail with \(r_m=\lceil\rho M_m\rceil\): on these sparse schedules,
an eventual \(1-\rho\) robust fraction holds for arbitrary sampled batch sizes,
with only the entropy-sensitive maximum remembering \(\log M_m\).
Corollary~\ref{cor:normalizer-tail-power-borel-window-vanishing-fractions}
lets this discarded fraction vary with \(m\), so \(\rho_m\to0\) is allowed
whenever \(\rho_m\ge e d_m^{-\tau}\) eventually for some \(\tau<B\).
Corollary~\ref{cor:normalizer-tail-empirical}
is the fixed-budget consequence \(M_a=O(n/\log^2(e+d))\), and carries the
random-sign tail, sign-averaged maximum, deterministic certificate, and scalar
edge-grid moment bounds; the same profile also gives a sampled-sign-list
entropy bound, gives binomial success counts for independent random search, and makes the
certificate robust under \(\Theta(\log d)\) sign flips with sampled robust-center
tails.  Corollary~\ref{cor:normalizer-tail-poly-confidence} packages this as a
joint matrix-and-search statement: \(O_B(1)\) sampled signs find such a robust
certificate with additional failure \(d^{-B}\).  In the larger
sample window \(n/\log^2(e+d)\to\infty\),
Corollary~\ref{cor:normalizer-tail-chernoff-density} gives the tunable
Chernoff tradeoff: \(p_d(a)\le\theta/\log^2(e+d)\) and a fixed empirical budget
\(R\ge1\), \(R>\theta\), cost
\(\exp\{-I(R,\theta)n/\log^2(e+d)\}\).
Corollary~\ref{cor:normalizer-tail-borel-all-dim-poly} makes the summable
polynomial regime concrete: if \(p_d(a)\le Cd^{-B}\) with \(B>1\), no
dimension sparsification is needed, and the almost-sure profile holds eventually
along the full sequence of integer dimensions.
Corollary~\ref{cor:normalizer-tail-search-rates} gives the same explicit
polylogarithmic, logarithmic-window, and polynomial failure rates for the full
random-search and robust-center certificate package.
Corollary~\ref{cor:normalizer-tail-abundance} packages the corresponding finite
abundance form: after the same non-Chernoff matrix exception, the whole sampled
list, the \(1-\rho\) good-sign fraction, and the \(1-\rho\) robust-center
fraction hold simultaneously outside an explicit binomial sign tail.
Corollary~\ref{cor:normalizer-tail-robust-fraction-confidence} turns this into
a constant-size randomized robust-fraction test with an added \(d^{-S}\)
confidence term, and Corollary~\ref{cor:normalizer-tail-poly-batch} gives the
finite polynomial-batch all-robust consequence \(M\le d^\beta\).
Corollary~\ref{cor:normalizer-tail-borel-search-stability} records that the same
Borel--Cantelli transfer also preserves the sampled random-search bounds and
logarithmic Hamming-stable robust-center certificates along the triangular array.
Corollary~\ref{cor:normalizer-tail-borel-counts} upgrades this to actual growing
sampled batches: after the sampled-list maximum tail and the binomial
failure-count tail are summable, all but the prescribed number of sampled signs
are eventually good and robust.  In the polynomial lower-tail regime, batches
\(M_m\le d_m^\beta\) are eventually entirely robust along polynomial dimension
schedules whenever \((A-\beta)\kappa>1\).
Corollary~\ref{cor:normalizer-tail-borel-log-summable-counts} isolates the
logarithmic-tail version in its sharp summability form: in logarithmic windows it
is enough that \(\sum_m(\log(e+d_m))^{\min\{\gamma,2\}-B}<\infty\), together
with the sampled sign count tails.  The later stretched, log-corrected, and
iterated-log schedule corollaries are explicit integral-test checks of this
single condition.
Corollary~\ref{cor:normalizer-tail-borel-log-schedule-counts} recasts this as a
dimension-schedule design rule: it is enough to choose any lower logarithmic
growth profile \(h_m\le\log(e+d_m)\) with
\(\sum_m h_m^{\min\{\gamma,2\}-B}<\infty\).
Corollary~\ref{cor:normalizer-tail-borel-log-vanishing-fractions} then lets the
discarded robust-center fraction vary with \(m\): arbitrary sampled batch sizes
have a \(1-\rho_m\) robust fraction whenever \(\rho_m\) stays above the natural
exponential binomial-tail threshold, so \(\rho_m\to0\) is allowed.
Corollary~\ref{cor:normalizer-tail-borel-all-dim-growing-fractions} takes
\(r_d=\lceil\rho M_d\rceil\) in the full integer polynomial-tail regime:
arbitrary sampled batch sizes have an eventual \(1-\rho\) robust-center
fraction once the count exponent satisfies \(A>1\).
Corollary~\ref{cor:normalizer-tail-borel-all-dim-vanishing-fractions} lets this
full-integer discarded fraction vary with \(d\): if
\(\rho_d\ge e d^{-\tau}\) for some \(\tau<A-1\), then \(\rho_d\to0\) is allowed
and the robust-center fraction tends to one with no upper bound on the batch
size.
Corollary~\ref{cor:normalizer-tail-borel-fractions} extracts the corresponding
constant-size robust-fraction theorem: once the matrix transfer is summable and
\(\sum_m d_m^{-S}<\infty\), \(O_{\rho,S}(1)\) sampled signs eventually contain
a \(1-\rho\) fraction of robust centers.
Corollary~\ref{cor:normalizer-tail-borel-poly-confidence} upgrades this from a
conditional sign-search statement to actual sampled-list success: if
\(\sum_m d_m^{-\alpha}<\infty\), then \(O_\alpha(1)\) sampled signs find robust
certificates for all but finitely many \(m\) almost surely over both the matrices
and the sign samples.
Corollary~\ref{cor:normalizer-tail-borel-sparse} then records the sparser
logarithmic-tail regime: polynomial lower-normalizer tails, and logarithmic
tails with exponent greater than \(3\), give the almost-sure profile along
exponentially sparse dimension schedules.
Corollary~\ref{cor:normalizer-tail-borel-sparse-sampled} upgrades this sparse
regime to actual sampled robust certificates: on exponentially sparse schedules
the sampled-list summand \(d_m^{-\alpha}\) is summable for every \(\alpha>0\).
Corollary~\ref{cor:normalizer-tail-borel-polynomial-scale} gives the polynomial
dimension-schedule counterpart: if \(p_d(a)\le C d^{-B}\) and
\(d_m\gtrsim m^\kappa\), then the same almost-sure profile follows whenever
\(B\kappa>1\), after choosing the auxiliary \(d^{-A}\) event with
\(A\kappa>1\).  Corollary~\ref{cor:normalizer-tail-borel-polynomial-sampled}
adds the sampled-list version on the same schedules, with the additional
summability condition \(\alpha\kappa>1\) for the finite random-sign list; the
robust-fraction form only needs \(S\kappa>1\) for the fraction confidence.  In
logarithmic sample windows \(n_m\lesssim\log^\gamma d_m\) along
stretched-exponential schedules \(d_m\ge\exp(c m^\kappa)\),
Corollary~\ref{cor:normalizer-tail-borel-window} improves this logarithmic
threshold to \(B>\min\{\gamma,2\}+1/\kappa\).
Corollary~\ref{cor:normalizer-tail-borel-window-log-corrected} records the
log-corrected endpoint: for
\(d_m\ge\exp(c m^\kappa(\log m)^\lambda)\), the same conclusion holds at
\(\kappa(B-\min\{\gamma,2\})=1\) provided
\(\lambda(B-\min\{\gamma,2\})>1\), with the sampled-certificate and
robust-fraction consequences on the same summable schedule.
Corollary~\ref{cor:normalizer-tail-borel-window-log-corrected-counts} adds the
matching sampled-count theorem at this endpoint, including all-robust
polynomial-in-\(d_m\) batches and subcritical log-corrected
stretched-exponential batches.
Corollary~\ref{cor:normalizer-tail-borel-window-iterlog-counts} pushes the same
sampled-count statement through one more integral-test layer:
if the first log-corrected endpoint is itself critical, an additional
\(\log(e+\log(e+m))^\mu\) factor with \(\mu(B-\min\{\gamma,2\})>1\) again gives
eventual sampled-count certificates, with the corresponding lexicographic
batch-growth rule.
Corollary~\ref{cor:normalizer-tail-borel-log-schedule-counts} is the schedule
design version of the same mechanism: any deterministic lower profile
\(h_m\le\log(e+d_m)\) with
\(\sum_m h_m^{\min\{\gamma,2\}-B}<\infty\) suffices, including the same
all-robust polynomial-batch and arbitrary-batch fraction consequences.
Corollary~\ref{cor:normalizer-tail-borel-log-vanishing-fractions} sharpens the
arbitrary-batch fraction conclusion: the discarded fraction may be a sequence
\(\rho_m\to0\), as long as \(\rho_m\ge e d_m^{-\tau}\) eventually for some
\(\tau<A\).
Corollary~\ref{cor:normalizer-tail-borel-window-sampled} adds the actual
sampled-list version of this logarithmic-window transfer; on these
stretched-exponential schedules, every \(\alpha>0\) gives a constant-size random
sign list which eventually contains robust certificates.
Corollary~\ref{cor:normalizer-tail-borel-window-fractions} gives the matching
robust-fraction version: for every \(S>0\) and \(0<\rho\le1\), a constant-size
sampled list eventually has \(O(\sqrt{n_m})\) maximum and a \(1-\rho\) fraction
of logarithmically robust centers.
Corollary~\ref{cor:normalizer-tail-borel-window-counts} adds the corresponding
growing-batch count form: on the same logarithmic-window sparse schedules,
summable sampled-list and binomial count tails force all but the prescribed
number of sampled signs to be good and robust eventually.
Corollary~\ref{cor:normalizer-tail-borel-window-growing-fractions} takes
\(r_m=\lceil\rho M_m\rceil\) in that count form, giving an eventual
\(1-\rho\) robust-center fraction for arbitrary sampled batch sizes.
Corollary~\ref{cor:normalizer-tail-large-sample}
is the specialization \(R=2L\), \(\theta=L\), while
Corollary~\ref{cor:normalizer-tail-relative-density} records the relative
excess form \(R=(1+\eta)L\).
Corollary~\ref{cor:normalizer-tail-chernoff-borel-density} records the summable
almost-sure profile/signing/moment version.
Corollary~\ref{cor:normalizer-tail-chernoff-search} carries these Chernoff and
relative-excess rates to the sampled random-search and robust-center certificate
package.
Corollary~\ref{cor:normalizer-tail-chernoff-joint-search} records the finite
joint matrix-and-sign form of that package: sampled-list maxima, good-sign
counts, and robust-center counts all inherit the Chernoff matrix exception plus
their deterministic edge-profile sign tails.
Corollary~\ref{cor:normalizer-tail-chernoff-abundance} packages the corresponding
finite abundance form: with the sampled-list maximum controlled, all but a
chosen \(\rho\)-fraction of sampled signs are good and robust outside the
matching binomial fraction tail.
Corollary~\ref{cor:normalizer-tail-chernoff-robust-fraction-confidence} turns
this into a constant-size randomized test: for prescribed \(\rho\) and
polynomial confidence exponent \(S\), \(O_{\rho,S}(1)\) sampled signs already
certify a \(1-\rho\) robust fraction.
Corollary~\ref{cor:normalizer-tail-chernoff-poly-batch} extracts the finite
polynomial-batch consequence: if \(M\le d^\beta\), then every sampled sign in
the batch is a robust center with only polynomial additional failure.
Corollary~\ref{cor:normalizer-tail-chernoff-poly-confidence} extracts the
finite actual sampled-list consequence: \(O_B(1)\) sampled signs find a robust
certificate with only the additional \(d^{-B}\) sign-search failure.
Corollary~\ref{cor:normalizer-tail-chernoff-borel-search} records the summable
almost-sure search version; together they include an all-dimensions consequence
at \(n(d)\gtrsim\log^3 d\).
Corollary~\ref{cor:normalizer-tail-chernoff-borel-sampled} upgrades the
summable Chernoff search transfer to actual sampled-list robust certificates
after adding the single summable sign-search term \(d_m^{-\alpha}\).
Corollary~\ref{cor:normalizer-tail-chernoff-borel-counts} upgrades the finite
joint tails to a Borel count theorem for growing sampled-sign batches: after a
summable sampled-list maximum tail and a summable binomial failure-count tail,
all but a prescribed number of sampled signs are eventually good and robust.
Corollary~\ref{cor:normalizer-tail-chernoff-borel-integer-vanishing-fractions}
specializes this count theorem to the full integer sequence in the critical
\(\log^3 d\) Chernoff window: arbitrary sampled batch sizes have robust-center
fractions tending to one whenever the discarded fraction satisfies
\(\rho_d\ge e d^{-\tau}\) eventually for some \(\tau<B-1\).
Corollary~\ref{cor:normalizer-tail-chernoff-borel-fractions} extracts the
constant-size robust-fraction consequence: if the \(d_m^{-S}\) confidence
series is summable, then \(O_{\rho,S}(1)\) sampled signs eventually contain a
\(1-\rho\) fraction of robust centers almost surely.
Corollary~\ref{cor:normalizer-tail-chernoff-borel-polynomial-fractions} makes
this explicit on polynomial dimension schedules \(d_m\gtrsim m^\kappa\), where
the confidence requirement becomes \(S\kappa>1\).
Corollary~\ref{cor:normalizer-tail-chernoff-borel-critical} extracts the
general critical-log schedule criterion.
Corollary~\ref{cor:normalizer-tail-chernoff-borel-window} makes the sparse
large-sample consequence explicit: along stretched-exponential schedules, any
sample window above the quadratic logarithmic scale is already summable.
Corollary~\ref{cor:normalizer-tail-chernoff-borel-window-sampled} adds the
actual sampled-list version on the same sparse schedules, with no lower bound on
\(\alpha\) beyond positivity.
Corollary~\ref{cor:normalizer-tail-chernoff-borel-window-fractions} adds the
matching robust-fraction version: on the same sparse schedules every \(S>0\)
gives a constant-size sampled list with an eventual \(1-\rho\) robust fraction.
Corollary~\ref{cor:normalizer-tail-chernoff-borel-window-counts} records the
corresponding growing-batch count form: on stretched-exponential schedules,
polynomial-in-\(d_m\) batches are eventually entirely robust as soon as the
sign-tail exponent beats the batch exponent.
Corollary~\ref{cor:normalizer-tail-chernoff-borel-window-growing-fractions}
again takes \(r_m=\lceil\rho M_m\rceil\), so arbitrary sampled batch sizes have
an eventual \(1-\rho\) robust-center fraction under the same sparse Chernoff
windows.
Corollary~\ref{cor:normalizer-tail-chernoff-borel-window-vanishing-fractions}
allows the Chernoff discarded fraction to vary with \(m\) and tend to zero,
provided \(\rho_m\ge e d_m^{-\tau}\) eventually for some \(\tau<B\).
Theorem~\ref{thm:convex-normalizers} verifies this input under convex
concentration, covering real or complex Gaussian and bounded coefficient laws
as well as Rademacher coefficients, and
Corollary~\ref{cor:strongly-log-concave-normalizers} adds the unbounded
strongly log-concave real and complex product laws.
Corollary~\ref{cor:convex-convolution-normalizers} further closes the class
under fixed finite independent linear combinations, giving examples such as
Gaussian plus Rademacher noise.
Corollary~\ref{cor:bounded-gaussian-scale-mixture-normalizers} adds centered
bounded Gaussian scale mixtures whose volatility is uniformly bounded above and
below; this covers variance mixtures by conditioning on the volatility field and
using Gaussian comparison plus concentration.
Theorem~\ref{thm:heterogeneous-normalizer-library} then removes identical
distribution from the conclusion: the \(n\) matrices may be drawn from different
laws in any uniform lower-normalizer library, including the convex,
strongly-log-concave, finite-convolution, and bounded scale-mixture classes just
listed.
Corollary~\ref{cor:matrix-level-radial-normalization} records the pathwise
invariance under arbitrary positive matrix-level radial scalings, so covered
realization-normalized models may be multiplied by heavy-tailed amplitudes
without changing the normalized array.
Theorem~\ref{thm:heterogeneous-average-normalizer-budget} further replaces the
uniform per-law lower-tail requirement by an aggregate lower-tail budget
\(\sum_i p_{d,i}(a)\min\{1,\log^2(e+d)/n\}=o(1)\), allowing sparse bad-law
contamination or nonuniform law quality.
Theorem~\ref{thm:heterogeneous-dyadic-normalizer-budget} replaces the cutoff
budget by an aggregate dyadic depth budget
\(\log(e+d)n^{-1}\sum_i\Exp Y_{d,i,a}^{\rm dy}=o(1)\), which distinguishes
moderately low normalizers from genuinely tiny ones.
Corollary~\ref{cor:heterogeneous-invmoment-normalizer-budget} makes this
criterion checkable from aggregate truncated inverse-normalizer moments.
Corollary~\ref{cor:heterogeneous-normalizer-operator-certificates} upgrades
these heterogeneous conclusions to sampled Toeplitz/Hankel operator certificates,
robust-fraction tests, polynomial batches, and almost-sure sampled-batch
schedules.
Corollary~\ref{cor:heterogeneous-invmoment-operator-schedules} records the
corresponding almost-sure sampled operator schedules under aggregate
inverse-moment summability.
Corollary~\ref{cor:heterogeneous-invmoment-operator-fractions} extracts the
fixed-fraction, vanishing-fraction, and polynomial-batch consequences for
arbitrary sampled batch sizes.
Corollary~\ref{cor:heterogeneous-invmoment-operator-batch-maxima} adds the
entropy-sensitive whole-list edge and operator-norm maximum to the same
heterogeneous inverse-moment sampled-batch conclusion.
Corollary~\ref{cor:heterogeneous-invmoment-operator-batch-maxima-schedules}
extracts the concrete all-integer, polynomial-schedule, stretched-exponential,
and Bernstein-window thresholds for these batch-maximal operator certificates.
Corollary~\ref{cor:heterogeneous-invp-operator-batch-maxima-schedules}
records the corresponding concrete aggregate inverse-\(p\) thresholds.
Corollary~\ref{cor:random-environment-invp-operator-batch-maxima-schedules}
extends those thresholds to random law environments whose empirical
inverse-\(p\) budget failures are summable.
Corollary~\ref{cor:random-environment-invp-expectation-schedules} gives
expectation and random-catalog criteria which imply that summability.
Corollary~\ref{cor:random-environment-invp-bernstein-catalog-schedules}
records the sharper independent bounded-catalog Bernstein check.
Corollary~\ref{cor:random-environment-invp-freedman-catalog-schedules} gives
the corresponding predictable bounded-catalog Freedman check for adaptive
environment selection.
Corollary~\ref{cor:random-environment-invp-truncated-freedman-catalog-schedules}
removes boundedness by adding a summable tail-excess budget.
Corollary~\ref{cor:random-environment-invp-moment-freedman-catalog-schedules}
verifies that excess from conditional moment envelopes.
Corollary~\ref{cor:random-environment-invp-auto-truncation-freedman-catalog-schedules}
chooses the truncation level automatically from a target tail exponent.
Corollary~\ref{cor:random-environment-invp-active-menu-auto-truncation}
makes the same automatic truncation simultaneous over finite active menus of
predictable catalog environments.
Corollary~\ref{cor:random-environment-invp-active-prefix-schedules} extracts
concrete active-prefix growth thresholds from that menu criterion.
Corollary~\ref{cor:random-environment-invp-random-active-prefix-counts} allows
the exposed active-prefix length itself to be random, after adding only a
summable overflow-tail condition.
Corollary~\ref{cor:random-environment-invp-subpoisson-active-counts} verifies
that condition from sub-Poisson proposal-count and occupancy envelopes.
Corollary~\ref{cor:random-environment-invp-predictable-active-intensity}
derives the same envelope from adaptive one-step activation intensities.
Corollary~\ref{cor:random-environment-invp-random-predictable-active-intensity}
allows those predictable intensity budgets to fail on a summable exceptional
set.
Corollary~\ref{cor:random-environment-invp-moment-predictable-active-intensity}
gives direct moment tests for that random intensity-envelope failure.
Corollary~\ref{cor:random-environment-invp-freedman-predictable-active-intensity}
instead certifies the random intensity envelope by scalar Freedman tails, and
Corollary~\ref{cor:random-environment-invp-truncated-freedman-predictable-active-intensity}
adds a truncation-and-excess variant for unbounded predictable loads.
Corollary~\ref{cor:random-environment-invp-auto-truncation-predictable-active-intensity}
chooses the truncation automatically from a conditional moment envelope and
records the resulting growth-rate tests.
Corollary~\ref{cor:random-environment-invp-auto-prefix-predictable-active-intensity}
then chooses the deterministic active prefix from the same intensity envelope.
Corollary~\ref{cor:random-environment-invp-end-to-end-auto-active-certificates}
combines this automatic prefix with active-menu catalog budgets to give the
sampled certificates in one check.
Corollary~\ref{cor:random-environment-invp-end-to-end-auto-active-schedules}
extracts the corresponding concrete rate schedules.
Corollary~\ref{cor:random-environment-invp-end-to-end-auto-active-sample-sizes}
turns their catalog moment series into explicit sample-size lower bounds under
aggregate catalog-growth envelopes.
Corollary~\ref{cor:random-environment-invp-end-to-end-auto-active-compatible-windows}
then extracts compatible sample-size windows from real catalog-growth exponents.
Corollary~\ref{cor:random-environment-invp-end-to-end-auto-active-per-stream-catalog-windows}
turns those aggregate exponents into per-stream catalog moment budgets.
Corollary~\ref{cor:random-environment-invp-end-to-end-auto-active-mean-tuned-windows}
also tunes the mean-budget exponent, leaving only per-stream first and
\(s\)-moment catalog rates.
Corollary~\ref{cor:convex-normalizers-sampled-operator} turns the same input
into constant-size sampled operator-norm robust certificates with logarithmic
Hamming radius, while
Corollary~\ref{cor:convex-normalizers-operator-abundance} gives the matching
operator-norm robust-fraction and polynomial-batch abundance forms.
Corollary~\ref{cor:convex-normalizers-borel-operator-batches} records the
almost-sure sampled-batch schedule version.  Finally
Corollary~\ref{cor:gaussian-realized-discrepancy} records the Gaussian
specialization.  Thus the remaining edge-profile gap is a normalizer-tail problem
for sub-Gaussian laws beyond these concentration, finite-convolution, and
bounded scale-mixture classes.
Equivalently, in the residual model one must control
$\max_{\omega}\abs{\sum_i\eps_iP_i(\omega)}$ over the $O(d)$ edge nodes $\omega$,
where $P_i$ is the normalized random symbol of $A_i$. The fluctuation of a single
such edge maximum is the near-circle edge problem of the companion paper
\cite{Shmalo-paper5}, whose local-limit problem for the unit-disk zero count
remains open and is driven by the near-circle shell. Question~\ref{q:main} is
the point at which that edge analysis and the deterministic banded discrepancy
theory meet.

We are explicit throughout about what is proved and what is not.
Theorem~\ref{thm:pathwise} (the pathwise banded bound),
Theorem~\ref{thm:bounded-aspect} (the bounded-aspect full-bandwidth endpoint),
Theorem~\ref{thm:expected-full} and
Corollaries~\ref{cor:expected-sign-list}--\ref{cor:expected-all-signs}
(the expected-normalized full random model and its deterministic/random
polynomial-list, sign-entropy, and logarithmic-window uniform-in-sign
strengthenings),
Theorem~\ref{thm:profile-normalized-fixed-sign} and
Corollaries~\ref{cor:profile-normalized-sign-entropy},
\ref{cor:profile-normalized-sign-entropy-borel},
\ref{cor:profile-normalized-sign-entropy-schedules}, and
\ref{cor:profile-normalized-coefficient-entropy}--\ref{cor:profile-normalized-coefficient-entropy-borel}
(profile-normalized independent non-identical lag arrays and their finite and
almost-sure sign/coefficient-family entropy bounds), and
Corollary~\ref{cor:expected-sampled-operator-certificates}
(expected-normalized sampled operator certificates),
Theorem~\ref{thm:normalizer-tail} and its quantitative and almost-sure
refinements, including
Corollary~\ref{cor:edge-to-operator-robust},
Corollary~\ref{cor:bounded-edge-to-operator-robust},
Corollary~\ref{cor:borel-edge-to-operator-batches},
Corollaries~\ref{cor:normalizer-tail-quant},
\ref{cor:normalizer-tail-search-stability},
\ref{cor:normalizer-tail-search-rates},
\ref{cor:normalizer-tail-chernoff-borel-density},
\ref{cor:normalizer-tail-chernoff-search},
\ref{cor:normalizer-tail-chernoff-joint-search},
\ref{cor:normalizer-tail-chernoff-abundance},
\ref{cor:normalizer-tail-chernoff-robust-fraction-confidence},
\ref{cor:normalizer-tail-chernoff-poly-batch},
\ref{cor:normalizer-tail-chernoff-poly-confidence},
\ref{cor:normalizer-tail-invmoment},
\ref{cor:normalizer-tail-invmoment-search},
\ref{cor:normalizer-tail-invmoment-abundance},
\ref{cor:normalizer-tail-invmoment-vanishing-fractions},
\ref{cor:normalizer-tail-invmoment-borel},
\ref{cor:normalizer-tail-invmoment-borel-abundance},
\ref{cor:normalizer-tail-invmoment-borel-vanishing-fractions},
\ref{cor:normalizer-tail-invmoment-borel-operator-batches},
\ref{cor:normalizer-tail-invmoment-borel-schedules},
\ref{cor:normalizer-tail-power-abundance},
\ref{cor:normalizer-tail-power-robust-fraction-confidence},
\ref{cor:normalizer-tail-power-poly-batch},
\ref{cor:normalizer-tail-power-borel-counts},
\ref{cor:normalizer-tail-power-borel-fractions},
\ref{cor:normalizer-tail-power-borel-poly-batch},
\ref{cor:normalizer-tail-power-borel-critical-integers},
\ref{cor:normalizer-tail-power-borel-critical-vanishing-fractions},
\ref{cor:normalizer-tail-power-borel-window},
\ref{cor:normalizer-tail-power-borel-window-sampled},
\ref{cor:normalizer-tail-power-borel-window-counts},
\ref{cor:normalizer-tail-power-borel-window-growing-fractions},
\ref{cor:normalizer-tail-power-borel-window-vanishing-fractions},
\ref{cor:normalizer-tail-chernoff-borel-search},
\ref{cor:normalizer-tail-chernoff-borel-sampled},
\ref{cor:normalizer-tail-chernoff-borel-counts},
\ref{cor:normalizer-tail-chernoff-borel-integer-vanishing-fractions},
\ref{cor:normalizer-tail-chernoff-borel-fractions},
\ref{cor:normalizer-tail-chernoff-borel-polynomial-fractions},
\ref{cor:normalizer-tail-chernoff-borel-critical},
\ref{cor:normalizer-tail-chernoff-borel-window},
\ref{cor:normalizer-tail-chernoff-borel-window-sampled},
\ref{cor:normalizer-tail-chernoff-borel-window-fractions},
\ref{cor:normalizer-tail-chernoff-borel-window-counts},
\ref{cor:normalizer-tail-chernoff-borel-window-growing-fractions},
\ref{cor:normalizer-tail-chernoff-borel-window-vanishing-fractions},
\ref{cor:normalizer-tail-borel-search-stability},
\ref{cor:normalizer-tail-abundance},
\ref{cor:normalizer-tail-robust-fraction-confidence},
\ref{cor:normalizer-tail-poly-batch},
\ref{cor:normalizer-tail-borel-counts},
\ref{cor:normalizer-tail-borel-log-summable-counts},
\ref{cor:normalizer-tail-borel-log-schedule-counts},
\ref{cor:normalizer-tail-borel-log-vanishing-fractions},
\ref{cor:normalizer-tail-borel-all-dim-growing-fractions},
\ref{cor:normalizer-tail-borel-all-dim-vanishing-fractions},
\ref{cor:normalizer-tail-borel-fractions},
\ref{cor:normalizer-tail-borel-all-dim-poly},
\ref{cor:normalizer-tail-borel-polynomial-scale},
\ref{cor:normalizer-tail-borel-window},
\ref{cor:normalizer-tail-borel-window-fractions},
\ref{cor:normalizer-tail-borel-window-counts},
\ref{cor:normalizer-tail-borel-window-growing-fractions},
\ref{cor:normalizer-tail-borel-window-log-corrected-counts},
\ref{cor:normalizer-tail-borel-window-iterlog-counts}, and
Theorem~\ref{thm:convex-normalizers} (the convex-concentrated realization-normalized
case), together with
Corollary~\ref{cor:strongly-log-concave-normalizers} (strongly log-concave
real and complex coefficient laws),
Corollary~\ref{cor:convex-convolution-normalizers} (finite independent linear
combinations of convex-concentrated laws),
Corollary~\ref{cor:bounded-gaussian-scale-mixture-normalizers} (bounded Gaussian
scale-mixture coefficient laws),
Theorem~\ref{thm:heterogeneous-normalizer-library} (heterogeneous
non-identically distributed lower-normalizer libraries),
Corollary~\ref{cor:matrix-level-radial-normalization} (matrix-level radial
scalings),
Theorem~\ref{thm:heterogeneous-average-normalizer-budget} (heterogeneous
average lower-tail budgets),
Theorem~\ref{thm:heterogeneous-dyadic-normalizer-budget} (heterogeneous dyadic
lower-tail budgets),
Corollary~\ref{cor:heterogeneous-invmoment-normalizer-budget} (heterogeneous
inverse-moment lower-normalizer budgets),
Corollary~\ref{cor:heterogeneous-normalizer-operator-certificates}
(heterogeneous sampled operator certificates and batches),
Corollary~\ref{cor:heterogeneous-invmoment-operator-schedules}
(heterogeneous inverse-moment sampled operator schedules),
Corollary~\ref{cor:heterogeneous-invmoment-operator-fractions}
(heterogeneous inverse-moment sampled operator robust fractions),
Corollary~\ref{cor:heterogeneous-invmoment-operator-batch-maxima}
(heterogeneous inverse-moment sampled-list operator maxima),
Corollary~\ref{cor:heterogeneous-invmoment-operator-batch-maxima-schedules}
(concrete heterogeneous inverse-moment sampled-list operator schedules),
Corollary~\ref{cor:heterogeneous-invp-operator-batch-maxima-schedules}
(concrete heterogeneous aggregate inverse-\(p\) sampled-list operator schedules),
Corollary~\ref{cor:random-environment-invp-operator-batch-maxima-schedules}
(random-environment aggregate inverse-\(p\) sampled-list operator schedules),
Corollary~\ref{cor:random-environment-invp-expectation-schedules}
(expectation and random-catalog criteria for those schedules),
Corollary~\ref{cor:random-environment-invp-bernstein-catalog-schedules}
(independent bounded-catalog Bernstein criteria for those schedules),
Corollary~\ref{cor:random-environment-invp-freedman-catalog-schedules}
(predictable bounded-catalog Freedman criteria for those schedules),
Corollary~\ref{cor:random-environment-invp-truncated-freedman-catalog-schedules}
(truncated Freedman criteria for unbounded predictable catalogs),
Corollary~\ref{cor:random-environment-invp-moment-freedman-catalog-schedules}
(conditional-moment criteria for the truncated predictable-catalog schedules),
Corollary~\ref{cor:random-environment-invp-auto-truncation-freedman-catalog-schedules}
(automatic truncation schedules for those moment criteria),
Corollary~\ref{cor:random-environment-invp-active-menu-auto-truncation}
(active-menu automatic truncation for predictable catalogs),
Corollary~\ref{cor:random-environment-invp-active-prefix-schedules}
(concrete active-prefix schedules for those menus),
Corollary~\ref{cor:random-environment-invp-random-active-prefix-counts}
(random active-prefix counts with summable overflow),
Corollary~\ref{cor:random-environment-invp-subpoisson-active-counts}
(sub-Poisson active-count overflow tests),
Corollary~\ref{cor:random-environment-invp-predictable-active-intensity}
(predictable active-intensity overflow tests),
Corollary~\ref{cor:random-environment-invp-random-predictable-active-intensity}
(random predictable-intensity overflow tests),
Corollary~\ref{cor:random-environment-invp-moment-predictable-active-intensity}
(moment tests for random predictable-intensity envelopes),
Corollary~\ref{cor:random-environment-invp-freedman-predictable-active-intensity}
(Freedman tests for random predictable-intensity envelopes),
Corollary~\ref{cor:random-environment-invp-truncated-freedman-predictable-active-intensity}
(truncated Freedman tests for unbounded predictable-intensity envelopes),
Corollary~\ref{cor:random-environment-invp-auto-truncation-predictable-active-intensity}
(automatic moment truncation for predictable active intensities),
Corollary~\ref{cor:random-environment-invp-auto-prefix-predictable-active-intensity}
(automatic deterministic prefixes from predictable active-intensity moments),
Corollary~\ref{cor:random-environment-invp-end-to-end-auto-active-certificates}
(end-to-end automatic active-prefix sampled certificates),
Corollary~\ref{cor:random-environment-invp-end-to-end-auto-active-schedules}
(concrete end-to-end automatic active-prefix schedules),
Corollary~\ref{cor:random-environment-invp-end-to-end-auto-active-sample-sizes}
(sample-size forms of those schedules),
Corollary~\ref{cor:random-environment-invp-end-to-end-auto-active-compatible-windows}
(compatible sample-size windows for catalog-growth exponents),
Corollary~\ref{cor:random-environment-invp-end-to-end-auto-active-per-stream-catalog-windows}
(per-stream catalog budgets for those compatible windows),
Corollary~\ref{cor:random-environment-invp-end-to-end-auto-active-mean-tuned-windows}
(first-moment-tuned per-stream compatible windows),
Corollary~\ref{cor:convex-normalizers-sampled-operator} (constant-size sampled
operator-norm robust certificates in the same lower-tail classes) and
Corollary~\ref{cor:convex-normalizers-operator-abundance} (robust-fraction and
polynomial-batch operator-certificate abundance), and
Corollary~\ref{cor:convex-normalizers-borel-operator-batches} (the corresponding
almost-sure sampled-batch schedules), are
theorems.  Together they imply
Corollary~\ref{cor:gaussian-realized-discrepancy} and the same
realization-normalized discrepancy conclusion for Rademacher, bounded real or
complex, strongly log-concave real or complex coefficient laws, bounded Gaussian
scale mixtures, and fixed finite independent linear combinations of those
convex-concentrated inputs, even when the matrices are sampled heterogeneously
from such a uniform library; the
heterogeneous sampled-certificate corollary gives the same randomized search and
robust-batch forms for those libraries, and the average-budget theorem allows
nonuniform law quality whenever the aggregate low-normalizer budget is small.
Question~\ref{q:main} remains open at
the full arbitrary sub-Gaussian level; Section~\ref{sec:conjecture} gives the
remaining lower-tail reduction and supporting heuristics. Section~\ref{sec:moments}
records proved fixed-edge fourth- and sixth-moment identities, uniform
phase-square cancellation and scalar even-moment estimates in the
expected-normalized model, packages these scalar moments into an
expected-normalized edge-maximum moment/signing and sampled robust-certificate
theorem, and gives a
moment-method template for the still-open all-frequency higher-moment problem.

\subsection*{Companions} This note sits between two companions. The
deterministic discrepancy engine, including the banded Toeplitz theorem used in
Section~\ref{sec:pathwise} and the circulant reduction used in
Section~\ref{sec:conjecture}, is \cite{Shmalo-paper12}. The fine behavior of the
random-Toeplitz spectral edge, including the local limit problem that controls
the obstruction in Section~\ref{sec:conjecture}, is \cite{Shmalo-paper5}.

\section{Setup}

We use the conventions of \cite{Shmalo-paper12}. For a finite lag sequence
$(a_k)_{\abs k<d}$ the Toeplitz matrix is $T_d(a)=(a_{r-s})_{0\le r,s<d}$, and if
$\varphi\in L^\infty(\T)$ has $\wh\varphi(k)=a_k$ for $\abs k<d$ then $\varphi$
induces $T_d(a)$. A Toeplitz matrix has \emph{half-bandwidth} $w$ if $a_k=0$ for
$\abs k>w$, equivalently if its symbol is a trigonometric polynomial of degree at
most $w$. The full-bandwidth case is $w=d-1$. The reversal matrix $J_d$,
$(J_dx)_r=x_{d-1-r}$, turns a Hankel matrix into a Toeplitz matrix of equal
operator norm, so every statement below transfers from Toeplitz to Hankel by the
Toeplitz--Hankel reversal argument of \cite{Shmalo-paper12}; we phrase the results for Toeplitz
matrices.

A \emph{random Toeplitz matrix} is $T_d(\xi)$ with i.i.d.\ entries
$\xi_{-(d-1)},\ldots,\xi_{d-1}$, mean zero and unit variance. Its operator norm
is of order $\sqrt{d\log d}$ \cite{Meckes}; we write
\[
N_d:=\Exp\norm{T_d(\xi)}_{\op}\asymp\sqrt{d\log d},
\]
and call \(A=T_d(\xi)/N_d\) an \emph{expected-normalized random Toeplitz matrix};
it satisfies \(\Exp\norm{A}_{\op}\asymp1\), but it is not a contraction almost surely.
When a statement needs honest contractions we instead assume the realized bound
\(\norm{A}_{\op}\le1\). A \emph{random banded Toeplitz contraction of half-bandwidth
\(w\)} means a banded Toeplitz random matrix normalized so that this realized
operator-norm bound holds; the pathwise results below use only that bound and make no
use of the law.

Throughout, $C,c$ are absolute constants that may change from line to line, and
$\eps$ ranges over $\{-1,1\}^n$.

\section{Pathwise deterministic bounds}\label{sec:pathwise}

The banded and bounded-aspect full-bandwidth regimes need no probability. We
restate the deterministic inputs from the companion paper and read off the
random corollaries.

\begin{theorem}[Banded Toeplitz signing; {\cite{Shmalo-paper12}}]
\label{thm:band-input}
There is an absolute constant $C$ such that the following holds. Let
$A_i=T_d(p_i)$ be arbitrary complex Toeplitz matrices whose symbols have degree
at most $w$, with $\norm{A_i}_{\op}\le1$, and suppose $d\ge4w$ and $w\ge1$. Then
\[
\discop(A_1,\ldots,A_n)\le C\sqrt{n\log\left(e+\frac{w}{n}\right)} .
\]
The multilevel version from \cite{Shmalo-paper12} replaces \(w\) by \(w^r\) for
\(r\)-level banded block Toeplitz or block Hankel contractions with per-coordinate
half-bandwidth \(w\).
\end{theorem}

This is a statement about every fixed tuple of matrices satisfying the
hypotheses. Since a random banded Toeplitz contraction satisfies them on every
realization, the conclusion holds almost surely with no probabilistic argument.

\begin{theorem}[Pathwise random banded discrepancy]\label{thm:pathwise}
Let $A_1,\ldots,A_n$ be random banded Toeplitz contractions of half-bandwidth
$w$, drawn from any joint law whatsoever, normalized so that
$\norm{A_i}_{\op}\le1$ almost surely, on a dimension $d\ge4w$. Then on every
realization
\[
\discop(A_1,\ldots,A_n)\le C\sqrt{n\log\left(e+\frac{w}{n}\right)} ,
\]
with $C$ absolute. In particular, for $w\le n$,
\[
\discop(A_1,\ldots,A_n)=O(\sqrt n)\qquad\text{almost surely,}
\]
with a constant independent of the dimension $d$. For $r$-level banded block
Toeplitz or block Hankel contractions of per-coordinate half-bandwidth $w$ the
same holds with $w$ replaced by $w^r$, so the bound is $O_{r}(\sqrt n)$ almost
surely whenever $w^r\le n$, again uniformly in the ambient dimension.
\end{theorem}

\begin{proof}
Fix a realization. The matrices $A_i$ are banded Toeplitz contractions of
half-bandwidth $w$ with $d\ge4w$, so Theorem~\ref{thm:band-input} applies and
yields the displayed bound; the bound is the same on every realization, hence
holds surely. For $w\le n$ the logarithm is at most $\log(e+1)$, a constant, so
the right side is $O(\sqrt n)$, with no dependence on $d$ because none appears in
the bound. The multilevel statement is identical, using the multilevel banded
theorem of \cite{Shmalo-paper12} and $w^r\le n$.
\end{proof}

So for bounded bandwidth the random discrepancy problem is fully resolved, and
the resolution is deterministic and dimension-free. Randomness contributes
nothing because the deterministic bound already covers the worst case in the
band. This is the baseline against which the full-bandwidth question is the real
content.

\begin{remark}
The dimension-freeness is the substantive point. A random banded Toeplitz
contraction lives in $\M_d(\C)$ with $d$ possibly enormous, yet its signed
superposition has operator norm $O(\sqrt n)$ as soon as $w\le n$, with no penalty
for the size of the ambient matrices. This is the random shadow of the entropy
mechanism of \cite{Shmalo-paper12}: a banded symbol carries $O(w)$ scalar degrees
of freedom regardless of $d$, so the relevant Fourier norming set has size
$O(w)$, not $O(d)$.
\end{remark}

The full-bandwidth scalar theorem of \cite{Shmalo-paper12} gives a second
pathwise endpoint.  It is not dimension-free for \(d/n\to\infty\), but it
settles the random question whenever the aspect ratio is bounded.

\begin{theorem}[Pathwise full bandwidth at bounded aspect ratio]
\label{thm:bounded-aspect}
There is an absolute constant \(C\) such that every collection of arbitrary
complex \(d\times d\) Toeplitz contractions satisfies
\[
\discop(A_1,\ldots,A_n)
\le C\sqrt{n\log\left(e+\frac dn\right)} .
\]
Consequently, for every fixed \(\Lambda<\infty\), every random full Toeplitz
contraction model with \(d\le \Lambda n\) satisfies
\[
\discop(A_1,\ldots,A_n)\le C_\Lambda\sqrt n
\]
on every realization.  The same statements hold for Hankel contractions.
\end{theorem}

\begin{proof}
The first display is the scalar Toeplitz theorem of \cite{Shmalo-paper12}.  If
\(d\le\Lambda n\), then \(\log(e+d/n)\le\log(e+\Lambda)\), giving the stated
constant \(C_\Lambda=C\sqrt{\log(e+\Lambda)}\).  The Hankel case follows from
the Toeplitz--Hankel reversal argument in the same companion paper.
\end{proof}

\section{The full-bandwidth split}\label{sec:full}

We now take $w=d-1$ and $d/n\to\infty$, and we ask when random signs or optimized
signs improve on the deterministic $\sqrt{n\log(e+d/n)}$ bound. Expected
normalization and realization normalization lead to different questions.

\subsection{Expected normalization is already solved by fixed signs}

We use two normalizations and keep them apart. The \emph{entry} normalization
takes $B=T_d(\xi)$ with unit-variance entries, so $\norm{B}_{\op}\asymp
\sqrt{d\log d}$ by Meckes \cite{Meckes} and $\norm{B}_{\HS}\asymp d$. The
\emph{expected} normalization divides by the operator-norm scale,
$A=T_d(\xi)/N_d$ with $N_d=\Exp\norm{T_d(\xi)}_{\op}\asymp\sqrt{d\log d}$, so that
$\Exp\norm{A}_{\op}\asymp1$. This normalization preserves the exact Toeplitz
superposition structure under any fixed choice of signs, and that observation
already gives the square-root bound.

\begin{theorem}[Expected-normalized full Toeplitz is fixed-sign]\label{thm:expected-full}
Let \(d\ge2\), and let \(B_i=T_d(\xi^{(i)})\), \(1\le i\le n\), be independent full Toeplitz
matrices whose lag variables are independent, mean zero, unit variance, and have
sub-Gaussian norm at most \(K\). Put
\[
  N_d:=\Exp\norm{T_d(\xi)}_{\op}\asymp_K \sqrt{d\log d},
  \qquad A_i:=B_i/N_d .
\]
Then for every \(L>0\) and every deterministic sign vector
\(\eps\in\{-1,1\}^n\),
\[
  \Prob\left\{\norm{\sum_{i=1}^n\eps_iA_i}_{\op}
  > C_{L,K}\sqrt n\right\}\le d^{-L}.
\]
Consequently, with probability at least \(1-d^{-L}\),
\[
  \discop(A_1,\ldots,A_n)\le C_{L,K}\sqrt n
\]
in the expected-normalized full Toeplitz model, without optimizing the signs.
The same conclusion holds for the corresponding expected-normalized Hankel
model.
\end{theorem}

\begin{proof}
Fix the signs.  The signed sum of the unnormalized matrices is again Toeplitz:
\[
  \sum_i\eps_iB_i=T_d(\eta),\qquad
  \eta_k:=\sum_{i=1}^n\eps_i\xi_k^{(i)} .
\]
The variables \(\eta_k\), \(\abs{k}<d\), are independent, centered, and
sub-Gaussian with norm at most \(C K\sqrt n\). Let
\[
  Q(e^{i\theta})=\sum_{\abs{k}<d}\eta_ke^{ik\theta}
\]
be the associated trigonometric symbol.  Since \(T_d(\eta)\) is a compression of
Laurent multiplication by \(Q\),
\[
  \norm{T_d(\eta)}_{\op}\le\norm{Q}_{L^\infty(\T)}.
\]
For each fixed \(\theta\), \(Q(e^{i\theta})\) has sub-Gaussian norm at most
\(C K\sqrt{nd}\). Choose a uniform grid \(\Gamma\subset\T\) of \(8d\) points.
The Bernstein inequality for trigonometric polynomials gives
\(\norm{Q'}_\infty\le d\norm{Q}_\infty\), hence
\(\norm{Q}_\infty\le2\max_{\Gamma}\abs Q\). Therefore, for a large enough
constant \(C_{L,K}\),
\[
  \Prob\{\norm{Q}_\infty>C_{L,K}\sqrt{nd\log d}\}
  \le 8d\cdot 2\exp(-c C_{L,K}^2\log d/K^2)
  \le d^{-L}.
\]
Meckes' two-sided estimate \cite{Meckes} gives
\(N_d\ge c_K\sqrt{d\log d}\). Dividing by \(N_d\) yields the displayed
\(C_{L,K}\sqrt n\) bound. The discrepancy bound follows by taking this one fixed
sign vector in the minimum.  Hankel matrices are reduced to Toeplitz matrices by
the reversal \(J_d\), which preserves operator norm.
\end{proof}

\begin{theorem}[Profile-normalized independent lags are fixed-sign]
\label{thm:profile-normalized-fixed-sign}
Let \(d\ge2\).  For \(1\le i\le n\) and \(|k|<d\), let
\(\zeta_{i,k}\) be independent centered real or complex random variables.  Let
\(\sigma_{i,k}\ge0\) be deterministic numbers such that \(\zeta_{i,k}\) has
sub-Gaussian norm at most \(K\sigma_{i,k}\), and put
\[
  S_i:=\left(\sum_{|k|<d}\sigma_{i,k}^2\right)^{1/2}.
\]
Let \(B_i=T_d((\zeta_{i,k})_{|k|<d})\), and define
\[
  A_i:=
  \begin{cases}
  B_i/(S_i\sqrt{\log(e+d)}), & S_i>0,\\
  0, & S_i=0.
  \end{cases}
\]
Then for every \(L>0\) and every deterministic sign vector
\(\eps\in\{-1,1\}^n\),
\[
  \Prob\left\{\norm{\sum_{i=1}^n\eps_iA_i}_{\op}
  > C_{L,K}\sqrt n\right\}\le d^{-L}.
\]
More generally, if \(\mathcal S\subset\{-1,1\}^n\) is deterministic and
\(\abs{\mathcal S}\le d^\alpha\), then with probability at least \(1-d^{-L}\),
\[
  \max_{\eps\in\mathcal S}
  \norm{\sum_{i=1}^n\eps_iA_i}_{\op}
  \le C_{L,\alpha,K}\sqrt n .
\]
The same conclusions hold for the corresponding profile-normalized full Hankel
matrices.
\end{theorem}

\begin{proof}
It is enough to prove the sign-list assertion; the first display is the case
\(\abs{\mathcal S}=1\).  Fix a sign vector \(\eps\), let
\(I=\{i:S_i>0\}\), and write the Toeplitz symbol of the signed normalized sum as
\[
  Q_\eps(e^{i\theta})
  =
  \sum_{i\in I}\eps_i\frac1{S_i\sqrt{\ell}}
  \sum_{|k|<d}\zeta_{i,k}e^{ik\theta},
  \qquad \ell=\log(e+d).
\]
For each fixed \(\theta\), the \(i\)-th inner normalized summand has
sub-Gaussian norm at most
\[
  C_K\frac{(\sum_{|k|<d}\sigma_{i,k}^2)^{1/2}}{S_i\sqrt{\ell}}
  =\frac{C_K}{\sqrt{\ell}}.
\]
Independence over \(i\) therefore gives \(Q_\eps(e^{i\theta})\) sub-Gaussian
norm at most \(C_K\sqrt{n/\ell}\), uniformly in \(\theta\) and \(\eps\).

Let \(\Gamma\subset\T\) be a uniform grid of \(8d\) points.  Since \(Q_\eps\) is
a trigonometric polynomial of degree at most \(d-1\), Bernstein's inequality
gives \(\norm{Q_\eps}_{L^\infty(\T)}\le2\max_{\theta\in\Gamma}
\abs{Q_\eps(e^{i\theta})}\).  The Toeplitz operator norm is bounded by
\(\norm{Q_\eps}_{L^\infty(\T)}\).  Hence, after increasing
\(C_{L,\alpha,K}\),
\[
\begin{aligned}
  \Prob\left\{\max_{\eps\in\mathcal S}
  \norm{\sum_i\eps_iA_i}_{\op}
  > C_{L,\alpha,K}\sqrt n\right\}
  &\le
  16d\,\abs{\mathcal S}\,
  \exp(-c_K C_{L,\alpha,K}^2\ell)  \\
  &\le d^{-L}.
\end{aligned}
\]
The Hankel case follows by multiplying the corresponding Toeplitz matrix by
the reversal \(J_d\), which preserves operator norm.
\end{proof}

\begin{corollary}[Profile-normalized sign entropy bound]
\label{cor:profile-normalized-sign-entropy}
In the setting of Theorem~\ref{thm:profile-normalized-fixed-sign}, fix
\(L>0\).  Let \(\mathcal S\subset\{-1,1\}^n\) be a deterministic sign family
with \(\abs{\mathcal S}\le M\), where \(M\ge1\).  Then with probability at least
\(1-d^{-L}\),
\[
  \max_{\eps\in\mathcal S}
  \norm{\sum_{i=1}^n\eps_iA_i}_{\op}
  \le
  C_K\sqrt{n\left(L+1+\frac{\log(eM)}{\log(e+d)}\right)} .
\]
The same statement holds if \(\mathcal S\) is random but independent of the
matrices and satisfies \(\Exp\abs{\mathcal S}\le M\).  In particular, with
probability at least \(1-d^{-L}\),
\[
  \max_{\eps\in\{-1,1\}^n}
  \norm{\sum_{i=1}^n\eps_iA_i}_{\op}
  \le
  C_K\sqrt{n\left(L+1+\frac{n}{\log(e+d)}\right)} ,
\]
and if \(n\le\lambda\log(e+d)\), all sign patterns are simultaneously
\(O_{L,K,\lambda}(\sqrt n)\).  The same conclusions hold for the corresponding
profile-normalized full Hankel matrices.
\end{corollary}

\begin{proof}
The proof of Theorem~\ref{thm:profile-normalized-fixed-sign}, before the final
constant is chosen, gives the fixed-sign tail
\[
  \Prob\left\{
  \norm{\sum_i\eps_iA_i}_{\op}>C_Ku\sqrt n
  \right\}
  \le (e+d)^{-u^2}
  \qquad(u\ge1)
\]
for every deterministic \(\eps\).  Take
\[
  u^2=L+1+\frac{\log(eM)}{\log(e+d)}.
\]
A union bound over deterministic \(\mathcal S\) gives the first assertion.
For a random independent \(\mathcal S\), condition on the realized family and
then average:
\[
  \Prob\left\{\max_{\eps\in\mathcal S}
  \norm{\sum_i\eps_iA_i}_{\op}>C_Ku\sqrt n\right\}
  \le
  \Exp\abs{\mathcal S}\,(e+d)^{-u^2}
  \le d^{-L}.
\]
Taking \(M=2^n\) gives the whole-cube estimate.  The Hankel statement follows
again by Toeplitz--Hankel reversal.
\end{proof}

\begin{corollary}[Profile-normalized Borel sign entropy]
\label{cor:profile-normalized-sign-entropy-borel}
For each \(m\), let \(d_m\ge2\), \(n_m\ge1\), and let
\((A_{m,i})_{i\le n_m}\) be a profile-normalized full Toeplitz array satisfying
the hypotheses of Theorem~\ref{thm:profile-normalized-fixed-sign} with the same
sub-Gaussian constant \(K\).  Let \(\mathcal S_m\subset\{-1,1\}^{n_m}\) be
nonempty and deterministic, or nonempty almost surely, random, and independent
of the \(m\)-th matrix array.  Assume \(\Exp\abs{\mathcal S_m}\le M_m\),
\(M_m\ge1\).  If \(L_m\ge1\) and
\[
  \sum_m d_m^{-L_m}<\infty ,
\]
then, almost surely, for all sufficiently large \(m\),
\[
  \max_{\eps\in\mathcal S_m}
  \norm{\sum_{i=1}^{n_m}\eps_iA_{m,i}}_{\op}
  \le
  C_K\sqrt{n_m\left(L_m+1+
  \frac{\log(eM_m)}{\log(e+d_m)}\right)} .
\]
In particular, if \(M_m=2^{n_m}\), then the same event gives the simultaneous
whole-cube bound
\[
  \max_{\eps\in\{-1,1\}^{n_m}}
  \norm{\sum_{i=1}^{n_m}\eps_iA_{m,i}}_{\op}
  \le
  C_K\sqrt{n_m\left(L_m+1+\frac{n_m}{\log(e+d_m)}\right)}
\]
eventually almost surely.  If additionally \(n_m\le\lambda\log(e+d_m)\) and
\(\sum_m d_m^{-L}<\infty\) for some fixed \(L\), every sign pattern is
eventually \(O_{K,L,\lambda}(\sqrt{n_m})\).  The same conclusions hold for the
corresponding profile-normalized full Hankel arrays.
\end{corollary}

\begin{proof}
Apply Corollary~\ref{cor:profile-normalized-sign-entropy} to the \(m\)-th array
with \(L=L_m\).  The exceptional probabilities are bounded by \(d_m^{-L_m}\),
which is summable.  The first Borel--Cantelli lemma gives the eventual Toeplitz
conclusion; independence across \(m\) is not needed.  The whole-cube and
logarithmic-window claims are the special cases \(M_m=2^{n_m}\) and
\(n_m\le\lambda\log(e+d_m)\).  The Hankel statement follows by applying the
same argument after Toeplitz--Hankel reversal.
\end{proof}

\begin{corollary}[Profile-normalized sign-entropy schedules]
\label{cor:profile-normalized-sign-entropy-schedules}
In the setting of Corollary~\ref{cor:profile-normalized-sign-entropy-borel},
the following concrete schedule consequences hold.

\begin{enumerate}[label=\textup{(\roman*)}]
\item If \(d_m\ge m^\kappa\) eventually for some \(\kappa>0\) and
\(\Exp\abs{\mathcal S_m}\le d_m^\alpha\) for some \(\alpha<\infty\), then
almost surely, for all sufficiently large \(m\),
\[
  \max_{\eps\in\mathcal S_m}
  \norm{\sum_{i=1}^{n_m}\eps_iA_{m,i}}_{\op}
  \le C_{K,\alpha,\kappa}\sqrt{n_m}.
\]
\item If \(d_m\ge \exp(c m^\sigma)\) eventually for some \(c,\sigma>0\) and
\(\Exp\abs{\mathcal S_m}\le d_m^\alpha\) for some \(\alpha<\infty\), then the
same conclusion holds with a constant \(C_{K,\alpha}\).
\item In either dimension-growth schedule, if
\(\mathcal S_m=\{-1,1\}^{n_m}\) and
\(n_m\le\lambda\log(e+d_m)\) eventually, then almost surely every sign pattern
is eventually \(O_{K,\lambda,\kappa}(\sqrt{n_m})\) in the polynomial schedule
and \(O_{K,\lambda}(\sqrt{n_m})\) in the stretched-exponential schedule.
\end{enumerate}
The same conclusions hold for the corresponding profile-normalized full Hankel
arrays.
\end{corollary}

\begin{proof}
For the polynomial schedule choose \(L>\max\{1,\kappa^{-1}\}\).  Then
\[
  \sum_m d_m^{-L}\le\sum_m m^{-\kappa L}<\infty .
\]
For the stretched-exponential schedule take \(L=1\), since
\[
  \sum_m d_m^{-1}\le\sum_m \exp(-c m^\sigma)<\infty .
\]
In both cases,
Corollary~\ref{cor:profile-normalized-sign-entropy-borel} applies with a fixed
summable tail exponent.  If \(\Exp\abs{\mathcal S_m}\le d_m^\alpha\), then
\[
  \frac{\log(eM_m)}{\log(e+d_m)}\le C_\alpha,
\]
uniformly for \(d_m\ge2\), giving the first two bounds.  For the whole cube,
take \(M_m=2^{n_m}\); the assumption \(n_m\le\lambda\log(e+d_m)\) gives
\[
  \frac{\log(eM_m)}{\log(e+d_m)}
  \le C_\lambda .
\]
The Hankel statement is again the Toeplitz statement after reversal.
\end{proof}

\begin{corollary}[Profile-normalized coefficient-family entropy]
\label{cor:profile-normalized-coefficient-entropy}
In the setting of Theorem~\ref{thm:profile-normalized-fixed-sign}, fix
\(L>0\).  Let \(\mathcal V\subset\mathbb C^n\) be a finite nonempty deterministic
coefficient family with \(\abs{\mathcal V}\le M\), \(M\ge1\), and
\[
  \norm{v}_2\le R\qquad(v\in\mathcal V).
\]
Then with probability at least \(1-d^{-L}\),
\[
  \max_{v\in\mathcal V}
  \norm{\sum_{i=1}^n v_iA_i}_{\op}
  \le
  C_K R\sqrt{L+1+\frac{\log(eM)}{\log(e+d)}} .
\]
The same statement holds if \(\mathcal V\) is finite and nonempty almost
surely, random but independent of the matrices, satisfies
\(\Exp\abs{\mathcal V}\le M\), and obeys \(\norm{v}_2\le R\) for all
\(v\in\mathcal V\) almost surely.  The same
conclusions hold for the corresponding profile-normalized full Hankel
matrices.
\end{corollary}

\begin{proof}
If \(R=0\), every displayed sum vanishes, so assume \(R>0\).  Fix a
deterministic \(v\in\mathbb C^n\), let \(I=\{i:S_i>0\}\), and write
\[
  Q_v(e^{i\theta})
  =
  \sum_{i\in I}v_i\frac1{S_i\sqrt{\ell}}
  \sum_{|k|<d}\zeta_{i,k}e^{ik\theta},
  \qquad \ell=\log(e+d).
\]
For each fixed \(\theta\), the \(i\)-th normalized summand has sub-Gaussian
norm at most \(C_K/\sqrt{\ell}\).  Independence over \(i\) gives
\[
  \norm{Q_v(e^{i\theta})}_{\psi_2}
  \le C_K\frac{\norm{v}_2}{\sqrt{\ell}}
  \le C_K\frac R{\sqrt{\ell}},
\]
uniformly in \(\theta\) and \(v\).  Using the same \(8d\)-point grid and
Bernstein inequality as in the proof of
Theorem~\ref{thm:profile-normalized-fixed-sign}, we obtain the fixed-vector
tail
\[
  \Prob\left\{
  \norm{\sum_i v_iA_i}_{\op}>C_KuR
  \right\}
  \le (e+d)^{-u^2}
  \qquad(u\ge1).
\]
Choose
\[
  u^2=L+1+\frac{\log(eM)}{\log(e+d)} .
\]
The union bound over deterministic \(\mathcal V\) gives the asserted
probability.  For a random independent \(\mathcal V\), condition on the
realized family and average, using \(\Exp\abs{\mathcal V}\le M\).  The Hankel
claim follows from Toeplitz--Hankel reversal.
\end{proof}

\begin{corollary}[Profile-normalized Borel coefficient entropy]
\label{cor:profile-normalized-coefficient-entropy-borel}
For each \(m\), let \(d_m\ge2\), \(n_m\ge1\), and let
\((A_{m,i})_{i\le n_m}\) be a profile-normalized full Toeplitz array satisfying
the hypotheses of Theorem~\ref{thm:profile-normalized-fixed-sign} with the same
sub-Gaussian constant \(K\).  Let
\(\mathcal V_m\subset\mathbb C^{n_m}\) be finite and nonempty deterministic, or
finite and nonempty almost surely, random, and independent of the \(m\)-th
matrix array.  Assume
\[
  \Exp\abs{\mathcal V_m}\le M_m,\qquad
  \sup_{v\in\mathcal V_m}\norm{v}_2\le R_m
\]
almost surely, where \(M_m\ge1\).  If \(L_m\ge1\) and
\[
  \sum_m d_m^{-L_m}<\infty ,
\]
then, almost surely, for all sufficiently large \(m\),
\[
  \max_{v\in\mathcal V_m}
  \norm{\sum_{i=1}^{n_m}v_iA_{m,i}}_{\op}
  \le
  C_K R_m\sqrt{L_m+1+
  \frac{\log(eM_m)}{\log(e+d_m)}} .
\]
In particular, if \(d_m\ge m^\kappa\) eventually and
\(M_m\le d_m^\alpha\) for some \(\kappa>0\), \(\alpha<\infty\), then the
right side is \(O_{K,\alpha,\kappa}(R_m)\) eventually almost surely.  If
\(d_m\ge \exp(c m^\sigma)\) eventually and \(M_m\le d_m^\alpha\), then it is
eventually \(O_{K,\alpha}(R_m)\).  The same conclusions hold for the
corresponding profile-normalized full Hankel arrays.
\end{corollary}

\begin{proof}
Apply Corollary~\ref{cor:profile-normalized-coefficient-entropy} to the
\(m\)-th array with \(L=L_m\).  The exceptional probabilities are bounded by
\(d_m^{-L_m}\), so Borel--Cantelli gives the eventual Toeplitz conclusion.  In
the polynomial schedule choose \(L>\max\{1,\kappa^{-1}\}\), and in the
stretched-exponential schedule choose \(L=1\).  Then
\(\sum_m d_m^{-L}<\infty\) in the corresponding schedule, while
\(M_m\le d_m^\alpha\) gives
\[
  \frac{\log(eM_m)}{\log(e+d_m)}\le C_\alpha .
\]
The Hankel statement again follows by Toeplitz--Hankel reversal.
\end{proof}

\begin{corollary}[Expected-normalized polynomial sign lists are harmless]
\label{cor:expected-sign-list}
In the setting of Theorem~\ref{thm:expected-full}, fix \(L>0\) and
\(\alpha<\infty\).  Let
\(\mathcal S\subset\{-1,1\}^n\) be a deterministic sign list with
\(\abs{\mathcal S}\le d^\alpha\).  Then with probability at least \(1-d^{-L}\),
\[
  \max_{\eps\in\mathcal S}
  \norm{\sum_{i=1}^n\eps_iA_i}_{\op}
  \le C_{L,\alpha,K}\sqrt n .
\]
The same simultaneous statement holds if \(\mathcal S\) is random but
independent of the matrices and either satisfies
\(\abs{\mathcal S}\le d^\alpha\) almost surely or only satisfies
\(\Exp\abs{\mathcal S}\le d^\alpha\).  It also holds for the corresponding
expected-normalized Hankel model.
\end{corollary}

\begin{proof}
Apply Theorem~\ref{thm:expected-full} with exponent \(L+\alpha+1\).  For each
fixed sign vector the exceptional probability is at most
\(d^{-L-\alpha-1}\).  A union bound over deterministic \(\mathcal S\) gives
exceptional probability at most \(d^{-L-1}\), hence at most \(d^{-L}\).  The
Hankel case again follows by Toeplitz--Hankel reversal.

If \(\mathcal S\) is random and independent of the matrices, condition on the
realized list and use the fixed-sign union bound:
\[
  \Prob\left\{\left.
  \max_{\eps\in\mathcal S}
  \norm{\sum_{i=1}^n\eps_iA_i}_{\op}
  > C_{L,\alpha,K}\sqrt n
  \ \right|\ \mathcal S\right\}
  \le \abs{\mathcal S}\,d^{-L-\alpha-1}.
\]
Taking expectations gives an exceptional probability at most
\(\Exp\abs{\mathcal S}\,d^{-L-\alpha-1}\le d^{-L-1}\le d^{-L}\).  The
almost-sure size assumption is the special case
\(\Exp\abs{\mathcal S}\le d^\alpha\).
\end{proof}

\begin{corollary}[Expected-normalized sign entropy bound]
\label{cor:expected-sign-entropy}
In the setting of Theorem~\ref{thm:expected-full}, fix \(L>0\).  Let
\(\mathcal S\subset\{-1,1\}^n\) be a deterministic sign family with
\(\abs{\mathcal S}\le M\), where \(M\ge1\).  Then with probability at least
\(1-d^{-L}\),
\[
  \max_{\eps\in\mathcal S}
  \norm{\sum_{i=1}^n\eps_iA_i}_{\op}
  \le
  C_K\sqrt{n\left(L+1+\frac{\log(eM)}{\log(e+d)}\right)} .
\]
The same statement holds if \(\mathcal S\) is random but independent of the
matrices and satisfies \(\Exp\abs{\mathcal S}\le M\).  In particular, with
probability at least \(1-d^{-L}\),
\[
  \max_{\eps\in\{-1,1\}^n}
  \norm{\sum_{i=1}^n\eps_iA_i}_{\op}
  \le
  C_K\sqrt{n\left(L+1+\frac{n}{\log(e+d)}\right)}
\]
in the expected-normalized Toeplitz and Hankel models.
\end{corollary}

\begin{proof}
The proof of Theorem~\ref{thm:expected-full}, before the exponent is fixed,
gives the following tail estimate: after enlarging \(C_K\), for every fixed
sign vector \(\eps\) and every \(u\ge1\),
\[
  \Prob\left\{
  \norm{\sum_{i=1}^n\eps_iA_i}_{\op}>C_Ku\sqrt n
  \right\}
  \le (e+d)^{-u^2}.
\]
Indeed, for the signed symbol \(Q\), the \(8d\)-point net gives
\[
  \Prob\{\norm{Q}_\infty>C_Ku\sqrt{nd\log(e+d)}\}
  \le (e+d)^{-u^2}.
\]
Dividing by \(N_d\ge c_K\sqrt{d\log(e+d)}\) gives the displayed tail.

Set
\[
  u^2=L+2+\frac{\log(eM)}{\log(e+d)} .
\]
For deterministic \(\mathcal S\), the union bound gives
\[
  \Prob\left\{
  \max_{\eps\in\mathcal S}
  \norm{\sum_{i=1}^n\eps_iA_i}_{\op}>C_Ku\sqrt n
  \right\}
  \le
  M(e+d)^{-L-2}(eM)^{-1}
  \le d^{-L},
\]
after adjusting constants for \(d\ge2\).  Absorbing the harmless \(+2\) into
the displayed \(+1\) changes only \(C_K\).

If \(\mathcal S\) is random and independent of the matrices, condition on
\(\mathcal S\) and take expectations in the same inequality; the factor
\(\abs{\mathcal S}\) is replaced by \(\Exp\abs{\mathcal S}\le M\).
Finally take \(M=2^n\). Since
\(\log(e2^n)\le C(1+n)\), the stated all-sign bound follows, with another
absolute change of \(C_K\).  The Hankel case is again obtained by reversal.
\end{proof}

\begin{corollary}[Expected-normalized signs are uniformly harmless at logarithmic size]
\label{cor:expected-all-signs}
In the setting of Theorem~\ref{thm:expected-full}, fix \(L>0\) and
\(\lambda<\infty\).  If \(n\le\lambda\log d\), then with probability at least
\(1-d^{-L}\),
\[
  \max_{\eps\in\{-1,1\}^n}
  \norm{\sum_{i=1}^n\eps_iA_i}_{\op}
  \le C_{L,\lambda,K}\sqrt n .
\]
The same simultaneous-in-sign statement holds for the corresponding
expected-normalized Hankel model.
\end{corollary}

\begin{proof}
This is the all-sign specialization in
Corollary~\ref{cor:expected-sign-entropy}.  If \(n\le\lambda\log d\), then
\(n/\log(e+d)\le C\lambda\), so the entropy factor is bounded in terms of
\(L,\lambda\) only.
\end{proof}

\begin{remark}\label{rem:edge-low-freq}
Theorem~\ref{thm:expected-full} is why the expected-normalized full random model
should not be listed as open.  It does not contradict the deterministic
worst-case lower bound \(\sqrt{n\log(e+d/n)}\): the latter is attained by
adversarial Toeplitz symbols arranged to peak together, whereas the random
expected-normalized sum is just another random Toeplitz matrix whose symbol has
variance \(n\).  What the theorem does not provide is a pathwise statement for
every tuple of full Toeplitz contractions, nor does it by itself settle the
realization-normalized model of Question~\ref{q:main}, where each matrix is
rescaled by its own random operator norm before the signing problem is posed.
\end{remark}

\subsection{Intrinsic freeness does not certify the square-root bound}
\label{rem:bbvh}

The sharp matrix-concentration theory of Bandeira, Boedihardjo and van Handel
\cite{BBvH} improves on the Khintchine logarithm precisely when the summands
behave as though freely independent. For a self-adjoint Gaussian model
$X=\sum_ig_iA_i$ with i.i.d.\ standard Gaussian $g_i$, their bound has the schematic
form
\begin{equation}\label{eq:bbvh}
\Exp\norm{X}_{\op}\le\norm{X_{\mathrm{free}}}+
C\,v(X)^{1/2}\,\sigma(X)^{1/2}\,(\log d)^{3/4},
\end{equation}
where $\sigma(X)^2=\norm{\Exp X^2}_{\op}$ is the variance parameter,
$\norm{X_{\mathrm{free}}}$ is the operator norm of the free (or matrix-semicircular)
model with the same covariance, and $v(X)$ is the free-deviation parameter, a
covariance norm that measures how far the summands are from being free. The
improvement over Khintchine occurs when $v(X)\ll\sigma(X)$, so that the
$(\log d)^{3/4}$ term is lower order; the free model then carries the main term
with no logarithm.

\begin{proposition}[Toeplitz matrices are maximally delocalized]
\label{prop:bbvh}
Let $A=T_d(\xi)/N_d$ be a random full Toeplitz matrix with unit-variance
entries, $N_d\asymp\sqrt{d\log d}$. Then
\[
\norm{A}_{\HS}\asymp\frac{\sqrt d}{\sqrt{\log d}}\gg1=\norm{A}_{\op}\ (\text{order}).
\]
Consequently, for the Gaussian Toeplitz model $X=\sum_ig_iA_i$ on $n\le d$ such
expected-normalized matrices, the free-deviation parameter and the variance
parameter are of the same order,
\[
v(X)\asymp\sigma(X),
\]
rather than $v(X)\ll\sigma(X)$. The intrinsic-freeness bound \eqref{eq:bbvh}
therefore has its $(\log d)^{3/4}$ correction term at least as large as its main
term, and reduces to the Khintchine scale; it does not certify
$\Exp\norm{X}_{\op}=O(\sqrt n)$.
\end{proposition}

\begin{proof}
The Hilbert--Schmidt norm of an unnormalized $d\times d$ Toeplitz matrix
$T_d(\xi)$ is
\[
\norm{T_d(\xi)}_{\HS}^2=\sum_{0\le r,s<d}\abs{\xi_{r-s}}^2
=\sum_{\abs k<d}(d-\abs k)\abs{\xi_k}^2\asymp d^2
\]
in expectation, since each lag $k$ appears $d-\abs k$ times and the $\xi_k$ have
unit variance; hence $\norm{T_d(\xi)}_{\HS}\asymp d$. Dividing by
$N_d\asymp\sqrt{d\log d}$,
\[
\norm{A}_{\HS}=\frac{\norm{T_d(\xi)}_{\HS}}{N_d}\asymp\frac{d}{\sqrt{d\log d}}
=\frac{\sqrt d}{\sqrt{\log d}}.
\]
An expected-normalized Toeplitz matrix with operator norm of order $1$ and Hilbert--Schmidt norm of order
$\sqrt d/\sqrt{\log d}$ has stable rank $\norm{A}_{\HS}^2/\norm{A}_{\op}^2\asymp
d/\log d$, i.e.\ it is delocalized over almost the full dimension. This is the
opposite of a low-rank perturbation, for which $v\ll\sigma$.

The deviation parameter $v(X)$ in \cite{BBvH} is, up to the normalization there,
controlled from below by the operator norm of the matrix of Hilbert--Schmidt
overlaps of the $A_i$, and for the variance parameter one has
$\sigma(X)^2=\norm{\sum_iA_i^2}_{\op}\asymp n$ since the \(A_i\) have operator
norm of order one and lie in nearly orthogonal Toeplitz directions. The overlap matrix has diagonal entries
$\norm{A_i}_{\HS}^2\asymp d/\log d$, which are far larger than $1$; tracking these
through the definition of $v$ in \cite{BBvH} gives $v(X)^2\gtrsim
\sigma(X)^2$, so $v(X)\asymp\sigma(X)$ rather than $v(X)\ll\sigma(X)$. With
$v\asymp\sigma\asymp\sqrt n$ the correction term in \eqref{eq:bbvh} is
\[
v(X)^{1/2}\sigma(X)^{1/2}(\log d)^{3/4}\asymp\sqrt n\,(\log d)^{3/4},
\]
which exceeds the main term $\norm{X_{\mathrm{free}}}\asymp\sqrt n$ once
$\log d\to\infty$. The bound then reads
$\Exp\norm{X}_{\op}\lesssim\sqrt n\,(\log d)^{3/4}$, and since
$(\log d)^{3/4}>(\log d)^{1/2}$ this is in fact weaker than the plain Khintchine
bound $\sqrt{n\log d}$; in particular it does not give $O(\sqrt n)$.
\end{proof}

\begin{remark}[Numerical check]
The ratio $v(X)/\sigma(X)$ was evaluated directly for signed superpositions of
independent random Toeplitz contractions across a range of dimensions and
sample sizes. It remained bounded away from zero throughout, taking values in the
band $2.5$ to $5.3$, with no decay as $d$ grew. This is consistent with
Proposition~\ref{prop:bbvh}: the parameter count forces $v$ to stay at the scale
of $\sigma$, so the intrinsic-freeness route cannot certify the square-root bound
for this class, and the question is genuinely about typical realizations rather
than about a soft concentration estimate.
\end{remark}

\section{The residual question and its reduction}\label{sec:conjecture}

The expected-normalized theorem of Section~\ref{sec:full} removes the easiest
random model. We now isolate the harder realization-normalized contraction
question and identify the precise analytic obstacle.

\subsection{Why the worst case should not be typical}

The deterministic bound $\sqrt{n\log(e+d/n)}$ is sharp in the worst case: the
matching lower bound of \cite[Thm.~5.30]{Shmalo-paper12} is realized by
circulants whose symbols are arranged, adversarially, so that for every sign
pattern some common frequency carries a large signed sum. Such a configuration is
a measure-zero conspiracy. For independent random symbols the values
$P_i(\omega)$ at a fixed frequency $\omega$ are independent mean-zero scalars of
bounded variance, and the obstruction to $O(\sqrt n)$ is whether, across the
$O(d)$ frequencies simultaneously, the best sign pattern can keep every
$\abs{\sum_i\eps_iP_i(\omega)}$ at the $\sqrt n$ scale. The deterministic theory
already does this for every band up to $w\le n$, i.e.\ for every symbol supported
on $O(n)$ Fourier modes; the residual question is whether the typical
realization-normalized full symbol, although supported on all $d$ modes, has no
worse a signed edge than a banded one.

\subsection{The circulant reduction}

We make the reduction explicit. Recall from \cite[\S5]{Shmalo-paper12} that
signing $d\times d$ circulants is, on the nose, complex vector discrepancy of
their $d$ Fourier symbols: if $C_i$ is the circulant with symbol vector
$v(C_i)=(P_i(\omega))_{\omega}\in\C^d$, where $\omega$ ranges over the $d$th
roots of unity, then
\[
\norm{\sum_{i=1}^n\eps_iC_i}_{\op}
=\max_{\omega}\abs{\sum_{i=1}^n\eps_iP_i(\omega)}
=\norm{\sum_{i=1}^n\eps_iv(C_i)}_\infty .
\]
A full Toeplitz contraction is approximated by a circulant of one larger size up
to a controlled boundary error, the de la Vall\'ee Poussin construction of
\cite[Rem.~5.12]{Shmalo-paper12}, so up to constants the discrepancy of random
Toeplitz contractions is the $\ell^\infty$ vector discrepancy of the random
sample vectors $v(C_i)\in\C^d$.

\begin{proposition}[Reduction to edge control]\label{prop:reduction}
Question~\ref{q:main} has an affirmative answer if the following is true. Let
$v_1,\ldots,v_n\in\C^d$ be the random symbol-sample vectors of independent random
realization-normalized Toeplitz contractions, \(n\le d\le e^{o(n)}\) and
\(d/n\to\infty\). Then with probability tending to one there exist signs
\(\eps\) with
\[
\max_{\omega}\abs{\sum_{i=1}^n\eps_iv_i(\omega)}\le C\sqrt n .
\]
\end{proposition}

\begin{proof}
By the circulant approximation,
\[
  \discop(A_1,\ldots,A_n)
  \le C\max_\omega\abs{\sum_i\eps_iv_i(\omega)}+E,
\]
where \(E\) is the boundary error of the de la Vall\'ee Poussin step, of order
\(o(\sqrt n)\) for the relevant parameters by
\cite[Rem.~5.12]{Shmalo-paper12}. The displayed hypothesis bounds the first
term by \(C\sqrt n\), hence \(\discop=O(\sqrt n)\) on the high-probability event.
\end{proof}

\begin{corollary}[Robust edge-to-operator transfer]
\label{cor:edge-to-operator-robust}
Let \((d_m,n_m)\) be a sequence with
\[
  n_m\le d_m,\qquad d_m/n_m\to\infty,\qquad \log d_m=o(n_m).
\]
For each \(m\), let \(B_{m,i}=T_{d_m}(\xi^{(m,i)})\),
\(1\le i\le n_m\), be full Toeplitz matrices with symbols \(P_{m,i}\), and put
\[
  A_{m,i}:=\frac{B_{m,i}}{\norm{B_{m,i}}_{\op}},
  \qquad
  v_{m,i}(\omega):=\frac{P_{m,i}(\omega)}{\norm{B_{m,i}}_{\op}},
  \qquad \omega^{d_m}=1.
\]
Let \(\mathcal F_m\subset\{-1,1\}^{n_m}\) be any sign family.  If, for some
constant \(L\), all sufficiently large \(m\) satisfy
\[
  \max_{\eps\in\mathcal F_m}\max_{\omega^{d_m}=1}
  \left|\sum_{i=1}^{n_m}\eps_i v_{m,i}(\omega)\right|
  \le L\sqrt{n_m},
\]
then there is \(L'<\infty\), depending only on \(L\) and the deterministic
Toeplitz--circulant reduction constants, such that all sufficiently large \(m\)
satisfy
\[
  \max_{\eps\in\mathcal F_m}
  \left\|\sum_{i=1}^{n_m}\eps_iA_{m,i}\right\|_{\op}
  \le L'\sqrt{n_m}.
\]
The same conclusion holds for the corresponding full Hankel contractions
\(\widetilde A_{m,i}=J_{d_m}A_{m,i}\).  In particular, any Hamming ball of
signs with \(O(\sqrt{n_m})\) edge maximum is an \(O(\sqrt{n_m})\)
operator-norm robust certificate for both the realization-normalized Toeplitz
and Hankel families.
\end{corollary}

\begin{proof}
The proof of Proposition~\ref{prop:reduction} gives, uniformly over sign
vectors,
\[
  \left\|\sum_{i=1}^{n_m}\eps_iA_{m,i}\right\|_{\op}
  \le
  C_{\rm red}\max_{\omega^{d_m}=1}
  \left|\sum_{i=1}^{n_m}\eps_i v_{m,i}(\omega)\right|
  +E_m,
\]
where the de la Vall\'ee Poussin boundary error satisfies
\(E_m=o(\sqrt{n_m})\) under \(d_m/n_m\to\infty\) and \(\log d_m=o(n_m)\).
Taking the maximum over \(\mathcal F_m\) and absorbing \(E_m\) into the
constant gives the Toeplitz claim.  Multiplication by the reversal \(J_{d_m}\)
preserves operator norm and turns full Toeplitz into full Hankel matrices, so
the Hankel claim follows identically.
\end{proof}

\begin{corollary}[Bounded-normalization robust edge-to-operator transfer]
\label{cor:bounded-edge-to-operator-robust}
Let \((d_m,n_m)\) be a sequence with
\[
  n_m\le d_m,\qquad d_m/n_m\to\infty,\qquad \log d_m=o(n_m).
\]
For each \(m\), let \(A_{m,i}\), \(1\le i\le n_m\), be full Toeplitz matrices
with symbols \(Q_{m,i}\), and assume
\[
  \max_{1\le i\le n_m}\norm{A_{m,i}}_{\op}\le K_0
\]
for all sufficiently large \(m\).  Let
\(\mathcal F_m\subset\{-1,1\}^{n_m}\) be any sign family.  If, for some
constant \(L\), all sufficiently large \(m\) satisfy
\[
  \max_{\eps\in\mathcal F_m}\max_{\omega^{d_m}=1}
  \left|\sum_{i=1}^{n_m}\eps_i Q_{m,i}(\omega)\right|
  \le L\sqrt{n_m},
\]
then there is \(L'<\infty\), depending only on \(L\), \(K_0\), and the
deterministic Toeplitz--circulant reduction constants, such that all
sufficiently large \(m\) satisfy
\[
  \max_{\eps\in\mathcal F_m}
  \left\|\sum_{i=1}^{n_m}\eps_iA_{m,i}\right\|_{\op}
  \le L'\sqrt{n_m}.
\]
The same conclusion holds for the corresponding full Hankel matrices
\(\widetilde A_{m,i}=J_{d_m}A_{m,i}\).  Thus every logarithmic Hamming ball
whose edge-grid maximum is \(O(\sqrt{n_m})\) is an actual operator-norm robust
certificate for any uniformly operator-bounded Toeplitz normalization.
\end{corollary}

\begin{proof}
The proof of Proposition~\ref{prop:reduction} is deterministic and uses only
the edge-grid bound and a uniform bound on the individual Toeplitz summands in
the de la Vall\'ee Poussin boundary term.  With
\(\max_i\norm{A_{m,i}}_{\op}\le K_0\), it gives, uniformly over all sign
vectors,
\[
  \left\|\sum_{i=1}^{n_m}\eps_iA_{m,i}\right\|_{\op}
  \le
  C_{\rm red}\max_{\omega^{d_m}=1}
  \left|\sum_{i=1}^{n_m}\eps_i Q_{m,i}(\omega)\right|
  +E_m(K_0),
\]
where \(E_m(K_0)=o_{K_0}(\sqrt{n_m})\) under
\(d_m/n_m\to\infty\) and \(\log d_m=o(n_m)\).  Taking the maximum over
\(\mathcal F_m\) and absorbing the boundary term proves the Toeplitz claim.
Multiplication by \(J_{d_m}\) preserves operator norm and turns full Toeplitz
matrices into full Hankel matrices, so the Hankel claim is identical.
\end{proof}

\begin{corollary}[Borel robust batches transfer to operator norm]
\label{cor:borel-edge-to-operator-batches}
Let \((d_m,n_m)\) satisfy
\[
  n_m\le d_m,\qquad d_m/n_m\to\infty,\qquad \log d_m=o(n_m).
\]
For each \(m\), let \(B_{m,i}=T_{d_m}(\xi^{(m,i)})\),
\(1\le i\le n_m\), be full Toeplitz matrices with symbols \(P_{m,i}\), and put
\[
  A_{m,i}:=\frac{B_{m,i}}{\norm{B_{m,i}}_{\op}},
  \qquad
  \widetilde A_{m,i}:=J_{d_m}A_{m,i},
  \qquad
  v_{m,i}(\omega):=\frac{P_{m,i}(\omega)}{\norm{B_{m,i}}_{\op}} .
\]
For each \(m\), let
\(\delta_m^{(1)},\ldots,\delta_m^{(M_m)}\) be sampled signs, and let
\(1\le r_m\le M_m\).  Suppose that on an event \(\Omega\), all but finitely
many \(m\) have a set
\(S_m\subset\{1,\ldots,M_m\}\) with \(\abs{S_m}\ge M_m-r_m+1\) such that, for
some constants \(L<\infty\) and radii \(q_m\),
\[
  \max_{\omega^{d_m}=1}
  \left|\sum_{i=1}^{n_m}\eps'_i v_{m,i}(\omega)\right|
  \le L\sqrt{n_m}
\]
for every \(s\in S_m\) and every sign vector \(\eps'\) with
\[
  \#\{i:\eps'_i\ne\delta_{m,i}^{(s)}\}\le q_m .
\]
Then, on the same event \(\Omega\), after increasing \(L\) by a deterministic
constant depending only on the Toeplitz--circulant reduction, all but finitely
many \(m\) have the same set \(S_m\) and the same radii \(q_m\) satisfying
\[
  \left\|\sum_{i=1}^{n_m}\eps'_i A_{m,i}\right\|_{\op}
  \le L'\sqrt{n_m},
  \qquad
  \left\|\sum_{i=1}^{n_m}\eps'_i \widetilde A_{m,i}\right\|_{\op}
  \le L'\sqrt{n_m}
\]
for every \(s\in S_m\) and every \(\eps'\) in the same Hamming ball.
Consequently, every almost-sure sampled-count or robust-fraction theorem in
this paper whose conclusion supplies logarithmic-radius \(O(\sqrt{n_m})\)
edge-grid robust centers also supplies actual realization-normalized full
Toeplitz and full Hankel operator-norm robust centers in the above
full-bandwidth range.
\end{corollary}

\begin{proof}
For each sufficiently large \(m\), let \(\mathcal F_m\) be the union of the
Hamming balls centered at \(\{\delta_m^{(s)}:s\in S_m\}\) with radius \(q_m\).
The hypothesis gives
\[
  \max_{\eps'\in\mathcal F_m}\max_{\omega^{d_m}=1}
  \left|\sum_{i=1}^{n_m}\eps'_i v_{m,i}(\omega)\right|
  \le L\sqrt{n_m}.
\]
Apply Corollary~\ref{cor:edge-to-operator-robust} to the sign family
\(\mathcal F_m\).  The Toeplitz estimate follows immediately, and the Hankel
estimate follows from the same reversal argument used there.  The final sentence
is just this deterministic implication applied on the almost-sure event supplied
by any of the sampled-count Borel corollaries.
\end{proof}

General deterministic vector-discrepancy theorems for \(d\) coordinates still
allow the growing \(\sqrt{n\log(e+d/n)}\) scale in this range. Thus
Proposition~\ref{prop:reduction} transfers the remaining problem to the joint
behavior of the \(n\) random edge vectors: one must use randomness of the symbols
to rule out the adversarial common-frequency peaking that makes the deterministic
bound sharp.

\begin{proposition}[Edge-variance profile criterion]\label{prop:variance-profile}
Let \(v_1,\ldots,v_n\in\C^d\) be arbitrary edge-sample vectors and set
\[
  V^2:=\max_{1\le j\le d}\sum_{i=1}^n \abs{v_i(j)}^2 .
\]
If
\[
  V^2\le C_0\,\frac{n}{\log(e+d)},
\]
then for every \(A>0\), independent Rademacher signs satisfy
\[
  \Prob_\eps\left\{
    \max_{1\le j\le d}\abs{\sum_{i=1}^n\eps_i v_i(j)}
    > C_{A,C_0}\sqrt n
  \right\}\le d^{-A}.
\]
Consequently, under this deterministic variance-profile bound, the edge-control
hypothesis in Proposition~\ref{prop:reduction} holds with high probability over
the signs.
\end{proposition}

\begin{proof}
For a fixed coordinate \(j\), the complex Hoeffding inequality gives
\[
  \Prob_\eps\left\{\abs{\sum_i\eps_i v_i(j)}>t\right\}
  \le 4\exp\left(-c\,\frac{t^2}{\sum_i\abs{v_i(j)}^2}\right)
  \le 4\exp\left(-c\,\frac{t^2}{V^2}\right),
\]
by applying the real inequality to the real and imaginary parts. A union bound
over the \(d\) coordinates yields
\[
  \Prob_\eps\left\{\max_j\abs{\sum_i\eps_i v_i(j)}>t\right\}
  \le 4d\exp\left(-c\,\frac{t^2}{V^2}\right).
\]
With \(t=C_{A,C_0}\sqrt n\) and \(V^2\le C_0 n/\log(e+d)\), the right-hand side
is at most \(d^{-A}\) after increasing \(C_{A,C_0}\).
\end{proof}

\begin{corollary}[Random search under an edge profile]
\label{cor:edge-profile-random-search}
With the notation of Proposition~\ref{prop:variance-profile}, let
\(\delta^{(1)},\ldots,\delta^{(M)}\) be independent uniform sign vectors in
\(\{-1,1\}^n\), independent of everything else, where \(M\ge1\) is an integer.
For every \(u\ge0\),
\[
  \Prob_\delta\left\{
  \max_{1\le s\le M}\max_{1\le j\le d}
  \left|\sum_{i=1}^n\delta_i^{(s)}v_i(j)\right|
  >
  C V\sqrt{\log(e+d)+\log(eM)+u}
  \right\}\le e^{-u}.
\]
Consequently, if \(V^2\le C_0n/\log(e+d)\), then with \(\delta\)-probability
at least \(1-e^{-u}\),
\[
  \max_{1\le s\le M}\max_{1\le j\le d}
  \left|\sum_{i=1}^n\delta_i^{(s)}v_i(j)\right|
  \le
  C_{C_0}\sqrt{
    n\left(1+\frac{\log(eM)+u}{\log(e+d)}\right)} .
\]
Moreover, under the same variance-profile hypothesis, for every \(A>0\) and
every integer \(M\ge1\),
\[
  \Prob_\delta\left\{
  \min_{1\le s\le M}\max_{1\le j\le d}
  \left|\sum_{i=1}^n\delta_i^{(s)}v_i(j)\right|
  > C_{A,C_0}\sqrt n
  \right\}\le d^{-AM}.
\]
More generally, if
\[
  G_M:=\#\left\{1\le s\le M:
  \max_{1\le j\le d}
  \left|\sum_{i=1}^n\delta_i^{(s)}v_i(j)\right|
  \le C_{A,C_0}\sqrt n\right\},
\]
then for every integer \(1\le r\le M\),
\[
  \Prob_\delta\{G_M\le M-r\}
  \le \binom{M}{r}\,d^{-Ar}.
\]
Thus, once the matrix realization has the edge profile, \(M\) independent
random-sign trials amplify the failure probability from \(d^{-A}\) to
\(d^{-AM}\), and the number of sampled certificates has the corresponding
binomial upper tail for failures.
\end{corollary}

\begin{proof}
The proof of Proposition~\ref{prop:variance-profile} gives, for one sign vector
\(\delta\),
\[
  \Prob_\delta\left\{
  \max_j\left|\sum_i\delta_i v_i(j)\right|>t
  \right\}
  \le 4d\exp(-ct^2/V^2).
\]
Taking a union bound over \(s=1,\ldots,M\) gives
\[
  \Prob_\delta\left\{
  \max_s\max_j\left|\sum_i\delta_i^{(s)}v_i(j)\right|>t
  \right\}
  \le 4dM\exp(-ct^2/V^2).
\]
With
\[
  t=C V\sqrt{\log(e+d)+\log(eM)+u},
\]
the last expression is at most \(e^{-u}\) after increasing the absolute
constant \(C\).  Substituting \(V^2\le C_0n/\log(e+d)\) gives the displayed
entropy-sensitive sampled-list bound.

For the final assertions, Proposition~\ref{prop:variance-profile} says that a
single random sign vector fails the \(C_{A,C_0}\sqrt n\) threshold with
probability at most \(d^{-A}\).  The \(M\) trials are independent, so the
probability that all \(M\) fail is at most \(d^{-AM}\).
If \(F_s\) is the failure event for trial \(s\), then for every fixed set
\(R\subset\{1,\ldots,M\}\) of size \(r\),
\[
  \Prob_\delta\left(\bigcap_{s\in R}F_s\right)\le d^{-Ar}.
\]
The event \(G_M\le M-r\) means that at least \(r\) of the failure events occur.
A union bound over the \(r\)-subsets \(R\) gives
\(\binom{M}{r}d^{-Ar}\).
\end{proof}

\begin{proposition}[Edge-profile scalar moment criterion]
\label{prop:edge-moment-profile}
Let \(v_1,\ldots,v_n\in\C^d\) and set
\[
  V^2:=\max_{1\le j\le d}\sum_{i=1}^n\abs{v_i(j)}^2 .
\]
Then, simultaneously for every integer \(k\ge1\),
\[
  \max_{1\le j\le d}
  \Exp_\eps\left|\sum_{i=1}^n\eps_i v_i(j)\right|^{2k}
  \le (CkV^2)^k .
\]
In particular, if \(V^2\le C_0n/\log(e+d)\), then
\[
  \max_{1\le j\le d}
  \Exp_\eps\left|\sum_{i=1}^n\eps_i v_i(j)\right|^{2k}
  \le
  \left(C_{C_0}k\,\frac n{\log(e+d)}\right)^k
  \qquad(k\ge1).
\]
\end{proposition}

\begin{proof}
For deterministic complex scalars \(a_i\), the standard Rademacher moment bound
gives
\[
  \left(\Exp_\eps\left|\sum_i\eps_i a_i\right|^{2k}\right)^{1/(2k)}
  \le C\sqrt{k}\left(\sum_i|a_i|^2\right)^{1/2}.
\]
Indeed, apply the real sub-Gaussian moment bound to the real and imaginary
parts and use
\(|z|\le|\operatorname{Re}z|+|\operatorname{Im}z|\).  Taking
\(a_i=v_i(j)\), maximizing over \(j\), and raising to the power \(2k\) proves
the first display; the second follows from the displayed bound on \(V^2\).
\end{proof}

\begin{corollary}[Edge-profile moment bound for the edge maximum]
\label{cor:edge-max-moment-profile}
With the notation of Proposition~\ref{prop:edge-moment-profile}, put
\[
  M_\eps:=\max_{1\le j\le d}\left|\sum_{i=1}^n\eps_i v_i(j)\right|.
\]
Then, for every integer \(k\ge1\),
\[
  \Exp_\eps M_\eps^{2k}\le d\,(CkV^2)^k,
\]
and hence
\[
  \left(\Exp_\eps M_\eps^{2k}\right)^{1/(2k)}
  \le C\,d^{1/(2k)}\sqrt{k}\,V .
\]
In particular, if \(V^2\le C_0n/\log(e+d)\) and
\(k\ge\log(e+d)\), then
\[
  \left(\Exp_\eps M_\eps^{2k}\right)^{1/(2k)}
  \le C_{C_0}\sqrt n\,\sqrt{\frac{k}{\log(e+d)}} .
\]
Taking \(k=\lceil(A+2)\log(e+d)\rceil\) and applying Markov's inequality
recovers the high-probability edge-control conclusion of
Proposition~\ref{prop:variance-profile}.
\end{corollary}

\begin{proof}
Since \(M_\eps^{2k}\le\sum_{j=1}^d|\sum_i\eps_i v_i(j)|^{2k}\), the first
display follows from Proposition~\ref{prop:edge-moment-profile}.  Taking
\(2k\)-th roots gives the second display.  If \(k\ge\log(e+d)\), then
\(d^{1/(2k)}\le e^{1/2}\), and the displayed bound on \(V^2\) gives the third
display.  For \(k=\lceil(A+2)\log(e+d)\rceil\), the third display is at most
\(C_{A,C_0}\sqrt n\); increasing the constant in the threshold and using
Markov's inequality gives a \(d^{-A}\) tail bound.
\end{proof}

\begin{corollary}[Edge-profile signing certificate]
\label{cor:edge-profile-certificate}
With the notation of Proposition~\ref{prop:edge-moment-profile},
\[
  \Exp_\eps\max_{1\le j\le d}\left|\sum_{i=1}^n\eps_i v_i(j)\right|
  \le C\,V\sqrt{\log(e+d)} .
\]
Consequently,
\[
  \min_{\eps\in\{-1,1\}^n}
  \max_{1\le j\le d}\left|\sum_{i=1}^n\eps_i v_i(j)\right|
  \le C\,V\sqrt{\log(e+d)} .
\]
In particular, if \(V^2\le C_0n/\log(e+d)\), then there is a sign vector
\(\eps\in\{-1,1\}^n\) for which
\[
  \max_{1\le j\le d}\left|\sum_{i=1}^n\eps_i v_i(j)\right|
  \le C_{C_0}\sqrt n .
\]
\end{corollary}

\begin{proof}
Let \(M_\eps\) be the edge maximum from
Corollary~\ref{cor:edge-max-moment-profile} and put
\(k_0=\lceil\log(e+d)\rceil\).  Since \(M_\eps\ge0\),
\[
  \Exp_\eps M_\eps\le
  \left(\Exp_\eps M_\eps^{2k_0}\right)^{1/(2k_0)}
  \le C\,d^{1/(2k_0)}\sqrt{k_0}\,V
  \le C\,V\sqrt{\log(e+d)} .
\]
The minimum over signs is bounded by the average over signs.  The final display
follows by substituting \(V^2\le C_0n/\log(e+d)\).
\end{proof}

\begin{corollary}[Hamming-stable edge-profile certificates]
\label{cor:edge-profile-hamming}
With the notation of Proposition~\ref{prop:edge-moment-profile}, suppose a sign
vector \(\eps\in\{-1,1\}^n\) satisfies
\[
  \max_{1\le j\le d}\left|\sum_{i=1}^n\eps_i v_i(j)\right|\le R .
\]
If \(\eps'\in\{-1,1\}^n\) differs from \(\eps\) in at most \(r\) coordinates,
then
\[
  \max_{1\le j\le d}\left|\sum_{i=1}^n\eps'_i v_i(j)\right|
  \le R+2\sqrt r\,V .
\]
Consequently, if \(V^2\le C_0n/\log(e+d)\), then for every \(A>0\), with
Rademacher probability at least \(1-d^{-A}\) a random sign vector \(\eps\) is
the center of a Hamming ball of radius
\[
  r_0=\left\lfloor c_{A,C_0}\log(e+d)\right\rfloor
\]
on which every sign vector \(\eps'\) satisfies
\[
  \max_{1\le j\le d}\left|\sum_{i=1}^n\eps'_i v_i(j)\right|
  \le C_{A,C_0}\sqrt n .
\]
Equivalently, at least a \(1-d^{-A}\) fraction of all sign vectors are centers
of such logarithmic-radius good Hamming balls.
Moreover, if \(\delta^{(1)},\ldots,\delta^{(M)}\) are independent uniform sign
vectors and \(H_M\) is the number of sampled signs that are centers of such
radius-\(r_0\) good Hamming balls, then for every integer \(1\le r\le M\),
\[
  \Prob_\delta\{H_M\le M-r\}\le \binom{M}{r}\,d^{-Ar}.
\]
\end{corollary}

\begin{proof}
Let \(F=\{i:\eps_i\ne\eps'_i\}\), so \(|F|\le r\).  For each coordinate \(j\),
\[
  \sum_i\eps'_iv_i(j)-\sum_i\eps_iv_i(j)
  =\sum_{i\in F}(\eps'_i-\eps_i)v_i(j).
\]
Since \(|\eps'_i-\eps_i|\le2\), Cauchy's inequality gives
\[
  \left|\sum_{i\in F}(\eps'_i-\eps_i)v_i(j)\right|
  \le 2\sum_{i\in F}|v_i(j)|
  \le 2\sqrt r\left(\sum_{i\in F}|v_i(j)|^2\right)^{1/2}
  \le 2\sqrt r\,V.
\]
Taking the maximum over \(j\) proves the deterministic stability estimate.

Under \(V^2\le C_0n/\log(e+d)\), Proposition~\ref{prop:variance-profile} gives
with sign probability at least \(1-d^{-A}\) a center \(\eps\) satisfying
\[
  \max_j\left|\sum_i\eps_i v_i(j)\right|\le C_{A,C_0}\sqrt n .
\]
Choose \(c_{A,C_0}>0\) small enough that
\[
  2\sqrt{r_0}\,V\le C'_{A,C_0}\sqrt n .
\]
The first part then gives the claimed bound on the whole Hamming ball after
enlarging \(C_{A,C_0}\).  Since Rademacher probability is uniform measure on the
sign cube, the final fraction statement is the same assertion.

For the sampled-center assertion, let \(E_s\) be the event that
\(\delta^{(s)}\) is not the center of such a good Hamming ball.  The preceding
paragraph gives \(\Prob(E_s)\le d^{-A}\).  For every fixed
\(R\subset\{1,\ldots,M\}\) of size \(r\), independence gives
\[
  \Prob_\delta\left(\bigcap_{s\in R}E_s\right)\le d^{-Ar}.
\]
If \(H_M\le M-r\), at least \(r\) of the events \(E_s\) occur.  A union bound
over the \(r\)-subsets \(R\) gives the displayed estimate.
\end{proof}

\begin{corollary}[Polynomial-confidence robust signing search]
\label{cor:edge-profile-poly-confidence}
With the notation of Proposition~\ref{prop:variance-profile}, assume
\[
  V^2\le C_0\frac n{\log(e+d)} .
\]
Then there are constants \(C_{C_0}\) and \(c_{C_0}>0\) such that the following
holds.  For every \(B>0\), if
\[
  M\ge \lceil B\rceil
\]
independent uniform sign vectors
\(\delta^{(1)},\ldots,\delta^{(M)}\in\{-1,1\}^n\) are sampled, then with
probability at least \(1-d^{-B}\) at least one sampled sign vector is the center
of a Hamming ball of radius
\[
  r_0=\left\lfloor c_{C_0}\log(e+d)\right\rfloor
\]
on which every sign vector \(\eps'\) satisfies
\[
  \max_{1\le j\le d}\left|\sum_{i=1}^n\eps'_i v_i(j)\right|
  \le C_{C_0}\sqrt n .
\]
Thus a constant-size random list gives a robust \(O(\sqrt n)\) signing
certificate with arbitrary polynomial confidence.
\end{corollary}

\begin{proof}
Apply Corollary~\ref{cor:edge-profile-hamming} with exponent \(A=1\).  A single
uniform sign fails to be the center of the displayed good Hamming ball with
probability at most \(d^{-1}\).  The \(M\) trials are independent, so the
probability that none of them is such a robust center is at most
\[
  d^{-M}\le d^{-B}.
\]
\end{proof}

\begin{remark}[The remaining random input]\label{rem:variance-profile}
For the realization-normalized Toeplitz contraction model in
Question~\ref{q:main}, the vectors in Proposition~\ref{prop:reduction} have the
form
\[
  v_i(\omega)=\frac{P_i(\omega)}{\norm{T_d(\xi^{(i)})}_{\op}},
  \qquad \omega^d=1 .
\]
Thus Proposition~\ref{prop:variance-profile} reduces the residual question to
the empirical edge-delocalization estimate
\[
  \max_{\omega^d=1}
  \sum_{i=1}^n
  \frac{\abs{P_i(\omega)}^2}
       {\norm{T_d(\xi^{(i)})}_{\op}^2}
  \;\lesssim\; \frac{n}{\log d}
\]
with high probability.  This is stronger and more targeted than asking a
generic vector-discrepancy theorem to beat the worst case: it says that, after
realization normalization, no edge node should carry more than its natural
average share of the \(n\) random symbols.  The expected-normalized theorem
above is the same heuristic with the random denominators replaced by the
deterministic scale \(N_d\).  Theorem~\ref{thm:gaussian-realized-profile}
removes the random-denominator difficulty for standard Gaussian coefficients by
combining Gaussian norm concentration with a bad-normalizer truncation.  The
general form of that argument is Theorem~\ref{thm:normalizer-tail}: apart from
standard sub-Gaussian root-value bounds and a crude Hilbert--Schmidt floor, the
only missing one-matrix input is
\[
  \Prob\{\norm{T_d(\xi)}_{\op}<a\,\Exp\norm{T_d(\xi)}_{\op}\}
  \min\{n,\log^2(e+d)\}=o(1).
\]
Thus the residual problem is no longer an arbitrary dependence question between
edge values and normalizers; it is enough to rule out a logarithmic-density
population of low-normalizer samples.
\end{remark}

\begin{proposition}[Expected edge profile]\label{prop:expected-profile}
In the expected-normalized model of Theorem~\ref{thm:expected-full}, let
\[
  v_i(\omega):=\frac{P_i(\omega)}{N_d},\qquad \omega^d=1,
\]
where \(P_i\) is the full Toeplitz symbol of \(B_i\).  For every \(A>0\), if
\(n\ge C_{A,K}\log(e+d)\), then
\[
  \Prob\left\{
    \max_{\omega^d=1}\sum_{i=1}^n \abs{v_i(\omega)}^2
    > C_{A,K}\frac{n}{\log(e+d)}
  \right\}\le d^{-A}.
\]
Thus the edge-variance profile required in
Proposition~\ref{prop:variance-profile} holds automatically in the
expected-normalized model, at the natural \(n/\log d\) scale.
\end{proposition}

\begin{proof}
Fix a \(d\)th root of unity \(\omega\).  The random variable
\[
  P_i(\omega)=\sum_{\abs{k}<d}\xi_k^{(i)}\omega^k
\]
is centered sub-Gaussian with norm at most \(C_K\sqrt d\), hence
\(\abs{P_i(\omega)}^2\) is sub-exponential with norm at most \(C_K d\) and mean
at most \(C_K d\).  Since \(N_d\ge c_K\sqrt{d\log(e+d)}\), the variables
\(\abs{v_i(\omega)}^2\) are independent nonnegative sub-exponential variables
with mean and sub-exponential norm \(O_K(1/\log(e+d))\). Bernstein's inequality
therefore gives
\[
  \Prob\left\{
    \sum_{i=1}^n\abs{v_i(\omega)}^2
    > C_{A,K}\frac{n}{\log(e+d)}
  \right\}\le \exp(-c_{A,K}n)
\]
after increasing \(C_{A,K}\).  If \(n\ge C_{A,K}\log(e+d)\), the right side is
at most \(d^{-A-1}\).  A union bound over the \(d\) roots gives the claim.
\end{proof}

\begin{theorem}[Normalizer-tail criterion]\label{thm:normalizer-tail}
Let \(B_i=T_d(\xi^{(i)})\), \(1\le i\le n\), be independent full Toeplitz
matrices whose lag variables are independent, mean zero, unit variance, and have
sub-Gaussian norm at most \(K\).  Let \(P_i\) be their symbols and put
\[
  N_d:=\Exp\norm{T_d(\xi)}_{\op}.
\]
Assume \(n\le d\), \(\log d=o(n)\), and that for some \(a>0\)
\[
  p_d(a):=\Prob\{\norm{T_d(\xi)}_{\op}<aN_d\}
  \quad\text{satisfies}\quad
  p_d(a)\min\{n,\log^2(e+d)\}=o(1).
\]
Then
\[
  \max_{\omega^d=1}
  \sum_{i=1}^n
  \frac{\abs{P_i(\omega)}^2}{\norm{B_i}_{\op}^2}
  \le C_{K,a}\,\frac{n}{\log(e+d)}
\]
with probability tending to one.  Consequently the realization-normalized edge
control in Proposition~\ref{prop:reduction} holds with high probability over
the matrices and independent signs.  On the same high-probability matrix event,
\[
  \Exp_\eps\max_{\omega^d=1}\left|\sum_{i=1}^n
  \eps_i\frac{P_i(\omega)}{\norm{B_i}_{\op}}\right|
  \le C_{K,a}\sqrt n,
\]
so in particular a deterministic sign vector with \(O_{K,a}(\sqrt n)\) edge
maximum exists.  Also on this event, simultaneously for every integer \(k\ge1\),
\[
  \max_{\omega^d=1}
  \Exp_\eps\left|\sum_{i=1}^n
  \eps_i\frac{P_i(\omega)}{\norm{B_i}_{\op}}\right|^{2k}
  \le
  \left(C_{K,a}k\,\frac n{\log(e+d)}\right)^k .
\]
If also \(d/n\to\infty\), then
\[
  \discop(B_1/\norm{B_1}_{\op},\ldots,B_n/\norm{B_n}_{\op})
  =O_{K,a}(\sqrt n)
\]
with probability tending to one.
\end{theorem}

\begin{proof}
Write \(\ell=\log(e+d)\), and split
\[
  G=\{i:\norm{B_i}_{\op}\ge aN_d\},\qquad G^c=\{1,\ldots,n\}\setminus G .
\]
On the event of Proposition~\ref{prop:expected-profile},
\[
  \max_{\omega^d=1}\sum_{i\in G}
  \frac{\abs{P_i(\omega)}^2}{\norm{B_i}_{\op}^2}
  \le a^{-2}\max_{\omega^d=1}\sum_{i=1}^n
  \frac{\abs{P_i(\omega)}^2}{N_d^2}
  \le C_{K,a}\frac n\ell ,
\]
with probability tending to one, since \(\ell=o(n)\).

The number \(M=\abs{G^c}\) of bad normalizers satisfies
\[
  \Prob\{M>n/\ell^2\}\le p_d(a)\min\{n,\ell^2\}=o(1).
\]
For \(n\ge\ell^2\) this is Markov's inequality; for \(n<\ell^2\) the event
forces at least one bad normalizer and is bounded by \(np_d(a)\).
It remains to bound the cost of those bad samples.  A union bound over
\(i\le n\le d\) and the \(d\) roots gives
\[
  \max_{i,\omega^d=1}\abs{P_i(\omega)}
  \le C_K\sqrt{d\ell}
\]
with probability tending to one, because each \(P_i(\omega)\) has
sub-Gaussian norm at most \(C_K\sqrt d\).

We also need a denominator floor which is much weaker than the desired
operator-norm scale.  Since \(\Exp\abs{\xi}^2=1\) and
\(\Exp\abs{\xi}^4\le C_K\), Paley--Zygmund gives constants \(c_K,\tau_K>0\)
such that \(\Prob\{\abs{\xi}^2\ge\tau_K\}\ge c_K\).  For \(\abs{k}\le d/2\)
the Toeplitz weight \(d-\abs{k}\) is at least \(d/2\), so a Chernoff bound gives
\[
  \norm{B_i}_{\HS}^2
  =\sum_{\abs k<d}(d-\abs k)\abs{\xi_k^{(i)}}^2
  \ge c_K d^2
\]
for all \(i\le n\), with probability tending to one.  Therefore
\[
  \norm{B_i}_{\op}\ge d^{-1/2}\norm{B_i}_{\HS}\ge c_K\sqrt d
\]
uniformly in \(i\le n\) on the same event.

On the intersection of these events, every bad sample contributes at most
\[
  \frac{\abs{P_i(\omega)}^2}{\norm{B_i}_{\op}^2}\le C_K\ell
\]
to each root \(\omega\).  Since \(M\le n/\ell^2\), the total bad contribution is
at most \(C_K n/\ell\), uniformly over the roots.  Combining good and bad
samples proves the displayed profile.

The signed edge-control statement is Proposition~\ref{prop:variance-profile}
applied on this high-probability matrix event.  The conditional expectation and
existence statement is Corollary~\ref{cor:edge-profile-certificate}.  The final
discrepancy bound follows from Proposition~\ref{prop:reduction}.
The scalar moment estimate is Proposition~\ref{prop:edge-moment-profile}
applied to the same profile.
\end{proof}

\begin{corollary}[Data-dependent bad-normalizer decomposition]
\label{cor:normalizer-tail-datadep}
Fix \(A>0\) and \(a>0\).  Let \(B_i=T_d(\xi^{(i)})\),
\(1\le i\le n\), be independent full Toeplitz matrices whose lag variables are
independent, mean zero, unit variance, and have sub-Gaussian norm at most \(K\).
Let \(P_i\) be their symbols, put
\(N_d=\Exp\norm{T_d(\xi)}_{\op}\), write \(\ell=\log(e+d)\), and set
\[
  M_a:=\#\{1\le i\le n:\norm{B_i}_{\op}<aN_d\}.
\]
There is a constant \(C_{A,K,a}\) such that, whenever
\[
  n\le d,\qquad n\ge C_{A,K,a}\ell ,
\]
one has, outside an event of probability at most \(C_{A,K,a}d^{-A}\),
\[
  \max_{\omega^d=1}
  \sum_{i=1}^n
  \frac{\abs{P_i(\omega)}^2}{\norm{B_i}_{\op}^2}
  \le
  C_{A,K,a}\left(\frac n\ell+M_a\ell\right).
\]
On the same matrix event, independent signs satisfy
\[
  \Prob_\eps\left\{
  \max_{\omega^d=1}\left|\sum_{i=1}^n
  \eps_i\frac{P_i(\omega)}{\norm{B_i}_{\op}}\right|
  > C_{A,K,a}\sqrt{\,n+M_a\ell^2\,}
  \right\}\le d^{-A},
\]
and
\[
  \Exp_\eps\max_{\omega^d=1}\left|\sum_{i=1}^n
  \eps_i\frac{P_i(\omega)}{\norm{B_i}_{\op}}\right|
  \le C_{A,K,a}\sqrt{\,n+M_a\ell^2\,}.
\]
Thus a deterministic sign vector with edge maximum
\(C_{A,K,a}\sqrt{n+M_a\ell^2}\) exists.  On the same event, simultaneously for
every integer \(k\ge1\),
\[
  \max_{\omega^d=1}
  \Exp_\eps\left|\sum_{i=1}^n
  \eps_i\frac{P_i(\omega)}{\norm{B_i}_{\op}}\right|^{2k}
  \le
  \left(C_{A,K,a}k\left(\frac n\ell+M_a\ell\right)\right)^k .
\]
\end{corollary}

\begin{proof}
Let
\[
  G=\{i:\norm{B_i}_{\op}\ge aN_d\}.
\]
The expected-normalized edge-profile event of
Proposition~\ref{prop:expected-profile} gives
\[
  \max_{\omega^d=1}\sum_{i\in G}
  \frac{\abs{P_i(\omega)}^2}{\norm{B_i}_{\op}^2}
  \le C_{A,K,a}\frac n\ell
\]
outside an event of probability at most \(d^{-A}\), after increasing the
constant and using \(n\ge C_{A,K,a}\ell\).

The same union bound and Paley--Zygmund/Chernoff Hilbert--Schmidt floor used in
the proof of Theorem~\ref{thm:normalizer-tail} give, outside another
\(C_{A,K}d^{-A}\) event,
\[
  \max_{i,\omega^d=1}\abs{P_i(\omega)}\le C_{A,K}\sqrt{d\ell},
  \qquad
  \min_i\norm{B_i}_{\op}\ge c_K\sqrt d .
\]
Hence every index outside \(G\) contributes at most \(C_{A,K}\ell\) to each
edge node.  Since there are exactly \(M_a\) such indices, the bad-normalizer
contribution is at most \(C_{A,K}M_a\ell\).  Combining the good and bad
contributions proves the data-dependent profile.

Condition on matrices satisfying that profile and put
\[
  V^2:=C_{A,K,a}\left(\frac n\ell+M_a\ell\right).
\]
The Hoeffding-and-union-bound proof of Proposition~\ref{prop:variance-profile}
gives
\[
  \Prob_\eps\left\{
  \max_{\omega^d=1}\left|\sum_i\eps_i
  \frac{P_i(\omega)}{\norm{B_i}_{\op}}\right|>t
  \right\}
  \le 4d\exp(-ct^2/V^2).
\]
Taking \(t=C_{A,K,a}\sqrt{V^2\ell}\) gives the displayed tail, since
\(V^2\ell\le C_{A,K,a}(n+M_a\ell^2)\).  The sign-averaged maximum and
deterministic certificate follow from
Corollary~\ref{cor:edge-profile-certificate}; the scalar moment estimates follow
from Proposition~\ref{prop:edge-moment-profile}.
\end{proof}

\begin{corollary}[Dyadic bad-normalizer decomposition]
\label{cor:normalizer-tail-dyadic}
Fix \(A>0\) and \(a>0\).  In the setting of
Corollary~\ref{cor:normalizer-tail-datadep}, put
\(\ell=\log(e+d)\) and
\[
  J_\ell:=\left\lceil \frac12\log_2(e+\ell)\right\rceil .
\]
For \(0\le r<J_\ell\), define
\[
  M_{a,r}:=
  \#\left\{1\le i\le n:
  2^{-(r+1)}aN_d\le \norm{B_i}_{\op}<2^{-r}aN_d\right\},
\]
and put
\[
  M_{a,\infty}:=
  \#\left\{1\le i\le n:
  \norm{B_i}_{\op}<2^{-J_\ell}aN_d\right\}.
\]
Set
\[
  W_a:=\sum_{0\le r<J_\ell}4^{r+1}M_{a,r}
       +\ell\,M_{a,\infty}.
\]
There is a constant \(C_{A,K,a}\) such that, whenever
\[
  n\le d,\qquad n\ge C_{A,K,a}\ell ,
\]
one has, outside an event of probability at most \(C_{A,K,a}d^{-A}\),
\[
  \max_{\omega^d=1}
  \sum_{i=1}^n
  \frac{\abs{P_i(\omega)}^2}{\norm{B_i}_{\op}^2}
  \le
  C_{A,K,a}\left(\frac n\ell+W_a\right).
\]
On the same matrix event, independent signs satisfy
\[
  \Prob_\eps\left\{
  \max_{\omega^d=1}\left|\sum_{i=1}^n
  \eps_i\frac{P_i(\omega)}{\norm{B_i}_{\op}}\right|
  >
  C_{A,K,a}\sqrt{\,n+W_a\ell\,}
  \right\}\le d^{-A},
\]
and
\[
  \Exp_\eps\max_{\omega^d=1}\left|\sum_{i=1}^n
  \eps_i\frac{P_i(\omega)}{\norm{B_i}_{\op}}\right|
  \le C_{A,K,a}\sqrt{\,n+W_a\ell\,}.
\]
Thus a deterministic sign vector with edge maximum
\(C_{A,K,a}\sqrt{n+W_a\ell}\) exists.  On the same event, simultaneously for
every integer \(k\ge1\),
\[
  \max_{\omega^d=1}
  \Exp_\eps\left|\sum_{i=1}^n
  \eps_i\frac{P_i(\omega)}{\norm{B_i}_{\op}}\right|^{2k}
  \le
  \left(C_{A,K,a}k\left(\frac n\ell+W_a\right)\right)^k .
\]
In particular, if the realized dyadic bad-normalizer budget satisfies
\[
  W_a\le R\frac n\ell
\]
for some \(R\ge1\), then the realization-normalized edge profile, random-sign
tail, sign-averaged maximum, deterministic certificate, and scalar moment
package all hold at the \(C_{A,K,a,R}\sqrt n\) scale.
\end{corollary}

\begin{proof}
Let
\[
  G=\{i:\norm{B_i}_{\op}\ge aN_d\}.
\]
As in Corollary~\ref{cor:normalizer-tail-datadep}, the expected-normalized
profile event gives
\[
  \max_{\omega^d=1}\sum_{i\in G}
  \frac{\abs{P_i(\omega)}^2}{\norm{B_i}_{\op}^2}
  \le C_{A,K,a}\frac n\ell
\]
outside a \(d^{-A}\) event.  Intersect this with the standard root-value and
Hilbert--Schmidt floor event from the proof of
Theorem~\ref{thm:normalizer-tail}:
\[
  \max_{i,\omega^d=1}\abs{P_i(\omega)}\le C_{A,K}\sqrt{d\ell},
  \qquad
  \min_i\norm{B_i}_{\op}\ge c_K\sqrt d .
\]
Using \(N_d\ge c_K\sqrt{d\ell}\), the first of these bounds implies
\[
  \max_{i,\omega^d=1}\frac{\abs{P_i(\omega)}^2}{N_d^2}
  \le C_{A,K}.
\]
If \(i\) lies in the \(r\)-th dyadic bin, then
\[
  \frac{\abs{P_i(\omega)}^2}{\norm{B_i}_{\op}^2}
  \le C_{A,K,a}\,4^{r+1}
  \qquad(\omega^d=1).
\]
The final bin is controlled by the Hilbert--Schmidt floor instead:
\[
  \frac{\abs{P_i(\omega)}^2}{\norm{B_i}_{\op}^2}
  \le C_{A,K}\ell .
\]
Summing over all dyadic bins gives the displayed profile bound with \(W_a\).

Condition on any matrix realization satisfying this profile and put
\[
  V^2:=C_{A,K,a}\left(\frac n\ell+W_a\right).
\]
The Hoeffding-and-union-bound proof of Proposition~\ref{prop:variance-profile}
gives
\[
  \Prob_\eps\left\{
  \max_{\omega^d=1}\left|\sum_i\eps_i
  \frac{P_i(\omega)}{\norm{B_i}_{\op}}\right|>t
  \right\}
  \le 4d\exp(-ct^2/V^2).
\]
Taking \(t=C_{A,K,a}\sqrt{V^2\ell}\) yields the displayed sign tail, since
\(V^2\ell\le C_{A,K,a}(n+W_a\ell)\).  The expectation and deterministic
certificate follow from Corollary~\ref{cor:edge-profile-certificate}, and the
scalar moment estimates follow from Proposition~\ref{prop:edge-moment-profile}.
The final assertion is the specialization \(W_a\le Rn/\ell\).
\end{proof}

\begin{corollary}[Probabilistic dyadic-tail criterion]
\label{cor:normalizer-tail-dyadic-prob}
Fix \(A>0\), \(a>0\), and \(R\ge1\).  In the setting of
Corollary~\ref{cor:normalizer-tail-dyadic}, let \(Y_{d,a}\) be the one-matrix
dyadic depth variable
\[
\begin{aligned}
  Y_{d,a}:={}&\sum_{0\le r<J_\ell}4^{r+1}
  \mathbf 1\!\left\{
  2^{-(r+1)}aN_d\le \norm{T_d(\xi)}_{\op}<2^{-r}aN_d
  \right\}  \\
  &+\ell\,
  \mathbf 1\!\left\{\norm{T_d(\xi)}_{\op}<2^{-J_\ell}aN_d\right\},
\end{aligned}
\]
and put
\[
  \mu_{d,a}^{\rm dy}:=\Exp Y_{d,a}.
\]
Then, whenever \(n\le d\) and \(n\ge C_{A,K,a}\ell\),
\[
  \Prob\left\{
  \max_{\omega^d=1}
  \sum_{i=1}^n
  \frac{\abs{P_i(\omega)}^2}{\norm{B_i}_{\op}^2}
  > C_{A,K,a,R}\frac n\ell
  \right\}
  \le
  C_{A,K,a,R}d^{-A}
  +\frac{\ell\,\mu_{d,a}^{\rm dy}}{R}.
\]
On the complement of this event, the random-sign tail, sign-averaged maximum,
deterministic certificate, and scalar edge-grid moment package of
Corollary~\ref{cor:normalizer-tail-dyadic} hold at the
\(C_{A,K,a,R}\sqrt n\) scale.  In particular, if
\(\mu_{d,a}^{\rm dy}=o(1/\ell)\), then the realization-normalized edge profile
holds with probability tending to one; if also \(d/n\to\infty\), the
corresponding full Toeplitz/Hankel contractions have \(O_{K,a,R}(\sqrt n)\)
discrepancy with probability tending to one.

There is also a large-sample exponential form.  There is an absolute
\(c_{R,\theta}>0\) such that, if \(0<\theta<R\) and
\[
  \mu_{d,a}^{\rm dy}\le \frac{\theta}{\ell},
\]
then the preceding exceptional probability can be replaced by
\[
  C_{A,K,a,R}d^{-A}
  +\exp\left(-c_{R,\theta}\frac n{\ell^2}\right).
\]
Thus the dyadic lower-tail input gives the same Chernoff-scale profile
probability as the single-cutoff theory, but with the low-normalizer mass
weighted by its actual dyadic depth.
\end{corollary}

\begin{proof}
For independent matrices \(B_i=T_d(\xi^{(i)})\), the realized budget in
Corollary~\ref{cor:normalizer-tail-dyadic} is
\[
  W_a=\sum_{i=1}^n Y_{d,a}^{(i)}
\]
with \(Y_{d,a}^{(i)}\) independent copies of \(Y_{d,a}\).  Hence
\[
  \Prob\left\{W_a>R\frac n\ell\right\}
  \le \frac{\ell\,\Exp W_a}{Rn}
  =\frac{\ell\,\mu_{d,a}^{\rm dy}}{R}
\]
by Markov's inequality.  On the event \(W_a\le Rn/\ell\),
Corollary~\ref{cor:normalizer-tail-dyadic} gives the displayed profile and the
signing, expectation, certificate, and scalar-moment conclusions outside only
the standard \(C_{A,K,a}d^{-A}\) matrix event.  Proposition~\ref{prop:reduction}
and the Toeplitz--Hankel reversal give the final discrepancy consequence.

For the exponential form, first note that
\[
  Y_{d,a}\le C\ell .
\]
Indeed \(J_\ell=\lceil\frac12\log_2(e+\ell)\rceil\), so
\(4^{J_\ell}\le C(e+\ell)\le C'\ell\), after increasing the absolute constant
to cover bounded \(\ell\).  If \(\mu_{d,a}^{\rm dy}\le\theta/\ell\), then
\[
  \Exp W_a\le \theta\frac n\ell .
\]
Bernstein's inequality for independent variables bounded by \(C\ell\), with
\[
  \sum_i\operatorname{Var}(Y_{d,a}^{(i)})
  \le C\ell\,\Exp W_a\le C\theta n,
\]
gives
\[
  \Prob\left\{W_a>R\frac n\ell\right\}
  \le
  \exp\left(
  -c\frac{(R-\theta)^2n^2/\ell^2}
          {\theta n+(R-\theta)n}
  \right)
  \le
  \exp\left(-c_{R,\theta}\frac n{\ell^2}\right).
\]
Combining this with the \(d^{-A}\) matrix event proves the last assertion.
\end{proof}

\begin{corollary}[Inverse-moment dyadic-tail criterion]
\label{cor:normalizer-tail-invmoment}
Fix \(A>0\), \(a>0\), and \(R\ge1\).  In the setting of
Corollary~\ref{cor:normalizer-tail-dyadic-prob}, define the one-matrix
truncated inverse-square normalizer mass
\[
  \nu_{d,a}^{(2)}
  :=
  \Exp\left[
  \mathbf 1_{\{\norm{T_d(\xi)}_{\op}<aN_d\}}
  \min\left\{
  \left(\frac{aN_d}{\norm{T_d(\xi)}_{\op}}\right)^2,\ell
  \right\}
  \right],
  \qquad \ell=\log(e+d),
\]
where the minimum is interpreted as \(\ell\) on the event
\(\norm{T_d(\xi)}_{\op}=0\).  Then
\[
  \mu_{d,a}^{\rm dy}\le C\,\nu_{d,a}^{(2)}
\]
with an absolute constant \(C\).  Consequently, whenever \(n\le d\) and
\(n\ge C_{A,K,a}\ell\),
\[
  \Prob\left\{
  \max_{\omega^d=1}
  \sum_{i=1}^n
  \frac{\abs{P_i(\omega)}^2}{\norm{B_i}_{\op}^2}
  > C_{A,K,a,R}\frac n\ell
  \right\}
  \le
  C_{A,K,a,R}d^{-A}
  +\frac{C\ell\,\nu_{d,a}^{(2)}}{R}.
\]
On the complement, the profile, random-sign tail, sign-averaged maximum,
deterministic certificate, and scalar edge-grid moment conclusions of
Corollary~\ref{cor:normalizer-tail-dyadic} hold at the
\(C_{A,K,a,R}\sqrt n\) scale.

There is also a large-sample form.  If \(0<\theta<R\) and
\[
  C\,\nu_{d,a}^{(2)}\le \frac{\theta}{\ell},
\]
then the preceding exceptional probability can be replaced by
\[
  C_{A,K,a,R}d^{-A}
  +\exp\left(-c_{R,\theta}\frac n{\ell^2}\right).
\]
Thus any estimate of the truncated inverse-square normalizer mass at order
\(O(1/\ell)\) is enough for the Chernoff-scale dyadic transfer.

More generally, for \(p>0\) put
\[
  I_{d,a}^{(p)}
  :=
  \Exp\left[
  \mathbf 1_{\{\norm{T_d(\xi)}_{\op}<aN_d\}}
  \left(\frac{aN_d}{\norm{T_d(\xi)}_{\op}}\right)^p
  \right].
\]
Then
\[
  \nu_{d,a}^{(2)}
  \le
  \ell^{(2-p)_+/2}\,I_{d,a}^{(p)}.
\]
In particular, any finite inverse \(p\)-moment with \(p\ge2\) controls the
dyadic criterion directly, while \(0<p<2\) only costs the explicit factor
\(\ell^{1-p/2}\).
\end{corollary}

\begin{proof}
Let \(X=\norm{T_d(\xi)}_{\op}\) and \(Z=aN_d/X\) on the event
\(\{X<aN_d\}\).  If \(X\) lies in the \(r\)-th finite dyadic bin from the
definition of \(Y_{d,a}\), then \(Z^2>4^r\), and hence
\[
  4^{r+1}\le4Z^2.
\]
At the same time \(r<J_\ell\), so
\[
  4^{r+1}\le 4^{J_\ell}\le C\ell .
\]
Therefore the finite-bin contribution is bounded by
\[
  4^{r+1}\le C\min\{Z^2,\ell\}.
\]
On the bottom bin \(X<2^{-J_\ell}aN_d\), one has
\(Z^2>4^{J_\ell}\ge\ell\), so the bottom weight \(\ell\) is exactly bounded by
\(\min\{Z^2,\ell\}\).  Thus
\[
  Y_{d,a}\le
  C\,\mathbf 1_{\{X<aN_d\}}\min\{Z^2,\ell\}
\]
pointwise, which gives
\(\mu_{d,a}^{\rm dy}\le C\nu_{d,a}^{(2)}\).  The probability and
large-sample conclusions are the corresponding Markov and Bernstein parts of
Corollary~\ref{cor:normalizer-tail-dyadic-prob}.

Finally, for \(u\ge1\),
\[
  \min\{u^2,\ell\}\le \ell^{(2-p)_+/2}u^p .
\]
Indeed, for \(p\ge2\) this is immediate from \(u^2\le u^p\), while for
\(0<p<2\) it follows by considering separately \(u\le\sqrt\ell\) and
\(u>\sqrt\ell\).  Applying this with \(u=Z\) proves the final inverse-moment
bound.
\end{proof}

\begin{corollary}[Inverse-moment random-search and robust certificates]
\label{cor:normalizer-tail-invmoment-search}
Fix \(A>0\), \(a>0\), and \(R\ge1\).  In the setting of
Corollary~\ref{cor:normalizer-tail-invmoment}, write
\[
  v_i(\omega):=\frac{P_i(\omega)}{\norm{B_i}_{\op}},
  \qquad \omega^d=1,
  \qquad \ell=\log(e+d).
\]
Then, whenever \(n\le d\) and \(n\ge C_{A,K,a,R}\ell\), outside a matrix event
of probability at most
\[
  C_{A,K,a,R}d^{-A}+\frac{C\ell\,\nu_{d,a}^{(2)}}{R},
\]
the following all hold conditionally on the matrices.

First, for every integer \(M\ge1\), if
\(\delta^{(1)},\ldots,\delta^{(M)}\) are independent uniform sign vectors, then
for every \(u\ge0\),
\[
  \Prob_\delta\left\{
  \max_{1\le s\le M}\max_{\omega^d=1}
  \left|\sum_{i=1}^n\delta_i^{(s)}v_i(\omega)\right|
  >
  C_{A,K,a,R}\sqrt{
    n\left(1+\frac{\log(eM)+u}{\ell}\right)}
  \right\}\le e^{-u}.
\]
Second, for every integer \(M\ge1\),
\[
  \Prob_\delta\left\{
  \min_{1\le s\le M}\max_{\omega^d=1}
  \left|\sum_{i=1}^n\delta_i^{(s)}v_i(\omega)\right|
  > C_{A,K,a,R}\sqrt n
  \right\}\le d^{-AM}.
\]
More generally, if \(G_M\) is the number of sampled sign vectors whose edge
maximum is at most \(C_{A,K,a,R}\sqrt n\), then for every integer
\(1\le r\le M\),
\[
  \Prob_\delta\{G_M\le M-r\}\le \binom Mr d^{-Ar}.
\]

Finally, with Rademacher probability at least \(1-d^{-A}\), a random sign
vector is the center of a Hamming ball of radius
\[
  r_0=\left\lfloor c_{A,K,a,R}\ell\right\rfloor
\]
on which every signing has edge maximum at most \(C_{A,K,a,R}\sqrt n\).  If
\(H_M\) counts the sampled signs that are centers of such good Hamming balls,
then for every integer \(1\le r\le M\),
\[
  \Prob_\delta\{H_M\le M-r\}\le \binom Mr d^{-Ar}.
\]

Consequently, for every \(B>0\), if
\[
  M\ge\lceil B\rceil
\]
independent uniform signs are sampled independently of the matrices, then the
joint matrix-and-sign probability that no sampled sign is such a robust center
is at most
\[
  C_{A,K,a,R}d^{-A}+\frac{C\ell\,\nu_{d,a}^{(2)}}{R}+d^{-B}.
\]

There is also a large-sample form.  If \(C_*\) denotes the absolute constant in
Corollary~\ref{cor:normalizer-tail-invmoment}, \(0<\theta<R\), and
\[
  C_*\nu_{d,a}^{(2)}\le\frac{\theta}{\ell},
\]
then in all preceding statements the matrix exceptional probability can be
replaced by
\[
  C_{A,K,a,R}d^{-A}
  +\exp\left(-c_{R,\theta}\frac n{\ell^2}\right),
\]
and the joint robust-search failure by the same expression plus \(d^{-B}\).
For every \(p>0\), all occurrences of \(\nu_{d,a}^{(2)}\) may be replaced by
\[
  \ell^{(2-p)_+/2}
  \Exp\left[
  \mathbf 1_{\{\norm{T_d(\xi)}_{\op}<aN_d\}}
  \left(\frac{aN_d}{\norm{T_d(\xi)}_{\op}}\right)^p
  \right].
\]
\end{corollary}

\begin{proof}
Corollary~\ref{cor:normalizer-tail-invmoment} gives the edge profile
\[
  \max_{\omega^d=1}\sum_{i=1}^n|v_i(\omega)|^2
  \le C_{A,K,a,R}\frac n\ell
\]
outside the displayed matrix event.  Condition on any realization satisfying
this profile.  The sampled-list maximum, \(d^{-AM}\) amplification, and
binomial good-sign count follow from
Corollary~\ref{cor:edge-profile-random-search}, after absorbing the profile
constant.  The Hamming radius and robust-center binomial tail follow from
Corollary~\ref{cor:edge-profile-hamming}.  The polynomial-confidence robust
search assertion follows from Corollary~\ref{cor:edge-profile-poly-confidence}
with \(B\) there equal to the displayed confidence exponent, followed by a union
bound with the matrix event.

The large-sample replacement is the large-sample part of
Corollary~\ref{cor:normalizer-tail-invmoment}.  The final inverse \(p\)-moment
replacement is also the final assertion of
Corollary~\ref{cor:normalizer-tail-invmoment}.
\end{proof}

\begin{corollary}[Inverse-moment sampled abundance]
\label{cor:normalizer-tail-invmoment-abundance}
Fix \(A>0\), \(a>0\), \(R\ge1\), \(B>0\), and a fraction \(0<\rho\le1\).
In the setting of Corollary~\ref{cor:normalizer-tail-invmoment}, write
\[
  \ell=\log(e+d),\qquad
  v_i(\omega):=\frac{P_i(\omega)}{\norm{B_i}_{\op}},
  \qquad \omega^d=1.
\]
Let \(\delta^{(1)},\ldots,\delta^{(M)}\) be independent uniform sign vectors,
independent of the matrices, where \(M\ge1\), and put
\[
  r_\rho=\lceil\rho M\rceil .
\]
There are constants \(C_{A,B,K,a,R}\) and \(c_{A,B,K,a,R}>0\) such that,
whenever
\[
  n\le d,\qquad n\ge C_{A,B,K,a,R}\ell ,
\]
then, for every \(u\ge0\), with probability at least
\[
  1-C_{A,B,K,a,R}d^{-A}
  -\frac{C\ell\,\nu_{d,a}^{(2)}}{R}
  -e^{-u}-2\binom{M}{r_\rho}d^{-B r_\rho},
\]
all three conclusions below hold simultaneously.

First, the whole sampled list obeys
\[
\begin{aligned}
  \max_{1\le s\le M}\max_{\omega^d=1}
  \left|\sum_{i=1}^n\delta_i^{(s)}
  \frac{P_i(\omega)}{\norm{B_i}_{\op}}\right|
  \le
  C_{A,B,K,a,R}\sqrt{
    n\left(1+\frac{\log(eM)+u}{\ell}\right)} .
\end{aligned}
\]
Second, at least \(M-r_\rho+1\) sampled signs have edge maximum at most
\[
  C_{A,B,K,a,R}\sqrt n .
\]
Third, at least \(M-r_\rho+1\) sampled signs are centers of Hamming balls of
radius
\[
  \left\lfloor c_{A,B,K,a,R}\ell\right\rfloor
\]
on which every signing \(\eps'\) has edge maximum at most
\[
  C_{A,B,K,a,R}\sqrt n .
\]
In particular, at least a \(1-\rho\) fraction of the sampled signs are good,
and at least a \(1-\rho\) fraction are robust centers.

The preceding display has two useful specializations.  First, fix \(S>0\), and
choose \(B>0\) and \(M\ge1\) so that
\[
  B\ge 2\log_2(e/\rho),
  \qquad
  M\ge \frac{2(S+1)}{B\rho}.
\]
If \(d\ge2\), then, taking \(u=S\ell\), with probability at least
\[
  1-C_{A,B,K,a,R}d^{-A}
  -\frac{C\ell\,\nu_{d,a}^{(2)}}{R}
  -2d^{-S},
\]
the whole sampled list has edge maximum \(C_{A,B,K,a,R,S,\rho}\sqrt n\), and
at least a \(1-\rho\) fraction of the sampled signs are robust centers of
radius \(\lfloor c_{A,B,K,a,R}\ell\rfloor\).

Second, fix \(\beta\ge0\), \(\sigma>0\), and suppose \(M\le d^\beta\).
Then, after enlarging constants depending on \(\beta,\sigma\), with
probability at least
\[
  1-C_{A,K,a,R,\beta,\sigma}d^{-A}
  -\frac{C\ell\,\nu_{d,a}^{(2)}}{R}
  -d^{-\sigma},
\]
every sampled sign is the center of a Hamming ball of radius
\[
  \left\lfloor c_{A,K,a,R,\beta,\sigma}\ell\right\rfloor
\]
on which every signing has edge maximum at most
\[
  C_{A,K,a,R,\beta,\sigma}\sqrt n .
\]
In particular, the entire polynomial-size sampled batch is robust and good.

There is also a large-sample form.  If \(C_*\) denotes the absolute constant in
Corollary~\ref{cor:normalizer-tail-invmoment}, \(0<\theta<R\), and
\[
  C_*\nu_{d,a}^{(2)}\le\frac{\theta}{\ell},
\]
then in every displayed probability bound above the matrix part
\[
  C_{A,K,a,R}d^{-A}+\frac{C\ell\,\nu_{d,a}^{(2)}}{R}
\]
may be replaced, after changing constants, by
\[
  C_{A,K,a,R}d^{-A}
  +\exp\left(-c_{R,\theta}\frac n{\ell^2}\right).
\]
For every \(p>0\), all occurrences of \(\nu_{d,a}^{(2)}\) may again be replaced
by the inverse \(p\)-moment expression from
Corollary~\ref{cor:normalizer-tail-invmoment-search}.
\end{corollary}

\begin{proof}
The general abundance statement is
Corollary~\ref{cor:normalizer-tail-invmoment-search} combined with a union
bound over the sampled-list event, the good-sign count event, and the
robust-center count event, using the sign-count exponent \(B\).  The matrix
failure is exactly the Markov-form matrix failure in
Corollary~\ref{cor:normalizer-tail-invmoment-search}.

For the robust-fraction specialization, use
\[
  \binom{M}{r_\rho}\le (e/\rho)^{r_\rho}
\]
and \(B\ge2\log_2(e/\rho)\) to get, for \(d\ge2\),
\[
  \binom{M}{r_\rho}d^{-B r_\rho}
  \le d^{-B r_\rho/2}.
\]
Since \(r_\rho\ge\rho M\), the lower bound on \(M\) gives
\(B r_\rho/2\ge S+1\), hence the two binomial failure terms are bounded by
\(d^{-S}\).  Taking \(u=S\ell\) also gives \(e^{-u}\le d^{-S}\), and the
sampled-list maximum is \(O_{A,B,K,a,R,S,\rho}(\sqrt n)\) because \(M\) is
fixed in terms of \(B,S,\rho\).

For the polynomial-batch specialization, apply the robust-center count with
sign exponent \(B_{\rm sign}=\beta+\sigma\) and \(r=1\).  The conditional sign
failure is at most
\[
  \binom M1 d^{-(\beta+\sigma)}\le d^{-\sigma}.
\]
On the complementary event all sampled signs are robust centers, and therefore
all have edge maximum at the same \(O(\sqrt n)\) scale.  The large-sample and
inverse \(p\)-moment replacements are the corresponding final assertions of
Corollary~\ref{cor:normalizer-tail-invmoment-search}.
\end{proof}

\begin{corollary}[Inverse-moment finite vanishing discarded fractions]
\label{cor:normalizer-tail-invmoment-vanishing-fractions}
Fix \(A>0\), \(a>0\), \(R\ge1\), a sign-count exponent \(B>0\), and
\(0<\tau<B\).  In the setting of
Corollary~\ref{cor:normalizer-tail-invmoment}, write
\[
  \ell=\log(e+d),\qquad
  v_i(\omega)=\frac{P_i(\omega)}{\norm{B_i}_{\op}},
  \qquad \omega^d=1 .
\]
Let \(\delta^{(1)},\ldots,\delta^{(M)}\) be independent uniform sign vectors,
independent of the matrices, where \(M\ge1\).  Let \(0<\rho\le1\) satisfy
\[
  \rho\ge e\,d^{-\tau},
\]
and put \(r_\rho=\lceil\rho M\rceil\).  There are constants
\(C_{A,B,K,a,R}\) and \(c_{A,B,K,a,R}>0\) such that, whenever
\[
  d\ge2,\qquad n\le d,\qquad n\ge C_{A,B,K,a,R}\ell ,
\]
then, for every \(u\ge0\), with probability at least
\[
  1-C_{A,B,K,a,R}d^{-A}
  -\frac{C\ell\,\nu_{d,a}^{(2)}}{R}
  -e^{-u}-2d^{-(B-\tau)},
\]
all three conclusions below hold simultaneously.

First, the whole sampled list obeys
\[
\begin{aligned}
  \max_{1\le s\le M}\max_{\omega^d=1}
  \left|\sum_{i=1}^n\delta_i^{(s)}v_i(\omega)\right|
  \le
  C_{A,B,K,a,R}
  \sqrt{n\left(1+\frac{\log(eM)+u}{\ell}\right)} .
\end{aligned}
\]
Second, at least \(M-r_\rho+1\) sampled signs have edge maximum at most
\[
  C_{A,B,K,a,R}\sqrt n .
\]
Third, at least \(M-r_\rho+1\) sampled signs are centers of Hamming balls of
radius
\[
  \left\lfloor c_{A,B,K,a,R}\ell\right\rfloor
\]
on which every signing has edge maximum at most
\[
  C_{A,B,K,a,R}\sqrt n .
\]
In particular, at least a \(1-\rho\) fraction of the sampled signs are robust
centers.  No upper bound on \(M\) is required; only the sampled-list maximum
retains the entropy factor \(\log(eM)\).

There is also a large-sample form.  If \(C_*\) denotes the absolute constant in
Corollary~\ref{cor:normalizer-tail-invmoment}, \(0<\theta<R\), and
\[
  C_*\nu_{d,a}^{(2)}\le\frac{\theta}{\ell},
\]
then the matrix part of the displayed probability bound may be replaced, after
changing constants, by
\[
  C_{A,K,a,R}d^{-A}
  +\exp\left(-c_{R,\theta}\frac n{\ell^2}\right).
\]
For every \(p>0\), all occurrences of \(\nu_{d,a}^{(2)}\) may again be replaced
by the inverse \(p\)-moment expression from
Corollary~\ref{cor:normalizer-tail-invmoment-search}.
\end{corollary}

\begin{proof}
Apply Corollary~\ref{cor:normalizer-tail-invmoment-abundance} with the same
fraction \(\rho\).  The only point is to simplify the binomial count tail.  Since
\[
  \binom{M}{r_\rho}
  \le \left(\frac{eM}{r_\rho}\right)^{r_\rho}
  \le \left(\frac e\rho\right)^{r_\rho},
\]
the lower bound \(\rho\ge e d^{-\tau}\) gives
\[
  \binom{M}{r_\rho}d^{-B r_\rho}
  \le d^{-(B-\tau)r_\rho}
  \le d^{-(B-\tau)} .
\]
Substituting this into the finite abundance theorem gives the displayed
probability and the sampled-list, good-count, and robust-center-count
conclusions.  Finally,
\[
  M-r_\rho+1\ge M-\rho M,
\]
so the robust-center count is at least a \(1-\rho\) fraction.  The large-sample
and inverse \(p\)-moment variants are the corresponding final assertions of
Corollary~\ref{cor:normalizer-tail-invmoment-abundance}.
\end{proof}

\begin{corollary}[Almost-sure inverse-moment dyadic-tail transfer]
\label{cor:normalizer-tail-invmoment-borel}
Fix \(A>0\), \(a>0\), and \(R\ge1\).  Let \((d_m,n_m)_{m\ge1}\) be a
triangular-array sequence with \(n_m\le d_m\), put
\(\ell_m=\log(e+d_m)\), and assume
\[
  n_m\ge C_{A,K,a,R}\ell_m
  \qquad(m\ge1).
\]
For each \(m\), let \(B_{m,i}=T_{d_m}(\xi^{(m,i)})\),
\(1\le i\le n_m\), be independent full Toeplitz matrices whose lag variables
are independent, mean zero, unit variance, and have sub-Gaussian norm at most
\(K\), and let \(P_{m,i}\) be their symbols.  Define
\[
  \nu_{m,a}^{(2)}
  :=
  \Exp\left[
  \mathbf 1_{\{\norm{T_{d_m}(\xi)}_{\op}<aN_{d_m}\}}
  \min\left\{
  \left(\frac{aN_{d_m}}{\norm{T_{d_m}(\xi)}_{\op}}\right)^2,\ell_m
  \right\}
  \right],
\]
where the minimum is interpreted as \(\ell_m\) on the event
\(\norm{T_{d_m}(\xi)}_{\op}=0\).  If
\[
  \sum_{m=1}^{\infty}
  \left(d_m^{-A}+\frac{\ell_m\nu_{m,a}^{(2)}}{R}\right)<\infty,
\]
then, with probability one over the matrix triangular array, all but finitely
many \(m\) satisfy
\[
  \max_{\omega^{d_m}=1}
  \sum_{i=1}^{n_m}
  \frac{\abs{P_{m,i}(\omega)}^2}{\norm{B_{m,i}}_{\op}^2}
  \le C_{A,K,a,R}\frac{n_m}{\ell_m}.
\]
For those \(m\), the signing, sign-averaged maximum,
deterministic-certificate, and scalar edge-grid moment conclusions of
Corollary~\ref{cor:normalizer-tail-dyadic} hold with
\((d,n,\ell)\) replaced by \((d_m,n_m,\ell_m)\) at scale
\(C_{A,K,a,R}\sqrt{n_m}\).

Moreover, put
\[
  v_{m,i}(\omega):=
  \frac{P_{m,i}(\omega)}{\norm{B_{m,i}}_{\op}},
  \qquad \omega^{d_m}=1.
\]
On the same almost-sure event, all but finitely many \(m\) also have the
conditional sampled-list maximum, \(d_m^{-AM}\) success amplification,
binomial good-sign counts, Hamming radius \(\lfloor c_{A,K,a,R}\ell_m\rfloor\),
and robust-center binomial tails of
Corollary~\ref{cor:normalizer-tail-borel-search-stability}, with \(v_{m,i}\)
as above and constants depending on \(A,K,a,R\).

There is also a large-sample form.  If, for the absolute constant \(C\) in
Corollary~\ref{cor:normalizer-tail-invmoment}, there is a
\(0<\theta<R\) such that
\[
  C\nu_{m,a}^{(2)}\le \frac{\theta}{\ell_m}
\]
for all sufficiently large \(m\), and if
\[
  \sum_{m=1}^{\infty}
  \left(d_m^{-A}+\exp\left(-c_{R,\theta}\frac{n_m}{\ell_m^2}\right)\right)
  <\infty,
\]
then all preceding almost-sure matrix and conditional sign-search conclusions
hold with this Bernstein summability condition in place of the Markov
summability condition.

Finally, suppose independent uniform sign lists
\[
  \delta_m^{(1)},\ldots,\delta_m^{(M_m)}\in\{-1,1\}^{n_m},
  \qquad M_m\ge\lceil\alpha\rceil,
\]
are sampled independently of the matrices.  If
\[
  \sum_{m=1}^{\infty}
  \left(d_m^{-A}+\frac{\ell_m\nu_{m,a}^{(2)}}{R}+d_m^{-\alpha}\right)<\infty,
\]
then, with probability one over the matrices and sampled signs, all but
finitely many \(m\) contain a sampled sign that is the center of a Hamming ball
of radius \(\lfloor c_{A,K,a,R}\ell_m\rfloor\) on which every signing
\(\eps'\) satisfies
\[
  \max_{\omega^{d_m}=1}
  \left|\sum_{i=1}^{n_m}
  \eps'_i\frac{P_{m,i}(\omega)}{\norm{B_{m,i}}_{\op}}\right|
  \le C_{A,K,a,R}\sqrt{n_m}.
\]
The same sampled-certificate conclusion holds under the large-sample
Bernstein summability condition after adding the term \(d_m^{-\alpha}\).

For every \(p>0\), all displayed summability conditions remain valid if
\(\nu_{m,a}^{(2)}\) is replaced by
\[
  \ell_m^{(2-p)_+/2}
  \Exp\left[
  \mathbf 1_{\{\norm{T_{d_m}(\xi)}_{\op}<aN_{d_m}\}}
  \left(\frac{aN_{d_m}}{\norm{T_{d_m}(\xi)}_{\op}}\right)^p
  \right].
\]
\end{corollary}

\begin{proof}
Apply Corollary~\ref{cor:normalizer-tail-invmoment} at each index \(m\).  The
Markov part gives a matrix exceptional probability bounded by a constant times
\[
  d_m^{-A}+\frac{\ell_m\nu_{m,a}^{(2)}}{R},
\]
so the first displayed summability condition and Borel--Cantelli give the
eventual edge profile.  The signing, sign-averaged maximum,
deterministic-certificate, and scalar moment statements are deterministic
consequences of that same good profile event in
Corollary~\ref{cor:normalizer-tail-dyadic}.

Once the edge profile holds, Corollaries~\ref{cor:edge-profile-random-search}
and \ref{cor:edge-profile-hamming}, applied with the profile constant
\(C_{A,K,a,R}\), give the conditional sampled-list, success-count, and robust
Hamming-center conclusions.  This is exactly the package displayed in
Corollary~\ref{cor:normalizer-tail-borel-search-stability}, with constants
depending also on \(R\).

If \(C\nu_{m,a}^{(2)}\le\theta/\ell_m\) eventually, the large-sample part of
Corollary~\ref{cor:normalizer-tail-invmoment} replaces the Markov matrix
exception by
\[
  C_{A,K,a,R}d_m^{-A}
  +\exp\left(-c_{R,\theta}\frac{n_m}{\ell_m^2}\right).
\]
The Bernstein summability assumption therefore gives the same almost-sure
eventual conclusions.

For the sampled-certificate assertion, condition on a matrix realization with
the eventual edge profile and apply
Corollary~\ref{cor:edge-profile-poly-confidence} with \(B=\alpha\).  The
failure probability for the sampled sign list is at most \(d_m^{-\alpha}\), so
the displayed joint summability condition and Borel--Cantelli give an eventual
sampled robust center.  The Bernstein variant is identical, with the matrix
exception replaced as above.  The final inverse \(p\)-moment replacement is the
last assertion of Corollary~\ref{cor:normalizer-tail-invmoment}.
\end{proof}

\begin{corollary}[Almost-sure inverse-moment sampled abundance]
\label{cor:normalizer-tail-invmoment-borel-abundance}
Fix \(A>0\), \(a>0\), \(R\ge1\), and a sign-count exponent
\(B_{\rm s}>0\).  Let \((d_m,n_m)_{m\ge1}\), \(\ell_m\),
\(B_{m,i}=T_{d_m}(\xi^{(m,i)})\), \(P_{m,i}\), and
\(\nu_{m,a}^{(2)}\) be as in
Corollary~\ref{cor:normalizer-tail-invmoment-borel}, with
\[
  n_m\le d_m,\qquad
  n_m\ge C_{A,B_{\rm s},K,a,R}\ell_m .
\]
At each \(m\), independently of the matrices, sample independent uniform sign
vectors
\[
  \delta_m^{(1)},\ldots,\delta_m^{(M_m)}\in\{-1,1\}^{n_m},
  \qquad M_m\ge1.
\]
Let \(u_m\ge0\) and \(1\le r_m\le M_m\) be deterministic parameters.  If
\[
  \sum_{m=1}^{\infty}
  \left(
  d_m^{-A}+\frac{\ell_m\nu_{m,a}^{(2)}}{R}
  +e^{-u_m}+\binom{M_m}{r_m}d_m^{-B_{\rm s}r_m}
  \right)<\infty,
\]
then, with probability one over the matrices and sampled signs, all but
finitely many \(m\) satisfy all three conclusions below, with
\[
  v_{m,i}(\omega):=
  \frac{P_{m,i}(\omega)}{\norm{B_{m,i}}_{\op}},
  \qquad \omega^{d_m}=1.
\]
First,
\[
\begin{aligned}
  \max_{1\le s\le M_m}\max_{\omega^{d_m}=1}
  \left|\sum_{i=1}^{n_m}\delta_{m,i}^{(s)}v_{m,i}(\omega)\right|
  \le
  C_{A,B_{\rm s},K,a,R}
  \sqrt{n_m\left(1+\frac{\log(eM_m)+u_m}{\ell_m}\right)} .
\end{aligned}
\]
Second, at least \(M_m-r_m+1\) sampled signs have edge maximum at most
\[
  C_{A,B_{\rm s},K,a,R}\sqrt{n_m}.
\]
Third, at least \(M_m-r_m+1\) sampled signs are centers of Hamming balls of
radius
\[
  \left\lfloor c_{A,B_{\rm s},K,a,R}\ell_m\right\rfloor
\]
on which every signing has edge maximum at most
\[
  C_{A,B_{\rm s},K,a,R}\sqrt{n_m}.
\]

There are two useful consequences.  First, fix \(0<\rho\le1\) and \(S>0\).
Choose \(B_{\rm s}\) and an integer \(M\ge1\) so that
\[
  B_{\rm s}\ge2\log_2(e/\rho),
  \qquad
  M\ge \frac{2(S+1)}{B_{\rm s}\rho}.
\]
If \(d_m\ge2\) and
\[
  \sum_{m=1}^{\infty}
  \left(d_m^{-A}+\frac{\ell_m\nu_{m,a}^{(2)}}{R}+d_m^{-S}\right)<\infty,
\]
then any fixed sampled list of this size \(M\) has, eventually almost surely,
edge maximum \(C_{A,B_{\rm s},K,a,R,S,\rho}\sqrt{n_m}\) on the whole list and a
\(1-\rho\) fraction of logarithmic Hamming-robust centers.

Second, fix \(\beta\ge0\) and \(\sigma>0\).  If \(M_m\le d_m^\beta\) and
\[
  \sum_{m=1}^{\infty}
  \left(d_m^{-A}+\frac{\ell_m\nu_{m,a}^{(2)}}{R}+d_m^{-\sigma}\right)<\infty,
\]
then, eventually almost surely, every sampled sign in the polynomial-size batch
is the center of a Hamming ball of radius
\[
  \left\lfloor c_{A,K,a,R,\beta,\sigma}\ell_m\right\rfloor
\]
on which every signing has edge maximum at most
\[
  C_{A,K,a,R,\beta,\sigma}\sqrt{n_m}.
\]
In particular, the whole sampled batch is robust and good.

There is also a large-sample form.  If \(C_*\) denotes the absolute constant in
Corollary~\ref{cor:normalizer-tail-invmoment}, \(0<\theta<R\), and
\[
  C_*\nu_{m,a}^{(2)}\le \frac{\theta}{\ell_m}
\]
for all sufficiently large \(m\), then each displayed summability condition
above remains valid with the matrix part
\[
  d_m^{-A}+\frac{\ell_m\nu_{m,a}^{(2)}}{R}
\]
replaced by
\[
  d_m^{-A}+\exp\left(-c_{R,\theta}\frac{n_m}{\ell_m^2}\right).
\]
For every \(p>0\), all occurrences of \(\nu_{m,a}^{(2)}\) may again be replaced
by the inverse \(p\)-moment expression from
Corollary~\ref{cor:normalizer-tail-invmoment-borel}.
\end{corollary}

\begin{proof}
For the general count statement, apply
Corollary~\ref{cor:normalizer-tail-invmoment-abundance} at index \(m\), with
\[
  \rho_m=\frac{r_m}{M_m}.
\]
Then \(r_{\rho_m}=r_m\), and the constants in the finite abundance statement do
not depend on \(\rho_m\).  The displayed summability condition dominates the
finite matrix-and-sign failure probabilities, so Borel--Cantelli gives the
eventual sampled-list maximum, good-sign count, and robust-center count.

For the fixed robust-fraction consequence, set
\[
  r_m=\lceil\rho M\rceil,\qquad u_m=S\ell_m .
\]
As in the proof of
Corollary~\ref{cor:normalizer-tail-invmoment-abundance},
\[
  \binom{M}{r_m}d_m^{-B_{\rm s}r_m}\le d_m^{-S}
  \qquad(d_m\ge2),
\]
after increasing constants, and \(e^{-u_m}\le d_m^{-S}\).  The whole sampled
list is \(O_{A,B_{\rm s},K,a,R,S,\rho}(\sqrt{n_m})\) because \(M\) is fixed, and
\[
  M-r_m+1\ge M-\rho M .
\]
Thus at least a \(1-\rho\) fraction of the sampled list consists of robust
centers.

For the polynomial-batch consequence, apply the polynomial-batch specialization
of Corollary~\ref{cor:normalizer-tail-invmoment-abundance} at each \(m\).  Its
additional sign failure is at most \(d_m^{-\sigma}\), so the displayed
summability condition again gives the eventual all-robust batch conclusion by
Borel--Cantelli.  The large-sample and inverse \(p\)-moment replacements are the
corresponding final assertions of
Corollary~\ref{cor:normalizer-tail-invmoment-abundance}.
\end{proof}

\begin{corollary}[Inverse-moment vanishing discarded fractions]
\label{cor:normalizer-tail-invmoment-borel-vanishing-fractions}
Fix \(A>0\), \(a>0\), \(R\ge1\), a sign-count exponent \(B_{\rm s}>0\),
and \(0<\tau<B_{\rm s}\).  In the setting of
Corollary~\ref{cor:normalizer-tail-invmoment-borel-abundance}, assume
\(d_m\ge2\) eventually, let \(M_m\ge1\) be arbitrary deterministic sampled-batch
sizes, let \(u_m\ge0\), and let \((\rho_m)_{m\ge1}\subset(0,1]\) satisfy
\[
  \rho_m\ge e\,d_m^{-\tau}
  \qquad\text{for all sufficiently large }m.
\]
Put
\[
  r_m=\lceil\rho_mM_m\rceil .
\]
If
\[
  \sum_{m=1}^{\infty}
  \left(
  d_m^{-A}+\frac{\ell_m\nu_{m,a}^{(2)}}{R}
  +e^{-u_m}+d_m^{-(B_{\rm s}-\tau)}
  \right)<\infty,
\]
then, with probability one over the matrices and sampled signs, all but
finitely many \(m\) satisfy the sampled-list maximum bound from
Corollary~\ref{cor:normalizer-tail-invmoment-borel-abundance} and at least a
\(1-\rho_m\) fraction of the \(M_m\) sampled signs have edge maximum
\[
  C_{A,B_{\rm s},K,a,R}\sqrt{n_m}
\]
and are centers of Hamming balls of radius
\[
  \left\lfloor c_{A,B_{\rm s},K,a,R}\ell_m\right\rfloor
\]
on which every signing has edge maximum at most
\[
  C_{A,B_{\rm s},K,a,R}\sqrt{n_m}.
\]
No upper bound on \(M_m\) is required.  In particular, if \(\rho_m\to0\), then
the robust-center fraction tends to one along the sampled batches, while the
whole-list maximum retains only the entropy factor \(\log(eM_m)\).

If the large-sample hypothesis
\[
  C_*\nu_{m,a}^{(2)}\le\frac{\theta}{\ell_m}
\]
holds eventually for some \(0<\theta<R\), then the same conclusion holds with
the matrix summand
\[
  d_m^{-A}+\frac{\ell_m\nu_{m,a}^{(2)}}{R}
\]
replaced by
\[
  d_m^{-A}+\exp\left(-c_{R,\theta}\frac{n_m}{\ell_m^2}\right).
\]
For every \(p>0\), the same inverse \(p\)-moment replacement as in
Corollary~\ref{cor:normalizer-tail-invmoment-borel-abundance} is valid.
\end{corollary}

\begin{proof}
Apply Corollary~\ref{cor:normalizer-tail-invmoment-borel-abundance} with the
displayed \(r_m\).  It remains only to check its binomial count-tail summability.
Since \(r_m=\lceil\rho_mM_m\rceil\),
\[
  \binom{M_m}{r_m}
  \le \left(\frac{eM_m}{r_m}\right)^{r_m}
  \le \left(\frac e{\rho_m}\right)^{r_m}.
\]
The lower bound on \(\rho_m\) gives, for all large \(m\),
\[
  \binom{M_m}{r_m}d_m^{-B_{\rm s}r_m}
  \le d_m^{-(B_{\rm s}-\tau)r_m}
  \le d_m^{-(B_{\rm s}-\tau)}.
\]
The displayed summability condition therefore implies the general sampled-count
hypothesis of Corollary~\ref{cor:normalizer-tail-invmoment-borel-abundance}.
Finally,
\[
  M_m-r_m+1\ge M_m-\rho_mM_m,
\]
so the robust-center count is at least a \(1-\rho_m\) fraction.  The
large-sample and inverse \(p\)-moment variants are the corresponding final
assertions of Corollary~\ref{cor:normalizer-tail-invmoment-borel-abundance}.
\end{proof}

\begin{corollary}[Inverse-moment operator-norm robust batches]
\label{cor:normalizer-tail-invmoment-borel-operator-batches}
In the setting of
Corollary~\ref{cor:normalizer-tail-invmoment-borel-abundance}, assume in
addition that
\[
  d_m/n_m\to\infty,
  \qquad
  \log d_m=o(n_m).
\]
Put
\[
  A_{m,i}:=\frac{B_{m,i}}{\norm{B_{m,i}}_{\op}},
  \qquad
  \widetilde A_{m,i}:=J_{d_m}A_{m,i}.
\]
If the summability hypothesis of
Corollary~\ref{cor:normalizer-tail-invmoment-borel-abundance} holds for
parameters \(u_m\) and \(r_m\), then, with probability one over the matrices and
sampled signs, all but finitely many \(m\) satisfy the sampled-list edge maximum
bound from that corollary and at least \(M_m-r_m+1\) sampled signs are centers
of Hamming balls of radius
\[
  \left\lfloor c_{A,B_{\rm s},K,a,R}\ell_m\right\rfloor
\]
on which every signing \(\eps'\) obeys both operator-norm estimates
\[
  \left\|\sum_{i=1}^{n_m}\eps'_iA_{m,i}\right\|_{\op}
  \le C_{A,B_{\rm s},K,a,R}\sqrt{n_m},
  \qquad
  \left\|\sum_{i=1}^{n_m}\eps'_i\widetilde A_{m,i}\right\|_{\op}
  \le C_{A,B_{\rm s},K,a,R}\sqrt{n_m},
\]
after increasing the constant.

Consequently, under the hypotheses of
Corollary~\ref{cor:normalizer-tail-invmoment-borel-vanishing-fractions}, the
same almost-sure event has the following stronger conclusion: all but finitely
many \(m\) have at least a \(1-\rho_m\) fraction of the sampled signs as
logarithmic-radius robust centers for the realization-normalized full Toeplitz
and full Hankel operator-norm discrepancy problems.  No upper bound on \(M_m\)
is required; only the sampled-list edge maximum keeps the entropy factor
\(\log(eM_m)\).

The large-sample Bernstein and inverse \(p\)-moment replacements from
Corollaries~\ref{cor:normalizer-tail-invmoment-borel-abundance} and
\ref{cor:normalizer-tail-invmoment-borel-vanishing-fractions} remain valid.
\end{corollary}

\begin{proof}
The abundance corollary gives, eventually almost surely, at least
\(M_m-r_m+1\) sampled signs whose Hamming balls have edge maximum at most
\(C_{A,B_{\rm s},K,a,R}\sqrt{n_m}\).  Apply
Corollary~\ref{cor:edge-to-operator-robust} to each of these Hamming balls,
or equivalently to their union, and absorb the deterministic reduction constant.
This gives the Toeplitz and Hankel operator-norm estimates on the same
almost-sure event.

For the vanishing-fraction statement, use
Corollary~\ref{cor:normalizer-tail-invmoment-borel-vanishing-fractions}; its
count estimate gives
\[
  M_m-r_m+1\ge M_m-\rho_mM_m.
\]
The large-sample and inverse \(p\)-moment variants are inherited from the same
two inverse-moment Borel corollaries, since the edge-to-operator transfer is
deterministic.
\end{proof}

\begin{corollary}[Concrete inverse-moment Borel schedules]
\label{cor:normalizer-tail-invmoment-borel-schedules}
Fix \(a>0\), \(R\ge1\), and \(C_0<\infty\).  For each \(d\), let
\(\ell_d=\log(e+d)\) and let \(\nu_{d,a}^{(2)}\) be the truncated
inverse-square normalizer mass from
Corollary~\ref{cor:normalizer-tail-invmoment}.  Let
\((d_m,n_m)_{m\ge1}\) be a triangular-array sequence with \(n_m\le d_m\), put
\(\ell_m=\log(e+d_m)\), and let \(B_{m,i}=T_{d_m}(\xi^{(m,i)})\),
\(1\le i\le n_m\), be independent full Toeplitz matrices whose lag variables
are independent, mean zero, unit variance, and have sub-Gaussian norm at most
\(K\), with symbols \(P_{m,i}\).

In each of the following Markov-form situations, with constants depending on the
displayed matrix-event exponent \(A\), the almost-sure matrix and conditional
sign-search conclusions of
Corollary~\ref{cor:normalizer-tail-invmoment-borel} hold:
\begin{enumerate}[label=\textup{(\roman*)},leftmargin=2.2em]
\item \(d_m=m\), \(A>1\),
\[
  C_{A,K,a,R}\ell_m\le n_m\le d_m,
  \qquad
  \nu_{d,a}^{(2)}\le C_0d^{-B}\quad(d\ge1)
\]
for some \(B>1\).
\item \(d_m\ge c_0m^\kappa\) for some \(c_0,\kappa>0\) and \(A>0\),
\[
  C_{A,K,a,R}\ell_m\le n_m\le d_m,
  \qquad
  \nu_{d,a}^{(2)}\le C_0d^{-B}\quad(d\ge1),
\]
and \(A\kappa>1\), \(B\kappa>1\).
\item \(d_m\ge\exp(c_0m^\kappa)\) for some \(c_0,\kappa>0\) and \(A>0\),
\[
  C_{A,K,a,R}\ell_m\le n_m\le d_m,
  \qquad
  \nu_{d,a}^{(2)}\le C_0\ell_d^{-B}\quad(d\ge1),
\]
and \(B>1+1/\kappa\).
\end{enumerate}
Thus, in each case, with probability one over the matrices, all but finitely
many \(m\) satisfy
\[
  \max_{\omega^{d_m}=1}
  \sum_{i=1}^{n_m}
  \frac{\abs{P_{m,i}(\omega)}^2}{\norm{B_{m,i}}_{\op}^2}
  \le C_{A,K,a,R}\frac{n_m}{\ell_m},
\]
together with the dyadic signing, sign-averaged maximum,
deterministic-certificate, scalar moment, sampled-list, success-count,
Hamming-ball, and robust-center conditional conclusions of
Corollary~\ref{cor:normalizer-tail-invmoment-borel}.

If independent sign lists
\[
  \delta_m^{(1)},\ldots,\delta_m^{(M_m)}\in\{-1,1\}^{n_m},
  \qquad M_m\ge\lceil\alpha\rceil,
\]
are also sampled, then the actual sampled robust-certificate conclusion of
Corollary~\ref{cor:normalizer-tail-invmoment-borel} also holds whenever
\(\sum_m d_m^{-\alpha}<\infty\).  In particular, this requires
\(\alpha>1\) in the all-integer case, \(\alpha\kappa>1\) in the polynomial
schedule case, and no lower bound beyond \(\alpha>0\) in the
stretched-exponential case.

The robust-fraction and polynomial-batch conclusions of
Corollary~\ref{cor:normalizer-tail-invmoment-borel-abundance} also hold on the
same Markov-form schedules after adding the corresponding sign summability
condition.  Thus a fixed \(O_{\rho,S}(1)\) sampled list has an eventual
\(1-\rho\) fraction of logarithmically robust centers whenever
\(\sum_m d_m^{-S}<\infty\), which means \(S>1\) in the all-integer case,
\(S\kappa>1\) in the polynomial schedule case, and any \(S>0\) in the
stretched-exponential case.  If \(M_m\le d_m^\beta\), then the whole sampled
batch is eventually robust whenever \(\sum_m d_m^{-\sigma}<\infty\), which means
\(\sigma>1\), \(\sigma\kappa>1\), and any \(\sigma>0\), respectively.
The vanishing discarded-fraction conclusion of
Corollary~\ref{cor:normalizer-tail-invmoment-borel-vanishing-fractions} also
holds with no upper bound on \(M_m\): for any sign-count exponent \(B_{\rm s}\)
and any \(0<\tau<B_{\rm s}\), arbitrary sampled batches have an eventual
\(1-\rho_m\) robust-center fraction whenever
\(\rho_m\ge e d_m^{-\tau}\) eventually, \(\sum_m e^{-u_m}<\infty\), and
\(\sum_m d_m^{-(B_{\rm s}-\tau)}<\infty\).  Thus the all-integer case requires
\(B_{\rm s}-\tau>1\), the polynomial schedule case requires
\((B_{\rm s}-\tau)\kappa>1\), and the stretched-exponential case needs only
\(\tau<B_{\rm s}\).

There is also a critical Bernstein form.  Let \(C_*\) be the absolute constant
in Corollary~\ref{cor:normalizer-tail-invmoment}.  Suppose
\[
  \nu_{d,a}^{(2)}\le \frac{\eta}{\ell_d}\qquad(d\ge1),
  \qquad C_*\eta<R.
\]
Choose \(C_*\eta<\theta<R\).  Then the same eventual matrix and conditional
sign-search conclusions hold in each triangular array for which
\[
  \sum_{m=1}^{\infty}
  \left(d_m^{-A}+\exp\left(-c_{R,\theta}\frac{n_m}{\ell_m^2}\right)\right)
  <\infty .
\]
Consequently, along all integer dimensions, if \(A>1\) and
\[
  n(d)\le d,\qquad n(d)\ge q\log^3(e+d)\quad\text{for all large }d,
  \qquad q>\frac1{c_{R,\theta}},
\]
then the conclusions hold almost surely for all sufficiently large \(d\).  On
stretched-exponential schedules \(d_m\ge\exp(c_0m^\kappa)\), they hold if
\[
  n_m\ge q\,\ell_m^2\log(e+m)
  \quad\text{eventually with}\quad q>\frac1{c_{R,\theta}},
\]
and also under the simpler superquadratic logarithmic window
\[
  n_m\ge Q\,\ell_m^{2+\tau}
  \qquad\text{eventually}
\]
for any \(Q,\tau>0\).  The corresponding actual sampled robust-certificate
claims again follow after adding any sign-list exponent \(\alpha\) with
\(\sum_m d_m^{-\alpha}<\infty\).
The same critical Bernstein schedules inherit the fixed robust-fraction
conclusion after adding \(\sum_m d_m^{-S}<\infty\), and the all-robust
polynomial-batch conclusion after adding \(\sum_m d_m^{-\sigma}<\infty\).  Thus
the all-integer \(\log^3d\) window requires \(S>1\) and \(\sigma>1\), while the
stretched-exponential critical windows allow every \(S,\sigma>0\).
They also inherit the vanishing discarded-fraction conclusion under the same
extra sign-series condition \(\sum_m d_m^{-(B_{\rm s}-\tau)}<\infty\): hence
\(B_{\rm s}-\tau>1\) in the all-integer critical \(\log^3d\) window and only
\(\tau<B_{\rm s}\) on the stretched-exponential critical windows.
\end{corollary}

\begin{proof}
It is enough to verify the summability hypotheses in
Corollary~\ref{cor:normalizer-tail-invmoment-borel}.  In the all-integer
polynomial-decay case,
\[
  \sum_m m^{-A}<\infty,
  \qquad
  \sum_m \ell_m\nu_{d_m,a}^{(2)}
  \le C_0\sum_m \log(e+m)m^{-B}<\infty
\]
because \(A,B>1\).  In the polynomial schedule case,
\[
  d_m^{-A}\le C m^{-A\kappa},
  \qquad
  \ell_m\nu_{d_m,a}^{(2)}
  \le C_{B,C_0,c_0}(\log(e+m))m^{-B\kappa},
\]
and both series converge under \(A\kappa>1\) and \(B\kappa>1\).  In the
stretched-exponential logarithmic-decay case, \(\sum_m d_m^{-A}<\infty\) and
\[
  \ell_m\nu_{d_m,a}^{(2)}
  \le C_0\ell_m^{1-B}
  \le C_{B,C_0,c_0}m^{\kappa(1-B)},
\]
which is summable precisely when \(B>1+1/\kappa\).  This proves the
Markov-form matrix and conditional sign-search conclusions.  The sampled-list
assertion is the sampled-certificate part of
Corollary~\ref{cor:normalizer-tail-invmoment-borel}, since the additional
summability condition is exactly \(\sum_m d_m^{-\alpha}<\infty\).
The robust-fraction and polynomial-batch addenda follow in the same way from
Corollary~\ref{cor:normalizer-tail-invmoment-borel-abundance}; the displayed
conditions \(S>1\), \(S\kappa>1\), \(\sigma>1\), and \(\sigma\kappa>1\) are
exactly the convergence tests for the extra series \(\sum_m d_m^{-S}\) and
\(\sum_m d_m^{-\sigma}\), while stretched-exponential growth makes both series
converge for every positive exponent.
The vanishing discarded-fraction addendum follows from
Corollary~\ref{cor:normalizer-tail-invmoment-borel-vanishing-fractions}; the
conditions \(B_{\rm s}-\tau>1\) and \((B_{\rm s}-\tau)\kappa>1\) are exactly the
convergence tests for \(\sum_m d_m^{-(B_{\rm s}-\tau)}\), while
stretched-exponential growth makes this series converge for every
\(B_{\rm s}-\tau>0\).

For the critical branch, the hypothesis on \(\eta\) gives
\[
  C_*\nu_{d_m,a}^{(2)}\le\frac{\theta}{\ell_m}
\]
for every \(m\), so the large-sample part of
Corollary~\ref{cor:normalizer-tail-invmoment-borel} applies under the displayed
Bernstein summability condition.  Along integer dimensions,
\[
  \exp\left(-c_{R,\theta}\frac{n(d)}{\log^2(e+d)}\right)
  \le (e+d)^{-c_{R,\theta}q},
\]
which is summable by the lower bound on \(q\); \(\sum_d d^{-A}\) is summable
because \(A>1\).  On a stretched-exponential schedule,
\[
  \exp\left(-c_{R,\theta}\frac{n_m}{\ell_m^2}\right)
  \le (e+m)^{-c_{R,\theta}q}
\]
under the first displayed lower bound on \(n_m\), and the exponent assumption
again gives summability.  Under \(n_m\ge Q\ell_m^{2+\tau}\), the same term is
bounded by \(\exp(-cQ\ell_m^\tau)\le\exp(-c' m^{\kappa\tau})\), which is
summable.  The sampled robust-certificate addendum is again the sampled part of
Corollary~\ref{cor:normalizer-tail-invmoment-borel}.
The robust-fraction and polynomial-batch critical addenda are the corresponding
parts of Corollary~\ref{cor:normalizer-tail-invmoment-borel-abundance}, with the
same sign-series convergence tests as above.
The critical vanishing discarded-fraction addendum is the corresponding
large-sample form of
Corollary~\ref{cor:normalizer-tail-invmoment-borel-vanishing-fractions}, again
with only the extra sign-series convergence test.
\end{proof}

\begin{corollary}[Power-curve dyadic-tail criterion]
\label{cor:normalizer-tail-power-curve}
Fix \(A>0\), \(a>0\), \(L>0\), \(\gamma\ge1\), and \(\beta>2\).  In the
setting of Corollary~\ref{cor:normalizer-tail-dyadic-prob}, put
\[
  F_{d,a}(t):=
  \Prob\left\{\norm{T_d(\xi)}_{\op}<t\,aN_d\right\},
  \qquad 0<t\le1 .
\]
Assume that, for the dimension \(d\),
\[
  F_{d,a}(t)\le L\,\ell^{-\gamma}t^\beta
  \qquad(0<t\le1),
  \qquad \ell=\log(e+d).
\]
Then
\[
  \mu_{d,a}^{\rm dy}\le C_\beta L\,\ell^{-\gamma}.
\]
Consequently, whenever \(n\le d\), \(n\ge C_{A,K,a}\ell\), and \(R\ge1\),
\[
  \Prob\left\{
  \max_{\omega^d=1}
  \sum_{i=1}^n
  \frac{\abs{P_i(\omega)}^2}{\norm{B_i}_{\op}^2}
  > C_{A,K,a,R}\frac n\ell
  \right\}
  \le
  C_{A,K,a,R}d^{-A}
  +\frac{C_\beta L}{R\,\ell^{\gamma-1}} .
\]
On the complement, the profile, random-sign tail, sign-averaged maximum,
deterministic certificate, and scalar edge-grid moment conclusions of
Corollary~\ref{cor:normalizer-tail-dyadic} hold at the
\(C_{A,K,a,R}\sqrt n\) scale.

In the critical-density case \(\gamma=1\), if \(R>C_\beta L\), the exceptional
probability can instead be bounded by
\[
  C_{A,K,a,R,\beta,L}d^{-A}
  +\exp\left(-c_{R,\beta,L}\frac n{\ell^2}\right).
\]
Thus a power lower-tail curve with exponent strictly larger than \(2\) permits
top-scale low-normalizer mass of order \(1/\log d\), rather than the
single-cutoff order \(1/\log^2 d\), as long as the deeper tail is charged by
the dyadic profile.
\end{corollary}

\begin{proof}
For each \(0\le r<J_\ell\), the \(r\)-th dyadic event in the definition of
\(Y_{d,a}\) has probability at most \(F_{d,a}(2^{-r})\).  Hence
\[
\begin{aligned}
  \mu_{d,a}^{\rm dy}
  &\le
  \sum_{0\le r<J_\ell}4^{r+1}F_{d,a}(2^{-r})
  +\ell F_{d,a}(2^{-J_\ell})  \\
  &\le
  \frac{4L}{\ell^\gamma}
  \sum_{r\ge0}2^{(2-\beta)r}
  +L\,\ell^{1-\gamma}2^{-\beta J_\ell}.
\end{aligned}
\]
Since \(\beta>2\), the geometric sum is finite.  Also
\(2^{J_\ell}\ge(e+\ell)^{1/2}\), and therefore
\[
  \ell\,2^{-\beta J_\ell}
  \le \ell(e+\ell)^{-\beta/2}
  \le C_\beta
\]
uniformly in \(d\).  This proves
\(\mu_{d,a}^{\rm dy}\le C_\beta L\ell^{-\gamma}\).
The Markov bound and the profile package are exactly
Corollary~\ref{cor:normalizer-tail-dyadic-prob}.

If \(\gamma=1\) and \(R>C_\beta L\), enlarge \(C_\beta\) in the preceding
display only once and choose \(\theta\) with
\(C_\beta L<\theta<R\).  The exponential form of
Corollary~\ref{cor:normalizer-tail-dyadic-prob} then gives the displayed
Bernstein bound.
\end{proof}

\begin{corollary}[Empirical bad-normalizer criterion]
\label{cor:normalizer-tail-empirical}
Fix \(A>0\), \(a>0\), and \(R\ge1\).  Let \(B_i=T_d(\xi^{(i)})\),
\(1\le i\le n\), be independent full Toeplitz matrices whose lag variables are
independent, mean zero, unit variance, and have sub-Gaussian norm at most \(K\).
Let \(P_i\) be their symbols, put \(N_d=\Exp\norm{T_d(\xi)}_{\op}\),
\(\ell=\log(e+d)\), and
\[
  M_a:=\#\{1\le i\le n:\norm{B_i}_{\op}<aN_d\}.
\]
There is a constant \(C_{A,K,a,R}\) such that, whenever
\[
  n\le d,\qquad n\ge C_{A,K,a,R}\ell ,
\]
one has
\[
  \Prob\left(
  \left\{
  \max_{\omega^d=1}
  \sum_{i=1}^n
  \frac{\abs{P_i(\omega)}^2}{\norm{B_i}_{\op}^2}
  > C_{A,K,a,R}\frac n\ell
  \right\}
  \cap
  \left\{M_a\le R\frac n{\ell^2}\right\}
  \right)
  \le C_{A,K,a,R}d^{-A}.
\]
Whenever the displayed profile bound holds, independent signs have
\[
  \Prob_\eps\left\{
  \max_{\omega^d=1}\left|\sum_{i=1}^n
  \eps_i\frac{P_i(\omega)}{\norm{B_i}_{\op}}\right|
  > C_{A,K,a,R}\sqrt n
  \right\}\le d^{-A},
\]
and
\[
  \Exp_\eps\max_{\omega^d=1}\left|\sum_{i=1}^n
  \eps_i\frac{P_i(\omega)}{\norm{B_i}_{\op}}\right|
  \le C_{A,K,a,R}\sqrt n .
\]
Thus a deterministic sign vector with edge maximum
\(C_{A,K,a,R}\sqrt n\) exists.  On the same profile event, simultaneously for
every integer \(k\ge1\),
\[
  \max_{\omega^d=1}
  \Exp_\eps\left|\sum_{i=1}^n
  \eps_i\frac{P_i(\omega)}{\norm{B_i}_{\op}}\right|^{2k}
  \le
  \left(C_{A,K,a,R}k\,\frac n\ell\right)^k .
\]
\end{corollary}

\begin{proof}
Let
\[
  G=\{i:\norm{B_i}_{\op}\ge aN_d\}.
\]
The expected-normalized edge-profile event of
Proposition~\ref{prop:expected-profile} gives
\[
  \max_{\omega^d=1}\sum_{i\in G}
  \frac{\abs{P_i(\omega)}^2}{\norm{B_i}_{\op}^2}
  \le C_{A,K,a}\frac n\ell
\]
outside an event of probability at most \(d^{-A}\), after increasing constants.
The same union bound and Paley--Zygmund/Chernoff Hilbert--Schmidt floor used in
the proof of Theorem~\ref{thm:normalizer-tail} give, outside another
\(C_{A,K}d^{-A}\) event,
\[
  \max_{i,\omega^d=1}\abs{P_i(\omega)}\le C_{A,K}\sqrt{d\ell},
  \qquad
  \min_i\norm{B_i}_{\op}\ge c_K\sqrt d .
\]
Hence every index outside \(G\) contributes at most \(C_{A,K}\ell\) to each
edge node.  On \(\{M_a\le Rn/\ell^2\}\), the full bad-normalizer contribution is
therefore at most \(C_{A,K,R}n/\ell\).  Combining the good and bad contributions
proves the displayed profile estimate.  The random-sign tail, conditional
expectation bound, deterministic certificate, and scalar moment estimates are
Proposition~\ref{prop:variance-profile},
Corollary~\ref{cor:edge-profile-certificate}, and
Proposition~\ref{prop:edge-moment-profile}.
\end{proof}

\begin{corollary}[Quantitative normalizer-tail bookkeeping]
\label{cor:normalizer-tail-quant}
Fix \(A>0\) and \(a>0\).  Let \(B_i=T_d(\xi^{(i)})\), \(1\le i\le n\), be
independent full Toeplitz matrices whose lag variables are independent, mean
zero, unit variance, and have sub-Gaussian norm at most \(K\).  Let \(P_i\) be
their symbols, put \(N_d=\Exp\norm{T_d(\xi)}_{\op}\),
\(\ell=\log(e+d)\), and
\[
  p_d(a)=\Prob\{\norm{T_d(\xi)}_{\op}<aN_d\}.
\]
There is a constant \(C_{A,K,a}\) such that, whenever
\[
  n\le d,\qquad n\ge C_{A,K,a}\ell ,
\]
one has
\[
  \Prob\left\{
  \max_{\omega^d=1}
  \sum_{i=1}^n
  \frac{\abs{P_i(\omega)}^2}{\norm{B_i}_{\op}^2}
  > C_{A,K,a}\frac n\ell
  \right\}
  \le C_{A,K,a}d^{-A}+p_d(a)\min\{n,\ell^2\}.
\]
On the complement of this exceptional event, independent signs satisfy
\[
  \Prob_\eps\left\{
  \max_{\omega^d=1}\left|\sum_{i=1}^n
  \eps_i\frac{P_i(\omega)}{\norm{B_i}_{\op}}\right|
  > C_{A,K,a}\sqrt n
  \right\}\le d^{-A},
\]
and
\[
  \Exp_\eps\max_{\omega^d=1}\left|\sum_{i=1}^n
  \eps_i\frac{P_i(\omega)}{\norm{B_i}_{\op}}\right|
  \le C_{A,K,a}\sqrt n .
\]
Thus a deterministic sign vector with edge maximum \(C_{A,K,a}\sqrt n\)
exists.  On the same event, simultaneously for every integer \(k\ge1\),
\[
  \max_{\omega^d=1}
  \Exp_\eps\left|\sum_{i=1}^n
  \eps_i\frac{P_i(\omega)}{\norm{B_i}_{\op}}\right|^{2k}
  \le
  \left(C_{A,K,a}k\,\frac n\ell\right)^k .
\]
\end{corollary}

\begin{proof}
Let
\[
  M_a:=\#\{1\le i\le n:\norm{B_i}_{\op}<aN_d\}.
\]
Since \(M_a\) is binomially dominated with mean \(np_d(a)\), the bad-count event
obeys
\[
  \Prob\{M_a>n/\ell^2\}\le p_d(a)\min\{n,\ell^2\}.
\]
Indeed, if \(n\ge\ell^2\) this is Markov's inequality, while if
\(n<\ell^2\) the event is \(\{M_a\ge1\}\), bounded by \(np_d(a)\).
Corollary~\ref{cor:normalizer-tail-empirical}, with \(R=1\), gives the profile,
signing, expectation, and moment statements outside an additional
\(C_{A,K,a}d^{-A}\) event.  Absorb constants to get the displayed bound.
\end{proof}

\begin{corollary}[Quantitative random-search and stable-center package]
\label{cor:normalizer-tail-search-stability}
Fix \(A>0\) and \(a>0\).  In the setting of
Corollary~\ref{cor:normalizer-tail-quant}, write
\[
  v_i(\omega):=\frac{P_i(\omega)}{\norm{B_i}_{\op}},
  \qquad \omega^d=1,
  \qquad
  \ell=\log(e+d).
\]
There is a constant \(C_{A,K,a}\) such that, whenever
\[
  n\le d,\qquad n\ge C_{A,K,a}\ell,
\]
outside a matrix event of probability at most
\[
  C_{A,K,a}d^{-A}+p_d(a)\min\{n,\ell^2\},
\]
the following all hold conditionally on the matrices.

First, for every integer \(M\ge1\), if
\(\delta^{(1)},\ldots,\delta^{(M)}\) are independent uniform sign vectors, then
for every \(u\ge0\),
\[
  \Prob_\delta\left\{
  \max_{1\le s\le M}\max_{\omega^d=1}
  \left|\sum_{i=1}^n\delta_i^{(s)}v_i(\omega)\right|
  >
  C_{A,K,a}\sqrt{
    n\left(1+\frac{\log(eM)+u}{\ell}\right)}
  \right\}
  \le e^{-u}.
\]
Second, for every integer \(M\ge1\),
\[
  \Prob_\delta\left\{
  \min_{1\le s\le M}\max_{\omega^d=1}
  \left|\sum_{i=1}^n\delta_i^{(s)}v_i(\omega)\right|
  > C_{A,K,a}\sqrt n
  \right\}\le d^{-AM}.
\]
More generally, if \(G_M\) denotes the number of sampled sign vectors whose edge
maximum is at most \(C_{A,K,a}\sqrt n\), then for every integer \(1\le r\le M\),
\[
  \Prob_\delta\{G_M\le M-r\}\le \binom Mr d^{-Ar}.
\]

Finally, with Rademacher probability at least \(1-d^{-A}\), a single random sign
vector is the center of a Hamming ball of radius
\[
  r_0=\left\lfloor c_{A,K,a}\ell\right\rfloor
\]
on which every sign vector in the ball has edge maximum at most
\(C_{A,K,a}\sqrt n\).  If \(H_M\) is the number of sampled signs among
\(\delta^{(1)},\ldots,\delta^{(M)}\) that are centers of such good Hamming balls,
then for every integer \(1\le r\le M\),
\[
  \Prob_\delta\{H_M\le M-r\}\le \binom Mr d^{-Ar}.
\]
\end{corollary}

\begin{proof}
Corollary~\ref{cor:normalizer-tail-quant} gives, outside the displayed matrix
exceptional event,
\[
  \max_{\omega^d=1}\sum_{i=1}^n |v_i(\omega)|^2
  \le C_{A,K,a}\frac n\ell .
\]
Condition on any matrix realization satisfying this profile.  The sampled-list
bound and the binomial success-count estimate are exactly
Corollary~\ref{cor:edge-profile-random-search}, after absorbing the profile
constant into \(C_{A,K,a}\).  The Hamming-ball assertion and the sampled
robust-center count are Corollary~\ref{cor:edge-profile-hamming}, again with the
same absorption of constants.
\end{proof}

\begin{corollary}[Polynomial-confidence normalized robust search]
\label{cor:normalizer-tail-poly-confidence}
Fix \(A>0\), \(B>0\), and \(a>0\).  In the setting of
Corollary~\ref{cor:normalizer-tail-quant}, let
\(\delta^{(1)},\ldots,\delta^{(M)}\) be independent uniform sign vectors,
independent of the matrices, with \(M\ge\lceil B\rceil\).  There is a constant
\(C_{A,K,a}\) and a constant \(c_{A,K,a}>0\) such that, whenever
\[
  n\le d,\qquad n\ge C_{A,K,a}\ell,\qquad \ell=\log(e+d),
\]
the joint matrix-and-sign probability that no sampled sign is a robust
certificate is at most
\[
  C_{A,K,a}d^{-A}+p_d(a)\min\{n,\ell^2\}+d^{-B}.
\]
Equivalently, outside that exceptional probability, some sampled
\(\delta^{(s)}\) is the center of a Hamming ball of radius
\[
  r_0=\left\lfloor c_{A,K,a}\ell\right\rfloor
\]
on which every sign vector \(\eps'\) satisfies
\[
  \max_{\omega^d=1}\left|\sum_{i=1}^n
  \eps'_i\frac{P_i(\omega)}{\norm{B_i}_{\op}}\right|
  \le C_{A,K,a}\sqrt n .
\]
Thus after the same normalizer-tail matrix event, a constant-size random list
finds a logarithmically stable \(O(\sqrt n)\) certificate at any prescribed
polynomial confidence level.
\end{corollary}

\begin{proof}
On the good matrix event of Corollary~\ref{cor:normalizer-tail-quant},
\[
  \max_{\omega^d=1}\sum_i
  \left|\frac{P_i(\omega)}{\norm{B_i}_{\op}}\right|^2
  \le C_{A,K,a}\frac n\ell .
\]
Condition on such a realization and apply
Corollary~\ref{cor:edge-profile-poly-confidence}, absorbing the displayed
profile constant into the constants.  The conditional sign-search failure is at
most \(d^{-B}\).  Adding the matrix exceptional probability from
Corollary~\ref{cor:normalizer-tail-quant} proves the claim.
\end{proof}

\begin{corollary}[Normalizer-tail sampled-list abundance]
\label{cor:normalizer-tail-abundance}
Fix \(A>0\), \(B>0\), \(a>0\), and a fraction \(0<\rho\le1\).  In the setting
of Corollary~\ref{cor:normalizer-tail-quant}, write
\[
  \ell=\log(e+d),\qquad
  v_i(\omega):=\frac{P_i(\omega)}{\norm{B_i}_{\op}},
  \qquad \omega^d=1.
\]
Let \(\delta^{(1)},\ldots,\delta^{(M)}\) be independent uniform sign vectors,
independent of the matrices, where \(M\ge1\), and put
\[
  r_\rho=\lceil \rho M\rceil .
\]
There are constants \(C_{A,B,K,a}\) and \(c_{A,B,K,a}>0\) such that, whenever
\[
  n\le d,\qquad n\ge C_{A,B,K,a}\ell ,
\]
then, for every \(u\ge0\), with probability at least
\[
  1-C_{A,B,K,a}\left(d^{-A}+p_d(a)\min\{n,\ell^2\}\right)
  -e^{-u}-2\binom{M}{r_\rho}d^{-B r_\rho},
\]
all three conclusions below hold simultaneously.

First, the whole sampled list obeys
\[
\begin{aligned}
  \max_{1\le s\le M}\max_{\omega^d=1}
  \left|\sum_{i=1}^n\delta_i^{(s)}
  \frac{P_i(\omega)}{\norm{B_i}_{\op}}\right|
  \le
  C_{A,B,K,a}\sqrt{
    n\left(1+\frac{\log(eM)+u}{\ell}\right)} .
\end{aligned}
\]
Second, at least \(M-r_\rho+1\) sampled signs have edge maximum at most
\[
  C_{A,B,K,a}\sqrt n .
\]
Third, at least \(M-r_\rho+1\) sampled signs are centers of Hamming balls of
radius
\[
  \left\lfloor c_{A,B,K,a}\ell\right\rfloor
\]
on which every signing \(\eps'\) has edge maximum at most
\[
  C_{A,B,K,a}\sqrt n .
\]
In particular, at least a \(1-\rho\) fraction of the sampled signs are good,
and at least a \(1-\rho\) fraction are robust centers.
\end{corollary}

\begin{proof}
By Corollary~\ref{cor:normalizer-tail-quant}, outside a matrix event of
probability at most
\[
  C_{A,K,a}d^{-A}+p_d(a)\min\{n,\ell^2\},
\]
one has the edge profile
\[
  \max_{\omega^d=1}\sum_i |v_i(\omega)|^2
  \le C_{A,K,a}\frac n\ell .
\]
Condition on a matrix realization satisfying this profile.  The sampled-list
maximum follows from Corollary~\ref{cor:edge-profile-random-search}.  Applying
the same corollary with exponent \(B\) gives the good-sign count tail
\(\binom{M}{r_\rho}d^{-B r_\rho}\), after absorbing the profile constant into
the displayed constants.  Corollary~\ref{cor:edge-profile-hamming}, again with
exponent \(B\), gives the same tail for the robust-center count.  A union bound
over the matrix event and the three sign failures proves the displayed
probability bound.  Finally \(M-r_\rho+1\ge M-\rho M\), which gives the
fraction statements.
\end{proof}

\begin{corollary}[Normalizer-tail polynomial-confidence robust fractions]
\label{cor:normalizer-tail-robust-fraction-confidence}
Fix \(A>0\), \(S>0\), \(a>0\), and a fraction \(0<\rho\le1\).  Choose \(B>0\)
and an integer \(M\ge1\) so that
\[
  B\ge 2\log_2(e/\rho),
  \qquad
  M\ge \frac{2(S+1)}{B\rho}.
\]
Write \(\ell=\log(e+d)\).  In the setting of
Corollary~\ref{cor:normalizer-tail-quant}, suppose \(d\ge2\), and let
\(\delta^{(1)},\ldots,\delta^{(M)}\) be independent uniform sign vectors,
independent of the matrices.  There are constants \(C_{A,B,K,a}\) and
\(c_{A,B,K,a}>0\) such that, whenever
\[
  n\le d,\qquad n\ge C_{A,B,K,a}\ell ,
\]
then, for every \(u\ge0\), with probability at least
\[
  1-C_{A,B,K,a}\left(d^{-A}+p_d(a)\min\{n,\ell^2\}\right)
  -e^{-u}-d^{-S},
\]
the whole sampled list obeys the entropy-sensitive maximum bound from
Corollary~\ref{cor:normalizer-tail-abundance}, and at least a \(1-\rho\)
fraction of the sampled signs are centers of radius
\[
  \left\lfloor c_{A,B,K,a}\ell\right\rfloor
\]
Hamming balls on which every signing \(\eps'\) has edge maximum at most
\[
  C_{A,B,K,a}\sqrt n .
\]
The same \(1-\rho\) fraction also has edge maximum at most
\(C_{A,B,K,a}\sqrt n\).

In particular, taking \(u=S\ell\), the same event has probability at least
\[
  1-C_{A,B,K,a}\left(d^{-A}+p_d(a)\min\{n,\ell^2\}\right)-2d^{-S},
\]
and the whole sampled list has edge maximum at most
\[
  C_{A,B,K,a,S,\rho}\sqrt n .
\]
\end{corollary}

\begin{proof}
Let \(r_\rho=\lceil\rho M\rceil\).  By
Corollary~\ref{cor:normalizer-tail-abundance}, it remains only to simplify the
binomial fraction tail.  Since \(r_\rho\ge\rho M\),
\[
  \binom{M}{r_\rho}\le \left(\frac{eM}{r_\rho}\right)^{r_\rho}
  \le (e/\rho)^{r_\rho}.
\]
For \(d\ge2\) and \(B\ge2\log_2(e/\rho)\),
\[
  \binom{M}{r_\rho}d^{-B r_\rho}
  \le d^{-B r_\rho/2}.
\]
The lower bound on \(M\) gives \(B r_\rho/2\ge S+1\), so
\[
  2\binom{M}{r_\rho}d^{-B r_\rho}\le 2d^{-(S+1)}\le d^{-S}.
\]
Substituting this into Corollary~\ref{cor:normalizer-tail-abundance} proves the
first assertion.  If \(u=S\ell\), then \(e^{-u}\le d^{-S}\), and
\[
  1+\frac{\log(eM)+u}{\ell}
  \le C_{S,\rho,B}
\]
because \(M\) is fixed once \(S,\rho,B\) are fixed.  This gives the
constant-scale sampled-list maximum.
\end{proof}

\begin{corollary}[Normalizer-tail polynomial-batch robust search]
\label{cor:normalizer-tail-poly-batch}
Fix \(A>0\), \(a>0\), \(\beta\ge0\), and \(\sigma>0\).  Write
\(\ell=\log(e+d)\).  In the setting of
Corollary~\ref{cor:normalizer-tail-quant}, let
\(\delta^{(1)},\ldots,\delta^{(M)}\) be independent uniform sign vectors,
independent of the matrices, with
\[
  1\le M\le d^\beta .
\]
There are constants \(C_{A,K,a,\beta,\sigma}\) and
\(c_{A,K,a,\beta,\sigma}>0\) such that, whenever
\[
  n\le d,\qquad n\ge C_{A,K,a,\beta,\sigma}\ell ,
\]
with probability at least
\[
  1-C_{A,K,a,\beta,\sigma}
  \left(d^{-A}+p_d(a)\min\{n,\ell^2\}\right)-d^{-\sigma},
\]
every sampled sign \(\delta^{(s)}\), \(1\le s\le M\), is the center of a
Hamming ball of radius
\[
  r_0=\left\lfloor c_{A,K,a,\beta,\sigma}\ell\right\rfloor
\]
on which every sign vector \(\eps'\) satisfies
\[
  \max_{\omega^d=1}\left|\sum_{i=1}^n
  \eps'_i\frac{P_i(\omega)}{\norm{B_i}_{\op}}\right|
  \le C_{A,K,a,\beta,\sigma}\sqrt n .
\]
In particular, the whole sampled batch itself has edge maximum at most
\[
  C_{A,K,a,\beta,\sigma}\sqrt n .
\]
\end{corollary}

\begin{proof}
By Corollary~\ref{cor:normalizer-tail-quant}, outside a matrix event of
probability at most
\[
  C_{A,K,a}d^{-A}+p_d(a)\min\{n,\ell^2\},
\]
one has the edge profile
\[
  \max_{\omega^d=1}\sum_i
  \left|\frac{P_i(\omega)}{\norm{B_i}_{\op}}\right|^2
  \le C_{A,K,a}\frac n\ell .
\]
Condition on a matrix realization satisfying this profile and apply
Corollary~\ref{cor:edge-profile-hamming} with sign exponent
\[
  B=\beta+\sigma .
\]
The robust-center count fails with conditional sign probability at most
\[
  \binom{M}{1}d^{-(\beta+\sigma)}
  \le d^{-\sigma}.
\]
On the complementary event \(H_M\ge M\), every sampled sign is a robust center,
and hence every sampled sign also has edge maximum at most the displayed
constant times \(\sqrt n\).  Absorb the dependence on \(\beta,\sigma\) into the
constants.
\end{proof}

\begin{corollary}[Power-curve random-search and stable-center criterion]
\label{cor:normalizer-tail-power-search}
Fix \(A>0\), \(a>0\), \(L>0\), \(\gamma\ge1\), \(\beta>2\), and \(R\ge1\).
In the setting of Corollary~\ref{cor:normalizer-tail-power-curve}, assume
\[
  F_{d,a}(t)\le L\,\ell^{-\gamma}t^\beta
  \qquad(0<t\le1),
  \qquad \ell=\log(e+d).
\]
Then there is a constant \(C_{A,K,a,R}\) such that, whenever
\[
  n\le d,\qquad n\ge C_{A,K,a,R}\ell ,
\]
outside a matrix event of probability at most
\[
  C_{A,K,a,R}d^{-A}
  +\frac{C_\beta L}{R\,\ell^{\gamma-1}},
\]
the following all hold conditionally on the matrices.  With
\[
  v_i(\omega):=\frac{P_i(\omega)}{\norm{B_i}_{\op}},
  \qquad \omega^d=1,
\]
every \(M\)-sample of independent uniform sign vectors
\(\delta^{(1)},\ldots,\delta^{(M)}\) satisfies, for every \(u\ge0\),
\[
  \Prob_\delta\left\{
  \max_{1\le s\le M}\max_{\omega^d=1}
  \left|\sum_{i=1}^n\delta_i^{(s)}v_i(\omega)\right|
  >
  C_{A,K,a,R}\sqrt{
    n\left(1+\frac{\log(eM)+u}{\ell}\right)}
  \right\}
  \le e^{-u}.
\]
Moreover,
\[
  \Prob_\delta\left\{
  \min_{1\le s\le M}\max_{\omega^d=1}
  \left|\sum_{i=1}^n\delta_i^{(s)}v_i(\omega)\right|
  > C_{A,K,a,R}\sqrt n
  \right\}\le d^{-AM},
\]
and if \(G_M\) counts sampled signs whose edge maximum is at most
\(C_{A,K,a,R}\sqrt n\), then
\[
  \Prob_\delta\{G_M\le M-r\}\le \binom Mr d^{-Ar}
  \qquad(1\le r\le M).
\]
Finally, with Rademacher probability at least \(1-d^{-A}\), a random sign vector
is the center of a Hamming ball of radius
\[
  r_0=\left\lfloor c_{A,K,a,R}\ell\right\rfloor
\]
on which every sign vector has edge maximum at most \(C_{A,K,a,R}\sqrt n\);
the sampled robust-center count obeys the same binomial tail
\[
  \Prob_\delta\{H_M\le M-r\}\le \binom Mr d^{-Ar}
  \qquad(1\le r\le M).
\]

If \(\gamma=1\) and \(R>C_\beta L\), the matrix exceptional probability can
instead be bounded by
\[
  C_{A,K,a,R,\beta,L}d^{-A}
  +\exp\left(-c_{R,\beta,L}\frac n{\ell^2}\right).
\]
\end{corollary}

\begin{proof}
Corollary~\ref{cor:normalizer-tail-power-curve} gives the edge profile
\[
  \max_{\omega^d=1}\sum_{i=1}^n |v_i(\omega)|^2
  \le C_{A,K,a,R}\frac n\ell
\]
outside the displayed matrix event, with the exponential replacement in the
critical-density case.  Condition on any realization satisfying this profile.
The sampled-list bound and success-count estimate follow from
Corollary~\ref{cor:edge-profile-random-search}, while the Hamming-ball and
sampled robust-center statements follow from
Corollary~\ref{cor:edge-profile-hamming}, after absorbing the profile constant
into \(C_{A,K,a,R}\) and \(c_{A,K,a,R}\).
\end{proof}

\begin{corollary}[Power-curve sampled-list abundance]
\label{cor:normalizer-tail-power-abundance}
Fix \(A>0\), \(B>0\), \(a>0\), \(L>0\), \(\gamma\ge1\), \(\beta>2\),
\(R\ge1\), and a fraction \(0<\rho\le1\).  In the setting of
Corollary~\ref{cor:normalizer-tail-power-curve}, assume
\[
  F_{d,a}(t)\le L\,\ell^{-\gamma}t^\beta
  \qquad(0<t\le1),
  \qquad \ell=\log(e+d).
\]
Let \(\delta^{(1)},\ldots,\delta^{(M)}\) be independent uniform sign vectors,
independent of the matrices, where \(M\ge1\), and put
\[
  r_\rho=\lceil\rho M\rceil .
\]
There are constants \(C_{A,B,K,a,R}\) and \(c_{A,B,K,a,R}>0\) such that,
whenever
\[
  n\le d,\qquad n\ge C_{A,B,K,a,R}\ell ,
\]
then, for every \(u\ge0\), with probability at least
\[
  1-C_{A,B,K,a,R}\left(d^{-A}+\frac{C_\beta L}{R\ell^{\gamma-1}}\right)
  -e^{-u}-2\binom{M}{r_\rho}d^{-B r_\rho},
\]
all three conclusions below hold simultaneously.

First, the whole sampled list obeys
\[
\begin{aligned}
  \max_{1\le s\le M}\max_{\omega^d=1}
  \left|\sum_{i=1}^n\delta_i^{(s)}
  \frac{P_i(\omega)}{\norm{B_i}_{\op}}\right|
  \le
  C_{A,B,K,a,R}\sqrt{
    n\left(1+\frac{\log(eM)+u}{\ell}\right)} .
\end{aligned}
\]
Second, at least \(M-r_\rho+1\) sampled signs have edge maximum at most
\[
  C_{A,B,K,a,R}\sqrt n .
\]
Third, at least \(M-r_\rho+1\) sampled signs are centers of Hamming balls of
radius
\[
  \left\lfloor c_{A,B,K,a,R}\ell\right\rfloor
\]
on which every signing \(\eps'\) has edge maximum at most
\[
  C_{A,B,K,a,R}\sqrt n .
\]
In particular, at least a \(1-\rho\) fraction of the sampled signs are good,
and at least a \(1-\rho\) fraction are robust centers.

If \(\gamma=1\) and \(R>C_\beta L\), then the same conclusion holds with the
matrix part of the exceptional probability replaced by
\[
  C_{A,B,K,a,R,\beta,L}d^{-A}
  +\exp\left(-c_{R,\beta,L}\frac n{\ell^2}\right).
\]
\end{corollary}

\begin{proof}
By Corollary~\ref{cor:normalizer-tail-power-curve}, outside a matrix event of
probability at most
\[
  C_{A,K,a,R}d^{-A}+\frac{C_\beta L}{R\ell^{\gamma-1}},
\]
one has the edge profile
\[
  \max_{\omega^d=1}\sum_i
  \left|\frac{P_i(\omega)}{\norm{B_i}_{\op}}\right|^2
  \le C_{A,K,a,R}\frac n\ell .
\]
Condition on a matrix realization satisfying this profile.  The sampled-list
maximum follows from Corollary~\ref{cor:edge-profile-random-search}.  Applying
the same corollary with exponent \(B\) gives the good-sign count tail
\(\binom{M}{r_\rho}d^{-B r_\rho}\), after absorbing the profile constant into
the displayed constants.  Corollary~\ref{cor:edge-profile-hamming}, again with
exponent \(B\), gives the same tail for the robust-center count.  A union bound
over the matrix event and the three sign failures proves the displayed
probability bound.

In the critical-density case, use the Bernstein replacement in
Corollary~\ref{cor:normalizer-tail-power-curve} for the matrix event.  Finally
\(M-r_\rho+1\ge M-\rho M\), which gives the fraction statements.
\end{proof}

\begin{corollary}[Power-curve polynomial-confidence robust fractions]
\label{cor:normalizer-tail-power-robust-fraction-confidence}
Fix \(A>0\), \(S>0\), \(a>0\), \(L>0\), \(\gamma\ge1\), \(\beta>2\),
\(R\ge1\), and a fraction \(0<\rho\le1\).  Choose \(B>0\) and an integer
\(M\ge1\) so that
\[
  B\ge 2\log_2(e/\rho),
  \qquad
  M\ge \frac{2(S+1)}{B\rho}.
\]
In the setting of Corollary~\ref{cor:normalizer-tail-power-curve}, suppose
\(d\ge2\) and
\[
  F_{d,a}(t)\le L\,\ell^{-\gamma}t^\beta
  \qquad(0<t\le1),
  \qquad \ell=\log(e+d).
\]
Let \(\delta^{(1)},\ldots,\delta^{(M)}\) be independent uniform sign vectors,
independent of the matrices.  There are constants \(C_{A,B,K,a,R}\) and
\(c_{A,B,K,a,R}>0\) such that, whenever
\[
  n\le d,\qquad n\ge C_{A,B,K,a,R}\ell ,
\]
then, for every \(u\ge0\), with probability at least
\[
  1-C_{A,B,K,a,R}\left(d^{-A}+\frac{C_\beta L}{R\ell^{\gamma-1}}\right)
  -e^{-u}-d^{-S},
\]
the whole sampled list obeys the entropy-sensitive maximum bound from
Corollary~\ref{cor:normalizer-tail-power-abundance}, and at least a
\(1-\rho\) fraction of the sampled signs are centers of radius
\[
  \left\lfloor c_{A,B,K,a,R}\ell\right\rfloor
\]
Hamming balls on which every signing \(\eps'\) has edge maximum at most
\[
  C_{A,B,K,a,R}\sqrt n .
\]
The same \(1-\rho\) fraction also has edge maximum at most
\(C_{A,B,K,a,R}\sqrt n\).

In particular, taking \(u=S\ell\), the same event has probability at least
\[
  1-C_{A,B,K,a,R}\left(d^{-A}+\frac{C_\beta L}{R\ell^{\gamma-1}}\right)
  -2d^{-S},
\]
and the whole sampled list has edge maximum at most
\[
  C_{A,B,K,a,R,S,\rho}\sqrt n .
\]

If \(\gamma=1\) and \(R>C_\beta L\), then both conclusions remain valid with the
matrix part of the exceptional probability replaced by
\[
  C_{A,B,K,a,R,\beta,L}d^{-A}
  +\exp\left(-c_{R,\beta,L}\frac n{\ell^2}\right).
\]
\end{corollary}

\begin{proof}
Let \(r_\rho=\lceil\rho M\rceil\).  By
Corollary~\ref{cor:normalizer-tail-power-abundance}, it remains only to
simplify the binomial fraction tail.  Since \(r_\rho\ge\rho M\),
\[
  \binom{M}{r_\rho}\le \left(\frac{eM}{r_\rho}\right)^{r_\rho}
  \le (e/\rho)^{r_\rho}.
\]
For \(d\ge2\) and \(B\ge2\log_2(e/\rho)\),
\[
  \binom{M}{r_\rho}d^{-B r_\rho}
  \le d^{-B r_\rho/2}.
\]
The lower bound on \(M\) gives \(B r_\rho/2\ge S+1\), hence
\[
  2\binom{M}{r_\rho}d^{-B r_\rho}\le 2d^{-(S+1)}\le d^{-S}.
\]
Substituting this into Corollary~\ref{cor:normalizer-tail-power-abundance}
proves the first assertion, including the critical-density replacement.  If
\(u=S\ell\), then \(e^{-u}\le d^{-S}\), and
\[
  1+\frac{\log(eM)+u}{\ell}
  \le C_{S,\rho,B}.
\]
This gives the constant-scale sampled-list maximum.
\end{proof}

\begin{corollary}[Power-curve polynomial-batch robust search]
\label{cor:normalizer-tail-power-poly-batch}
Fix \(A>0\), \(a>0\), \(L>0\), \(\gamma\ge1\), \(\beta_0>2\), \(R\ge1\),
\(\beta_{\rm b}\ge0\), and \(\sigma>0\).  In the setting of
Corollary~\ref{cor:normalizer-tail-power-curve}, assume
\[
  F_{d,a}(t)\le L\,\ell^{-\gamma}t^{\beta_0}
  \qquad(0<t\le1),
  \qquad \ell=\log(e+d).
\]
Let \(\delta^{(1)},\ldots,\delta^{(M)}\) be independent uniform sign vectors,
independent of the matrices, with
\[
  1\le M\le d^{\beta_{\rm b}} .
\]
There are constants \(C_{A,K,a,R,\beta_{\rm b},\sigma}\) and
\(c_{A,K,a,R,\beta_{\rm b},\sigma}>0\) such that, whenever
\[
  n\le d,\qquad n\ge C_{A,K,a,R,\beta_{\rm b},\sigma}\ell ,
\]
with probability at least
\[
  1-C_{A,K,a,R,\beta_{\rm b},\sigma}
  \left(d^{-A}+\frac{C_{\beta_0}L}{R\ell^{\gamma-1}}\right)-d^{-\sigma},
\]
every sampled sign \(\delta^{(s)}\), \(1\le s\le M\), is the center of a
Hamming ball of radius
\[
  r_0=\left\lfloor c_{A,K,a,R,\beta_{\rm b},\sigma}\ell\right\rfloor
\]
on which every sign vector \(\eps'\) satisfies
\[
  \max_{\omega^d=1}\left|\sum_{i=1}^n
  \eps'_i\frac{P_i(\omega)}{\norm{B_i}_{\op}}\right|
  \le C_{A,K,a,R,\beta_{\rm b},\sigma}\sqrt n .
\]
In particular, the whole sampled batch itself has edge maximum at most
\[
  C_{A,K,a,R,\beta_{\rm b},\sigma}\sqrt n .
\]

If \(\gamma=1\) and \(R>C_{\beta_0}L\), then the same conclusion holds with the
matrix part of the exceptional probability replaced by
\[
  C_{A,K,a,R,\beta_0,L,\beta_{\rm b},\sigma}d^{-A}
  +\exp\left(-c_{R,\beta_0,L}\frac n{\ell^2}\right).
\]
\end{corollary}

\begin{proof}
By Corollary~\ref{cor:normalizer-tail-power-curve}, outside the displayed
power-curve matrix event one has the edge profile
\[
  \max_{\omega^d=1}\sum_i
  \left|\frac{P_i(\omega)}{\norm{B_i}_{\op}}\right|^2
  \le C_{A,K,a,R}\frac n\ell .
\]
Condition on a matrix realization satisfying this profile and apply
Corollary~\ref{cor:edge-profile-hamming} with sign exponent
\[
  B=\beta_{\rm b}+\sigma .
\]
The robust-center count fails with conditional sign probability at most
\[
  \binom{M}{1}d^{-(\beta_{\rm b}+\sigma)}
  \le d^{-\sigma}.
\]
On the complementary event every sampled sign is a robust center, and therefore
also has edge maximum at most the displayed constant times \(\sqrt n\).  The
critical-density assertion again uses the Bernstein matrix-event replacement in
Corollary~\ref{cor:normalizer-tail-power-curve}.
\end{proof}

\begin{corollary}[Almost-sure power-curve search transfer]
\label{cor:normalizer-tail-power-borel-search}
Fix \(A>0\), \(a>0\), \(L>0\), \(\gamma\ge1\), \(\beta>2\), and \(R\ge1\).
There is a constant \(C_{A,K,a,R}\) with the following property.  Let
\((d_m,n_m)_{m\ge1}\) be a triangular-array sequence with \(n_m\le d_m\), put
\(\ell_m=\log(e+d_m)\), and assume
\[
  n_m\ge C_{A,K,a,R}\ell_m \qquad(m\ge1).
\]
At each \(m\), let \(B_{m,i}=T_{d_m}(\xi^{(m,i)})\), \(1\le i\le n_m\), be
independent full Toeplitz matrices with independent mean-zero, unit-variance
lag variables of sub-Gaussian norm at most \(K\), and let \(P_{m,i}\) be their
symbols.  Suppose that, for every \(m\),
\[
  F_{d_m,a}(t):=
  \Prob\left\{\norm{T_{d_m}(\xi)}_{\op}
  <t\,a\,\Exp\norm{T_{d_m}(\xi)}_{\op}\right\}
  \le L\,\ell_m^{-\gamma}t^\beta
  \qquad(0<t\le1).
\]
If
\[
  \sum_{m=1}^{\infty}
  \left(d_m^{-A}+\ell_m^{1-\gamma}\right)<\infty,
\]
then, with probability one over the matrix triangular array, all but finitely
many \(m\) satisfy the full random-search and stable-center package of
Corollary~\ref{cor:normalizer-tail-power-search}, with
\(d,n,\ell\) replaced by \(d_m,n_m,\ell_m\).

There is also a critical-density Bernstein form.  If \(\gamma=1\) and
\(R>C_\beta L\), then the same almost-sure conclusion holds under the weaker
summability condition
\[
  \sum_{m=1}^{\infty}
  \left(d_m^{-A}
  +\exp\left(-c_{R,\beta,L}\frac{n_m}{\ell_m^2}\right)\right)<\infty .
\]
Consequently, along the full integer dimension sequence, if \(A>1\),
\[
  F_{d,a}(t)\le L\,\log^{-1}(e+d)\,t^\beta
  \qquad(0<t\le1,\ d\ge1),
\]
\(\beta>2\), \(R>C_\beta L\), and if \(n(d)\le d\) while, for all sufficiently
large \(d\),
\[
  n(d)\ge q\,\log^3(e+d)
  \qquad\text{for some } q>\frac1{c_{R,\beta,L}},
\]
then almost surely all sufficiently large integer dimensions satisfy the same
random-search and stable-center package.
\end{corollary}

\begin{proof}
Apply Corollary~\ref{cor:normalizer-tail-power-search} at each \(m\), enlarging
\(C_{A,K,a,R}\) if necessary.  In the Markov form, the matrix exceptional
probability is at most
\[
  C_{A,K,a,R}d_m^{-A}+C_{\beta}L R^{-1}\ell_m^{1-\gamma},
\]
which is summable by hypothesis.  The first Borel--Cantelli lemma gives an
almost-sure event on which all but finitely many matrix arrays satisfy the
finite power-curve search package.

In the critical-density case, the exponential replacement in
Corollary~\ref{cor:normalizer-tail-power-search} gives summable exceptional
probabilities under the displayed Bernstein condition, so the same
Borel--Cantelli argument applies.  For the full-dimension consequence take
\(d_m=m\).  Then \(\sum_m m^{-A}<\infty\), and for all large \(m\),
\[
  \exp\left(-c_{R,\beta,L}\frac{n(m)}{\log^2(e+m)}\right)
  \le (e+m)^{-c_{R,\beta,L}q},
\]
which is summable because \(c_{R,\beta,L}q>1\).  The finitely many initial
dimensions do not affect the eventual almost-sure conclusion.
\end{proof}

\begin{corollary}[Power-curve sampled robust certificates]
\label{cor:normalizer-tail-power-borel-sampled}
Fix \(A>0\), \(\alpha>0\), \(a>0\), \(L>0\), \(\gamma\ge1\), \(\beta>2\), and
\(R\ge1\).  There is a constant \(C_{A,K,a,R}\) with the following property.
Let \((d_m,n_m)_{m\ge1}\) be as in
Corollary~\ref{cor:normalizer-tail-power-borel-search}, with
\(\ell_m=\log(e+d_m)\), and suppose
\[
  n_m\ge C_{A,K,a,R}\ell_m\qquad(m\ge1).
\]
Assume the same power-curve lower-tail bound
\[
  F_{d_m,a}(t)\le L\,\ell_m^{-\gamma}t^\beta
  \qquad(0<t\le1,\ m\ge1).
\]
For each \(m\), sample \(M_m\) independent uniform sign vectors in
\(\{-1,1\}^{n_m}\), independently of the matrices, with
\[
  M_m\ge\lceil\alpha\rceil .
\]
If
\[
  \sum_{m=1}^{\infty}
  \left(d_m^{-A}+\ell_m^{1-\gamma}+d_m^{-\alpha}\right)<\infty,
\]
then, with probability one over the matrices and the sampled signs, all but
finitely many \(m\) have a sampled sign which is the center of a Hamming ball of
radius
\[
  r_m=\left\lfloor c_{A,K,a,R}\ell_m\right\rfloor
\]
on which every signing \(\eps'\) satisfies
\[
  \max_{\omega^{d_m}=1}\left|\sum_{i=1}^{n_m}
  \eps'_i\frac{P_{m,i}(\omega)}{\norm{B_{m,i}}_{\op}}\right|
  \le C_{A,K,a,R}\sqrt{n_m}.
\]

In the critical-density case \(\gamma=1\), if \(R>C_\beta L\), the same actual
sampled-certificate conclusion holds under the weaker summability condition
\[
  \sum_{m=1}^{\infty}
  \left(d_m^{-A}
  +\exp\left(-c_{R,\beta,L}\frac{n_m}{\ell_m^2}\right)
  +d_m^{-\alpha}\right)<\infty .
\]
Consequently, along the full integer dimension sequence, if \(A>1\),
\(\alpha>1\),
\[
  F_{d,a}(t)\le L\,\log^{-1}(e+d)\,t^\beta
  \qquad(0<t\le1,\ d\ge1),
\]
\(\beta>2\), \(R>C_\beta L\), and
\[
  n(d)\ge q\,\log^3(e+d)
  \qquad\text{for all large }d,\qquad
  q>\frac1{c_{R,\beta,L}},
\]
then any fixed list size \(M\ge\lceil\alpha\rceil\) gives, almost surely,
sampled logarithmically robust \(O(\sqrt{n(d)})\) certificates for all
sufficiently large integer dimensions.
\end{corollary}

\begin{proof}
Apply Corollary~\ref{cor:normalizer-tail-power-search} at each \(m\).  In the
Markov form, outside a matrix event of probability at most
\[
  C_{A,K,a,R}d_m^{-A}+C_\beta L R^{-1}\ell_m^{1-\gamma},
\]
the edge profile
\[
  \max_{\omega^{d_m}=1}
  \sum_{i=1}^{n_m}\left|
  \frac{P_{m,i}(\omega)}{\norm{B_{m,i}}_{\op}}\right|^2
  \le C_{A,K,a,R}\frac{n_m}{\ell_m}
\]
holds.  Condition on such a realization and apply
Corollary~\ref{cor:edge-profile-poly-confidence}, absorbing the profile constant.
The sampled-sign failure is at most \(d_m^{-\alpha}\).  Therefore the total
failure probability at index \(m\) is bounded by a constant times
\[
  d_m^{-A}+\ell_m^{1-\gamma}+d_m^{-\alpha},
\]
which is summable by hypothesis.  Borel--Cantelli gives the first assertion.

For \(\gamma=1\) and \(R>C_\beta L\), use the Bernstein exceptional probability
from Corollary~\ref{cor:normalizer-tail-power-search} in place of the Markov
term.  This gives the second displayed summability condition.  The
full-dimension consequence is obtained by taking \(d_m=m\), \(n_m=n(m)\), and
\(M_m=M\).  The series \(\sum_m m^{-A}\) and \(\sum_m m^{-\alpha}\) are finite,
and the displayed lower bound on \(n(m)\) makes
\[
  \exp\left(-c_{R,\beta,L}\frac{n(m)}{\log^2(e+m)}\right)
  \le (e+m)^{-c_{R,\beta,L}q}
\]
summable.
\end{proof}

\begin{corollary}[Almost-sure power-curve sampled-search counts]
\label{cor:normalizer-tail-power-borel-counts}
Fix \(A>0\), \(B>0\), \(a>0\), \(L>0\), \(\gamma\ge1\), \(\beta>2\), and
\(R\ge1\).  There are constants \(C_{A,B,K,a,R}\) and \(c_{A,B,K,a,R}>0\) with
the following property.  Let \((d_m,n_m)_{m\ge1}\) be a triangular-array sequence
with \(n_m\le d_m\), put \(\ell_m=\log(e+d_m)\), and assume
\[
  n_m\ge C_{A,B,K,a,R}\ell_m \qquad(m\ge1).
\]
At each \(m\), let \(B_{m,i}=T_{d_m}(\xi^{(m,i)})\), \(1\le i\le n_m\), be
independent full Toeplitz matrices with independent mean-zero, unit-variance lag
variables of sub-Gaussian norm at most \(K\), and let \(P_{m,i}\) be their
symbols.  Suppose that, for every \(m\),
\[
  F_{d_m,a}(t):=
  \Prob\left\{\norm{T_{d_m}(\xi)}_{\op}
  <t\,a\,\Exp\norm{T_{d_m}(\xi)}_{\op}\right\}
  \le L\,\ell_m^{-\gamma}t^\beta
  \qquad(0<t\le1).
\]
Independently of the matrices, sample independent uniform sign vectors
\[
  \delta_m^{(1)},\ldots,\delta_m^{(M_m)}\in\{-1,1\}^{n_m},
  \qquad M_m\ge1 .
\]
Let \(u_m\ge0\) and \(1\le r_m\le M_m\) be deterministic parameters.  If
\[
  \sum_{m=1}^{\infty}
  \left(d_m^{-A}+\ell_m^{1-\gamma}+e^{-u_m}
  +\binom{M_m}{r_m}d_m^{-B r_m}\right)<\infty,
\]
then, with probability one over matrices and sampled signs, all but finitely many
\(m\) satisfy all of the following.

First, the whole sampled list obeys
\[
\begin{aligned}
  \max_{1\le s\le M_m}\max_{\omega^{d_m}=1}
  \left|\sum_{i=1}^{n_m}\delta_{m,i}^{(s)}
  \frac{P_{m,i}(\omega)}{\norm{B_{m,i}}_{\op}}\right|
  \le
  C_{A,B,K,a,R}\sqrt{
    n_m\left(1+\frac{\log(eM_m)+u_m}{\ell_m}\right)} .
\end{aligned}
\]
Second, at least \(M_m-r_m+1\) sampled signs have edge maximum at most
\[
  C_{A,B,K,a,R}\sqrt{n_m}.
\]
Third, at least \(M_m-r_m+1\) sampled signs are centers of Hamming balls of
radius
\[
  \left\lfloor c_{A,B,K,a,R}\ell_m\right\rfloor
\]
on which every signing \(\eps'\) has edge maximum at most
\[
  C_{A,B,K,a,R}\sqrt{n_m}.
\]

In the critical-density case \(\gamma=1\), if \(R>C_\beta L\), the same
conclusion holds under the weaker summability condition
\[
  \sum_{m=1}^{\infty}
  \left(d_m^{-A}
  +\exp\left(-c_{R,\beta,L}\frac{n_m}{\ell_m^2}\right)
  +e^{-u_m}
  +\binom{M_m}{r_m}d_m^{-B r_m}\right)<\infty .
\]
Consequently, along the full integer dimension sequence, suppose \(A>1\),
\(s>1\),
\[
  F_{d,a}(t)\le L\,\log^{-1}(e+d)t^\beta
  \qquad(0<t\le1,\ d\ge1),
\]
\(R>C_\beta L\), and
\[
  n(d)\le d,\qquad
  n(d)\ge q\,\log^3(e+d)
  \qquad\text{for all large }d,\qquad
  q>\frac1{c_{R,\beta,L}} .
\]
If \(M_d\le d^{\beta_{\rm b}}\) for some \(\beta_{\rm b}\ge0\) and the fixed
sign-count exponent satisfies \(B>\beta_{\rm b}+1\), then taking \(r_d=1\) and
\[
  u_d=s\log(e+d)
\]
gives an almost-sure event on which, for all sufficiently large integer
dimensions, the whole sampled batch has edge maximum
\[
  O_{A,B,K,a,R,s,\beta_{\rm b}}(\sqrt{n(d)})
\]
and every sampled sign in the batch is a robust center.
\end{corollary}

\begin{proof}
By Corollary~\ref{cor:normalizer-tail-power-curve}, outside a matrix event of
probability at most
\[
  C_{A,K,a,R}d_m^{-A}+C_\beta L R^{-1}\ell_m^{1-\gamma},
\]
one has the edge profile
\[
  \max_{\omega^{d_m}=1}\sum_{i=1}^{n_m}
  \left|\frac{P_{m,i}(\omega)}{\norm{B_{m,i}}_{\op}}\right|^2
  \le C_{A,K,a,R}\frac{n_m}{\ell_m}.
\]
Condition on a matrix realization satisfying this profile.  The sampled-list
maximum follows from Corollary~\ref{cor:edge-profile-random-search}.  Applying
that corollary with exponent \(B\) gives the good-sign count failure
\[
  \binom{M_m}{r_m}d_m^{-B r_m},
\]
after absorbing the profile constant into \(C_{A,B,K,a,R}\).  The sampled
robust-center count follows from Corollary~\ref{cor:edge-profile-hamming}, again
with exponent \(B\), with the same binomial tail.  Hence the total failure
probability at index \(m\) is bounded, up to a fixed constant, by the summand in
the displayed Markov-form series.  Borel--Cantelli gives the first assertion.

For \(\gamma=1\) and \(R>C_\beta L\), use the critical-density replacement in
Corollary~\ref{cor:normalizer-tail-power-curve} for the matrix event.  The same
Borel--Cantelli argument applies under the second displayed summability
condition.

For the full integer-dimensional consequence take \(d_m=m\), \(n_m=n(m)\), and
\(u_m=s\log(e+m)\).  The series \(\sum_m m^{-A}\) and
\(\sum_m e^{-u_m}\) are finite because \(A>1\) and \(s>1\).  The sample-size
lower bound gives
\[
  \exp\left(-c_{R,\beta,L}\frac{n(m)}{\log^2(e+m)}\right)
  \le (e+m)^{-c_{R,\beta,L}q},
\]
which is summable by the assumption on \(q\).  Finally, with \(r_m=1\),
\[
  \binom{M_m}{1}m^{-B}\le m^{\beta_{\rm b}-B},
\]
and this is summable because \(B>\beta_{\rm b}+1\).  The sampled-list maximum is
\(O(\sqrt{n(m)})\) since
\[
  \frac{\log(eM_m)+u_m}{\log(e+m)}
  \le C_{s,\beta_{\rm b}}
\]
for all large \(m\).  With \(r_m=1\), the count conclusions say that every
sampled sign is good and robust.
\end{proof}

\begin{corollary}[Almost-sure power-curve robust fractions]
\label{cor:normalizer-tail-power-borel-fractions}
Fix \(A>0\), \(S>0\), \(a>0\), \(L>0\), \(\gamma\ge1\), \(\beta>2\),
\(R\ge1\), and a fraction \(0<\rho\le1\).  Choose \(B>0\) and an integer
\(M\ge1\) so that
\[
  B\ge 2\log_2(e/\rho),
  \qquad
  M\ge \frac{2(S+1)}{B\rho}.
\]
There are constants \(C_{A,B,K,a,R}\) and \(c_{A,B,K,a,R}>0\) with the
following property.  Let \((d_m,n_m)_{m\ge1}\) be a triangular-array sequence
with
\[
  d_m\ge2,\qquad n_m\le d_m,\qquad \ell_m=\log(e+d_m),
\]
and assume
\[
  n_m\ge C_{A,B,K,a,R}\ell_m \qquad(m\ge1).
\]
At each \(m\), let \(B_{m,i}=T_{d_m}(\xi^{(m,i)})\), \(1\le i\le n_m\), be
independent full Toeplitz matrices with independent mean-zero, unit-variance lag
variables of sub-Gaussian norm at most \(K\), and let \(P_{m,i}\) be their
symbols.  Independently of the matrices, sample independent uniform sign vectors
\[
  \delta_m^{(1)},\ldots,\delta_m^{(M)}\in\{-1,1\}^{n_m}.
\]
Assume that, for every \(m\),
\[
  F_{d_m,a}(t):=
  \Prob\left\{\norm{T_{d_m}(\xi)}_{\op}
  <t\,a\,\Exp\norm{T_{d_m}(\xi)}_{\op}\right\}
  \le L\,\ell_m^{-\gamma}t^\beta
  \qquad(0<t\le1),
\]
and suppose
\[
  \sum_{m=1}^{\infty}\left(d_m^{-A}+\ell_m^{1-\gamma}+d_m^{-S}\right)<\infty .
\]
Then, with probability one over the matrices and sampled signs, all but
finitely many \(m\) satisfy both conclusions below.

First, the whole sampled list obeys
\[
  \max_{1\le s\le M}\max_{\omega^{d_m}=1}
  \left|\sum_{i=1}^{n_m}\delta_{m,i}^{(s)}
  \frac{P_{m,i}(\omega)}{\norm{B_{m,i}}_{\op}}\right|
  \le C_{A,B,K,a,R,S,\rho}\sqrt{n_m}.
\]
Second, at least a \(1-\rho\) fraction of the \(M\) sampled signs are centers of
Hamming balls of radius
\[
  \left\lfloor c_{A,B,K,a,R}\ell_m\right\rfloor
\]
on which every signing \(\eps'\) has edge maximum at most
\[
  C_{A,B,K,a,R}\sqrt{n_m}.
\]
The same \(1-\rho\) fraction also has edge maximum at most
\(C_{A,B,K,a,R}\sqrt{n_m}\).

In the critical-density case \(\gamma=1\), if \(R>C_\beta L\), the same
conclusion holds under the weaker summability condition
\[
  \sum_{m=1}^{\infty}
  \left(d_m^{-A}
  +\exp\left(-c_{R,\beta,L}\frac{n_m}{\ell_m^2}\right)
  +d_m^{-S}\right)<\infty .
\]
Consequently, on any stretched-exponential dimension schedule
\[
  d_m\ge\exp(c_0m^\kappa)
  \qquad(c_0,\kappa>0),
\]
the Markov-form conclusion holds for every \(S>0\) whenever
\(\gamma>1+1/\kappa\).  In the critical-density case \(\gamma=1\), the same
conclusion holds on such schedules if either
\[
  n_m\ge q\,\ell_m^2\log(e+m)
  \qquad\text{for all large }m,\qquad
  q>\frac1{c_{R,\beta,L}},
\]
or, more simply, if for some \(\tau>0\) and \(Q>0\),
\[
  n_m\ge Q\,\ell_m^{2+\tau}
  \qquad\text{for all large }m .
\]
\end{corollary}

\begin{proof}
Put \(u_m=S\ell_m\).  Apply
Corollary~\ref{cor:normalizer-tail-power-robust-fraction-confidence} at each
dimension \(m\).  In the Markov form, the failure probability at \(m\) is bounded
by a constant times
\[
  d_m^{-A}+\ell_m^{1-\gamma}+e^{-u_m}+d_m^{-S}.
\]
Since \(e^{-u_m}=(e+d_m)^{-S}\le d_m^{-S}\), these probabilities are summable by
hypothesis.  Borel--Cantelli gives the eventual sampled-list and robust-fraction
conclusions.  The sampled-list bound is \(O(\sqrt{n_m})\) because \(M\) is fixed
and
\[
  1+\frac{\log(eM)+u_m}{\ell_m}
  \le C_{S,\rho,B}.
\]

For \(\gamma=1\) and \(R>C_\beta L\), use the critical-density replacement in
Corollary~\ref{cor:normalizer-tail-power-robust-fraction-confidence}; the same
Borel--Cantelli argument applies with the displayed exponential summand.

On a stretched-exponential dimension schedule, \(\sum_m d_m^{-A}\) and
\(\sum_m d_m^{-S}\) are finite for every \(A,S>0\), while
\(\ell_m\ge c m^\kappa\) for all large \(m\).  Thus
\[
  \sum_m \ell_m^{1-\gamma}
  \le C\sum_m m^{-\kappa(\gamma-1)}<\infty
\]
whenever \(\gamma>1+1/\kappa\).  In the critical-density case, the lower bound
\(n_m\ge q\ell_m^2\log(e+m)\) gives
\[
  \exp\left(-c_{R,\beta,L}\frac{n_m}{\ell_m^2}\right)
  \le (e+m)^{-c_{R,\beta,L}q},
\]
which is summable when \(q>1/c_{R,\beta,L}\).  The alternative
\(n_m\ge Q\ell_m^{2+\tau}\) gives a stretched-exponential summand
\(\exp(-c'm^{\kappa\tau})\), also summable.
\end{proof}

\begin{corollary}[Almost-sure power-curve polynomial batches]
\label{cor:normalizer-tail-power-borel-poly-batch}
Fix \(A>0\), \(a>0\), \(L>0\), \(\gamma\ge1\), \(\beta_0>2\), \(R\ge1\),
\(\beta_{\rm b}\ge0\), and \(\sigma>0\).  There are constants
\(C_{A,K,a,R,\beta_{\rm b},\sigma}\) and
\(c_{A,K,a,R,\beta_{\rm b},\sigma}>0\) with the following property.  Let
\((d_m,n_m)_{m\ge1}\) be a triangular-array sequence with
\[
  d_m\ge2,\qquad n_m\le d_m,\qquad \ell_m=\log(e+d_m),
\]
and assume
\[
  n_m\ge C_{A,K,a,R,\beta_{\rm b},\sigma}\ell_m \qquad(m\ge1).
\]
At each \(m\), let \(B_{m,i}=T_{d_m}(\xi^{(m,i)})\), \(1\le i\le n_m\), be
independent full Toeplitz matrices with independent mean-zero, unit-variance lag
variables of sub-Gaussian norm at most \(K\), and let \(P_{m,i}\) be their
symbols.  Independently of the matrices, sample \(M_m\) independent uniform sign
vectors in \(\{-1,1\}^{n_m}\), where
\[
  1\le M_m\le d_m^{\beta_{\rm b}} .
\]
Assume that, for every \(m\),
\[
  F_{d_m,a}(t)\le L\,\ell_m^{-\gamma}t^{\beta_0}
  \qquad(0<t\le1),
\]
and suppose
\[
  \sum_{m=1}^{\infty}
  \left(d_m^{-A}+\ell_m^{1-\gamma}+d_m^{-\sigma}\right)<\infty .
\]
Then, with probability one over the matrices and sampled signs, all but
finitely many \(m\) have the following property: every sampled sign
\(\delta_m^{(s)}\), \(1\le s\le M_m\), is the center of a Hamming ball of radius
\[
  \left\lfloor c_{A,K,a,R,\beta_{\rm b},\sigma}\ell_m\right\rfloor
\]
on which every signing \(\eps'\) satisfies
\[
  \max_{\omega^{d_m}=1}\left|\sum_{i=1}^{n_m}
  \eps'_i\frac{P_{m,i}(\omega)}{\norm{B_{m,i}}_{\op}}\right|
  \le C_{A,K,a,R,\beta_{\rm b},\sigma}\sqrt{n_m}.
\]
In particular, the whole sampled batch itself has edge maximum
\[
  O_{A,K,a,R,\beta_{\rm b},\sigma}(\sqrt{n_m})
\]
eventually almost surely.

In the critical-density case \(\gamma=1\), if \(R>C_{\beta_0}L\), the same
conclusion holds under the weaker summability condition
\[
  \sum_{m=1}^{\infty}
  \left(d_m^{-A}
  +\exp\left(-c_{R,\beta_0,L}\frac{n_m}{\ell_m^2}\right)
  +d_m^{-\sigma}\right)<\infty .
\]
Consequently, on any stretched-exponential dimension schedule
\[
  d_m\ge\exp(c_0m^\kappa)
  \qquad(c_0,\kappa>0),
\]
the Markov-form conclusion holds for every \(\sigma>0\) whenever
\(\gamma>1+1/\kappa\).  In the critical-density case \(\gamma=1\), the same
conclusion holds on such schedules under either of the sample-size alternatives
\[
  n_m\ge q\,\ell_m^2\log(e+m),\qquad q>\frac1{c_{R,\beta_0,L}},
\]
for all large \(m\), or
\[
  n_m\ge Q\,\ell_m^{2+\tau}
\]
for all large \(m\), with some \(\tau>0\) and \(Q>0\).
\end{corollary}

\begin{proof}
Apply Corollary~\ref{cor:normalizer-tail-power-poly-batch} at each dimension.
In the Markov form, the failure probability at \(m\) is bounded by a constant
times
\[
  d_m^{-A}+\ell_m^{1-\gamma}+d_m^{-\sigma},
\]
which is summable by hypothesis.  Borel--Cantelli gives the eventual all-robust
batch conclusion.

For \(\gamma=1\) and \(R>C_{\beta_0}L\), use the critical-density replacement in
Corollary~\ref{cor:normalizer-tail-power-poly-batch}.  The stated exponential
summability then gives the same almost-sure conclusion.  The
stretched-exponential consequences are verified exactly as in
Corollary~\ref{cor:normalizer-tail-power-borel-fractions}: powers
\(d_m^{-A}\) and \(d_m^{-\sigma}\) are summable, \(\ell_m^{1-\gamma}\) is
summable when \(\gamma>1+1/\kappa\), and the two critical sample-size lower
bounds make \(\exp(-c n_m/\ell_m^2)\) summable.
\end{proof}

\begin{corollary}[Critical power-curve abundance on all integer dimensions]
\label{cor:normalizer-tail-power-borel-critical-integers}
Fix \(a>0\), \(L>0\), \(\beta>2\), \(R>C_\beta L\), and \(A>1\).  Let
\(d=1,2,3,\ldots\), put \(\ell_d=\log(e+d)\), and let \(n(d)\le d\).  At
dimension \(d\), let
\(B_{d,i}=T_d(\xi^{(d,i)})\), \(1\le i\le n(d)\), be independent full Toeplitz
matrices with independent mean-zero, unit-variance lag variables of
sub-Gaussian norm at most \(K\), and let \(P_{d,i}\) be their symbols.  Assume
the critical power-curve lower-tail bound
\[
  F_{d,a}(t):=
  \Prob\left\{\norm{T_d(\xi)}_{\op}
  <t\,a\,\Exp\norm{T_d(\xi)}_{\op}\right\}
  \le L\,\ell_d^{-1}t^\beta
  \qquad(0<t\le1,\ d\ge1),
\]
and suppose that, for all sufficiently large \(d\),
\[
  n(d)\ge q\,\ell_d^3
  \qquad\text{with}\qquad
  q>\frac1{c_{R,\beta,L}}.
\]
Then the following two all-integer conclusions hold.

First, fix \(0<\rho\le1\), \(S>1\), choose \(B>0\) and an integer \(M\ge1\)
such that
\[
  B\ge 2\log_2(e/\rho),
  \qquad
  M\ge \frac{2(S+1)}{B\rho},
\]
and at every \(d\) sample \(M\) independent uniform sign vectors, independently
of the matrices.  With probability one over matrices and signs, all sufficiently
large integer dimensions have sampled-list maximum
\[
  \max_{1\le s\le M}\max_{\omega^d=1}
  \left|\sum_{i=1}^{n(d)}\delta_{d,i}^{(s)}
  \frac{P_{d,i}(\omega)}{\norm{B_{d,i}}_{\op}}\right|
  \le C_{A,B,K,a,R,S,\rho}\sqrt{n(d)},
\]
and at least a \(1-\rho\) fraction of the sampled signs are centers of Hamming
balls of radius \(\lfloor c_{A,B,K,a,R}\ell_d\rfloor\) on which every signing
has edge maximum at most \(C_{A,B,K,a,R}\sqrt{n(d)}\).  The same fraction also
has edge maximum at most \(C_{A,B,K,a,R}\sqrt{n(d)}\).

Second, fix \(\beta_{\rm b}\ge0\) and \(\sigma>1\), and at dimension \(d\)
sample any independent batch of signs with
\[
  1\le M_d\le d^{\beta_{\rm b}} .
\]
Then, with probability one over matrices and signs, all sufficiently large
integer dimensions have every sampled sign as the center of a Hamming ball of
radius \(\lfloor c_{A,K,a,R,\beta_{\rm b},\sigma}\ell_d\rfloor\) on which every
signing has edge maximum at most
\[
  C_{A,K,a,R,\beta_{\rm b},\sigma}\sqrt{n(d)}.
\]
In particular, the whole sampled batch itself has edge maximum
\[
  O_{A,K,a,R,\beta_{\rm b},\sigma}(\sqrt{n(d)})
\]
eventually almost surely.
\end{corollary}

\begin{proof}
The lower bound \(n(d)\ge q\ell_d^3\) eventually dominates every fixed multiple
of \(\ell_d\), so the admissibility hypotheses in the two source corollaries
hold after discarding finitely many dimensions.

Apply Corollary~\ref{cor:normalizer-tail-power-borel-fractions} with
\(d_m=m\), \(n_m=n(m)\), \(\ell_m=\log(e+m)\), and \(\gamma=1\).  The series
\(\sum_m m^{-A}\) and \(\sum_m m^{-S}\) are finite because \(A>1\) and \(S>1\),
while the lower bound on \(n(m)\) gives
\[
  \exp\left(-c_{R,\beta,L}\frac{n(m)}{\ell_m^2}\right)
  \le (e+m)^{-c_{R,\beta,L}q}
\]
for all large \(m\), which is summable by the assumption on \(q\).  This proves
the robust-fraction conclusion.

The all-robust polynomial-batch conclusion follows in the same way from
Corollary~\ref{cor:normalizer-tail-power-borel-poly-batch}: the only additional
sampled-sign summand is \(\sum_m m^{-\sigma}\), which is finite because
\(\sigma>1\).
\end{proof}

\begin{corollary}[Critical power-curve vanishing fractions on all integer dimensions]
\label{cor:normalizer-tail-power-borel-critical-vanishing-fractions}
Fix \(a>0\), \(L>0\), \(\beta>2\), \(R>C_\beta L\), \(A>1\), and
\(B>1\).  Fix \(0<\tau_0<B-1\).  Let \(d=1,2,3,\ldots\), put
\(\ell_d=\log(e+d)\), and let \(n(d)\le d\).  At dimension \(d\), let
\(B_{d,i}=T_d(\xi^{(d,i)})\), \(1\le i\le n(d)\), be independent full Toeplitz
matrices with independent mean-zero, unit-variance lag variables of
sub-Gaussian norm at most \(K\), and let \(P_{d,i}\) be their symbols.  Assume
the critical power-curve lower-tail bound
\[
  F_{d,a}(t):=
  \Prob\left\{\norm{T_d(\xi)}_{\op}
  <t\,a\,\Exp\norm{T_d(\xi)}_{\op}\right\}
  \le L\,\ell_d^{-1}t^\beta
  \qquad(0<t\le1,\ d\ge1),
\]
and suppose that, for all sufficiently large \(d\),
\[
  n(d)\ge q\,\ell_d^3
  \qquad\text{with}\qquad
  q>\frac1{c_{R,\beta,L}}.
\]
At each \(d\), independently sample an arbitrary deterministic number \(M_d\ge1\)
of uniform sign vectors in \(\{-1,1\}^{n(d)}\).  Let \(u_d\ge0\) satisfy
\[
  \sum_{d=1}^{\infty}e^{-u_d}<\infty ,
\]
and let \((\rho_d)_{d\ge1}\subset(0,1]\) satisfy
\[
  \rho_d\ge e\,d^{-\tau_0}
\]
for all sufficiently large \(d\).  Then, with probability one over the matrices
and sampled signs, all sufficiently large integer dimensions satisfy both
conclusions below.

First, the whole sampled list obeys
\[
\begin{aligned}
  \max_{1\le s\le M_d}\max_{\omega^d=1}
  \left|\sum_{i=1}^{n(d)}\delta_{d,i}^{(s)}
  \frac{P_{d,i}(\omega)}{\norm{B_{d,i}}_{\op}}\right|
  \le
  C_{A,B,K,a,R}\sqrt{
    n(d)\left(1+\frac{\log(eM_d)+u_d}{\ell_d}\right)} .
\end{aligned}
\]
Second, at least a \(1-\rho_d\) fraction of the \(M_d\) sampled signs have edge
maximum at most \(C_{A,B,K,a,R}\sqrt{n(d)}\) and are centers of Hamming balls of
radius
\[
  \left\lfloor c_{A,B,K,a,R}\ell_d\right\rfloor
\]
on which every signing \(\eps'\) has edge maximum at most
\[
  C_{A,B,K,a,R}\sqrt{n(d)}.
\]
No upper bound on \(M_d\) is required.  In particular, if \(\rho_d\to0\), then
the robust-center fraction tends to one along the full integer dimension
sequence.
\end{corollary}

\begin{proof}
Apply Corollary~\ref{cor:normalizer-tail-power-borel-counts} with
\[
  d_m=m,\qquad n_m=n(m),\qquad r_m=\lceil\rho_mM_m\rceil ,
\]
and with the critical-density replacement.  The eventual lower bound
\(n(m)\ge q\ell_m^3\) implies the admissibility condition
\(n_m\ge C_{A,B,K,a,R}\ell_m\) after discarding finitely many dimensions, and
\[
  \exp\left(-c_{R,\beta,L}\frac{n(m)}{\ell_m^2}\right)
  \le (e+m)^{-c_{R,\beta,L}q}
\]
is summable by the assumption on \(q\).  Also \(\sum_m m^{-A}<\infty\) because
\(A>1\), and the sampled-list tail is summable by the hypothesis on \(u_m\).

It remains to check the binomial count tail.  Since
\(r_m=\lceil\rho_mM_m\rceil\),
\[
  \binom{M_m}{r_m}
  \le \left(\frac e{\rho_m}\right)^{r_m}.
\]
The lower bound \(\rho_m\ge e\,m^{-\tau_0}\) gives, for all large \(m\),
\[
  \binom{M_m}{r_m}m^{-B r_m}
  \le m^{-(B-\tau_0)r_m}\le m^{-(B-\tau_0)}.
\]
This is summable because \(B-\tau_0>1\).  The sampled-count theorem gives the
displayed sampled-list maximum and at least \(M_m-r_m+1\) robust centers
eventually.  Finally,
\[
  M_m-r_m+1\ge M_m-\rho_mM_m,
\]
so the lower bound is at least a \(1-\rho_m\) fraction.
\end{proof}

\begin{corollary}[Sparse power-curve Borel windows]
\label{cor:normalizer-tail-power-borel-window}
Fix \(A>0\), \(a>0\), \(L>0\), \(\beta>2\), and \(R\ge1\).  There is a
constant \(C_{A,K,a,R}\) with the following property.  Let \((d_m,n_m)_{m\ge1}\)
be a triangular-array sequence with \(n_m\le d_m\), \(\ell_m=\log(e+d_m)\), and
\[
  d_m\ge \exp(c_0m^\kappa)
  \qquad(m\ge1)
\]
for some \(c_0,\kappa>0\).  At each \(m\), let the matrices be as in
Corollary~\ref{cor:normalizer-tail-power-borel-search}, and assume
\[
  n_m\ge C_{A,K,a,R}\ell_m
  \qquad(m\ge1).
\]
If, for some \(\gamma>1+1/\kappa\),
\[
  F_{d_m,a}(t)\le L\,\ell_m^{-\gamma}t^\beta
  \qquad(0<t\le1,\ m\ge1),
\]
then almost surely all but finitely many \(m\) satisfy the full random-search
and stable-center package of
Corollary~\ref{cor:normalizer-tail-power-search}.

In the critical-density case \(\gamma=1\), assume \(R>C_\beta L\) and
\[
  F_{d_m,a}(t)\le L\,\ell_m^{-1}t^\beta
  \qquad(0<t\le1,\ m\ge1).
\]
Then the same eventual conclusion holds if either
\[
  n_m\ge q\,\ell_m^2\log(e+m)
  \qquad\text{for all large }m,\qquad
  q>\frac1{c_{R,\beta,L}},
\]
or, more simply, if for some \(\rho>0\) and \(Q>0\),
\[
  n_m\ge Q\,\ell_m^{2+\rho}
  \qquad\text{for all large }m .
\]
Thus on stretched-exponential dimension schedules the critical power-curve
density \(F_{d,a}(t)\lesssim \ell^{-1}t^\beta\) needs only a superquadratic
logarithmic sample window to give eventual randomized certificates.
\end{corollary}

\begin{proof}
Since \(d_m\ge\exp(c_0m^\kappa)\), the series \(\sum_m d_m^{-A}\) is finite.
Also \(\ell_m\ge c m^\kappa\) for all large \(m\).  Hence
\[
  \sum_m \ell_m^{1-\gamma}
  \le C\sum_m m^{-\kappa(\gamma-1)}<\infty
\]
when \(\gamma>1+1/\kappa\).  The Markov form of
Corollary~\ref{cor:normalizer-tail-power-borel-search} gives the first
assertion.

For the critical-density case, the condition
\(n_m\ge q\ell_m^2\log(e+m)\) gives
\[
  \exp\left(-c_{R,\beta,L}\frac{n_m}{\ell_m^2}\right)
  \le (e+m)^{-c_{R,\beta,L}q},
\]
which is summable when \(q>1/c_{R,\beta,L}\).  If instead
\(n_m\ge Q\ell_m^{2+\rho}\), then for all large \(m\),
\[
  \exp\left(-c_{R,\beta,L}\frac{n_m}{\ell_m^2}\right)
  \le \exp(-c' m^{\kappa\rho}),
\]
which is summable for every \(\rho>0\).  The Bernstein form of
Corollary~\ref{cor:normalizer-tail-power-borel-search} completes the proof.
\end{proof}

\begin{corollary}[Sparse sampled power-curve Borel windows]
\label{cor:normalizer-tail-power-borel-window-sampled}
Fix \(A>0\), \(\alpha>0\), \(a>0\), \(L>0\), \(\beta>2\), and \(R\ge1\).
There is a constant \(C_{A,K,a,R}\) with the following property.  Let
\((d_m,n_m)_{m\ge1}\) be a triangular-array sequence with \(n_m\le d_m\),
\(\ell_m=\log(e+d_m)\), and
\[
  d_m\ge \exp(c_0m^\kappa)
  \qquad(m\ge1)
\]
for some \(c_0,\kappa>0\).  At each \(m\), let the matrices be as in
Corollary~\ref{cor:normalizer-tail-power-borel-sampled}, assume
\[
  n_m\ge C_{A,K,a,R}\ell_m
  \qquad(m\ge1),
\]
and sample \(M_m\) independent uniform sign vectors in \(\{-1,1\}^{n_m}\),
independently of the matrices, with
\[
  M_m\ge\lceil\alpha\rceil .
\]
If, for some \(\gamma>1+1/\kappa\),
\[
  F_{d_m,a}(t)\le L\,\ell_m^{-\gamma}t^\beta
  \qquad(0<t\le1,\ m\ge1),
\]
then, with probability one over the matrices and sampled signs, all but
finitely many \(m\) have a sampled sign which is the center of a Hamming ball of
radius
\[
  r_m=\left\lfloor c_{A,K,a,R}\ell_m\right\rfloor
\]
on which every signing \(\eps'\) satisfies
\[
  \max_{\omega^{d_m}=1}\left|\sum_{i=1}^{n_m}
  \eps'_i\frac{P_{m,i}(\omega)}{\norm{B_{m,i}}_{\op}}\right|
  \le C_{A,K,a,R}\sqrt{n_m}.
\]

In the critical-density case \(\gamma=1\), assume \(R>C_\beta L\) and
\[
  F_{d_m,a}(t)\le L\,\ell_m^{-1}t^\beta
  \qquad(0<t\le1,\ m\ge1).
\]
Then the same actual sampled-certificate conclusion holds if either
\[
  n_m\ge q\,\ell_m^2\log(e+m)
  \qquad\text{for all large }m,\qquad
  q>\frac1{c_{R,\beta,L}},
\]
or, more simply, if for some \(\rho>0\) and \(Q>0\),
\[
  n_m\ge Q\,\ell_m^{2+\rho}
  \qquad\text{for all large }m .
\]
Thus on stretched-exponential dimension schedules, every \(\alpha>0\) and every
fixed list size \(M\ge\lceil\alpha\rceil\) give eventual sampled robust
certificates under the same sparse windows as
Corollary~\ref{cor:normalizer-tail-power-borel-window}.
\end{corollary}

\begin{proof}
Apply Corollary~\ref{cor:normalizer-tail-power-borel-sampled}.  The
stretched-exponential lower bound on \(d_m\) makes both \(\sum_m d_m^{-A}\) and
\(\sum_m d_m^{-\alpha}\) finite.  Also \(\ell_m\ge c m^\kappa\) for all large
\(m\), so
\[
  \sum_m \ell_m^{1-\gamma}
  \le C\sum_m m^{-\kappa(\gamma-1)}<\infty
\]
when \(\gamma>1+1/\kappa\).  This verifies the Markov-form summability condition
of Corollary~\ref{cor:normalizer-tail-power-borel-sampled}.

In the critical-density case, the lower bound
\(n_m\ge q\ell_m^2\log(e+m)\) gives
\[
  \exp\left(-c_{R,\beta,L}\frac{n_m}{\ell_m^2}\right)
  \le (e+m)^{-c_{R,\beta,L}q},
\]
which is summable when \(q>1/c_{R,\beta,L}\).  The alternative
\(n_m\ge Q\ell_m^{2+\rho}\) gives
\[
  \exp\left(-c_{R,\beta,L}\frac{n_m}{\ell_m^2}\right)
  \le \exp(-c' m^{\kappa\rho})
\]
for all large \(m\), again summable.  Together with the summability of
\(\sum_m d_m^{-A}\) and \(\sum_m d_m^{-\alpha}\), this verifies the critical
condition in Corollary~\ref{cor:normalizer-tail-power-borel-sampled}.
\end{proof}

\begin{corollary}[Sparse sampled-count power-curve Borel windows]
\label{cor:normalizer-tail-power-borel-window-counts}
Fix \(A>0\), \(B>0\), \(a>0\), \(L>0\), \(\beta>2\), \(R\ge1\), \(c_0>0\),
and \(\kappa>0\).  There are constants \(C_{A,B,K,a,R}\) and
\(c_{A,B,K,a,R}>0\) with the following property.  Let
\((d_m,n_m)_{m\ge1}\) be a triangular-array sequence with \(n_m\le d_m\),
\(\ell_m=\log(e+d_m)\), and
\[
  d_m\ge \exp(c_0m^\kappa),\qquad
  n_m\ge C_{A,B,K,a,R}\ell_m
  \qquad(m\ge1).
\]
At each \(m\), take the independent full Toeplitz triangular array at dimension
\(d_m\), and independently sample \(M_m\) uniform sign vectors in
\(\{-1,1\}^{n_m}\), with \(M_m\ge1\).  Let \(u_m\ge0\) and
\(1\le r_m\le M_m\) be deterministic parameters.  If, for some
\(\gamma>1+1/\kappa\),
\[
  F_{d_m,a}(t)\le L\,\ell_m^{-\gamma}t^\beta
  \qquad(0<t\le1,\ m\ge1),
\]
and
\[
  \sum_{m=1}^{\infty}
  \left(e^{-u_m}+\binom{M_m}{r_m}d_m^{-B r_m}\right)<\infty,
\]
then, with probability one over the matrices and sampled signs, all but
finitely many \(m\) satisfy the sampled-list maximum bound, good-sign count, and
robust-center count conclusions of
Corollary~\ref{cor:normalizer-tail-power-borel-counts}.

In the critical-density case \(\gamma=1\), assume \(R>C_\beta L\) and
\[
  F_{d_m,a}(t)\le L\,\ell_m^{-1}t^\beta
  \qquad(0<t\le1,\ m\ge1).
\]
Then the same eventual sampled-count conclusion holds, under the same
sign-summability condition, if either
\[
  n_m\ge q\,\ell_m^2\log(e+m)
  \qquad\text{for all large }m,\qquad
  q>\frac1{c_{R,\beta,L}},
\]
or, more simply, if for some \(\rho>0\) and \(Q>0\),
\[
  n_m\ge Q\,\ell_m^{2+\rho}
  \qquad\text{for all large }m .
\]

In particular, let \(s>1\), take \(u_m=s\log(e+m)\) and \(r_m=1\), and suppose
\[
  M_m\le d_m^{\beta_{\rm b}}
  \qquad(m\ge1)
\]
for some \(\beta_{\rm b}\ge0\).  Then every fixed sign-count exponent
\(B>\beta_{\rm b}\) gives an almost-sure event on which, for all sufficiently
large \(m\), every sampled sign in the batch is an \(O(\sqrt{n_m})\)-good robust
center and the whole batch satisfies the displayed sampled-list maximum bound.
More generally, if
\[
  M_m\le \exp(bm^\sigma)
\]
for some \(b>0\), \(0<\sigma\le\kappa\), then the same eventual all-sampled-sign
conclusion holds whenever either \(\sigma<\kappa\), or \(\sigma=\kappa\) and
\(b<Bc_0\).
\end{corollary}

\begin{proof}
Apply Corollary~\ref{cor:normalizer-tail-power-borel-counts}.  The
stretched-exponential lower bound on \(d_m\) makes \(\sum_m d_m^{-A}\) finite.
Also \(\ell_m\ge c m^\kappa\) for all large \(m\), so
\[
  \sum_m \ell_m^{1-\gamma}
  \le C\sum_m m^{-\kappa(\gamma-1)}<\infty
\]
when \(\gamma>1+1/\kappa\).  Together with the displayed sign summability, this
verifies the Markov-form hypothesis of the sampled-count Borel theorem.

In the critical-density case, the lower bound
\(n_m\ge q\ell_m^2\log(e+m)\) gives
\[
  \exp\left(-c_{R,\beta,L}\frac{n_m}{\ell_m^2}\right)
  \le (e+m)^{-c_{R,\beta,L}q},
\]
which is summable when \(q>1/c_{R,\beta,L}\).  The alternative
\(n_m\ge Q\ell_m^{2+\rho}\) gives
\[
  \exp\left(-c_{R,\beta,L}\frac{n_m}{\ell_m^2}\right)
  \le \exp(-c' m^{\kappa\rho})
\]
for all large \(m\), again summable.  With the summability of \(\sum_m d_m^{-A}\)
and the displayed sign terms, this verifies the critical hypothesis in
Corollary~\ref{cor:normalizer-tail-power-borel-counts}.

For the polynomial-in-\(d_m\) batch consequence,
\[
  \binom{M_m}{1}d_m^{-B}\le d_m^{\beta_{\rm b}-B},
\]
which is summable on the stretched-exponential schedule when
\(B>\beta_{\rm b}\).  For the \(\exp(bm^\sigma)\) batch consequence,
\[
  \binom{M_m}{1}d_m^{-B}
  \le \exp\left(bm^\sigma-Bc_0m^\kappa\right),
\]
which is summable if \(\sigma<\kappa\), or if \(\sigma=\kappa\) and
\(b<Bc_0\).  Finally, \(\sum_m e^{-u_m}<\infty\) because \(s>1\).  With
\(r_m=1\), the count conclusion says all sampled signs are good robust centers
eventually, and the comparison
\[
  \frac{\log(eM_m)+u_m}{\ell_m}=O(1)
\]
in both displayed batch regimes gives the \(O(\sqrt{n_m})\) sampled-list maximum.
\end{proof}

\begin{corollary}[Sparse power-curve growing robust fractions]
\label{cor:normalizer-tail-power-borel-window-growing-fractions}
Fix \(A>0\), \(a>0\), \(L>0\), \(\beta>2\), \(R\ge1\), \(c_0>0\),
\(\kappa>0\), and \(0<\rho\le1\).  There are constants
\(C_{A,K,a,R}\) and \(c_{A,K,a,R}>0\) with the following property.  Let
\((d_m,n_m)_{m\ge1}\) be a triangular-array sequence with \(n_m\le d_m\),
\(\ell_m=\log(e+d_m)\), and
\[
  d_m\ge \exp(c_0m^\kappa),\qquad
  n_m\ge C_{A,K,a,R}\ell_m
  \qquad(m\ge1).
\]
At each \(m\), take the independent full Toeplitz triangular array at dimension
\(d_m\), writing \(B_{m,i}=T_{d_m}(\xi^{(m,i)})\), \(1\le i\le n_m\), and
\(P_{m,i}\) for its symbol.  Independently sample an arbitrary deterministic
number \(M_m\ge1\) of uniform sign vectors
\[
  \delta_m^{(1)},\ldots,\delta_m^{(M_m)}\in\{-1,1\}^{n_m}.
\]
Let \(u_m\ge0\) be deterministic and assume
\[
  \sum_{m=1}^{\infty}e^{-u_m}<\infty .
\]
If, for some \(\gamma>1+1/\kappa\),
\[
  F_{d_m,a}(t)\le L\,\ell_m^{-\gamma}t^\beta
  \qquad(0<t\le1,\ m\ge1),
\]
then, with probability one over the matrices and sampled signs, all but
finitely many \(m\) satisfy both conclusions below.

First, the whole sampled list obeys the entropy-sensitive maximum bound
\[
\begin{aligned}
  \max_{1\le s\le M_m}\max_{\omega^{d_m}=1}
  \left|\sum_{i=1}^{n_m}\delta_{m,i}^{(s)}
  \frac{P_{m,i}(\omega)}{\norm{B_{m,i}}_{\op}}\right|
  \le
  C_{A,K,a,R}\sqrt{
    n_m\left(1+\frac{\log(eM_m)+u_m}{\ell_m}\right)} .
\end{aligned}
\]
Second, at least a \(1-\rho\) fraction of the \(M_m\) sampled signs have edge
maximum at most \(C_{A,K,a,R}\sqrt{n_m}\) and are centers of Hamming balls of
radius
\[
  \left\lfloor c_{A,K,a,R}\ell_m\right\rfloor
\]
on which every signing \(\eps'\) has edge maximum at most
\[
  C_{A,K,a,R}\sqrt{n_m}.
\]
No upper bound on \(M_m\) is required for this robust-fraction conclusion.

In the critical-density case \(\gamma=1\), assume \(R>C_\beta L\) and
\[
  F_{d_m,a}(t)\le L\,\ell_m^{-1}t^\beta
  \qquad(0<t\le1,\ m\ge1).
\]
Then the same eventual conclusions hold if either
\[
  n_m\ge q\,\ell_m^2\log(e+m)
  \qquad\text{for all large }m,\qquad
  q>\frac1{c_{R,\beta,L}},
\]
or, more simply, if for some \(\tau>0\) and \(Q>0\),
\[
  n_m\ge Q\,\ell_m^{2+\tau}
  \qquad\text{for all large }m .
\]
\end{corollary}

\begin{proof}
Apply Corollary~\ref{cor:normalizer-tail-power-borel-window-counts} with
sign-count exponent \(B=1\) and
\[
  r_m=\lceil\rho M_m\rceil .
\]
The only additional summability condition to check is the binomial count tail.
Since \(r_m\ge\rho M_m\),
\[
  \binom{M_m}{r_m}d_m^{-r_m}
  \le \left(\frac{eM_m}{r_m}\right)^{r_m}d_m^{-r_m}
  \le \left(\frac{e}{\rho}d_m^{-1}\right)^{r_m}.
\]
For all large \(m\), the last display is at most \(d_m^{-r_m/2}\), and hence at
most \(d_m^{-1/2}\), which is summable because
\(d_m\ge\exp(c_0m^\kappa)\).  The sampled-list tail is summable by the
hypothesis on \(u_m\).  Corollary~\ref{cor:normalizer-tail-power-borel-window-counts}
therefore gives the entropy-sensitive list maximum and at least
\[
  M_m-r_m+1\ge (1-\rho)M_m
\]
good robust centers eventually.  The same argument applies in the critical
density case, because the two displayed lower sample-window alternatives are
exactly the critical alternatives in
Corollary~\ref{cor:normalizer-tail-power-borel-window-counts}.
\end{proof}

\begin{corollary}[Vanishing sparse power-curve discarded fractions]
\label{cor:normalizer-tail-power-borel-window-vanishing-fractions}
Fix \(A>0\), \(B>0\), \(a>0\), \(L>0\), \(\beta>2\), \(R\ge1\),
\(c_0>0\), \(\kappa>0\), and \(0<\tau_0<B\).  There are constants
\(C_{A,B,K,a,R}\) and \(c_{A,B,K,a,R}>0\) with the following property.  Let
\((d_m,n_m)_{m\ge1}\) be a triangular-array sequence with \(n_m\le d_m\),
\(\ell_m=\log(e+d_m)\), and
\[
  d_m\ge \exp(c_0m^\kappa),\qquad
  n_m\ge C_{A,B,K,a,R}\ell_m
  \qquad(m\ge1).
\]
At each \(m\), take the independent full Toeplitz triangular array at dimension
\(d_m\), writing \(B_{m,i}=T_{d_m}(\xi^{(m,i)})\), \(1\le i\le n_m\), and
\(P_{m,i}\) for its symbol.  Independently sample an arbitrary deterministic
number \(M_m\ge1\) of uniform sign vectors
\[
  \delta_m^{(1)},\ldots,\delta_m^{(M_m)}\in\{-1,1\}^{n_m}.
\]
Let \(u_m\ge0\) be deterministic and assume
\[
  \sum_{m=1}^{\infty}e^{-u_m}<\infty .
\]
Let \((\rho_m)_{m\ge1}\subset(0,1]\) satisfy
\[
  \rho_m\ge e\,d_m^{-\tau_0}
\]
for all sufficiently large \(m\).  If, for some \(\gamma>1+1/\kappa\),
\[
  F_{d_m,a}(t)\le L\,\ell_m^{-\gamma}t^\beta
  \qquad(0<t\le1,\ m\ge1),
\]
then, with probability one over the matrices and sampled signs, all but
finitely many \(m\) satisfy both conclusions below.

First, the whole sampled list obeys
\[
\begin{aligned}
  \max_{1\le s\le M_m}\max_{\omega^{d_m}=1}
  \left|\sum_{i=1}^{n_m}\delta_{m,i}^{(s)}
  \frac{P_{m,i}(\omega)}{\norm{B_{m,i}}_{\op}}\right|
  \le
  C_{A,B,K,a,R}\sqrt{
    n_m\left(1+\frac{\log(eM_m)+u_m}{\ell_m}\right)} .
\end{aligned}
\]
Second, at least a \(1-\rho_m\) fraction of the \(M_m\) sampled signs have edge
maximum at most \(C_{A,B,K,a,R}\sqrt{n_m}\) and are centers of Hamming balls of
radius
\[
  \left\lfloor c_{A,B,K,a,R}\ell_m\right\rfloor
\]
on which every signing \(\eps'\) has edge maximum at most
\[
  C_{A,B,K,a,R}\sqrt{n_m}.
\]
No upper bound on \(M_m\) is required.  In particular, if \(\rho_m\to0\), then
the robust-center fraction tends to one along the sampled batches.

In the critical-density case \(\gamma=1\), assume \(R>C_\beta L\) and
\[
  F_{d_m,a}(t)\le L\,\ell_m^{-1}t^\beta
  \qquad(0<t\le1,\ m\ge1).
\]
Then the same eventual conclusions hold if either
\[
  n_m\ge q\,\ell_m^2\log(e+m)
  \qquad\text{for all large }m,\qquad
  q>\frac1{c_{R,\beta,L}},
\]
or, more simply, if for some \(\tau>0\) and \(Q>0\),
\[
  n_m\ge Q\,\ell_m^{2+\tau}
  \qquad\text{for all large }m .
\]
\end{corollary}

\begin{proof}
Apply Corollary~\ref{cor:normalizer-tail-power-borel-window-counts} with
sign-count exponent \(B\) and
\[
  r_m=\lceil\rho_mM_m\rceil .
\]
The matrix summability hypotheses are exactly those verified in the proof of
Corollary~\ref{cor:normalizer-tail-power-borel-window-growing-fractions}, and
the sampled-list tail is summable by the hypothesis on \(u_m\).  It remains only
to check the binomial count tail.  Since \(r_m=\lceil\rho_mM_m\rceil\),
\[
  \binom{M_m}{r_m}
  \le \left(\frac{eM_m}{r_m}\right)^{r_m}
  \le \left(\frac e{\rho_m}\right)^{r_m}.
\]
The lower bound on \(\rho_m\) gives, for all large \(m\),
\[
  \binom{M_m}{r_m}d_m^{-B r_m}
  \le d_m^{-(B-\tau_0)r_m}\le d_m^{-(B-\tau_0)}.
\]
This is summable on the stretched-exponential schedule.  Hence the sampled-count
theorem applies.  Finally,
\[
  M_m-r_m+1\ge M_m-\rho_mM_m,
\]
so the good robust-center count is at least a \(1-\rho_m\) fraction.  The
critical-density alternatives are exactly the alternatives already built into
Corollary~\ref{cor:normalizer-tail-power-borel-window-counts}.
\end{proof}

\begin{corollary}[Explicit lower-tail rates]
\label{cor:normalizer-tail-rates}
Fix \(A>0\), \(a>0\), and write \(\ell=\log(e+d)\).
In the setting of Corollary~\ref{cor:normalizer-tail-quant}, suppose
\[
  p_d(a):=\Prob\{\norm{T_d(\xi)}_{\op}<aN_d\}\le R_d .
\]
Then, whenever \(n\le d\) and \(n\ge C_{A,K,a}\ell\), the realization-normalized
edge profile fails with probability at most
\[
  C_{A,K,a}d^{-A}+R_d\min\{n,\ell^2\}.
\]
In particular:
\begin{enumerate}[label=\textup{(\roman*)}]
\item if \(R_d\le C_0\ell^{-B}\) with \(B>2\), then the exceptional probability
is at most
\[
  C_{A,K,a,C_0,B}\bigl(d^{-A}+\ell^{2-B}\bigr);
\]
\item more generally, if \(n\le C_1\ell^\gamma\) and
\(R_d\le C_0\ell^{-B}\), then the exceptional probability is at most
\[
  C_{A,K,a,C_0,C_1,B,\gamma}
  \bigl(d^{-A}+\ell^{\min\{\gamma,2\}-B}\bigr);
\]
\item if \(R_d\le C_0d^{-B}\) with \(B>0\), then for every \(0<B'<B\)
the exceptional probability is at most
\[
  C_{A,K,a,C_0,B,B'}\bigl(d^{-A}+d^{-B'}\bigr).
\]
\end{enumerate}
On the complement of the same exceptional event, independent signs have the
\(C_{A,K,a}\sqrt n\) edge bound with conditional probability at least
\(1-d^{-A}\), sign-averaged edge maximum \(C_{A,K,a}\sqrt n\), the same
deterministic certificate, and scalar edge-grid moments bounded by
\((C_{A,K,a}k\,n/\ell)^k\) for every \(2k\)-th moment.
\end{corollary}

\begin{proof}
The first display is exactly Corollary~\ref{cor:normalizer-tail-quant}.  The
first polylogarithmic case follows by substituting
\(R_d\le C_0\ell^{-B}\) and using \(\min\{n,\ell^2\}\le\ell^2\).  The
sample-window refinement uses
\(\min\{n,\ell^2\}\le C_1\ell^{\min\{\gamma,2\}}\).  For the polynomial case,
enlarge the constant to cover bounded \(d\); for large \(d\),
\(\ell^2\le C_{B,B'}d^{B-B'}\), so
\(R_d\min\{n,\ell^2\}\le C_{C_0,B,B'}d^{-B'}\).
\end{proof}

\begin{corollary}[Explicit search-stability rates]
\label{cor:normalizer-tail-search-rates}
Fix \(A>0\), \(a>0\), and write \(\ell=\log(e+d)\).  In the setting of
Corollary~\ref{cor:normalizer-tail-search-stability}, suppose
\[
  p_d(a):=\Prob\{\norm{T_d(\xi)}_{\op}<aN_d\}\le R_d .
\]
Then, whenever \(n\le d\) and \(n\ge C_{A,K,a}\ell\), outside a matrix event of
probability at most
\[
  C_{A,K,a}d^{-A}+R_d\min\{n,\ell^2\},
\]
the full random-search and stable-center package of
Corollary~\ref{cor:normalizer-tail-search-stability} holds.  In particular:
\begin{enumerate}[label=\textup{(\roman*)}]
\item if \(R_d\le C_0\ell^{-B}\) with \(B>2\), this package fails with probability
at most
\[
  C_{A,K,a,C_0,B}\bigl(d^{-A}+\ell^{2-B}\bigr);
\]
\item if \(n\le C_1\ell^\gamma\) and \(R_d\le C_0\ell^{-B}\), it fails with
probability at most
\[
  C_{A,K,a,C_0,C_1,B,\gamma}
  \bigl(d^{-A}+\ell^{\min\{\gamma,2\}-B}\bigr);
\]
\item if \(R_d\le C_0d^{-B}\) with \(B>0\), then for every \(0<B'<B\) it fails
with probability at most
\[
  C_{A,K,a,C_0,B,B'}\bigl(d^{-A}+d^{-B'}\bigr).
\]
\end{enumerate}
\end{corollary}

\begin{proof}
Corollary~\ref{cor:normalizer-tail-search-stability} gives the first exceptional
probability bound.  The three displayed specializations are the same elementary
substitutions as in Corollary~\ref{cor:normalizer-tail-rates}.
\end{proof}

\begin{corollary}[Almost-sure normalizer-tail transfer]
\label{cor:normalizer-tail-borel}
Fix \(A>0\) and \(a>0\).  Let \((d_m,n_m)_{m\ge1}\) be a sequence with
\(n_m\le d_m\), put \(\ell_m=\log(e+d_m)\), and assume
\[
  n_m\ge C_{A,K,a}\ell_m
  \qquad(m\ge1).
\]
For each \(m\), let \(B_{m,i}=T_{d_m}(\xi^{(m,i)})\), \(1\le i\le n_m\), be
independent full Toeplitz matrices whose lag variables are independent, mean
zero, unit variance, and have sub-Gaussian norm at most \(K\), and let
\(P_{m,i}\) be their symbols.  Let
\[
  p_{d_m}(a)
  =\Prob\{\norm{T_{d_m}(\xi)}_{\op}<a\,\Exp\norm{T_{d_m}(\xi)}_{\op}\}.
\]
If
\[
  \sum_{m=1}^{\infty}
  \left(d_m^{-A}+p_{d_m}(a)\min\{n_m,\ell_m^2\}\right)<\infty,
\]
then, with probability one, all but finitely many \(m\) satisfy
\[
  \max_{\omega^{d_m}=1}
  \sum_{i=1}^{n_m}
  \frac{\abs{P_{m,i}(\omega)}^2}{\norm{B_{m,i}}_{\op}^2}
  \le C_{A,K,a}\frac{n_m}{\ell_m}.
\]
For those \(m\), the signing, sign-averaged maximum, deterministic-certificate,
and scalar moment conclusions of Corollary~\ref{cor:normalizer-tail-quant} hold
with \(d,n,\ell\) replaced by \(d_m,n_m,\ell_m\).
\end{corollary}

\begin{proof}
Apply Corollary~\ref{cor:normalizer-tail-quant} at each index \(m\).  The
displayed summability assumption makes the profile exceptional probabilities
summable.  The first Borel--Cantelli lemma gives the almost-sure eventual
profile bound, and the remaining conclusions are deterministic consequences of
that same good profile event in Corollary~\ref{cor:normalizer-tail-quant}.
\end{proof}

\begin{corollary}[Almost-sure random-search and stable-center transfer]
\label{cor:normalizer-tail-borel-search-stability}
In the setting of Corollary~\ref{cor:normalizer-tail-borel}, assume the same
summability condition
\[
  \sum_{m=1}^{\infty}
  \left(d_m^{-A}+p_{d_m}(a)\min\{n_m,\ell_m^2\}\right)<\infty .
\]
For \(\omega^{d_m}=1\), put
\[
  v_{m,i}(\omega):=
  \frac{P_{m,i}(\omega)}{\norm{B_{m,i}}_{\op}} .
\]
Then, with probability one over the matrix triangular array, all but finitely
many \(m\) have the following conditional sign-search properties.

For every integer \(M\ge1\), if
\(\delta^{(1)},\ldots,\delta^{(M)}\) are independent uniform signs in
\(\{-1,1\}^{n_m}\), then for every \(u\ge0\),
\[
  \Prob_\delta\left\{
  \max_{1\le s\le M}\max_{\omega^{d_m}=1}
  \left|\sum_{i=1}^{n_m}\delta_i^{(s)}v_{m,i}(\omega)\right|
  >
  C_{A,K,a}\sqrt{
    n_m\left(1+\frac{\log(eM)+u}{\ell_m}\right)}
  \right\}
  \le e^{-u}.
\]
Also,
\[
  \Prob_\delta\left\{
  \min_{1\le s\le M}\max_{\omega^{d_m}=1}
  \left|\sum_{i=1}^{n_m}\delta_i^{(s)}v_{m,i}(\omega)\right|
  > C_{A,K,a}\sqrt {n_m}
  \right\}\le d_m^{-AM}.
\]
If \(G_{m,M}\) is the number of sampled sign vectors whose edge maximum is at
most \(C_{A,K,a}\sqrt{n_m}\), then for every integer \(1\le r\le M\),
\[
  \Prob_\delta\{G_{m,M}\le M-r\}\le \binom Mr d_m^{-Ar}.
\]

Finally, with Rademacher probability at least \(1-d_m^{-A}\), a random sign
vector is the center of a Hamming ball of radius
\[
  r_m=\left\lfloor c_{A,K,a}\ell_m\right\rfloor
\]
on which every signing has edge maximum at most \(C_{A,K,a}\sqrt{n_m}\).  If
\(H_{m,M}\) counts the sampled signs that are centers of such good balls, then
for every integer \(1\le r\le M\),
\[
  \Prob_\delta\{H_{m,M}\le M-r\}\le \binom Mr d_m^{-Ar}.
\]
\end{corollary}

\begin{proof}
Apply Corollary~\ref{cor:normalizer-tail-search-stability} at each \(m\).  The
same summability assumption makes its matrix exceptional probabilities
summable.  Borel--Cantelli gives an almost-sure event on which all but finitely
many matrix realizations satisfy the finite-size search-stability package.
Substituting \(d,n,\ell\) by \(d_m,n_m,\ell_m\) gives the displayed conditional
statements.
\end{proof}

\begin{corollary}[Almost-sure normalized sampled-search counts]
\label{cor:normalizer-tail-borel-counts}
Fix \(A>0\) and \(a>0\).  Let \((d_m,n_m)_{m\ge1}\) be a triangular-array
sequence with
\[
  d_m\ge2,\qquad n_m\le d_m,\qquad \ell_m=\log(e+d_m),
\]
and assume
\[
  n_m\ge C_{A,K,a}\ell_m\qquad(m\ge1).
\]
At each \(m\), let \(B_{m,i}=T_{d_m}(\xi^{(m,i)})\), \(1\le i\le n_m\), be
independent full Toeplitz matrices with independent mean-zero, unit-variance
lag variables of sub-Gaussian norm at most \(K\), and let \(P_{m,i}\) be their
symbols.  Independently of the matrices, sample independent uniform sign
vectors
\[
  \delta_m^{(1)},\ldots,\delta_m^{(M_m)}\in\{-1,1\}^{n_m},
  \qquad M_m\ge1 .
\]
Let \(u_m\ge0\) and \(1\le r_m\le M_m\) be deterministic parameters, and put
\[
  p_{d_m}(a)
  =\Prob\{\norm{T_{d_m}(\xi)}_{\op}
  <a\,\Exp\norm{T_{d_m}(\xi)}_{\op}\}.
\]
If
\[
  \sum_{m=1}^\infty
  \left[
  d_m^{-A}+p_{d_m}(a)\min\{n_m,\ell_m^2\}
  +e^{-u_m}+\binom{M_m}{r_m}d_m^{-A r_m}
  \right]<\infty ,
\]
then, with probability one over the matrices and sampled signs, all but
finitely many \(m\) satisfy all of the following.

First, the entire sampled list obeys
\[
  \max_{1\le s\le M_m}\max_{\omega^{d_m}=1}
  \left|\sum_{i=1}^{n_m}\delta_{m,i}^{(s)}
  \frac{P_{m,i}(\omega)}{\norm{B_{m,i}}_{\op}}\right|
  \le
  C_{A,K,a}\sqrt{
    n_m\left(1+\frac{\log(eM_m)+u_m}{\ell_m}\right)} .
\]
Second, at least \(M_m-r_m+1\) sampled signs have edge maximum at most
\(C_{A,K,a}\sqrt{n_m}\).  Third, at least \(M_m-r_m+1\) sampled signs are
centers of Hamming balls of radius
\[
  \left\lfloor c_{A,K,a}\ell_m\right\rfloor
\]
on which every signing \(\eps'\) has edge maximum at most
\(C_{A,K,a}\sqrt{n_m}\).

Consequently, along the full integer dimension sequence, suppose \(s>1\),
\(p_d(a)\le C_0d^{-B_0}\) with \(B_0>1\), and \(M_d\le d^\beta\) for some
\(\beta\ge0\).  If \(A>\beta+1\), then taking \(u_d=s\log(e+d)\) and
\(r_d=1\) gives an almost-sure event on which, for all sufficiently large
integer dimensions and every admissible window
\[
  C_{A,K,a}\log(e+d)\le n(d)\le d,
\]
every sampled sign in the batch is a robust center and the whole batch satisfies
the displayed sampled-list maximum bound.

More generally, on a polynomial dimension schedule \(d_m\ge c_0m^\kappa\),
suppose \(p_d(a)\le C_0d^{-B_0}\), \(M_m\le d_m^\beta\), and
\[
  B_0\kappa>1,\qquad (A-\beta)\kappa>1 .
\]
Then taking \(u_m=s\log(e+m)\), \(s>1\), and \(r_m=1\) gives the same eventual
all-good, all-robust sampled-batch conclusion for every admissible window
\[
  C_{A,K,a}\log(e+d_m)\le n_m\le d_m.
\]
The sampled-list maximum is then \(O_{A,K,a,s,\beta,\kappa,c_0}(\sqrt{n_m})\).
\end{corollary}

\begin{proof}
Apply Corollary~\ref{cor:normalizer-tail-search-stability} at each \(m\), with
\[
  (d,n,\ell,M,u,r)=(d_m,n_m,\ell_m,M_m,u_m,r_m).
\]
The matrix event contributes
\[
  C_{A,K,a}d_m^{-A}+p_{d_m}(a)\min\{n_m,\ell_m^2\}.
\]
Conditioned on the complementary matrix event, the sampled-list maximum fails
with probability at most \(e^{-u_m}\), while the good-sign count and
robust-center count each fail with probability at most
\[
  \binom{M_m}{r_m}d_m^{-A r_m}.
\]
Thus the joint matrix-and-sign failure probabilities are summable, up to a
constant factor.  The first Borel--Cantelli lemma gives an almost-sure event on
which all three conclusions hold for all but finitely many \(m\).

For the full-dimension consequence take \(d_m=m\), \(u_m=s\log(e+m)\), and
\(r_m=1\).  Then \(\sum_m m^{-A}<\infty\), \(\sum_m e^{-u_m}<\infty\), and
\[
  \binom{M_m}{1}m^{-A}\le m^{\beta-A},
\]
which is summable because \(A>\beta+1\).  The lower-normalizer contribution is
bounded by a constant times \(m^{-B_0}\log^2(e+m)\), hence is summable when
\(B_0>1\).  With \(r_m=1\), the count lower bounds say every sampled sign is
good and robust.

For the polynomial schedule, again set \(u_m=s\log(e+m)\) and \(r_m=1\).  The
matrix tail \(d_m^{-A}\) is summable because \((A-\beta)\kappa>1\) and
\(\beta\ge0\).  The lower-normalizer contribution is bounded by a constant times
\[
  m^{-B_0\kappa}\log^2(e+m),
\]
which is summable under \(B_0\kappa>1\).  The sampled-list tail is summable since
\(s>1\), and the binomial count term satisfies
\[
  \binom{M_m}{1}d_m^{-A}\le d_m^{\beta-A}
  \le C_{\beta,A,c_0}\,m^{-(A-\beta)\kappa},
\]
which is summable by hypothesis.  Finally
\[
  \frac{\log(eM_m)+u_m}{\ell_m}
  \le C_{s,\beta,\kappa,c_0}
\]
for all large \(m\), so the entropy-sensitive sampled-list maximum reduces to
the stated \(O(\sqrt{n_m})\) bound.  Finitely many initial indices are irrelevant.
\end{proof}

\begin{corollary}[Summable logarithmic-tail sampled counts]
\label{cor:normalizer-tail-borel-log-summable-counts}
Fix \(A>0\), \(a>0\), \(C_0<\infty\), \(C_1<\infty\), and \(\gamma\ge0\).  Put
\[
  \theta=\min\{\gamma,2\}.
\]
Suppose
\[
  p_d(a):=\Prob\{\norm{T_d(\xi)}_{\op}<a\,\Exp\norm{T_d(\xi)}_{\op}\}
  \le C_0\bigl(\log(e+d)\bigr)^{-B}
  \qquad(d\ge1).
\]
Let \((d_m,n_m)_{m\ge1}\) satisfy, with \(\ell_m=\log(e+d_m)\),
\[
  d_m\ge2,\qquad
  C_{A,K,a}\ell_m\le n_m\le \min\{d_m,C_1\ell_m^\gamma\}.
\]
At each \(m\), take independent full Toeplitz samples at dimension \(d_m\), and
independently sample \(M_m\) uniform sign vectors in \(\{-1,1\}^{n_m}\), with
\(M_m\ge1\).  Let \(u_m\ge0\) and \(1\le r_m\le M_m\) be deterministic
parameters.  If
\[
  \sum_{m=1}^{\infty}
  \left(
    d_m^{-A}+\ell_m^{\theta-B}
    +e^{-u_m}+\binom{M_m}{r_m}d_m^{-A r_m}
  \right)<\infty,
\]
then, with probability one over the matrices and sampled signs, all but
finitely many \(m\) satisfy
\[
\begin{aligned}
  \max_{1\le s\le M_m}\max_{\omega^{d_m}=1}
  \left|\sum_{i=1}^{n_m}\delta_{m,i}^{(s)}
  \frac{P_{m,i}(\omega)}{\norm{B_{m,i}}_{\op}}\right|
  \le
  C_{A,K,a}\sqrt{
    n_m\left(1+\frac{\log(eM_m)+u_m}{\ell_m}\right)} .
\end{aligned}
\]
Moreover, at least \(M_m-r_m+1\) sampled signs have edge maximum at most
\(C_{A,K,a}\sqrt{n_m}\), and at least \(M_m-r_m+1\) sampled signs are centers of
Hamming balls of radius \(\lfloor c_{A,K,a}\ell_m\rfloor\) on which every
signing has edge maximum at most \(C_{A,K,a}\sqrt{n_m}\).

In particular, if \(r_m=1\), \(M_m\le d_m^\beta\) for some \(0\le\beta<A\),
\[
  \sum_m\left(\ell_m^{\theta-B}+d_m^{\beta-A}+e^{-u_m}\right)<\infty,
  \qquad
  \sup_m\frac{u_m}{\ell_m}<\infty,
\]
then eventually every sampled sign in the batch is an \(O(\sqrt{n_m})\)-good
robust center and the whole batch satisfies an \(O(\sqrt{n_m})\) sampled-list
maximum bound.  If \(0<\rho\le1\), \(0<\tau_0<A\),
\(r_m=\lceil\rho M_m\rceil\), and
\[
  \sum_m\left(\ell_m^{\theta-B}+d_m^{-(A-\tau_0)}+e^{-u_m}\right)<\infty,
\]
then no upper bound on \(M_m\) is needed to get, eventually almost surely, a
\(1-\rho\) fraction of sampled signs which are logarithmically robust centers;
only the displayed sampled-list maximum keeps its entropy factor
\(\log(eM_m)\).
\end{corollary}

\begin{proof}
The logarithmic sample window gives
\[
  \min\{n_m,\ell_m^2\}\le C_{C_1}\ell_m^\theta .
\]
Hence
\[
  p_{d_m}(a)\min\{n_m,\ell_m^2\}
  \le C_{C_0,C_1}\ell_m^{\theta-B}.
\]
The displayed summability hypothesis is therefore stronger than the hypothesis
of Corollary~\ref{cor:normalizer-tail-borel-counts}, which gives the sampled-list
maximum and count assertions.

For the polynomial-in-\(d_m\) batch consequence, \(d_m^{-A}\le d_m^{\beta-A}\)
and
\[
  \binom{M_m}{1}d_m^{-A}\le d_m^{\beta-A}.
\]
The assumed boundedness of \(u_m/\ell_m\) also gives
\((\log(eM_m)+u_m)/\ell_m=O(1)\).  Thus the entropy-sensitive sampled-list
maximum has the \(O(\sqrt{n_m})\) form, and the count conclusion with \(r_m=1\)
says every sampled sign is good and robust.

For \(r_m=\lceil\rho M_m\rceil\),
\[
  \binom{M_m}{r_m}d_m^{-A r_m}
  \le
  \left(\frac{e}{\rho}d_m^{-A}\right)^{r_m}.
\]
For all large \(m\) this is at most
\[
  d_m^{-(A-\tau_0)r_m}\le d_m^{-(A-\tau_0)},
\]
which is summable by hypothesis.  Applying the first part gives the eventual
\(1-\rho\) robust-center fraction.
\end{proof}

\begin{corollary}[Logarithmic-tail sampled-count schedule design]
\label{cor:normalizer-tail-borel-log-schedule-counts}
Fix \(A>0\), \(a>0\), \(C_0<\infty\), \(C_1<\infty\), and \(\gamma\ge0\).  Put
\[
  \theta=\min\{\gamma,2\}.
\]
Suppose
\[
  p_d(a):=\Prob\{\norm{T_d(\xi)}_{\op}<a\,\Exp\norm{T_d(\xi)}_{\op}\}
  \le C_0\bigl(\log(e+d)\bigr)^{-B}
  \qquad(d\ge1).
\]
Let \(h_m\ge1\) be deterministic numbers such that
\[
  \sum_{m=1}^\infty h_m^{\theta-B}<\infty .
\]
Let \((d_m,n_m)_{m\ge1}\) satisfy, with \(\ell_m=\log(e+d_m)\),
\[
  d_m\ge2,\qquad \ell_m\ge h_m,\qquad
  C_{A,K,a}\ell_m\le n_m\le \min\{d_m,C_1\ell_m^\gamma\}.
\]
At each \(m\), take independent full Toeplitz samples at dimension \(d_m\), and
independently sample \(M_m\) uniform sign vectors in \(\{-1,1\}^{n_m}\), with
\(M_m\ge1\).  Let \(u_m\ge0\) and \(1\le r_m\le M_m\) be deterministic
parameters.  If
\[
  \sum_{m=1}^{\infty}
  \left(e^{-u_m}+\binom{M_m}{r_m}d_m^{-A r_m}\right)<\infty,
\]
then the sampled-list maximum, good-sign count, and robust-center count
conclusions of Corollary~\ref{cor:normalizer-tail-borel-log-summable-counts}
hold eventually almost surely.

In particular, if \(r_m=1\), \(M_m\le d_m^\beta\) for some \(0\le\beta<A\),
\(\sum_m e^{-u_m}<\infty\), and
\[
  \sup_m\frac{u_m}{\ell_m}<\infty,
\]
then eventually every sampled sign in the batch is an \(O(\sqrt{n_m})\)-good
robust center and the whole batch satisfies an \(O(\sqrt{n_m})\) sampled-list
maximum bound.  If \(0<\rho\le1\), \(r_m=\lceil\rho M_m\rceil\), and
\(\sum_m e^{-u_m}<\infty\), then no upper bound on \(M_m\) is needed to get,
eventually almost surely, a \(1-\rho\) fraction of sampled signs which are
logarithmically robust centers; only the displayed sampled-list maximum keeps
its entropy factor \(\log(eM_m)\).
\end{corollary}

\begin{proof}
Since the displayed series has terms at least one unless \(B>\theta\), the
hypothesis forces \(B>\theta\) and \(h_m\to\infty\).  Therefore, for every
\(c>0\), there is \(m_0(c)\) such that
\[
  e^{-c h_m}\le h_m^{\theta-B}\qquad(m\ge m_0(c)).
\]
The lower bound \(\ell_m\ge h_m\) gives
\[
  \ell_m^{\theta-B}\le h_m^{\theta-B},
  \qquad
  d_m^{-c}\le C_c e^{-c h_m}\le C_c h_m^{\theta-B}
\]
for all large \(m\).  Thus the matrix summands in
Corollary~\ref{cor:normalizer-tail-borel-log-summable-counts} are summable, and
the general sampled-count conclusion follows from the assumed sign-tail
summability.

If \(M_m\le d_m^\beta\) and \(r_m=1\), then
\[
  \binom{M_m}{1}d_m^{-A}\le d_m^{\beta-A}
  \le e^{-(A-\beta)h_m},
\]
which is summable by the preceding comparison.  Also
\((\log(eM_m)+u_m)/\ell_m=O(1)\), so the sampled-list maximum is
\(O(\sqrt{n_m})\), and \(r_m=1\) makes every sampled sign good and robust.

For \(r_m=\lceil\rho M_m\rceil\),
\[
  \binom{M_m}{r_m}d_m^{-A r_m}
  \le
  \left(\frac{e}{\rho}d_m^{-A}\right)^{r_m}
\]
for all large \(m\).  Hence, for any fixed \(0<\tau_0<A\), the last display is at
most
\[
  d_m^{-(A-\tau_0)}
  \le C_{A,\tau_0}e^{-(A-\tau_0)h_m}.
\]
This is summable by the first paragraph, so the
sampled-count theorem gives the robust-center fraction.
\end{proof}

\begin{corollary}[Vanishing discarded fractions on logarithmic-tail schedules]
\label{cor:normalizer-tail-borel-log-vanishing-fractions}
In the setting of
Corollary~\ref{cor:normalizer-tail-borel-log-schedule-counts}, fix
\(0<\tau_0<A\), and suppose
\((\rho_m)_{m\ge1}\subset(0,1]\) satisfies
\[
  \rho_m\ge e\,d_m^{-\tau_0}
  \qquad\text{for all sufficiently large }m,
\]
and set
\[
  r_m=\lceil\rho_mM_m\rceil .
\]
If
\[
  \sum_m e^{-u_m}<\infty ,
\]
then, with probability one over the matrices and sampled signs, all but finitely
many \(m\) satisfy the sampled-list maximum bound from
Corollary~\ref{cor:normalizer-tail-borel-log-schedule-counts}, and at least a
\(1-\rho_m\) fraction of the \(M_m\) sampled signs are logarithmically robust
centers.  No upper bound on \(M_m\) is required.  In particular, if
\(\rho_m\to0\), then the robust-center fraction tends to one along the sampled
batches, while the whole-list maximum retains only the entropy factor
\(\log(eM_m)\).
\end{corollary}

\begin{proof}
By Corollary~\ref{cor:normalizer-tail-borel-log-schedule-counts}, it is enough
to check the binomial count-tail summability.  Since \(r_m=\lceil\rho_mM_m\rceil\),
\[
  \binom{M_m}{r_m}
  \le \left(\frac{eM_m}{r_m}\right)^{r_m}
  \le \left(\frac e{\rho_m}\right)^{r_m}.
\]
The lower bound on \(\rho_m\) gives, for all large \(m\),
\[
  \binom{M_m}{r_m}d_m^{-A r_m}
  \le d_m^{-(A-\tau_0) r_m}\le d_m^{-(A-\tau_0)}.
\]
The proof of Corollary~\ref{cor:normalizer-tail-borel-log-schedule-counts} shows
\(\sum_m d_m^{-(A-\tau_0)}<\infty\) under the lower-profile hypothesis.  The
sampled-count theorem therefore applies.  Finally,
\[
  M_m-r_m+1\ge M_m-\rho_mM_m,
\]
so the count lower bound is at least a \(1-\rho_m\) fraction.
\end{proof}

\begin{corollary}[Full-integer arbitrary-batch robust fractions]
\label{cor:normalizer-tail-borel-all-dim-growing-fractions}
Fix \(A>1\), \(a>0\), \(s>1\), and \(0<\rho\le1\).  Suppose that
\[
  p_d(a):=\Prob\{\norm{T_d(\xi)}_{\op}<a\,\Exp\norm{T_d(\xi)}_{\op}\}
  \le C_0d^{-B_0}\qquad(d\ge1)
\]
for some \(C_0<\infty\) and \(B_0>1\).  Let \(n(d)\) be any admissible full
integer sample window with
\[
  C_{A,K,a}\log(e+d)\le n(d)\le d
  \qquad\text{for all large }d .
\]
At each integer dimension \(d\), let \(B_{d,i}=T_d(\xi^{(d,i)})\),
\(1\le i\le n(d)\), be independent full Toeplitz matrices with independent
mean-zero, unit-variance lag variables of sub-Gaussian norm at most \(K\), and
let \(P_{d,i}\) be their symbols.  Independently sample an arbitrary deterministic
number \(M_d\ge1\) of uniform sign vectors
\[
  \delta_d^{(1)},\ldots,\delta_d^{(M_d)}\in\{-1,1\}^{n(d)} .
\]
Then, with probability one over the matrices and sampled signs, all sufficiently
large integer dimensions satisfy both conclusions below.

First, the whole sampled list obeys
\[
\begin{aligned}
  \max_{1\le t\le M_d}\max_{\omega^d=1}
  \left|\sum_{i=1}^{n(d)}\delta_{d,i}^{(t)}
  \frac{P_{d,i}(\omega)}{\norm{B_{d,i}}_{\op}}\right|
  \le
  C_{A,K,a}\sqrt{
    n(d)\left(1+\frac{\log(eM_d)+s\log(e+d)}{\log(e+d)}\right)} .
\end{aligned}
\]
Second, at least a \(1-\rho\) fraction of the \(M_d\) sampled signs have edge
maximum at most \(C_{A,K,a}\sqrt{n(d)}\) and are centers of Hamming balls of
radius
\[
  \left\lfloor c_{A,K,a}\log(e+d)\right\rfloor
\]
on which every signing \(\eps'\) has edge maximum at most
\[
  C_{A,K,a}\sqrt{n(d)} .
\]
No upper bound on \(M_d\) is required for this robust-fraction conclusion.
\end{corollary}

\begin{proof}
Apply Corollary~\ref{cor:normalizer-tail-borel-counts} with \(d_m=m\),
\(n_m=n(m)\),
\[
  u_m=s\log(e+m),\qquad r_m=\lceil\rho M_m\rceil ,
\]
where \(M_m\) denotes the sampled batch size at dimension \(m\).  The matrix
summand \(\sum_m m^{-A}\) is finite because \(A>1\), the sampled-list summand is
finite because \(s>1\), and
\[
  p_m(a)\min\{n(m),\log^2(e+m)\}
  \le C_0m^{-B_0}\log^2(e+m)
\]
is summable because \(B_0>1\).  It remains to check the binomial count tail.
Since \(r_m\ge\rho M_m\),
\[
  \binom{M_m}{r_m}m^{-A r_m}
  \le \left(\frac{eM_m}{r_m}\right)^{r_m}m^{-A r_m}
  \le \left(\frac{e}{\rho}m^{-A}\right)^{r_m}.
\]
Choose \(0<\tau_0<A-1\).  For all large \(m\), the last display is at most
\[
  m^{-(A-\tau_0)r_m}\le m^{-(A-\tau_0)},
\]
which is summable because \(A-\tau_0>1\).  The sampled-count theorem gives the
entropy-sensitive list maximum and at least
\[
  M_m-r_m+1\ge (1-\rho)M_m
\]
good robust centers eventually.
\end{proof}

\begin{corollary}[Full-integer power-vanishing robust fractions]
\label{cor:normalizer-tail-borel-all-dim-vanishing-fractions}
Fix \(A>1\), \(a>0\), and \(0<\tau<A-1\).  Suppose that
\[
  p_d(a):=\Prob\{\norm{T_d(\xi)}_{\op}<a\,\Exp\norm{T_d(\xi)}_{\op}\}
  \le C_0d^{-B_0}\qquad(d\ge1)
\]
for some \(C_0<\infty\) and \(B_0>1\).  Let \(n(d)\) be any admissible full
integer sample window with
\[
  C_{A,K,a}\log(e+d)\le n(d)\le d
  \qquad\text{for all large }d .
\]
At each integer dimension \(d\), let \(B_{d,i}=T_d(\xi^{(d,i)})\),
\(1\le i\le n(d)\), be independent full Toeplitz matrices with independent
mean-zero, unit-variance lag variables of sub-Gaussian norm at most \(K\), and
let \(P_{d,i}\) be their symbols.  Independently sample an arbitrary
deterministic number \(M_d\ge1\) of uniform sign vectors
\[
  \delta_d^{(1)},\ldots,\delta_d^{(M_d)}\in\{-1,1\}^{n(d)} .
\]
Let \(u_d\ge0\) satisfy
\[
  \sum_{d=1}^{\infty}e^{-u_d}<\infty ,
\]
and let \((\rho_d)_{d\ge1}\subset(0,1]\) satisfy
\[
  \rho_d\ge e\,d^{-\tau}
  \qquad\text{for all sufficiently large }d .
\]
Then, with probability one over the matrices and sampled signs, all
sufficiently large integer dimensions satisfy both conclusions below.

First, the whole sampled list obeys
\[
\begin{aligned}
  \max_{1\le t\le M_d}\max_{\omega^d=1}
  \left|\sum_{i=1}^{n(d)}\delta_{d,i}^{(t)}
  \frac{P_{d,i}(\omega)}{\norm{B_{d,i}}_{\op}}\right|
  \le
  C_{A,K,a}\sqrt{
    n(d)\left(1+\frac{\log(eM_d)+u_d}{\log(e+d)}\right)} .
\end{aligned}
\]
Second, at least a \(1-\rho_d\) fraction of the \(M_d\) sampled signs have edge
maximum at most \(C_{A,K,a}\sqrt{n(d)}\) and are centers of Hamming balls of
radius
\[
  \left\lfloor c_{A,K,a}\log(e+d)\right\rfloor
\]
on which every signing \(\eps'\) has edge maximum at most
\[
  C_{A,K,a}\sqrt{n(d)} .
\]
No upper bound on \(M_d\) is required.  In particular, if \(\rho_d\to0\), then
the robust-center fraction tends to one along the full integer dimension
sequence.
\end{corollary}

\begin{proof}
Apply Corollary~\ref{cor:normalizer-tail-borel-counts} with \(d_m=m\),
\(n_m=n(m)\), and
\[
  r_m=\lceil\rho_mM_m\rceil ,
\]
where \(M_m\), \(u_m\), and \(\rho_m\) are the parameters at dimension \(m\).
The matrix summand \(\sum_m m^{-A}\) is finite because \(A>1\), the sampled-list
summand is finite by hypothesis, and
\[
  p_m(a)\min\{n(m),\log^2(e+m)\}
  \le C_0m^{-B_0}\log^2(e+m)
\]
is summable because \(B_0>1\).  For the binomial count tail,
\[
  \binom{M_m}{r_m}
  \le \left(\frac{eM_m}{r_m}\right)^{r_m}
  \le \left(\frac e{\rho_m}\right)^{r_m}.
\]
The lower bound \(\rho_m\ge e\,m^{-\tau}\) gives, for all large \(m\),
\[
  \binom{M_m}{r_m}m^{-A r_m}
  \le m^{-(A-\tau) r_m}\le m^{-(A-\tau)},
\]
which is summable because \(A-\tau>1\).  Hence the sampled-count theorem applies.
Finally,
\[
  M_m-r_m+1\ge M_m-\rho_mM_m,
\]
so the robust-center count is at least a \(1-\rho_m\) fraction.
\end{proof}

\begin{corollary}[Almost-sure normalized robust fractions]
\label{cor:normalizer-tail-borel-fractions}
Fix \(A>0\), \(S>0\), \(a>0\), and a fraction \(0<\rho\le1\).  Suppose
\[
  A\ge 2\log_2(e/\rho),
\]
and choose an integer \(M\ge1\) so that
\[
  M\ge \frac{2(S+1)}{A\rho}.
\]
Let \((d_m,n_m)_{m\ge1}\) be a triangular-array sequence with
\[
  d_m\ge2,\qquad n_m\le d_m,\qquad \ell_m=\log(e+d_m),
\]
and assume
\[
  n_m\ge C_{A,K,a}\ell_m \qquad(m\ge1).
\]
At each \(m\), let \(B_{m,i}=T_{d_m}(\xi^{(m,i)})\), \(1\le i\le n_m\), be
independent full Toeplitz matrices with independent mean-zero, unit-variance
lag variables of sub-Gaussian norm at most \(K\), and let \(P_{m,i}\) be their
symbols.  Independently of the matrices, sample independent uniform sign
vectors
\[
  \delta_m^{(1)},\ldots,\delta_m^{(M)}\in\{-1,1\}^{n_m}.
\]
Put
\[
  p_{d_m}(a)
  =\Prob\{\norm{T_{d_m}(\xi)}_{\op}
  <a\,\Exp\norm{T_{d_m}(\xi)}_{\op}\}.
\]
If
\[
  \sum_{m=1}^\infty
  \left(d_m^{-A}
  +p_{d_m}(a)\min\{n_m,\ell_m^2\}
  +d_m^{-S}\right)<\infty ,
\]
then, with probability one over the matrices and sampled signs, all but
finitely many \(m\) satisfy both conclusions below.

First, the whole sampled list obeys
\[
  \max_{1\le s\le M}\max_{\omega^{d_m}=1}
  \left|\sum_{i=1}^{n_m}\delta_{m,i}^{(s)}
  \frac{P_{m,i}(\omega)}{\norm{B_{m,i}}_{\op}}\right|
  \le C_{A,K,a,S,\rho}\sqrt{n_m}.
\]
Second, at least a \(1-\rho\) fraction of the \(M\) sampled signs are centers of
Hamming balls of radius
\[
  \left\lfloor c_{A,K,a}\ell_m\right\rfloor
\]
on which every signing \(\eps'\) has edge maximum at most
\[
  C_{A,K,a}\sqrt{n_m}.
\]
The same \(1-\rho\) fraction also has edge maximum at most
\(C_{A,K,a}\sqrt{n_m}\).

In particular, along the full integer dimension sequence, suppose \(A>1\),
\(S>1\), and
\[
  p_d(a)\le C_0d^{-B_0}\qquad(d\ge1)
\]
for some \(B_0>1\).  Then any admissible window
\[
  C_{A,K,a}\log(e+d)\le n(d)\le d
\]
has the same eventual robust-fraction conclusion almost surely.  More generally,
on polynomial dimension schedules \(d_m\ge c_0m^\kappa\), the same conclusion
holds under the three summability inequalities
\[
  A\kappa>1,\qquad B_0\kappa>1,\qquad S\kappa>1 .
\]
\end{corollary}

\begin{proof}
Let \(r_\rho=\lceil\rho M\rceil\), and put \(u_m=S\ell_m\).  Apply the finite
Corollary~\ref{cor:normalizer-tail-search-stability} at dimension \(m\).  The
matrix exceptional probability contributes
\[
  C_{A,K,a}d_m^{-A}+p_{d_m}(a)\min\{n_m,\ell_m^2\}.
\]
Outside that matrix event, the sampled-list maximum fails with conditional
probability at most \(e^{-u_m}\), while the good-sign count and robust-center
count each fail with conditional probability at most
\[
  \binom{M}{r_\rho}d_m^{-A r_\rho}.
\]
It remains to bound these two sign terms by the displayed \(d_m^{-S}\) summand.
Since \(r_\rho\ge\rho M\),
\[
  \binom{M}{r_\rho}\le \left(\frac{eM}{r_\rho}\right)^{r_\rho}
  \le (e/\rho)^{r_\rho}.
\]
For \(d_m\ge2\) and \(A\ge2\log_2(e/\rho)\),
\[
  \binom{M}{r_\rho}d_m^{-A r_\rho}
  \le d_m^{-A r_\rho/2}.
\]
The lower bound on \(M\) gives \(A r_\rho/2\ge S+1\), hence
\[
  2\binom{M}{r_\rho}d_m^{-A r_\rho}\le 2d_m^{-(S+1)}\le d_m^{-S}.
\]
Also \(e^{-u_m}=(e+d_m)^{-S}\le d_m^{-S}\).  Therefore the total
matrix-and-sign failure probability is summable.  Borel--Cantelli gives the
eventual sampled-list maximum and the \(1-\rho\) fraction of good and robust
sampled signs.  Since
\[
  1+\frac{\log(eM)+u_m}{\ell_m}
  \le C_{S,\rho,A},
\]
the sampled-list maximum has the displayed \(O(\sqrt{n_m})\) form.

The full-dimension polynomial-tail consequence follows by taking \(d_m=m\):
\(\sum_m m^{-A}\), \(\sum_m m^{-S}\), and
\(\sum_m m^{-B_0}\log^2(e+m)\) are finite under \(A,S,B_0>1\).  On polynomial
dimension schedules the same estimates are bounded by constants times
\[
  m^{-A\kappa},\qquad
  m^{-S\kappa},\qquad
  m^{-B_0\kappa}\log^2(e+m),
\]
which are summable under the three displayed inequalities.
\end{proof}

\begin{corollary}[Almost-sure sampled robust certificates]
\label{cor:normalizer-tail-borel-poly-confidence}
Fix \(A>0\), \(\alpha>0\), and \(a>0\).  Let
\((d_m,n_m)_{m\ge1}\) be a sequence with \(n_m\le d_m\), put
\(\ell_m=\log(e+d_m)\), and assume
\[
  n_m\ge C_{A,K,a}\ell_m
  \qquad(m\ge1).
\]
At each \(m\), let \(B_{m,i}=T_{d_m}(\xi^{(m,i)})\), \(1\le i\le n_m\), be
independent full Toeplitz matrices with independent mean-zero, unit-variance
lag variables of sub-Gaussian norm at most \(K\), and let \(P_{m,i}\) be their
symbols.  Let
\[
  p_{d_m}(a)
  =\Prob\{\norm{T_{d_m}(\xi)}_{\op}
  <a\,\Exp\norm{T_{d_m}(\xi)}_{\op}\}.
\]
For each \(m\), sample \(M_m\) independent uniform sign vectors
\[
  \delta_m^{(1)},\ldots,\delta_m^{(M_m)}\in\{-1,1\}^{n_m},
  \qquad M_m\ge\lceil\alpha\rceil ,
\]
independently of the matrices.  If
\[
  \sum_{m=1}^\infty
  \left(d_m^{-A}
  +p_{d_m}(a)\min\{n_m,\ell_m^2\}
  +d_m^{-\alpha}\right)<\infty ,
\]
then, with probability one over the matrices and the sampled signs, all but
finitely many \(m\) have a sampled sign \(\delta_m^{(s)}\) which is the center
of a Hamming ball of radius
\[
  r_m=\left\lfloor c_{A,K,a}\ell_m\right\rfloor
\]
on which every signing \(\eps'\) satisfies
\[
  \max_{\omega^{d_m}=1}\left|\sum_{i=1}^{n_m}
  \eps'_i\frac{P_{m,i}(\omega)}{\norm{B_{m,i}}_{\op}}\right|
  \le C_{A,K,a}\sqrt{n_m}.
\]
Thus the finite random search itself eventually returns logarithmically robust
\(O(\sqrt{n_m})\) certificates, rather than only a conditional certificate
probability.

In particular, along the full integer dimension sequence, if \(A>1\),
\(\alpha>1\), and
\[
  p_d(a)\le C_0d^{-B_0}\qquad(d\ge1)
\]
for some \(B_0>1\), then any admissible window
\[
  C_{A,K,a}\log(e+d)\le n(d)\le d
\]
and any fixed list size \(M\ge\lceil\alpha\rceil\) give, almost surely, robust
sampled certificates for all sufficiently large integer dimensions.
\end{corollary}

\begin{proof}
Apply Corollary~\ref{cor:normalizer-tail-poly-confidence} at each \(m\), with
\(B\) there replaced by \(\alpha\).  The probability that the \(m\)-th matrix
array and its sampled sign list fail to contain the displayed robust certificate
is at most
\[
  C_{A,K,a}d_m^{-A}
  +p_{d_m}(a)\min\{n_m,\ell_m^2\}
  +d_m^{-\alpha}.
\]
The series of these probabilities is finite by hypothesis, so the first
Borel--Cantelli lemma gives the eventual sampled-certificate conclusion.

For the full-dimension consequence take \(d_m=m\), \(n_m=n(m)\), and
\(M_m=M\).  Then \(\sum_m m^{-A}\) and \(\sum_m m^{-\alpha}\) are finite, while
\[
  p_m(a)\min\{n(m),\log^2(e+m)\}
  \le C_0m^{-B_0}\log^2(e+m)
\]
is summable because \(B_0>1\).  The general assertion applies.
\end{proof}

\begin{corollary}[All-dimensions almost-sure transfer from polynomial lower tails]
\label{cor:normalizer-tail-borel-all-dim-poly}
Fix \(a>0\), \(C_0<\infty\), and \(A>1\).  Suppose
\[
  p_d(a):=\Prob\{\norm{T_d(\xi)}_{\op}<a\,\Exp\norm{T_d(\xi)}_{\op}\}
  \le C_0d^{-B}
  \qquad(d\ge1)
\]
for some \(B>1\).  For each \(d\ge1\), let \(n(d)\) satisfy
\[
  C_{A,K,a}\log(e+d)\le n(d)\le d.
\]
Then, for independent full Toeplitz samples \(B_{d,i}=T_d(\xi^{(d,i)})\),
\(1\le i\le n(d)\), at every integer dimension \(d\), with symbols
\(P_{d,i}\), with probability one all but finitely many dimensions satisfy
\[
  \max_{\omega^d=1}
  \sum_{i=1}^{n(d)}
  \frac{\abs{P_{d,i}(\omega)}^2}{\norm{B_{d,i}}_{\op}^2}
  \le C_{A,K,a}\frac{n(d)}{\log(e+d)}.
\]
For those \(d\), the signing, sign-averaged maximum, deterministic-certificate,
and scalar moment conclusions of Corollary~\ref{cor:normalizer-tail-quant} hold
with \(n,\ell\) replaced by \(n(d),\log(e+d)\).
\end{corollary}

\begin{proof}
Apply Corollary~\ref{cor:normalizer-tail-borel} with \(d_m=m\) and
\(n_m=n(m)\).  The matrix-event part is summable because \(A>1\).  Also
\[
  p_m(a)\min\{n_m,\log^2(e+m)\}
  \le C_0m^{-B}\log^2(e+m),
\]
which is summable because \(B>1\).  Hence the general Borel--Cantelli transfer
gives the eventual profile, and the deterministic consequences again come from
Corollary~\ref{cor:normalizer-tail-quant}.
\end{proof}

\begin{corollary}[Sparse-scale almost-sure transfer]
\label{cor:normalizer-tail-borel-sparse}
Fix \(A>0\), \(a>0\), \(C_0<\infty\), and \(\lambda>1\).  Suppose either the
polynomial tail
\[
  p_d(a):=\Prob\{\norm{T_d(\xi)}_{\op}<a\,\Exp\norm{T_d(\xi)}_{\op}\}
  \le C_0d^{-B}
  \qquad(d\ge1).
\]
holds for some \(B>0\), or the logarithmic tail
\[
  p_d(a)\le C_0\bigl(\log(e+d)\bigr)^{-B}
  \qquad(d\ge1)
\]
holds for some \(B>3\).
Let \((d_m,n_m)_{m\ge1}\) satisfy
\[
  d_m\ge \lambda^m,\qquad
  C_{A,K,a}\log(e+d_m)\le n_m\le d_m .
\]
Then, for independent full Toeplitz samples at the dimensions \(d_m\), with
probability one all but finitely many \(m\) satisfy
\[
  \max_{\omega^{d_m}=1}
  \sum_{i=1}^{n_m}
  \frac{\abs{P_{m,i}(\omega)}^2}{\norm{B_{m,i}}_{\op}^2}
  \le C_{A,K,a}\frac{n_m}{\log(e+d_m)}.
\]
For those \(m\), the signing, sign-averaged maximum, deterministic-certificate,
and scalar moment conclusions of Corollary~\ref{cor:normalizer-tail-quant} hold
with \(d,n,\ell\) replaced by \(d_m,n_m,\log(e+d_m)\).
\end{corollary}

\begin{proof}
By Corollary~\ref{cor:normalizer-tail-borel}, it is enough to verify
summability.  Put \(\ell_m=\log(e+d_m)\).  We use
\(\min\{n_m,\ell_m^2\}\le \ell_m^2\).  In the polynomial-tail case, for
\(x\ge1\),
\[
  x^{-B}\bigl(\log(e+x)\bigr)^2\le C_Bx^{-B/2}.
\]
Then
\[
  p_{d_m}(a)\min\{n_m,\ell_m^2\}
  \le C_{B,C_0}d_m^{-B/2}
  \le C_{B,C_0,\lambda}\lambda^{-Bm/2}.
\]
In the logarithmic-tail case,
\[
  p_{d_m}(a)\min\{n_m,\ell_m^2\}
  \le C_0\ell_m^{2-B}
  \le C_{B,C_0,\lambda}m^{2-B},
\]
which is summable because \(B>3\).
Also \(d_m^{-A}\le \lambda^{-Am}\).  The required series is therefore finite,
and the Borel--Cantelli transfer applies.
\end{proof}

\begin{corollary}[Sparse-scale sampled robust certificates]
\label{cor:normalizer-tail-borel-sparse-sampled}
Fix \(A>0\), \(\alpha>0\), \(a>0\), \(C_0<\infty\), and \(\lambda>1\).
Suppose either the polynomial tail
\[
  p_d(a):=\Prob\{\norm{T_d(\xi)}_{\op}<a\,\Exp\norm{T_d(\xi)}_{\op}\}
  \le C_0d^{-B}
  \qquad(d\ge1)
\]
holds for some \(B>0\), or the logarithmic tail
\[
  p_d(a)\le C_0\bigl(\log(e+d)\bigr)^{-B}
  \qquad(d\ge1)
\]
holds for some \(B>3\).
Let \((d_m,n_m)_{m\ge1}\) satisfy
\[
  d_m\ge \lambda^m,\qquad
  C_{A,K,a}\log(e+d_m)\le n_m\le d_m .
\]
For each \(m\), sample \(M_m\) independent uniform sign vectors in
\(\{-1,1\}^{n_m}\), independently of the matrices, with
\[
  M_m\ge\lceil\alpha\rceil .
\]
Then, with probability one over the matrices and sampled signs, all but
finitely many \(m\) have a sampled sign which is the center of a Hamming ball of
radius
\[
  r_m=\left\lfloor c_{A,K,a}\log(e+d_m)\right\rfloor
\]
on which every signing \(\eps'\) satisfies
\[
  \max_{\omega^{d_m}=1}\left|\sum_{i=1}^{n_m}
  \eps'_i\frac{P_{m,i}(\omega)}{\norm{B_{m,i}}_{\op}}\right|
  \le C_{A,K,a}\sqrt{n_m}.
\]
Thus exponentially sparse dimensions need no lower bound on the sampled-list
confidence exponent beyond \(\alpha>0\).
\end{corollary}

\begin{proof}
Apply Corollary~\ref{cor:normalizer-tail-borel-poly-confidence}.  It remains to
check
\[
  \sum_m\left(d_m^{-A}
  +p_{d_m}(a)\min\{n_m,\ell_m^2\}
  +d_m^{-\alpha}\right)<\infty,
  \qquad \ell_m=\log(e+d_m).
\]
The exponential growth of \(d_m\) makes \(\sum_m d_m^{-A}\) and
\(\sum_m d_m^{-\alpha}\) finite.  The same estimates used in
Corollary~\ref{cor:normalizer-tail-borel-sparse} show that the middle summand is
summable: in the polynomial-tail case it is at most
\(C_{B,C_0,\lambda}\lambda^{-Bm/2}\), while in the logarithmic-tail case it is
at most \(C_{B,C_0,\lambda}m^{2-B}\), summable because \(B>3\).  The sampled
Borel--Cantelli transfer applies.
\end{proof}

\begin{corollary}[Polynomial-scale almost-sure transfer from polynomial lower tails]
\label{cor:normalizer-tail-borel-polynomial-scale}
Fix \(a>0\), \(C_0<\infty\), \(c_0>0\), and \(\kappa>0\).  Suppose
\[
  p_d(a):=\Prob\{\norm{T_d(\xi)}_{\op}<a\,\Exp\norm{T_d(\xi)}_{\op}\}
  \le C_0d^{-B}
  \qquad(d\ge1)
\]
for some \(B>0\).  Let \(A>0\) satisfy
\[
  A\kappa>1,\qquad B\kappa>1.
\]
Let \((d_m,n_m)_{m\ge1}\) satisfy
\[
  d_m\ge c_0m^\kappa,\qquad
  C_{A,K,a}\log(e+d_m)\le n_m\le d_m .
\]
Then, for independent full Toeplitz samples at the dimensions \(d_m\), with
probability one all but finitely many \(m\) satisfy
\[
  \max_{\omega^{d_m}=1}
  \sum_{i=1}^{n_m}
  \frac{\abs{P_{m,i}(\omega)}^2}{\norm{B_{m,i}}_{\op}^2}
  \le C_{A,K,a}\frac{n_m}{\log(e+d_m)}.
\]
For those \(m\), the signing, sign-averaged maximum, deterministic-certificate,
and scalar moment conclusions of Corollary~\ref{cor:normalizer-tail-quant} hold
with \(d,n,\ell\) replaced by \(d_m,n_m,\log(e+d_m)\).
\end{corollary}

\begin{proof}
By Corollary~\ref{cor:normalizer-tail-borel}, it is enough to check
\[
  \sum_m\left(d_m^{-A}+p_{d_m}(a)\min\{n_m,\ell_m^2\}\right)<\infty,
  \qquad \ell_m:=\log(e+d_m).
\]
The first summand is at most \(C_{A,c_0}m^{-A\kappa}\), which is summable by
\(A\kappa>1\).  For the second summand, \(\min\{n_m,\ell_m^2\}\le\ell_m^2\),
and
\[
  p_{d_m}(a)\ell_m^2
  \le C_{B,C_0,c_0}m^{-B\kappa}(\log(e+m))^2.
\]
This is summable by \(B\kappa>1\).  The Borel--Cantelli transfer gives the
eventual profile, and the remaining conclusions are again deterministic
consequences of that profile in Corollary~\ref{cor:normalizer-tail-quant}.
\end{proof}

\begin{corollary}[Polynomial-scale sampled robust certificates]
\label{cor:normalizer-tail-borel-polynomial-sampled}
Fix \(a>0\), \(C_0<\infty\), \(c_0>0\), \(\kappa>0\), and \(A,\alpha>0\).
Suppose
\[
  p_d(a):=\Prob\{\norm{T_d(\xi)}_{\op}<a\,\Exp\norm{T_d(\xi)}_{\op}\}
  \le C_0d^{-B}
  \qquad(d\ge1)
\]
for some \(B>0\).  Assume
\[
  A\kappa>1,\qquad B\kappa>1,\qquad \alpha\kappa>1.
\]
Let \((d_m,n_m)_{m\ge1}\) satisfy
\[
  d_m\ge c_0m^\kappa,\qquad
  C_{A,K,a}\log(e+d_m)\le n_m\le d_m .
\]
At each \(m\), sample \(M_m\) independent uniform sign vectors in
\(\{-1,1\}^{n_m}\), independently of the matrices, with
\[
  M_m\ge\lceil\alpha\rceil .
\]
Then, with probability one over the matrices and sampled signs, all but
finitely many \(m\) have a sampled sign which is the center of a Hamming ball of
radius
\[
  r_m=\left\lfloor c_{A,K,a}\log(e+d_m)\right\rfloor
\]
on which every signing \(\eps'\) satisfies
\[
  \max_{\omega^{d_m}=1}\left|\sum_{i=1}^{n_m}
  \eps'_i\frac{P_{m,i}(\omega)}{\norm{B_{m,i}}_{\op}}\right|
  \le C_{A,K,a}\sqrt{n_m}.
\]
Thus polynomially sparse dimension schedules need only
\(\alpha>1/\kappa\) sampled-list confidence, rather than the all-integer
condition \(\alpha>1\).
\end{corollary}

\begin{proof}
Apply Corollary~\ref{cor:normalizer-tail-borel-poly-confidence}.  It is enough
to verify
\[
  \sum_m\left(d_m^{-A}
  +p_{d_m}(a)\min\{n_m,\ell_m^2\}
  +d_m^{-\alpha}\right)<\infty,
  \qquad \ell_m=\log(e+d_m).
\]
The first and third summands are bounded by constants times
\(m^{-A\kappa}\) and \(m^{-\alpha\kappa}\), respectively, and are summable by
hypothesis.  For the middle term, \(\min\{n_m,\ell_m^2\}\le\ell_m^2\), hence
\[
  p_{d_m}(a)\ell_m^2
  \le C_{B,C_0,c_0}m^{-B\kappa}(\log(e+m))^2,
\]
which is summable because \(B\kappa>1\).  The general sampled-list Borel
corollary gives the conclusion.
\end{proof}

\begin{corollary}[Sparse almost-sure transfer in logarithmic sample windows]
\label{cor:normalizer-tail-borel-window}
Fix \(A>0\), \(a>0\), \(C_0<\infty\), \(C_1<\infty\), \(c_0>0\),
\(\kappa>0\), and \(\gamma\ge0\).  Suppose the logarithmic lower-normalizer
tail satisfies
\[
  p_d(a):=\Prob\{\norm{T_d(\xi)}_{\op}<a\,\Exp\norm{T_d(\xi)}_{\op}\}
  \le C_0\bigl(\log(e+d)\bigr)^{-B}
  \qquad(d\ge1)
\]
for some
\[
  B>\min\{\gamma,2\}+\frac1\kappa .
\]
Let \((d_m,n_m)_{m\ge1}\) satisfy, with \(\ell_m=\log(e+d_m)\),
\[
  d_m\ge\exp(c_0m^\kappa),\qquad
  C_{A,K,a}\ell_m\le n_m\le \min\{d_m,C_1\ell_m^\gamma\}.
\]
Then, for independent full Toeplitz samples at the dimensions \(d_m\), with
probability one all but finitely many \(m\) satisfy
\[
  \max_{\omega^{d_m}=1}
  \sum_{i=1}^{n_m}
  \frac{\abs{P_{m,i}(\omega)}^2}{\norm{B_{m,i}}_{\op}^2}
  \le C_{A,K,a}\frac{n_m}{\ell_m}.
\]
For those \(m\), the signing, sign-averaged maximum, deterministic-certificate,
and scalar moment conclusions of Corollary~\ref{cor:normalizer-tail-quant} hold
with \(d,n,\ell\) replaced by \(d_m,n_m,\ell_m\).
\end{corollary}

\begin{proof}
By Corollary~\ref{cor:normalizer-tail-borel}, it is enough to prove
\[
  \sum_m\left(d_m^{-A}+p_{d_m}(a)\min\{n_m,\ell_m^2\}\right)<\infty .
\]
The first term is summable because \(d_m\ge\exp(c_0m^\kappa)\).  For the
second term,
\[
  \min\{n_m,\ell_m^2\}
  \le C_{C_1}\ell_m^{\min\{\gamma,2\}},
\]
and hence
\[
  p_{d_m}(a)\min\{n_m,\ell_m^2\}
  \le C_{C_0,C_1}\ell_m^{\min\{\gamma,2\}-B}
  \le C_{C_0,C_1,c_0,\kappa}
        m^{\kappa(\min\{\gamma,2\}-B)}.
\]
This is summable precisely under the displayed condition
\(B>\min\{\gamma,2\}+1/\kappa\).  The Borel--Cantelli transfer then gives the
eventual profile, and the remaining statements are the deterministic
consequences of that profile in Corollary~\ref{cor:normalizer-tail-quant}.
\end{proof}

\begin{corollary}[Log-corrected sparse logarithmic-window transfer]
\label{cor:normalizer-tail-borel-window-log-corrected}
Fix \(A>0\), \(a>0\), \(C_0<\infty\), \(C_1<\infty\), \(c_0>0\),
\(\kappa>0\), \(\lambda\in\mathbb R\), and \(\gamma\ge0\).  Put
\[
  \theta=\min\{\gamma,2\}.
\]
Suppose the logarithmic lower-normalizer tail satisfies
\[
  p_d(a):=\Prob\{\norm{T_d(\xi)}_{\op}<a\,\Exp\norm{T_d(\xi)}_{\op}\}
  \le C_0\bigl(\log(e+d)\bigr)^{-B}
  \qquad(d\ge1),
\]
where
\[
  \kappa(B-\theta)>1
  \qquad\text{or}\qquad
  \kappa(B-\theta)=1\ \text{ and }\ \lambda(B-\theta)>1 .
\]
Let \((d_m,n_m)_{m\ge1}\) satisfy, with
\[
  \ell_m=\log(e+d_m),\qquad L_m=\log(e+m),
\]
the log-corrected sparse growth and logarithmic sample-window bounds
\[
  d_m\ge\exp(c_0m^\kappa L_m^\lambda),
  \qquad
  C_{A,K,a}\ell_m\le n_m\le \min\{d_m,C_1\ell_m^\gamma\}.
\]
Then, for independent full Toeplitz samples at the dimensions \(d_m\), with
probability one all but finitely many \(m\) satisfy
\[
  \max_{\omega^{d_m}=1}
  \sum_{i=1}^{n_m}
  \frac{\abs{P_{m,i}(\omega)}^2}{\norm{B_{m,i}}_{\op}^2}
  \le C_{A,K,a}\frac{n_m}{\ell_m}.
\]
For those \(m\), the signing, sign-averaged maximum, deterministic-certificate,
and scalar moment conclusions of Corollary~\ref{cor:normalizer-tail-quant} hold
with \(d,n,\ell\) replaced by \(d_m,n_m,\ell_m\).

Moreover, for every \(\alpha>0\), if \(M_m\ge\lceil\alpha\rceil\) independent
uniform sign vectors are sampled at dimension \(m\), then almost surely all but
finitely many \(m\) contain a sampled sign which is the center of a Hamming ball
of radius \(\lfloor c_{A,K,a}\ell_m\rfloor\) on which every signing has edge
maximum at most \(C_{A,K,a}\sqrt{n_m}\).  Finally, if \(0<\rho\le1\),
\(S>0\), \(A\ge2\log_2(e/\rho)\), and
\[
  M\ge \frac{2(S+1)}{A\rho},
\]
then a fixed list of \(M\) sampled signs has, eventually almost surely, a
\(1-\rho\) fraction of such logarithmically robust centers and the whole sampled
list has \(O_{A,K,a,S,\rho}(\sqrt{n_m})\) edge maximum.
\end{corollary}

\begin{proof}
By the logarithmic sample-window hypothesis,
\[
  \min\{n_m,\ell_m^2\}\le C_{C_1}\ell_m^\theta .
\]
The growth assumption gives \(\ell_m\ge c_0m^\kappa L_m^\lambda\) for all large
\(m\).  Hence
\[
\begin{aligned}
  p_{d_m}(a)\min\{n_m,\ell_m^2\}
  &\le C_{C_0,C_1}\ell_m^{\theta-B}  \\
  &\le
  C\,m^{-\kappa(B-\theta)}L_m^{-\lambda(B-\theta)} .
\end{aligned}
\]
The displayed condition on \(B\) is exactly the integral-test summability
condition for this last series.  Also, for every \(r>0\),
\[
  \sum_m d_m^{-r}
  \le
  \sum_m \exp\{-rc_0m^\kappa L_m^\lambda\}
  <\infty .
\]
Therefore
\[
  \sum_m\left(d_m^{-A}
  +p_{d_m}(a)\min\{n_m,\ell_m^2\}\right)<\infty .
\]
Corollary~\ref{cor:normalizer-tail-borel} gives the almost-sure eventual edge
profile and its deterministic signing, average, certificate, and moment
consequences.

Adding \(\sum_m d_m^{-\alpha}<\infty\) and applying
Corollary~\ref{cor:normalizer-tail-borel-poly-confidence} gives the sampled
robust-certificate conclusion.  Adding \(\sum_m d_m^{-S}<\infty\) and applying
Corollary~\ref{cor:normalizer-tail-borel-fractions} gives the fixed-list
robust-fraction conclusion under the displayed choices of \(A\) and \(M\).
\end{proof}

\begin{corollary}[Log-corrected sparse logarithmic-window sampled counts]
\label{cor:normalizer-tail-borel-window-log-corrected-counts}
Fix \(A>0\), \(a>0\), \(C_0<\infty\), \(C_1<\infty\), \(c_0>0\),
\(\kappa>0\), \(\lambda\in\mathbb R\), and \(\gamma\ge0\).  Put
\[
  \theta=\min\{\gamma,2\}.
\]
Suppose
\[
  p_d(a):=\Prob\{\norm{T_d(\xi)}_{\op}<a\,\Exp\norm{T_d(\xi)}_{\op}\}
  \le C_0\bigl(\log(e+d)\bigr)^{-B}
  \qquad(d\ge1),
\]
where
\[
  \kappa(B-\theta)>1
  \qquad\text{or}\qquad
  \kappa(B-\theta)=1\ \text{ and }\ \lambda(B-\theta)>1 .
\]
Let \((d_m,n_m)_{m\ge1}\) satisfy, with
\[
  \ell_m=\log(e+d_m),\qquad L_m=\log(e+m),
\]
\[
  d_m\ge\exp(c_0m^\kappa L_m^\lambda),
  \qquad
  C_{A,K,a}\ell_m\le n_m\le \min\{d_m,C_1\ell_m^\gamma\}.
\]
At each \(m\), take independent full Toeplitz samples at dimension \(d_m\), and
independently sample \(M_m\) uniform sign vectors in \(\{-1,1\}^{n_m}\), with
\(M_m\ge1\).  Let \(u_m\ge0\) and \(1\le r_m\le M_m\) be deterministic
parameters.  If
\[
  \sum_{m=1}^{\infty}
  \left(e^{-u_m}+\binom{M_m}{r_m}d_m^{-A r_m}\right)<\infty,
\]
then, with probability one over the matrices and sampled signs, all but
finitely many \(m\) satisfy
\[
\begin{aligned}
  \max_{1\le s\le M_m}\max_{\omega^{d_m}=1}
  \left|\sum_{i=1}^{n_m}\delta_{m,i}^{(s)}
  \frac{P_{m,i}(\omega)}{\norm{B_{m,i}}_{\op}}\right|
  \le
  C_{A,K,a}\sqrt{
    n_m\left(1+\frac{\log(eM_m)+u_m}{\ell_m}\right)} .
\end{aligned}
\]
Moreover, at least \(M_m-r_m+1\) sampled signs have edge maximum at most
\(C_{A,K,a}\sqrt{n_m}\), and at least \(M_m-r_m+1\) sampled signs are centers
of Hamming balls of radius \(\lfloor c_{A,K,a}\ell_m\rfloor\) on which every
signing has edge maximum at most \(C_{A,K,a}\sqrt{n_m}\).

In particular, let \(s>1\), take \(u_m=s\log(e+m)\), and take \(r_m=1\).  If
\[
  M_m\le d_m^\beta
\]
for some \(0\le\beta<A\), then eventually every sampled sign in the batch is an
\(O(\sqrt{n_m})\)-good robust center and the whole batch satisfies an
\(O(\sqrt{n_m})\) sampled-list maximum bound.  More generally, if
\[
  M_m\le \exp(bm^\sigma L_m^\nu)
\]
for some \(b>0\), \(0<\sigma\le\kappa\), and \(\nu\in\mathbb R\), then the same
eventual all-sampled-sign conclusion holds whenever
\[
  \sigma<\kappa,
  \qquad\text{or}\qquad
  \sigma=\kappa,\ \nu<\lambda,
  \qquad\text{or}\qquad
  \sigma=\kappa,\ \nu=\lambda,\ b<Ac_0 .
\]
Finally, for any \(0<\rho\le1\), if \(\sum_m e^{-u_m}<\infty\) and
\(r_m=\lceil\rho M_m\rceil\), then no upper bound on \(M_m\) is needed to get,
eventually almost surely, a \(1-\rho\) fraction of sampled signs which are
logarithmically robust centers; only the displayed sampled-list maximum keeps
its entropy factor \(\log(eM_m)\).
\end{corollary}

\begin{proof}
As in the proof of
Corollary~\ref{cor:normalizer-tail-borel-window-log-corrected}, the hypotheses
give
\[
  \sum_m\left(d_m^{-A}
  +p_{d_m}(a)\min\{n_m,\ell_m^2\}\right)<\infty .
\]
Therefore Corollary~\ref{cor:normalizer-tail-borel-counts} applies once the two
displayed sign-tail summands are summable.  This proves the general sampled-list
maximum, good-sign count, and robust-center count assertions.

For the polynomial-in-\(d_m\) batch consequence,
\[
  \binom{M_m}{1}d_m^{-A}\le d_m^{\beta-A},
\]
which is summable because \(\beta<A\) and
\(d_m\ge\exp(c_0m^\kappa L_m^\lambda)\).  Also
\[
  \frac{\log(eM_m)+u_m}{\ell_m}\le C_{s,\beta,c_0,\kappa,\lambda}
\]
for all large \(m\), so the sampled-list maximum is \(O(\sqrt{n_m})\).

For \(M_m\le\exp(bm^\sigma L_m^\nu)\),
\[
  \binom{M_m}{1}d_m^{-A}
  \le
  \exp\{bm^\sigma L_m^\nu-Ac_0m^\kappa L_m^\lambda\}.
\]
The displayed alternatives make this last series summable.  The same comparison
shows \((\log(eM_m)+u_m)/\ell_m=O(1)\), so the sampled-list maximum again has
the \(O(\sqrt{n_m})\) form.

Finally, if \(r_m=\lceil\rho M_m\rceil\), then
\[
  \binom{M_m}{r_m}d_m^{-A r_m}
  \le
  \left(\frac{e}{\rho}d_m^{-A}\right)^{r_m}.
\]
For every fixed \(0<\tau_0<A\) and all large \(m\), this is at most
\[
  d_m^{-(A-\tau_0)r_m}\le d_m^{-(A-\tau_0)},
\]
which is summable on the log-corrected sparse schedule.  The sampled-count
theorem gives the stated \(1-\rho\) robust-center fraction.
\end{proof}

\begin{corollary}[Iterated-log sparse logarithmic-window sampled counts]
\label{cor:normalizer-tail-borel-window-iterlog-counts}
Fix \(A>0\), \(a>0\), \(C_0<\infty\), \(C_1<\infty\), \(c_0>0\),
\(\kappa>0\), \(\lambda,\mu\in\mathbb R\), and \(\gamma\ge0\).  Put
\[
  \theta=\min\{\gamma,2\}.
\]
Suppose
\[
  p_d(a):=\Prob\{\norm{T_d(\xi)}_{\op}<a\,\Exp\norm{T_d(\xi)}_{\op}\}
  \le C_0\bigl(\log(e+d)\bigr)^{-B}
  \qquad(d\ge1),
\]
and write \(q=B-\theta\).  Assume
\[
\begin{gathered}
  \kappa q>1,\\
  \text{or}\quad \kappa q=1\ \text{ and }\ \lambda q>1,\\
  \text{or}\quad \kappa q=1,\ \lambda q=1,\ \text{ and }\ \mu q>1 .
\end{gathered}
\]
Let \((d_m,n_m)_{m\ge1}\) satisfy, with
\[
  \ell_m=\log(e+d_m),\qquad
  L_m=\log(e+m),\qquad
  H_m=\log(e+L_m),
\]
\[
  d_m\ge\exp(c_0m^\kappa L_m^\lambda H_m^\mu),
  \qquad
  C_{A,K,a}\ell_m\le n_m\le \min\{d_m,C_1\ell_m^\gamma\}.
\]
At each \(m\), take independent full Toeplitz samples at dimension \(d_m\), and
independently sample \(M_m\) uniform sign vectors in \(\{-1,1\}^{n_m}\), with
\(M_m\ge1\).  Let \(u_m\ge0\) and \(1\le r_m\le M_m\) be deterministic
parameters.  If
\[
  \sum_{m=1}^{\infty}
  \left(e^{-u_m}+\binom{M_m}{r_m}d_m^{-A r_m}\right)<\infty,
\]
then, with probability one over the matrices and sampled signs, all but
finitely many \(m\) satisfy
\[
\begin{aligned}
  \max_{1\le s\le M_m}\max_{\omega^{d_m}=1}
  \left|\sum_{i=1}^{n_m}\delta_{m,i}^{(s)}
  \frac{P_{m,i}(\omega)}{\norm{B_{m,i}}_{\op}}\right|
  \le
  C_{A,K,a}\sqrt{
    n_m\left(1+\frac{\log(eM_m)+u_m}{\ell_m}\right)} .
\end{aligned}
\]
Moreover, at least \(M_m-r_m+1\) sampled signs have edge maximum at most
\(C_{A,K,a}\sqrt{n_m}\), and at least \(M_m-r_m+1\) sampled signs are centers
of Hamming balls of radius \(\lfloor c_{A,K,a}\ell_m\rfloor\) on which every
signing has edge maximum at most \(C_{A,K,a}\sqrt{n_m}\).

In particular, let \(s>1\), take \(u_m=s\log(e+m)\), and take \(r_m=1\).  If
\[
  M_m\le d_m^\beta
\]
for some \(0\le\beta<A\), then eventually every sampled sign in the batch is an
\(O(\sqrt{n_m})\)-good robust center and the whole batch satisfies an
\(O(\sqrt{n_m})\) sampled-list maximum bound.  More generally, if
\[
  M_m\le \exp(bm^\sigma L_m^\nu H_m^\eta)
\]
for some \(b>0\), \(0<\sigma\le\kappa\), and \(\nu,\eta\in\mathbb R\), then the
same eventual all-sampled-sign conclusion holds whenever
\[
  \sigma<\kappa,
  \qquad\text{or}\qquad
  \sigma=\kappa,\ \nu<\lambda,
  \qquad\text{or}\qquad
  \sigma=\kappa,\ \nu=\lambda,\ \eta<\mu,
\]
or
\[
  \sigma=\kappa,\qquad \nu=\lambda,\qquad \eta=\mu,\qquad b<Ac_0 .
\]
Finally, for any \(0<\rho\le1\), if \(\sum_m e^{-u_m}<\infty\) and
\(r_m=\lceil\rho M_m\rceil\), then no upper bound on \(M_m\) is needed to get,
eventually almost surely, a \(1-\rho\) fraction of sampled signs which are
logarithmically robust centers; only the displayed sampled-list maximum keeps
its entropy factor \(\log(eM_m)\).
\end{corollary}

\begin{proof}
The matrix summability is the only new point.  Since
\(\min\{n_m,\ell_m^2\}\le C_1'\ell_m^\theta\), the lower bound on \(d_m\) gives
\[
  p_{d_m}(a)\min\{n_m,\ell_m^2\}
  \le C
  m^{-\kappa q}L_m^{-\lambda q}H_m^{-\mu q}.
\]
The displayed alternatives are exactly the integral-test conditions making this
series summable.  Also \(d_m^{-A}\le \exp(-c m^\kappa L_m^\lambda H_m^\mu)\) is
summable.  Hence Corollary~\ref{cor:normalizer-tail-borel-counts} applies once
the two sign-tail summands are summable, proving the general sampled-list
maximum and count assertions.

If \(M_m\le d_m^\beta\), then
\[
  \binom{M_m}{1}d_m^{-A}\le d_m^{\beta-A},
\]
which is summable for \(\beta<A\), and
\((\log(eM_m)+u_m)/\ell_m=O(1)\).  If
\(M_m\le\exp(bm^\sigma L_m^\nu H_m^\eta)\), then
\[
  \binom{M_m}{1}d_m^{-A}
  \le
  \exp\{bm^\sigma L_m^\nu H_m^\eta
        -Ac_0m^\kappa L_m^\lambda H_m^\mu\}.
\]
The lexicographic alternatives make this series summable, and give the same
boundedness of \((\log(eM_m)+u_m)/\ell_m\).  Finally, with
\(r_m=\lceil\rho M_m\rceil\),
\[
  \binom{M_m}{r_m}d_m^{-A r_m}
  \le
  \left(\frac{e}{\rho}d_m^{-A}\right)^{r_m}
\]
for all large \(m\).  For every fixed \(0<\tau_0<A\), this is eventually at most
\[
  d_m^{-(A-\tau_0)},
\]
which is summable.  The sampled-count theorem gives the claimed robust-center
fraction.
\end{proof}

\begin{corollary}[Sampled logarithmic-window sparse transfer]
\label{cor:normalizer-tail-borel-window-sampled}
Fix \(A>0\), \(\alpha>0\), \(a>0\), \(C_0<\infty\), \(C_1<\infty\),
\(c_0>0\), \(\kappa>0\), and \(\gamma\ge0\).  Suppose
\[
  p_d(a):=\Prob\{\norm{T_d(\xi)}_{\op}<a\,\Exp\norm{T_d(\xi)}_{\op}\}
  \le C_0\bigl(\log(e+d)\bigr)^{-B}
  \qquad(d\ge1)
\]
for some
\[
  B>\min\{\gamma,2\}+\frac1\kappa .
\]
Let \((d_m,n_m)_{m\ge1}\) satisfy, with \(\ell_m=\log(e+d_m)\),
\[
  d_m\ge\exp(c_0m^\kappa),\qquad
  C_{A,K,a}\ell_m\le n_m\le \min\{d_m,C_1\ell_m^\gamma\}.
\]
For each \(m\), sample \(M_m\) independent uniform sign vectors in
\(\{-1,1\}^{n_m}\), independently of the matrices, with
\[
  M_m\ge\lceil\alpha\rceil .
\]
Then, with probability one over the matrices and sampled signs, all but
finitely many \(m\) have a sampled sign which is the center of a Hamming ball of
radius
\[
  r_m=\left\lfloor c_{A,K,a}\ell_m\right\rfloor
\]
on which every signing \(\eps'\) satisfies
\[
  \max_{\omega^{d_m}=1}\left|\sum_{i=1}^{n_m}
  \eps'_i\frac{P_{m,i}(\omega)}{\norm{B_{m,i}}_{\op}}\right|
  \le C_{A,K,a}\sqrt{n_m}.
\]
Thus the logarithmic-window threshold of
Corollary~\ref{cor:normalizer-tail-borel-window} also gives actual sampled
robust certificates; on stretched-exponential schedules no restriction beyond
\(\alpha>0\) is needed for the sampled-list confidence exponent.
\end{corollary}

\begin{proof}
Apply Corollary~\ref{cor:normalizer-tail-borel-poly-confidence}.  It is enough
to prove
\[
  \sum_m\left(d_m^{-A}
  +p_{d_m}(a)\min\{n_m,\ell_m^2\}
  +d_m^{-\alpha}\right)<\infty .
\]
The first and third terms are summable because
\(d_m\ge\exp(c_0m^\kappa)\).  The middle term is summable by the same estimate
used in Corollary~\ref{cor:normalizer-tail-borel-window}:
\[
  p_{d_m}(a)\min\{n_m,\ell_m^2\}
  \le C\,m^{\kappa(\min\{\gamma,2\}-B)},
\]
and the exponent is smaller than \(-1\) by the displayed hypothesis on \(B\).
The sampled Borel--Cantelli transfer gives the conclusion.
\end{proof}

\begin{corollary}[Sparse logarithmic-window robust fractions]
\label{cor:normalizer-tail-borel-window-fractions}
Fix \(A>0\), \(S>0\), \(a>0\), \(C_0<\infty\), \(C_1<\infty\),
\(c_0>0\), \(\kappa>0\), \(\gamma\ge0\), and a fraction \(0<\rho\le1\).
Assume
\[
  A\ge 2\log_2(e/\rho),
\]
and choose an integer \(M\ge1\) so that
\[
  M\ge \frac{2(S+1)}{A\rho}.
\]
Suppose
\[
  p_d(a):=\Prob\{\norm{T_d(\xi)}_{\op}<a\,\Exp\norm{T_d(\xi)}_{\op}\}
  \le C_0\bigl(\log(e+d)\bigr)^{-B}
  \qquad(d\ge1)
\]
for some
\[
  B>\min\{\gamma,2\}+\frac1\kappa .
\]
Let \((d_m,n_m)_{m\ge1}\) satisfy, with \(\ell_m=\log(e+d_m)\),
\[
  d_m\ge\exp(c_0m^\kappa),\qquad
  C_{A,K,a}\ell_m\le n_m\le \min\{d_m,C_1\ell_m^\gamma\}.
\]
At each \(m\), take independent full Toeplitz samples at dimension \(d_m\), and
independently sample \(M\) uniform sign vectors in \(\{-1,1\}^{n_m}\).
Then, with probability one over the matrices and sampled signs, all but
finitely many \(m\) satisfy both conclusions below.

First, the whole sampled list obeys
\[
  \max_{1\le s\le M}\max_{\omega^{d_m}=1}
  \left|\sum_{i=1}^{n_m}\delta_{m,i}^{(s)}
  \frac{P_{m,i}(\omega)}{\norm{B_{m,i}}_{\op}}\right|
  \le C_{A,K,a,S,\rho}\sqrt{n_m}.
\]
Second, at least a \(1-\rho\) fraction of the \(M\) sampled signs are centers of
Hamming balls of radius
\[
  \left\lfloor c_{A,K,a}\ell_m\right\rfloor
\]
on which every signing \(\eps'\) has edge maximum at most
\[
  C_{A,K,a}\sqrt{n_m}.
\]
The same \(1-\rho\) fraction also has edge maximum at most
\(C_{A,K,a}\sqrt{n_m}\).  Thus, on stretched-exponential logarithmic-window
schedules, the robust-fraction confidence exponent has no lower summability
restriction beyond \(S>0\).
\end{corollary}

\begin{proof}
Apply Corollary~\ref{cor:normalizer-tail-borel-fractions}.  It is enough to
check
\[
  \sum_m\left(d_m^{-A}
  +p_{d_m}(a)\min\{n_m,\ell_m^2\}
  +d_m^{-S}\right)<\infty .
\]
The stretched-exponential lower bound on \(d_m\) makes both
\(\sum_m d_m^{-A}\) and \(\sum_m d_m^{-S}\) finite, for every positive \(A\)
and \(S\).  The middle term is controlled by the same estimate used in
Corollary~\ref{cor:normalizer-tail-borel-window}:
\[
  p_{d_m}(a)\min\{n_m,\ell_m^2\}
  \le C\,m^{\kappa(\min\{\gamma,2\}-B)}.
\]
The displayed hypothesis on \(B\) makes the exponent smaller than \(-1\), so
this series is finite.  The general robust-fraction Borel theorem gives the
eventual sampled-list maximum and the \(1-\rho\) robust-center fraction.
\end{proof}

\begin{corollary}[Sparse logarithmic-window sampled counts]
\label{cor:normalizer-tail-borel-window-counts}
Fix \(A>0\), \(a>0\), \(C_0<\infty\), \(C_1<\infty\), \(c_0>0\),
\(\kappa>0\), and \(\gamma\ge0\).  Suppose
\[
  p_d(a):=\Prob\{\norm{T_d(\xi)}_{\op}<a\,\Exp\norm{T_d(\xi)}_{\op}\}
  \le C_0\bigl(\log(e+d)\bigr)^{-B}
  \qquad(d\ge1)
\]
for some
\[
  B>\min\{\gamma,2\}+\frac1\kappa .
\]
Let \((d_m,n_m)_{m\ge1}\) satisfy, with \(\ell_m=\log(e+d_m)\),
\[
  d_m\ge\exp(c_0m^\kappa),\qquad
  C_{A,K,a}\ell_m\le n_m\le \min\{d_m,C_1\ell_m^\gamma\}.
\]
At each \(m\), take independent full Toeplitz samples at dimension \(d_m\), and
independently sample \(M_m\) uniform sign vectors in \(\{-1,1\}^{n_m}\), with
\(M_m\ge1\).  Let \(u_m\ge0\) and \(1\le r_m\le M_m\) be deterministic
parameters.  If
\[
  \sum_{m=1}^{\infty}
  \left(e^{-u_m}+\binom{M_m}{r_m}d_m^{-A r_m}\right)<\infty,
\]
then, with probability one over the matrices and sampled signs, all but
finitely many \(m\) satisfy the sampled-list maximum bound, good-sign count, and
robust-center count conclusions of
Corollary~\ref{cor:normalizer-tail-borel-counts}.

In particular, let \(s>1\), take \(u_m=s\log(e+m)\) and \(r_m=1\), and suppose
\[
  M_m\le d_m^\beta
  \qquad(m\ge1)
\]
for some \(0\le\beta<A\).  Then eventually every sampled sign in the batch is
an \(O(\sqrt{n_m})\)-good robust center and the whole batch satisfies the
displayed sampled-list maximum bound.  More generally, if
\[
  M_m\le \exp(bm^\sigma)
\]
for some \(b>0\), \(0<\sigma\le\kappa\), then the same eventual all-sampled-sign
conclusion holds whenever either \(\sigma<\kappa\), or \(\sigma=\kappa\) and
\(b<Ac_0\).
\end{corollary}

\begin{proof}
Apply Corollary~\ref{cor:normalizer-tail-borel-counts}.  The matrix summand
\(\sum_m d_m^{-A}\) is finite because \(d_m\ge\exp(c_0m^\kappa)\).  The
normalizer-tail summand is finite by the estimate from
Corollary~\ref{cor:normalizer-tail-borel-window}:
\[
  p_{d_m}(a)\min\{n_m,\ell_m^2\}
  \le C\,m^{\kappa(\min\{\gamma,2\}-B)},
\]
and the exponent is less than \(-1\).  The two sign-tail summands are finite by
the displayed hypothesis, so the sampled-count Borel theorem applies.

For the polynomial-in-\(d_m\) batch consequence,
\[
  \binom{M_m}{1}d_m^{-A}\le d_m^{\beta-A},
\]
which is summable on the stretched-exponential schedule when \(\beta<A\).  Also
\[
  \frac{\log(eM_m)+u_m}{\ell_m}\le C_{s,\beta,c_0,\kappa}
\]
for all large \(m\), so the sampled-list maximum is \(O(\sqrt{n_m})\).  For the
\(\exp(bm^\sigma)\) batch consequence,
\[
  \binom{M_m}{1}d_m^{-A}
  \le \exp\left(bm^\sigma-Ac_0m^\kappa\right),
\]
which is summable if \(\sigma<\kappa\), or if \(\sigma=\kappa\) and
\(b<Ac_0\).  The same comparison with \(\ell_m\) gives the \(O(\sqrt{n_m})\)
sampled-list maximum.
\end{proof}

\begin{corollary}[Sparse logarithmic-window growing robust fractions]
\label{cor:normalizer-tail-borel-window-growing-fractions}
Fix \(A>0\), \(a>0\), \(C_0<\infty\), \(C_1<\infty\), \(c_0>0\),
\(\kappa>0\), \(\gamma\ge0\), and \(0<\rho\le1\).  Suppose
\[
  p_d(a):=\Prob\{\norm{T_d(\xi)}_{\op}<a\,\Exp\norm{T_d(\xi)}_{\op}\}
  \le C_0\bigl(\log(e+d)\bigr)^{-B}
  \qquad(d\ge1)
\]
for some
\[
  B>\min\{\gamma,2\}+\frac1\kappa .
\]
Let \((d_m,n_m)_{m\ge1}\) satisfy, with \(\ell_m=\log(e+d_m)\),
\[
  d_m\ge\exp(c_0m^\kappa),\qquad
  C_{A,K,a}\ell_m\le n_m\le \min\{d_m,C_1\ell_m^\gamma\}.
\]
At each \(m\), let \(B_{m,i}=T_{d_m}(\xi^{(m,i)})\), \(1\le i\le n_m\), be
independent full Toeplitz matrices with independent mean-zero, unit-variance lag
variables of sub-Gaussian norm at most \(K\), and let \(P_{m,i}\) be their
symbols.  Independently sample an arbitrary deterministic number \(M_m\ge1\) of
uniform sign vectors
\[
  \delta_m^{(1)},\ldots,\delta_m^{(M_m)}\in\{-1,1\}^{n_m}.
\]
Let \(u_m\ge0\) be deterministic and assume
\[
  \sum_{m=1}^{\infty}e^{-u_m}<\infty .
\]
Then, with probability one over the matrices and sampled signs, all but
finitely many \(m\) satisfy both conclusions below.

First, the whole sampled list obeys
\[
\begin{aligned}
  \max_{1\le s\le M_m}\max_{\omega^{d_m}=1}
  \left|\sum_{i=1}^{n_m}\delta_{m,i}^{(s)}
  \frac{P_{m,i}(\omega)}{\norm{B_{m,i}}_{\op}}\right|
  \le
  C_{A,K,a}\sqrt{
    n_m\left(1+\frac{\log(eM_m)+u_m}{\ell_m}\right)} .
\end{aligned}
\]
Second, at least a \(1-\rho\) fraction of the \(M_m\) sampled signs have edge
maximum at most \(C_{A,K,a}\sqrt{n_m}\) and are centers of Hamming balls of
radius
\[
  \left\lfloor c_{A,K,a}\ell_m\right\rfloor
\]
on which every signing \(\eps'\) has edge maximum at most
\[
  C_{A,K,a}\sqrt{n_m}.
\]
No upper bound on \(M_m\) is required for this robust-fraction conclusion.
\end{corollary}

\begin{proof}
Apply Corollary~\ref{cor:normalizer-tail-borel-window-counts} with
\[
  r_m=\lceil\rho M_m\rceil .
\]
The only additional summability condition to check is the binomial count tail.
Since \(r_m\ge\rho M_m\),
\[
  \binom{M_m}{r_m}d_m^{-A r_m}
  \le \left(\frac{eM_m}{r_m}\right)^{r_m}d_m^{-A r_m}
  \le \left(\frac{e}{\rho}d_m^{-A}\right)^{r_m}.
\]
For every fixed \(0<\tau_0<A\) and all large \(m\), the last display is at most
\[
  d_m^{-(A-\tau_0)r_m}\le d_m^{-(A-\tau_0)},
\]
which is summable on the stretched-exponential schedule.
The sampled-list tail is summable by the hypothesis on \(u_m\).  The sampled
count theorem therefore gives the entropy-sensitive sampled-list maximum and at
least
\[
  M_m-r_m+1\ge (1-\rho)M_m
\]
robust centers eventually.
\end{proof}

\begin{corollary}[Tunable bounded-density Chernoff transfer]
\label{cor:normalizer-tail-chernoff-density}
Fix \(A>0\), \(a>0\), and constants \(R\ge1\) and \(0<\theta<R\).  Write
\(\ell=\log(e+d)\).  In the setting of
Corollary~\ref{cor:normalizer-tail-quant}, suppose
\[
  p_d(a):=\Prob\{\norm{T_d(\xi)}_{\op}<aN_d\}\le \frac{\theta}{\ell^2}.
\]
Set
\[
  I(R,\theta):=R\log\frac R\theta-R+\theta>0 .
\]
Then there is a constant \(C_{A,K,a,R}\) such that, whenever
\[
  n\le d,\qquad n\ge C_{A,K,a,R}\ell ,
\]
one has
\[
  \Prob\left\{
  \max_{\omega^d=1}
  \sum_{i=1}^n
  \frac{\abs{P_i(\omega)}^2}{\norm{B_i}_{\op}^2}
  > C_{A,K,a,R}\frac n\ell
  \right\}
  \le C_{A,K,a,R}d^{-A}
     +\exp\left(-I(R,\theta)\frac n{\ell^2}\right).
\]
On the complement of this exceptional event, independent signs have the
\(C_{A,K,a,R}\sqrt n\) edge bound with conditional probability at least
\(1-d^{-A}\), sign-averaged edge maximum \(C_{A,K,a,R}\sqrt n\), a
deterministic signing certificate, and scalar edge-grid moments bounded by
\((C_{A,K,a,R}k\,n/\ell)^k\) for every \(2k\)-th moment.
\end{corollary}

\begin{proof}
Let
\[
  M_a:=\#\{1\le i\le n:\norm{B_i}_{\op}<aN_d\}.
\]
Corollary~\ref{cor:normalizer-tail-empirical} gives the full profile and
signing package on \(\{M_a\le Rn/\ell^2\}\), outside a
\(C_{A,K,a,R}d^{-A}\) event.  It remains only to bound the bad-count excess.
The random variable \(M_a\) is binomially dominated by a binomial variable with
success probability \(p_d(a)\).  If \(R/\ell^2\ge1\), then
\(\{M_a>Rn/\ell^2\}\) is empty.  Otherwise \(\theta/\ell^2<R/\ell^2<1\), and
monotonicity in the success probability plus Chernoff's bound for
\(X\sim\operatorname{Bin}(n,\theta/\ell^2)\) gives
\[
  \Prob\{M_a>Rn/\ell^2\}
  \le \Prob\{X>Rn/\ell^2\}
  \le
  \exp\left(-\frac n{\ell^2}
  \left(R\log\frac R\theta-R+\theta\right)\right).
\]
Combining this with the empirical criterion proves the claim.
\end{proof}

\begin{corollary}[Almost-sure Chernoff profile transfer]
\label{cor:normalizer-tail-chernoff-borel-density}
Fix \(A>0\), \(a>0\), and constants \(R\ge1\) and \(0<\theta<R\).  There is a
constant \(C_{A,K,a,R}\) with the following property.  Let \((d_m,n_m)_{m\ge1}\)
be a triangular-array sequence with \(n_m\le d_m\), put
\(\ell_m=\log(e+d_m)\), and assume
\[
  n_m\ge C_{A,K,a,R}\ell_m \qquad(m\ge1).
\]
At each \(m\), let \(B_{m,i}=T_{d_m}(\xi^{(m,i)})\), \(1\le i\le n_m\), be
independent full Toeplitz matrices whose lag variables are independent, mean
zero, unit variance, and have sub-Gaussian norm at most \(K\), and let
\(P_{m,i}\) be their symbols.  Suppose
\[
  p_{d_m}(a):=
  \Prob\{\norm{T_{d_m}(\xi)}_{\op}
      <a\,\Exp\norm{T_{d_m}(\xi)}_{\op}\}
  \le \frac{\theta}{\ell_m^2}
  \qquad(m\ge1),
\]
and set \(I(R,\theta)=R\log(R/\theta)-R+\theta\).  If
\[
  \sum_{m=1}^\infty
  \left(d_m^{-A}+
  \exp\left(-I(R,\theta)\frac{n_m}{\ell_m^2}\right)\right)<\infty,
\]
then, with probability one over the matrix triangular array, all but finitely
many \(m\) satisfy
\[
  \max_{\omega^{d_m}=1}
  \sum_{i=1}^{n_m}
  \frac{\abs{P_{m,i}(\omega)}^2}{\norm{B_{m,i}}_{\op}^2}
  \le C_{A,K,a,R}\frac{n_m}{\ell_m}.
\]
For those \(m\), the independent-random-sign tail, sign-averaged edge maximum,
deterministic signing certificate, and scalar edge-grid moment conclusions of
Corollary~\ref{cor:normalizer-tail-chernoff-density} hold with
\(d,n,\ell\) replaced by \(d_m,n_m,\ell_m\).

In particular, if \(L\ge1\), \(\eta>0\), and
\(p_{d_m}(a)\le L/\ell_m^2\), then the same conclusion holds under
\[
  \sum_{m=1}^\infty
  \left(d_m^{-A}+
  \exp\left(-Lh(\eta)\frac{n_m}{\ell_m^2}\right)\right)<\infty,
  \qquad h(\eta)=(1+\eta)\log(1+\eta)-\eta .
\]
Consequently, along the full integer dimension sequence, if \(A>1\),
\[
  p_d(a)\le \frac{L}{\log^2(e+d)}
  \qquad(d\ge1),
\]
and if \(n(d)\le d\) while, for all sufficiently large \(d\),
\[
  n(d)\ge q\,\log^3(e+d)
  \qquad\text{for some } q>\frac1{Lh(\eta)},
\]
then almost surely all sufficiently large integer dimensions satisfy the same
edge-profile, signing, and scalar-moment package.  For \(\eta=1\), the threshold
is \(q>1/(L(2\log2-1))\).
\end{corollary}

\begin{proof}
Apply Corollary~\ref{cor:normalizer-tail-chernoff-density} at each \(m\), with
the constant \(C_{A,K,a,R}\) enlarged if necessary.  The displayed summability
condition makes the matrix exceptional probabilities summable.  The first
Borel--Cantelli lemma therefore gives an almost-sure event on which all but
finitely many \(m\) satisfy the finite Chernoff profile and signing package.

The relative-excess version is obtained by taking \(R=(1+\eta)L\) and
\(\theta=L\), so \(I((1+\eta)L,L)=Lh(\eta)\).  For the all-dimensions
consequence take \(d_m=m\).  Then \(\sum_m m^{-A}<\infty\), and for all large
\(m\),
\[
  \exp\left(-Lh(\eta)\frac{n(m)}{\log^2(e+m)}\right)
  \le (e+m)^{-Lh(\eta)q},
\]
which is summable because \(Lh(\eta)q>1\).  The finitely many initial dimensions
are irrelevant to the eventual almost-sure conclusion.
\end{proof}

\begin{corollary}[Tunable Chernoff search-stability transfer]
\label{cor:normalizer-tail-chernoff-search}
Fix \(A>0\), \(a>0\), and constants \(R\ge1\) and \(0<\theta<R\).  Write
\(\ell=\log(e+d)\).  In the setting of
Corollary~\ref{cor:normalizer-tail-chernoff-density}, suppose
\[
  p_d(a):=\Prob\{\norm{T_d(\xi)}_{\op}<aN_d\}\le \frac{\theta}{\ell^2}.
\]
Set \(I(R,\theta)=R\log(R/\theta)-R+\theta\).  Then there is a constant
\(C_{A,K,a,R}\) such that, whenever
\[
  n\le d,\qquad n\ge C_{A,K,a,R}\ell ,
\]
outside a matrix event of probability at most
\[
  C_{A,K,a,R}d^{-A}
  +\exp\left(-I(R,\theta)\frac n{\ell^2}\right),
\]
the full random-search and stable-center package of
Corollary~\ref{cor:normalizer-tail-search-stability} holds, with
\(C_{A,K,a}\) and \(c_{A,K,a}\) replaced by constants depending on
\((A,K,a,R)\).

In particular, if \(p_d(a)\le L/\ell^2\) with \(L\ge1\), then for every
\(\eta>0\) the same package holds outside probability
\[
  C_{A,K,a,L,\eta}d^{-A}
  +\exp\left(-Lh(\eta)\frac n{\ell^2}\right),
  \qquad h(\eta)=(1+\eta)\log(1+\eta)-\eta .
\]
Taking \(\eta=1\), there are constants \(C_{A,K,a,L},c_L>0\) such that the
package holds outside probability
\[
  C_{A,K,a,L}d^{-A}+\exp\left(-c_L\frac n{\ell^2}\right).
\]
\end{corollary}

\begin{proof}
Let \(M_a=\#\{i:\norm{B_i}_{\op}<aN_d\}\).  On the event
\(\{M_a\le Rn/\ell^2\}\), Corollary~\ref{cor:normalizer-tail-empirical} gives the
edge profile
\[
  \max_{\omega^d=1}\sum_{i=1}^n
  \frac{|P_i(\omega)|^2}{\norm{B_i}_{\op}^2}
  \le C_{A,K,a,R}\frac n\ell
\]
outside a \(C_{A,K,a,R}d^{-A}\) matrix event.  Conditioning on any realization
with this profile, Corollaries~\ref{cor:edge-profile-random-search} and
\ref{cor:edge-profile-hamming} give the sampled-list, success-count, Hamming-ball,
and sampled robust-center statements exactly as in
Corollary~\ref{cor:normalizer-tail-search-stability}, with constants depending
on \(R\).

The same Chernoff calculation as in the proof of
Corollary~\ref{cor:normalizer-tail-chernoff-density} gives
\[
  \Prob\{M_a>Rn/\ell^2\}
  \le \exp\left(-I(R,\theta)\frac n{\ell^2}\right).
\]
This proves the first assertion.  The relative-excess specialization is obtained
by taking \(R=(1+\eta)L\) and \(\theta=L\), so
\(I((1+\eta)L,L)=Lh(\eta)\).  The final display is the case \(\eta=1\), with
\(c_L=L(2\log2-1)\), after absorbing constants.
\end{proof}

\begin{corollary}[Chernoff joint sampled-search package]
\label{cor:normalizer-tail-chernoff-joint-search}
Fix \(A>0\), \(B>0\), \(a>0\), and constants \(R\ge1\) and \(0<\theta<R\).
Write \(\ell=\log(e+d)\).  In the setting of
Corollary~\ref{cor:normalizer-tail-chernoff-search}, suppose
\[
  p_d(a):=\Prob\{\norm{T_d(\xi)}_{\op}<aN_d\}\le \frac{\theta}{\ell^2}.
\]
Let \(\delta^{(1)},\ldots,\delta^{(M)}\) be independent uniform sign vectors,
independent of the matrices, where \(M\ge1\).  Put
\[
  v_i(\omega):=\frac{P_i(\omega)}{\norm{B_i}_{\op}},
  \qquad \omega^d=1,
\]
and set \(I(R,\theta)=R\log(R/\theta)-R+\theta\).  There are constants
\(C_{A,B,K,a,R}\) and \(c_{A,B,K,a,R}>0\) such that, whenever
\[
  n\le d,\qquad n\ge C_{A,B,K,a,R}\ell ,
\]
the following joint matrix-and-sign estimates hold.

First, for every \(u\ge0\),
\[
\begin{aligned}
&\Prob\Bigg\{
  \max_{1\le s\le M}\max_{\omega^d=1}
  \left|\sum_{i=1}^n\delta_i^{(s)}v_i(\omega)\right|
\\
&\hspace{35mm}>
  C_{A,B,K,a,R}\sqrt{
    n\left(1+\frac{\log(eM)+u}{\ell}\right)}
  \Bigg\}
\\
&\qquad\le
  C_{A,B,K,a,R}d^{-A}
  +\exp\left(-I(R,\theta)\frac n{\ell^2}\right)
  +e^{-u}.
\end{aligned}
\]

Second, let
\[
  G_M:=\#\left\{1\le s\le M:
  \max_{\omega^d=1}\left|\sum_{i=1}^n
  \delta_i^{(s)}v_i(\omega)\right|
  \le C_{A,B,K,a,R}\sqrt n\right\}.
\]
Then, for every integer \(1\le r\le M\),
\[
  \Prob\{G_M\le M-r\}
  \le
  C_{A,B,K,a,R}d^{-A}
  +\exp\left(-I(R,\theta)\frac n{\ell^2}\right)
  +\binom Mr d^{-Br}.
\]

Finally, call a sampled sign a robust center if it is the center of a Hamming
ball of radius
\[
  r_0=\left\lfloor c_{A,B,K,a,R}\ell\right\rfloor
\]
on which every sign vector \(\eps'\) satisfies
\[
  \max_{\omega^d=1}\left|\sum_{i=1}^n
  \eps'_i v_i(\omega)\right|
  \le C_{A,B,K,a,R}\sqrt n .
\]
If \(H_M\) is the number of robust centers among
\(\delta^{(1)},\ldots,\delta^{(M)}\), then for every integer \(1\le r\le M\),
\[
  \Prob\{H_M\le M-r\}
  \le
  C_{A,B,K,a,R}d^{-A}
  +\exp\left(-I(R,\theta)\frac n{\ell^2}\right)
  +\binom Mr d^{-Br}.
\]

In particular, if \(p_d(a)\le L/\ell^2\) with \(L\ge1\), then for every
\(\eta>0\) all three estimates remain true with the Chernoff term replaced by
\[
  \exp\left(-Lh(\eta)\frac n{\ell^2}\right),
  \qquad h(\eta)=(1+\eta)\log(1+\eta)-\eta ,
\]
after allowing the constants to depend on \(L\) and \(\eta\).
\end{corollary}

\begin{proof}
On the good matrix event in the proof of
Corollary~\ref{cor:normalizer-tail-chernoff-search},
\[
  \max_{\omega^d=1}\sum_{i=1}^n
  |v_i(\omega)|^2
  \le C_{A,K,a,R}\frac n\ell .
\]
That good event fails with probability at most
\[
  C_{A,K,a,R}d^{-A}
  +\exp\left(-I(R,\theta)\frac n{\ell^2}\right).
\]
Condition on any matrix realization satisfying the profile.  The sampled-list
maximum follows from Corollary~\ref{cor:edge-profile-random-search}, and the
good-sign count follows from the same corollary with exponent \(B\).  The
robust-center count follows from Corollary~\ref{cor:edge-profile-hamming}, again
with exponent \(B\).  Absorbing the profile constant into the displayed
constants gives the three conditional sign estimates; adding the matrix
exceptional probability gives the joint estimates.  The relative-excess
specialization is \(R=(1+\eta)L\) and \(\theta=L\).
\end{proof}

\begin{corollary}[Chernoff sampled-list abundance]
\label{cor:normalizer-tail-chernoff-abundance}
Fix \(A>0\), \(B>0\), \(a>0\), constants \(R\ge1\) and \(0<\theta<R\), and a
fraction \(0<\rho\le1\).  Write \(\ell=\log(e+d)\).  In the setting of
Corollary~\ref{cor:normalizer-tail-chernoff-search}, suppose
\[
  p_d(a):=\Prob\{\norm{T_d(\xi)}_{\op}<aN_d\}\le \frac{\theta}{\ell^2}.
\]
Let \(\delta^{(1)},\ldots,\delta^{(M)}\) be independent uniform sign vectors,
independent of the matrices, where \(M\ge1\), and put
\[
  r_\rho=\lceil \rho M\rceil,\qquad
  I(R,\theta)=R\log(R/\theta)-R+\theta .
\]
There are constants \(C_{A,B,K,a,R}\) and \(c_{A,B,K,a,R}>0\) such that,
whenever
\[
  n\le d,\qquad n\ge C_{A,B,K,a,R}\ell ,
\]
then, for every \(u\ge0\), with probability at least
\[
  1-
  C_{A,B,K,a,R}\left[
    d^{-A}+\exp\left(-I(R,\theta)\frac n{\ell^2}\right)\right]
  -e^{-u}-2\binom{M}{r_\rho}d^{-B r_\rho},
\]
all three conclusions below hold simultaneously.

First, the whole sampled list obeys
\[
\begin{aligned}
  \max_{1\le s\le M}\max_{\omega^d=1}
  \left|\sum_{i=1}^n\delta_i^{(s)}
  \frac{P_i(\omega)}{\norm{B_i}_{\op}}\right|
  \le
  C_{A,B,K,a,R}\sqrt{
    n\left(1+\frac{\log(eM)+u}{\ell}\right)} .
\end{aligned}
\]
Second, at least \(M-r_\rho+1\) sampled signs have edge maximum at most
\(C_{A,B,K,a,R}\sqrt n\).  Third, at least \(M-r_\rho+1\) sampled signs are
centers of Hamming balls of radius
\[
  \left\lfloor c_{A,B,K,a,R}\ell\right\rfloor
\]
on which every signing \(\eps'\) has edge maximum at most
\(C_{A,B,K,a,R}\sqrt n\).  In particular, at least a \(1-\rho\) fraction of the
sampled signs are good, and at least a \(1-\rho\) fraction are robust centers.

If \(p_d(a)\le L/\ell^2\) with \(L\ge1\), then for every \(\eta>0\) the same
conclusion holds with the Chernoff term replaced by
\[
  \exp\left(-Lh(\eta)\frac n{\ell^2}\right),
  \qquad h(\eta)=(1+\eta)\log(1+\eta)-\eta ,
\]
after allowing the constants to depend on \(L\) and \(\eta\).
\end{corollary}

\begin{proof}
Apply Corollary~\ref{cor:normalizer-tail-chernoff-joint-search} with
\[
  (M,r)=(M,r_\rho).
\]
The sampled-list maximum fails with probability at most the matrix exception
plus \(e^{-u}\).  The good-sign and robust-center count conclusions each fail
with probability at most the same matrix exception plus
\(\binom{M}{r_\rho}d^{-B r_\rho}\).  A union bound gives the displayed
simultaneous probability after absorbing the three copies of the matrix
exception into \(C_{A,B,K,a,R}\).  Since
\[
  M-r_\rho+1\ge M-\rho M,
\]
the count lower bounds imply the stated \(1-\rho\) fraction conclusions.  The
relative-excess form again takes \(R=(1+\eta)L\) and \(\theta=L\).
\end{proof}

\begin{corollary}[Chernoff polynomial-confidence robust fractions]
\label{cor:normalizer-tail-chernoff-robust-fraction-confidence}
Fix \(A>0\), \(S>0\), \(a>0\), constants \(R\ge1\) and \(0<\theta<R\), and a
fraction \(0<\rho\le1\).  Choose \(B>0\) and an integer \(M\ge1\) so that
\[
  B\ge 2\log_2(e/\rho),
  \qquad
  M\ge \frac{2(S+1)}{B\rho}.
\]
Write \(\ell=\log(e+d)\).  In the setting of
Corollary~\ref{cor:normalizer-tail-chernoff-search}, suppose \(d\ge2\) and
\[
  p_d(a):=\Prob\{\norm{T_d(\xi)}_{\op}<aN_d\}\le \frac{\theta}{\ell^2}.
\]
Let \(\delta^{(1)},\ldots,\delta^{(M)}\) be independent uniform sign vectors,
independent of the matrices.  Set \(I(R,\theta)=R\log(R/\theta)-R+\theta\).
There are constants \(C_{A,B,K,a,R}\) and \(c_{A,B,K,a,R}>0\) such that,
whenever
\[
  n\le d,\qquad n\ge C_{A,B,K,a,R}\ell ,
\]
then, for every \(u\ge0\), with probability at least
\[
  1-
  C_{A,B,K,a,R}\left[
    d^{-A}+\exp\left(-I(R,\theta)\frac n{\ell^2}\right)\right]
  -e^{-u}-d^{-S},
\]
the whole sampled list obeys the entropy-sensitive maximum bound from
Corollary~\ref{cor:normalizer-tail-chernoff-abundance}, and at least a
\(1-\rho\) fraction of the sampled signs are centers of radius
\[
  \left\lfloor c_{A,B,K,a,R}\ell\right\rfloor
\]
Hamming balls on which every signing \(\eps'\) has edge maximum at most
\[
  C_{A,B,K,a,R}\sqrt n .
\]
In particular, taking \(u=S\ell\), the same event has probability at least
\[
  1-
  C_{A,B,K,a,R}\left[
    d^{-A}+\exp\left(-I(R,\theta)\frac n{\ell^2}\right)\right]
  -2d^{-S},
\]
and the whole sampled list has edge maximum at most
\(C_{A,B,K,a,R,S,\rho}\sqrt n\).

If \(p_d(a)\le L/\ell^2\) with \(L\ge1\), then for every \(\eta>0\) the same
conclusions hold with the Chernoff term replaced by
\[
  \exp\left(-Lh(\eta)\frac n{\ell^2}\right),
  \qquad h(\eta)=(1+\eta)\log(1+\eta)-\eta ,
\]
after allowing the constants to depend on \(L\) and \(\eta\).
\end{corollary}

\begin{proof}
Let \(r_\rho=\lceil\rho M\rceil\).  The only point not already contained in
Corollary~\ref{cor:normalizer-tail-chernoff-abundance} is the simplification of
the binomial fraction tail.  Since \(r_\rho\ge\rho M\),
\[
  \binom{M}{r_\rho}\le \left(\frac{eM}{r_\rho}\right)^{r_\rho}
  \le (e/\rho)^{r_\rho}.
\]
For \(d\ge2\) and \(B\ge2\log_2(e/\rho)\), this gives
\[
  \binom{M}{r_\rho}d^{-B r_\rho}
  \le d^{-B r_\rho/2}.
\]
The lower bound on \(M\) implies \(B r_\rho/2\ge S+1\), hence
\[
  2\binom{M}{r_\rho}d^{-B r_\rho}\le 2d^{-(S+1)}\le d^{-S}.
\]
Substituting this into
Corollary~\ref{cor:normalizer-tail-chernoff-abundance} proves the first
assertion.  If \(u=S\ell\), then \(e^{-u}=(e+d)^{-S}\le d^{-S}\), while
\[
  1+\frac{\log(eM)+u}{\ell}
  \le C_{S,\rho,B}
\]
because \(M\) is fixed once \((S,\rho,B)\) are fixed.  This gives the
constant-scale sampled-list bound.  The relative-excess statement is the same
specialization \(R=(1+\eta)L\), \(\theta=L\).
\end{proof}

\begin{corollary}[Chernoff polynomial-batch robust search]
\label{cor:normalizer-tail-chernoff-poly-batch}
Fix \(A>0\), \(a>0\), constants \(R\ge1\) and \(0<\theta<R\), and parameters
\(\beta\ge0\) and \(\sigma>0\).  Write \(\ell=\log(e+d)\).  In the setting of
Corollary~\ref{cor:normalizer-tail-chernoff-search}, suppose
\[
  p_d(a):=\Prob\{\norm{T_d(\xi)}_{\op}<aN_d\}\le \frac{\theta}{\ell^2}.
\]
Let \(\delta^{(1)},\ldots,\delta^{(M)}\) be independent uniform sign vectors,
independent of the matrices, with
\[
  1\le M\le d^\beta .
\]
Set \(I(R,\theta)=R\log(R/\theta)-R+\theta\).  There are constants
\(C_{A,K,a,R,\beta,\sigma}\) and \(c_{A,K,a,R,\beta,\sigma}>0\) such that,
whenever
\[
  n\le d,\qquad n\ge C_{A,K,a,R,\beta,\sigma}\ell ,
\]
with probability at least
\[
  1-C_{A,K,a,R,\beta,\sigma}d^{-A}
  -\exp\left(-I(R,\theta)\frac n{\ell^2}\right)
  -d^{-\sigma},
\]
every sampled sign \(\delta^{(s)}\), \(1\le s\le M\), is the center of a
Hamming ball of radius
\[
  r_0=\left\lfloor c_{A,K,a,R,\beta,\sigma}\ell\right\rfloor
\]
on which every sign vector \(\eps'\) satisfies
\[
  \max_{\omega^d=1}\left|\sum_{i=1}^n
  \eps'_i\frac{P_i(\omega)}{\norm{B_i}_{\op}}\right|
  \le C_{A,K,a,R,\beta,\sigma}\sqrt n .
\]
In particular, the whole sampled batch itself has edge maximum at most
\(C_{A,K,a,R,\beta,\sigma}\sqrt n\).

If \(p_d(a)\le L/\ell^2\) with \(L\ge1\), then for every \(\eta>0\) the same
conclusion holds with the displayed exceptional probability replaced by
\[
  C_{A,K,a,L,\eta,\beta,\sigma}d^{-A}
  +\exp\left(-Lh(\eta)\frac n{\ell^2}\right)
  +d^{-\sigma},
  \qquad h(\eta)=(1+\eta)\log(1+\eta)-\eta .
\]
\end{corollary}

\begin{proof}
Apply Corollary~\ref{cor:normalizer-tail-chernoff-joint-search} with sign-count
exponent
\[
  B=\beta+\sigma
\]
and with \(r=1\) in the robust-center count.  The probability that at least one
of the \(M\) sampled signs is not a robust center is at most
\[
  C_{A,B,K,a,R}d^{-A}
  +\exp\left(-I(R,\theta)\frac n{\ell^2}\right)
  +M d^{-B}.
\]
Since \(M\le d^\beta\) and \(B=\beta+\sigma\), the last term is at most
\(d^{-\sigma}\).  Relabeling the constants gives the first assertion.  Since
the center itself is contained in its Hamming ball, the robust-center conclusion
also gives the displayed bound for every sampled sign.  The relative-excess
version is again the specialization \(R=(1+\eta)L\) and \(\theta=L\).
\end{proof}

\begin{corollary}[Chernoff polynomial-confidence robust search]
\label{cor:normalizer-tail-chernoff-poly-confidence}
Fix \(A>0\), \(B>0\), \(a>0\), and constants \(R\ge1\) and \(0<\theta<R\).
Write \(\ell=\log(e+d)\).  In the setting of
Corollary~\ref{cor:normalizer-tail-chernoff-search}, suppose
\[
  p_d(a):=\Prob\{\norm{T_d(\xi)}_{\op}<aN_d\}\le \frac{\theta}{\ell^2}.
\]
Let \(\delta^{(1)},\ldots,\delta^{(M)}\) be independent uniform sign vectors,
independent of the matrices, with \(M\ge\lceil B\rceil\).  Set
\[
  I(R,\theta)=R\log(R/\theta)-R+\theta .
\]
There are constants \(C_{A,K,a,R}\) and \(c_{A,K,a,R}>0\) such that, whenever
\[
  n\le d,\qquad n\ge C_{A,K,a,R}\ell ,
\]
the joint matrix-and-sign probability that no sampled sign is a robust
certificate is at most
\[
  C_{A,K,a,R}d^{-A}
  +\exp\left(-I(R,\theta)\frac n{\ell^2}\right)
  +d^{-B}.
\]
Equivalently, outside that exceptional probability, some sampled
\(\delta^{(s)}\) is the center of a Hamming ball of radius
\[
  r_0=\left\lfloor c_{A,K,a,R}\ell\right\rfloor
\]
on which every sign vector \(\eps'\) satisfies
\[
  \max_{\omega^d=1}\left|\sum_{i=1}^n
  \eps'_i\frac{P_i(\omega)}{\norm{B_i}_{\op}}\right|
  \le C_{A,K,a,R}\sqrt n .
\]

In particular, if \(p_d(a)\le L/\ell^2\) with \(L\ge1\), then for every
\(\eta>0\) the same conclusion holds with the displayed exceptional probability
replaced by
\[
  C_{A,K,a,L,\eta}d^{-A}
  +\exp\left(-Lh(\eta)\frac n{\ell^2}\right)
  +d^{-B},
  \qquad h(\eta)=(1+\eta)\log(1+\eta)-\eta .
\]
\end{corollary}

\begin{proof}
On the good matrix event in the proof of
Corollary~\ref{cor:normalizer-tail-chernoff-search},
\[
  \max_{\omega^d=1}\sum_{i=1}^n
  \frac{|P_i(\omega)|^2}{\norm{B_i}_{\op}^2}
  \le C_{A,K,a,R}\frac n\ell .
\]
That good event fails with probability at most
\[
  C_{A,K,a,R}d^{-A}
  +\exp\left(-I(R,\theta)\frac n{\ell^2}\right).
\]
Condition on any matrix realization satisfying the profile.  Applying
Corollary~\ref{cor:edge-profile-poly-confidence}, and absorbing the profile
constant into \(C_{A,K,a,R}\) and \(c_{A,K,a,R}\), shows that the sampled signs
fail to contain a robust Hamming-ball center with conditional probability at
most \(d^{-B}\).  Adding the matrix exceptional probability proves the first
claim.  The relative-excess specialization is \(R=(1+\eta)L\) and
\(\theta=L\), so \(I((1+\eta)L,L)=Lh(\eta)\).
\end{proof}

\begin{corollary}[Almost-sure Chernoff search-stability transfer]
\label{cor:normalizer-tail-chernoff-borel-search}
Fix \(A>0\), \(a>0\), and constants \(R\ge1\) and \(0<\theta<R\).  There is a
constant \(C_{A,K,a,R}\) with the following property.  Let \((d_m,n_m)_{m\ge1}\)
be a triangular-array sequence with \(n_m\le d_m\), put
\(\ell_m=\log(e+d_m)\), and assume
\[
  n_m\ge C_{A,K,a,R}\ell_m \qquad(m\ge1).
\]
At each \(m\), let \(B_{m,i}=T_{d_m}(\xi^{(m,i)})\), \(1\le i\le n_m\), be
independent full Toeplitz matrices with independent mean-zero, unit-variance
lag variables of sub-Gaussian norm at most \(K\), and let \(P_{m,i}\) be their
symbols.  Suppose
\[
  p_{d_m}(a):=
  \Prob\{\norm{T_{d_m}(\xi)}_{\op}
      <a\,\Exp\norm{T_{d_m}(\xi)}_{\op}\}
  \le \frac{\theta}{\ell_m^2}
  \qquad(m\ge1),
\]
and set \(I(R,\theta)=R\log(R/\theta)-R+\theta\).  If
\[
  \sum_{m=1}^\infty
  \left(d_m^{-A}+
  \exp\left(-I(R,\theta)\frac{n_m}{\ell_m^2}\right)\right)<\infty,
\]
then, with probability one over the matrix triangular array, all but finitely
many \(m\) satisfy the full random-search and stable-center package of
Corollary~\ref{cor:normalizer-tail-search-stability}, with
\(d,n,\ell\) replaced by \(d_m,n_m,\ell_m\) and constants depending on
\((A,K,a,R)\).

In particular, if \(L\ge1\), \(\eta>0\), and
\(p_{d_m}(a)\le L/\ell_m^2\), then the same conclusion holds under
\[
  \sum_{m=1}^\infty
  \left(d_m^{-A}+
  \exp\left(-Lh(\eta)\frac{n_m}{\ell_m^2}\right)\right)<\infty,
  \qquad h(\eta)=(1+\eta)\log(1+\eta)-\eta .
\]
Consequently, along the full integer dimension sequence, if \(A>1\),
\[
  p_d(a)\le \frac{L}{\log^2(e+d)}
  \qquad(d\ge1),
\]
and if \(n(d)\le d\) while, for all sufficiently large \(d\),
\[
  n(d)\ge q\,\log^3(e+d)
  \qquad\text{for some } q>\frac1{Lh(\eta)},
\]
then almost surely all sufficiently large integer dimensions satisfy the same
random-search and stable-center package.  For \(\eta=1\), the threshold is
\(q>1/(L(2\log2-1))\).
\end{corollary}

\begin{proof}
Apply Corollary~\ref{cor:normalizer-tail-chernoff-search} at each \(m\), with
the constant \(C_{A,K,a,R}\) enlarged if necessary.  The displayed summability
condition makes the matrix exceptional probabilities summable, so the first
Borel--Cantelli lemma gives an almost-sure event on which all but finitely many
matrix realizations satisfy the finite Chernoff search package.  Substituting
\(d,n,\ell\) by \(d_m,n_m,\ell_m\) gives the first assertion.

The relative-excess version is the specialization \(R=(1+\eta)L\) and
\(\theta=L\), exactly as in Corollary~\ref{cor:normalizer-tail-chernoff-search}.
For the full-dimension consequence take \(d_m=m\).  Then
\(\sum_m m^{-A}<\infty\) because \(A>1\), and the sample-size assumption gives
\[
  \exp\left(-Lh(\eta)\frac{n(m)}{\log^2(e+m)}\right)
  \le (e+m)^{-Lh(\eta)q}
\]
for all large \(m\), which is summable because \(Lh(\eta)q>1\).  Finite initial
dimensions do not affect the almost-sure eventual conclusion.
\end{proof}

\begin{corollary}[Almost-sure sampled Chernoff robust certificates]
\label{cor:normalizer-tail-chernoff-borel-sampled}
Fix \(A>0\), \(\alpha>0\), \(a>0\), and constants \(R\ge1\) and
\(0<\theta<R\).  There are constants \(C_{A,K,a,R}\) and \(c_{A,K,a,R}>0\)
with the following property.  Let \((d_m,n_m)_{m\ge1}\) be a triangular-array
sequence with \(n_m\le d_m\), put \(\ell_m=\log(e+d_m)\), and assume
\[
  n_m\ge C_{A,K,a,R}\ell_m \qquad(m\ge1).
\]
At each \(m\), let \(B_{m,i}=T_{d_m}(\xi^{(m,i)})\), \(1\le i\le n_m\), be
independent full Toeplitz matrices with independent mean-zero, unit-variance
lag variables of sub-Gaussian norm at most \(K\), and let \(P_{m,i}\) be their
symbols.  Independently of the matrices, sample independent uniform sign
vectors
\[
  \delta_m^{(1)},\ldots,\delta_m^{(M_m)}\in\{-1,1\}^{n_m},
  \qquad M_m\ge \lceil\alpha\rceil .
\]
Suppose
\[
  p_{d_m}(a):=
  \Prob\{\norm{T_{d_m}(\xi)}_{\op}
      <a\,\Exp\norm{T_{d_m}(\xi)}_{\op}\}
  \le \frac{\theta}{\ell_m^2}
  \qquad(m\ge1),
\]
and set \(I(R,\theta)=R\log(R/\theta)-R+\theta\).  If
\[
  \sum_{m=1}^\infty
  \left(d_m^{-A}
  +\exp\left(-I(R,\theta)\frac{n_m}{\ell_m^2}\right)
  +d_m^{-\alpha}\right)<\infty,
\]
then, with probability one over the matrices and sampled signs, all but
finitely many \(m\) have some sampled sign \(\delta_m^{(s_m)}\) which is the
center of a Hamming ball of radius
\[
  r_m=\left\lfloor c_{A,K,a,R}\ell_m\right\rfloor
\]
on which every signing \(\eps'\) satisfies
\[
  \max_{\omega^{d_m}=1}
  \left|\sum_{i=1}^{n_m}\eps'_i
  \frac{P_{m,i}(\omega)}{\norm{B_{m,i}}_{\op}}\right|
  \le C_{A,K,a,R}\sqrt{n_m}.
\]

In particular, if \(L\ge1\), \(\eta>0\), and
\(p_{d_m}(a)\le L/\ell_m^2\), then the same conclusion holds under
\[
  \sum_{m=1}^\infty
  \left(d_m^{-A}
  +\exp\left(-Lh(\eta)\frac{n_m}{\ell_m^2}\right)
  +d_m^{-\alpha}\right)<\infty,
  \qquad h(\eta)=(1+\eta)\log(1+\eta)-\eta .
\]
Consequently, along the full integer dimension sequence, if \(A>1\),
\(\alpha>1\),
\[
  p_d(a)\le \frac{L}{\log^2(e+d)}
  \qquad(d\ge1),
\]
and if \(n(d)\le d\) while, for all sufficiently large \(d\),
\[
  n(d)\ge q\,\log^3(e+d)
  \qquad\text{for some } q>\frac1{Lh(\eta)},
\]
then \(M_d\ge\lceil\alpha\rceil\) sampled signs per dimension eventually contain
such a robust certificate almost surely.  For \(\eta=1\), the threshold is
\(q>1/(L(2\log2-1))\).
\end{corollary}

\begin{proof}
For each \(m\), apply
Corollary~\ref{cor:normalizer-tail-chernoff-search}.  Outside a matrix event of
probability at most
\[
  C_{A,K,a,R}d_m^{-A}
  +\exp\left(-I(R,\theta)\frac{n_m}{\ell_m^2}\right),
\]
its proof gives the edge profile
\[
  \max_{\omega^{d_m}=1}\sum_{i=1}^{n_m}
  \frac{\abs{P_{m,i}(\omega)}^2}{\norm{B_{m,i}}_{\op}^2}
  \le C_{A,K,a,R}\frac{n_m}{\ell_m}.
\]
Conditioned on any matrix realization with this profile,
Corollary~\ref{cor:edge-profile-poly-confidence}, with \(B=\alpha\), says that
the sampled signs fail to contain a robust Hamming-ball center with probability
at most \(d_m^{-\alpha}\), after absorbing the profile constant into
\(C_{A,K,a,R}\) and \(c_{A,K,a,R}\).  Hence the joint matrix-and-sign failure
probability at dimension \(m\) is bounded by the \(m\)-th summand in the first
displayed series, up to constants.  Borel--Cantelli gives the eventual
sampled-certificate conclusion.

The relative-excess specialization is \(R=(1+\eta)L\) and \(\theta=L\), so
\(I((1+\eta)L,L)=Lh(\eta)\).  For the full-dimension consequence take
\(d_m=m\).  Then \(\sum_m m^{-A}<\infty\) and
\(\sum_m m^{-\alpha}<\infty\), while the sample-size lower bound gives
\[
  \exp\left(-Lh(\eta)\frac{n(m)}{\log^2(e+m)}\right)
  \le (e+m)^{-Lh(\eta)q}
\]
for all large \(m\), which is summable because \(Lh(\eta)q>1\).  The finitely
many initial dimensions are irrelevant.
\end{proof}

\begin{corollary}[Almost-sure Chernoff sampled-search counts]
\label{cor:normalizer-tail-chernoff-borel-counts}
Fix \(A>0\), \(B>0\), \(a>0\), and constants \(R\ge1\) and \(0<\theta<R\).
There are constants \(C_{A,B,K,a,R}\) and \(c_{A,B,K,a,R}>0\) with the following
property.  Let \((d_m,n_m)_{m\ge1}\) be a triangular-array sequence with
\(n_m\le d_m\), put \(\ell_m=\log(e+d_m)\), and assume
\[
  n_m\ge C_{A,B,K,a,R}\ell_m \qquad(m\ge1).
\]
At each \(m\), let \(B_{m,i}=T_{d_m}(\xi^{(m,i)})\), \(1\le i\le n_m\), be
independent full Toeplitz matrices with independent mean-zero, unit-variance
lag variables of sub-Gaussian norm at most \(K\), and let \(P_{m,i}\) be their
symbols.  Independently of the matrices, sample independent uniform sign
vectors
\[
  \delta_m^{(1)},\ldots,\delta_m^{(M_m)}\in\{-1,1\}^{n_m},
  \qquad M_m\ge1 .
\]
Let \(u_m\ge0\) and \(1\le r_m\le M_m\) be deterministic parameters.  Suppose
\[
  p_{d_m}(a):=
  \Prob\{\norm{T_{d_m}(\xi)}_{\op}
      <a\,\Exp\norm{T_{d_m}(\xi)}_{\op}\}
  \le \frac{\theta}{\ell_m^2}
  \qquad(m\ge1),
\]
set \(I(R,\theta)=R\log(R/\theta)-R+\theta\), and assume
\[
  \sum_{m=1}^\infty
  \left[
  d_m^{-A}
  +\exp\left(-I(R,\theta)\frac{n_m}{\ell_m^2}\right)
  +e^{-u_m}
  +\binom{M_m}{r_m}d_m^{-B r_m}
  \right]<\infty .
\]
Then, with probability one over the matrices and sampled signs, all but
finitely many \(m\) satisfy all of the following.

First, the entire sampled list obeys
\[
  \max_{1\le s\le M_m}\max_{\omega^{d_m}=1}
  \left|\sum_{i=1}^{n_m}\delta_{m,i}^{(s)}
  \frac{P_{m,i}(\omega)}{\norm{B_{m,i}}_{\op}}\right|
  \le
  C_{A,B,K,a,R}\sqrt{
    n_m\left(1+\frac{\log(eM_m)+u_m}{\ell_m}\right)} .
\]
Second, at least \(M_m-r_m+1\) of the sampled signs have edge maximum at most
\(C_{A,B,K,a,R}\sqrt{n_m}\).  Third, at least \(M_m-r_m+1\) sampled signs are
centers of Hamming balls of radius
\[
  \left\lfloor c_{A,B,K,a,R}\ell_m\right\rfloor
\]
on which every signing \(\eps'\) has edge maximum at most
\(C_{A,B,K,a,R}\sqrt{n_m}\).

In particular, if \(L\ge1\), \(\eta>0\), and
\(p_{d_m}(a)\le L/\ell_m^2\), then the same conclusion holds with the Chernoff
summand replaced by
\[
  \exp\left(-Lh(\eta)\frac{n_m}{\ell_m^2}\right),
  \qquad h(\eta)=(1+\eta)\log(1+\eta)-\eta .
\]
Consequently, along the full integer dimension sequence, suppose \(A>1\),
\(s>1\), \(L\ge1\), \(\eta>0\), and \(p_d(a)\le L/\log^2(e+d)\).  If
\[
  n(d)\ge q\,\log^3(e+d)
  \quad\text{for all large }d,\qquad
  q>\frac1{Lh(\eta)},
\]
and if a batch size satisfies \(M_d\le d^\beta\) for some \(\beta\ge0\), then
for every \(B>\beta+1\), taking \(r_d=1\) and
\(u_d=s\log(e+d)\) gives an almost-sure event on which, for all sufficiently
large integer dimensions, every sampled sign in the batch is a robust center and
the whole batch satisfies the displayed sampled-list maximum bound.
\end{corollary}

\begin{proof}
Apply Corollary~\ref{cor:normalizer-tail-chernoff-joint-search} at each \(m\).
In that application use
\[
  (d,n,\ell,M,u,r)=(d_m,n_m,\ell_m,M_m,u_m,r_m).
\]
The failure probabilities for the sampled-list
maximum, the good-sign count, and the robust-center count are bounded by a
constant times the summand in the displayed series, after absorbing the two
identical binomial count terms into the constant.  The first Borel--Cantelli
lemma therefore gives an almost-sure event on which all three conclusions hold
for all but finitely many \(m\).

The relative-excess form is obtained from the same specialization
\(R=(1+\eta)L\), \(\theta=L\), so \(I((1+\eta)L,L)=Lh(\eta)\).  For the
full-dimension consequence take \(d_m=m\), \(u_m=s\log(e+m)\), and \(r_m=1\).
Then \(\sum_m m^{-A}<\infty\), \(\sum_m e^{-u_m}<\infty\), and
\[
  \binom{M_m}{1}m^{-B}\le m^{\beta-B},
\]
which is summable when \(B>\beta+1\).  The lower bound on \(n(m)\) gives
\[
  \exp\left(-Lh(\eta)\frac{n(m)}{\log^2(e+m)}\right)
  \le (e+m)^{-Lh(\eta)q}
\]
for all large \(m\), and this is summable by the assumption on \(q\).  Thus the
general Borel statement applies.  With \(r_m=1\), the bounds
\(G_{M_m}\ge M_m\) and \(H_{M_m}\ge M_m\) say that every sampled sign in the
batch is good and robust.
\end{proof}

\begin{corollary}[All-integer Chernoff vanishing robust fractions]
\label{cor:normalizer-tail-chernoff-borel-integer-vanishing-fractions}
Fix \(A>1\), \(B>1\), \(a>0\), \(L\ge1\), \(\eta>0\), and \(0<\tau_0<B-1\), and put
\[
  h(\eta)=(1+\eta)\log(1+\eta)-\eta .
\]
There are constants \(C_{A,B,K,a,L,\eta}\) and \(c_{A,B,K,a,L,\eta}>0\) with the
following property.  Suppose
\[
  p_d(a):=\Prob\{\norm{T_d(\xi)}_{\op}
      <a\,\Exp\norm{T_d(\xi)}_{\op}\}
  \le \frac{L}{\log^2(e+d)}
  \qquad(d\ge1).
\]
Let \(n(d)\le d\) and assume that, for all sufficiently large \(d\),
\[
  n(d)\ge q\,\log^3(e+d)
  \qquad\text{for some } q>\frac1{Lh(\eta)} .
\]
At dimension \(d\), let \(B_{d,i}=T_d(\xi^{(d,i)})\), \(1\le i\le n(d)\), be
independent full Toeplitz matrices with independent mean-zero, unit-variance
lag variables of sub-Gaussian norm at most \(K\), and let \(P_{d,i}\) be their
symbols.  Independently of the matrices, sample an arbitrary deterministic
number \(M_d\ge1\) of uniform sign vectors in \(\{-1,1\}^{n(d)}\).  Let
\(u_d\ge0\) satisfy
\[
  \sum_{d=1}^{\infty}e^{-u_d}<\infty ,
\]
and let \((\rho_d)_{d\ge1}\subset(0,1]\) satisfy
\[
  \rho_d\ge e\,d^{-\tau_0}
\]
for all sufficiently large \(d\).  Then, with probability one over the matrices
and sampled signs, all sufficiently large integer dimensions satisfy both
conclusions below.

First, the whole sampled list obeys
\[
\begin{aligned}
  \max_{1\le s\le M_d}\max_{\omega^d=1}
  \left|\sum_{i=1}^{n(d)}\delta_{d,i}^{(s)}
  \frac{P_{d,i}(\omega)}{\norm{B_{d,i}}_{\op}}\right|
  \le
  C_{A,B,K,a,L,\eta}\sqrt{
    n(d)\left(1+\frac{\log(eM_d)+u_d}{\log(e+d)}\right)} .
\end{aligned}
\]
Second, at least a \(1-\rho_d\) fraction of the \(M_d\) sampled signs have edge
maximum at most \(C_{A,B,K,a,L,\eta}\sqrt{n(d)}\) and are centers of Hamming
balls of radius
\[
  \left\lfloor c_{A,B,K,a,L,\eta}\log(e+d)\right\rfloor
\]
on which every signing \(\eps'\) has edge maximum at most
\[
  C_{A,B,K,a,L,\eta}\sqrt{n(d)}.
\]
No upper bound on \(M_d\) is required.  In particular, if \(\rho_d\to0\), then
the robust-center fraction tends to one along all integer dimensions.
\end{corollary}

\begin{proof}
Apply Corollary~\ref{cor:normalizer-tail-chernoff-borel-counts} with
\[
  d_m=m,\qquad n_m=n(m),\qquad
  R=(1+\eta)L,\qquad \theta=L,\qquad
  r_m=\lceil\rho_mM_m\rceil .
\]
The eventual lower bound \(n(m)\ge q\log^3(e+m)\) implies
\(n_m\ge C_{A,B,K,a,L,\eta}\ell_m\) after discarding finitely many dimensions.
The matrix summands are summable because \(A>1\) and
\[
  \exp\left(-Lh(\eta)\frac{n(m)}{\log^2(e+m)}\right)
  \le (e+m)^{-Lh(\eta)q}
\]
for all large \(m\), with \(Lh(\eta)q>1\).  The sampled-list tail is summable by
the hypothesis on \(u_m\).

For the binomial count tail,
\[
  \binom{M_m}{r_m}
  \le \left(\frac e{\rho_m}\right)^{r_m}.
\]
The lower bound \(\rho_m\ge e\,m^{-\tau_0}\) gives, for all large \(m\),
\[
  \binom{M_m}{r_m}m^{-B r_m}
  \le m^{-(B-\tau_0)r_m}\le m^{-(B-\tau_0)},
\]
which is summable because \(B-\tau_0>1\).  The sampled-count theorem therefore gives
the sampled-list maximum and at least \(M_m-r_m+1\) good robust centers
eventually.  Since
\[
  M_m-r_m+1\ge M_m-\rho_mM_m,
\]
this is at least a \(1-\rho_m\) fraction.
\end{proof}

\begin{corollary}[Almost-sure Chernoff robust fractions]
\label{cor:normalizer-tail-chernoff-borel-fractions}
Fix \(A>0\), \(S>0\), \(a>0\), constants \(R\ge1\) and \(0<\theta<R\), and a
fraction \(0<\rho\le1\).  Choose \(B>0\) and an integer \(M\ge1\) so that
\[
  B\ge 2\log_2(e/\rho),
  \qquad
  M\ge \frac{2(S+1)}{B\rho}.
\]
There are constants \(C_{A,B,K,a,R}\) and \(c_{A,B,K,a,R}>0\) with the following
property.  Let \((d_m,n_m)_{m\ge1}\) be a triangular-array sequence with
\[
  d_m\ge2,\qquad n_m\le d_m,\qquad
  \ell_m=\log(e+d_m),
\]
and assume
\[
  n_m\ge C_{A,B,K,a,R}\ell_m \qquad(m\ge1).
\]
At each \(m\), let \(B_{m,i}=T_{d_m}(\xi^{(m,i)})\), \(1\le i\le n_m\), be
independent full Toeplitz matrices with independent mean-zero, unit-variance
lag variables of sub-Gaussian norm at most \(K\), and let \(P_{m,i}\) be their
symbols.  Independently of the matrices, sample independent uniform sign
vectors
\[
  \delta_m^{(1)},\ldots,\delta_m^{(M)}\in\{-1,1\}^{n_m}.
\]
Suppose
\[
  p_{d_m}(a):=
  \Prob\{\norm{T_{d_m}(\xi)}_{\op}
      <a\,\Exp\norm{T_{d_m}(\xi)}_{\op}\}
  \le \frac{\theta}{\ell_m^2}
  \qquad(m\ge1),
\]
set \(I(R,\theta)=R\log(R/\theta)-R+\theta\), and assume
\[
  \sum_{m=1}^\infty
  \left[
  d_m^{-A}
  +\exp\left(-I(R,\theta)\frac{n_m}{\ell_m^2}\right)
  +d_m^{-S}
  \right]<\infty .
\]
Then, with probability one over the matrices and sampled signs, all but
finitely many \(m\) satisfy both conclusions below.

First, the whole sampled list obeys
\[
  \max_{1\le s\le M}\max_{\omega^{d_m}=1}
  \left|\sum_{i=1}^{n_m}\delta_{m,i}^{(s)}
  \frac{P_{m,i}(\omega)}{\norm{B_{m,i}}_{\op}}\right|
  \le C_{A,B,K,a,R,S,\rho}\sqrt{n_m}.
\]
Second, at least a \(1-\rho\) fraction of the \(M\) sampled signs are centers of
Hamming balls of radius
\[
  \left\lfloor c_{A,B,K,a,R}\ell_m\right\rfloor
\]
on which every signing \(\eps'\) has edge maximum at most
\[
  C_{A,B,K,a,R}\sqrt{n_m}.
\]
The same \(1-\rho\) fraction also has edge maximum at most
\(C_{A,B,K,a,R}\sqrt{n_m}\).

In particular, if \(L\ge1\), \(\eta>0\), and
\(p_{d_m}(a)\le L/\ell_m^2\), then the same conclusion holds with the Chernoff
summand replaced by
\[
  \exp\left(-Lh(\eta)\frac{n_m}{\ell_m^2}\right),
  \qquad h(\eta)=(1+\eta)\log(1+\eta)-\eta .
\]
Consequently, along the full integer dimension sequence, suppose \(A>1\),
\(S>1\), \(L\ge1\), \(\eta>0\), and \(p_d(a)\le L/\log^2(e+d)\).  If
\[
  n(d)\ge q\,\log^3(e+d)
  \quad\text{for all large }d,\qquad
  q>\frac1{Lh(\eta)},
\]
then these \(M=O_{\rho,S,B}(1)\) sampled signs per dimension eventually have a
\(1-\rho\) fraction of robust centers almost surely.
\end{corollary}

\begin{proof}
Let \(r_\rho=\lceil\rho M\rceil\) and \(u_m=S\ell_m\).  We apply
Corollary~\ref{cor:normalizer-tail-chernoff-borel-counts} with the constant
batch size \(M_m=M\), count parameter \(r_m=r_\rho\), and sampled-list parameter
\(u_m\).  It remains only to check the summability of the two sign terms.

As in Corollary~\ref{cor:normalizer-tail-chernoff-robust-fraction-confidence},
\[
  \binom{M}{r_\rho}\le (e/\rho)^{r_\rho}.
\]
Since \(d_m\ge2\) and \(B\ge2\log_2(e/\rho)\),
\[
  \binom{M}{r_\rho}d_m^{-B r_\rho}
  \le d_m^{-B r_\rho/2}.
\]
The lower bound on \(M\) gives \(B r_\rho/2\ge S+1\), hence
\[
  \binom{M}{r_\rho}d_m^{-B r_\rho}\le d_m^{-(S+1)}\le d_m^{-S}.
\]
Also \(e^{-u_m}=(e+d_m)^{-S}\le d_m^{-S}\).  Therefore the summability
hypothesis of Corollary~\ref{cor:normalizer-tail-chernoff-borel-counts} is
satisfied.

That corollary gives the sampled-list maximum with
\[
  1+\frac{\log(eM)+u_m}{\ell_m}
  =1+S+\frac{\log(eM)}{\ell_m}
  \le C_{S,\rho,B},
\]
because \(M\) is fixed once the displayed parameters are fixed.  It also gives
at least \(M-r_\rho+1\) good sampled signs and at least \(M-r_\rho+1\) robust
centers.  Since \(M-r_\rho+1\ge M-\rho M\), both lower bounds are at least a
\(1-\rho\) fraction of the sampled list.

The relative-excess assertion follows from the same specialization
\(R=(1+\eta)L\), \(\theta=L\).  For the full-dimension consequence take
\(d_m=m\).  Then \(\sum_m m^{-A}<\infty\) and \(\sum_m m^{-S}<\infty\), while
the sample-size lower bound implies
\[
  \exp\left(-Lh(\eta)\frac{n(m)}{\log^2(e+m)}\right)
  \le (e+m)^{-Lh(\eta)q}
\]
for all large \(m\), which is summable because \(Lh(\eta)q>1\).  The finitely
many initial dimensions do not affect the almost-sure eventual conclusion.
\end{proof}

\begin{corollary}[Polynomial-scale Chernoff robust fractions]
\label{cor:normalizer-tail-chernoff-borel-polynomial-fractions}
Fix \(A>0\), \(S>0\), \(a>0\), constants \(R\ge1\) and \(0<\theta<R\), a
polynomial dimension exponent \(\kappa>0\), a constant \(c_0>0\), and a
fraction \(0<\rho\le1\).  Choose \(B>0\) and an integer \(M\ge1\) so that
\[
  B\ge 2\log_2(e/\rho),
  \qquad
  M\ge \frac{2(S+1)}{B\rho}.
\]
Assume
\[
  A\kappa>1,\qquad S\kappa>1.
\]
There are constants \(C_{A,B,K,a,R}\) and \(c_{A,B,K,a,R}>0\) with the following
property.  Let \((d_m,n_m)_{m\ge1}\) be a triangular-array sequence with
\[
  d_m\ge2,\qquad d_m\ge c_0m^\kappa,\qquad
  n_m\le d_m,\qquad \ell_m=\log(e+d_m),
\]
and assume
\[
  n_m\ge C_{A,B,K,a,R}\ell_m \qquad(m\ge1).
\]
At each \(m\), let \(B_{m,i}=T_{d_m}(\xi^{(m,i)})\), \(1\le i\le n_m\), be
independent full Toeplitz matrices with independent mean-zero, unit-variance
lag variables of sub-Gaussian norm at most \(K\), and let \(P_{m,i}\) be their
symbols.  Independently of the matrices, sample independent uniform sign
vectors
\[
  \delta_m^{(1)},\ldots,\delta_m^{(M)}\in\{-1,1\}^{n_m}.
\]
Suppose
\[
  p_{d_m}(a):=
  \Prob\{\norm{T_{d_m}(\xi)}_{\op}
      <a\,\Exp\norm{T_{d_m}(\xi)}_{\op}\}
  \le \frac{\theta}{\ell_m^2}
  \qquad(m\ge1),
\]
put \(I(R,\theta)=R\log(R/\theta)-R+\theta\), and assume that, for all
sufficiently large \(m\),
\[
  n_m\ge q\,\ell_m^2\log(e+m)
  \qquad\text{for some } q>\frac1{I(R,\theta)} .
\]
Then, with probability one over the matrices and sampled signs, all but
finitely many \(m\) satisfy both conclusions below.

First, the whole sampled list obeys
\[
  \max_{1\le s\le M}\max_{\omega^{d_m}=1}
  \left|\sum_{i=1}^{n_m}\delta_{m,i}^{(s)}
  \frac{P_{m,i}(\omega)}{\norm{B_{m,i}}_{\op}}\right|
  \le C_{A,B,K,a,R,S,\rho}\sqrt{n_m}.
\]
Second, at least a \(1-\rho\) fraction of the sampled signs are centers of
Hamming balls of radius
\[
  \left\lfloor c_{A,B,K,a,R}\ell_m\right\rfloor
\]
on which every signing \(\eps'\) has edge maximum at most
\[
  C_{A,B,K,a,R}\sqrt{n_m}.
\]
The same \(1-\rho\) fraction also has edge maximum at most
\(C_{A,B,K,a,R}\sqrt{n_m}\).

In particular, if \(L\ge1\), \(\eta>0\), and
\(p_{d_m}(a)\le L/\ell_m^2\), then the same conclusion holds under
\[
  n_m\ge q\,\ell_m^2\log(e+m)
  \qquad\text{for all large }m,
  \qquad
  q>\frac1{Lh(\eta)},
\]
where \(h(\eta)=(1+\eta)\log(1+\eta)-\eta\),
after allowing the constants to depend on \(L\) and \(\eta\).  Thus polynomial
dimension schedules \(d_m\gtrsim m^\kappa\) need only \(S>1/\kappa\) robust
fraction confidence, rather than the all-integer \(S>1\) condition.
\end{corollary}

\begin{proof}
Apply Corollary~\ref{cor:normalizer-tail-chernoff-borel-fractions}.  It is
enough to verify
\[
  \sum_m\left[
  d_m^{-A}
  +\exp\left(-I(R,\theta)\frac{n_m}{\ell_m^2}\right)
  +d_m^{-S}
  \right]<\infty .
\]
The polynomial lower bound on \(d_m\) gives
\[
  d_m^{-A}\le C_{A,c_0}m^{-A\kappa},
  \qquad
  d_m^{-S}\le C_{S,c_0}m^{-S\kappa},
\]
and these two series are finite by \(A\kappa>1\) and \(S\kappa>1\).  For all
large \(m\), the sample-window assumption gives
\[
  \exp\left(-I(R,\theta)\frac{n_m}{\ell_m^2}\right)
  \le (e+m)^{-I(R,\theta)q},
\]
which is summable because \(I(R,\theta)q>1\).  The general Borel-fraction
corollary gives the eventual sampled-list maximum and robust-center fraction.
The relative-excess assertion is the specialization \(R=(1+\eta)L\),
\(\theta=L\), for which \(I(R,\theta)=Lh(\eta)\).
\end{proof}

\begin{corollary}[Critical-log Chernoff Borel criterion]
\label{cor:normalizer-tail-chernoff-borel-critical}
Fix \(A>0\), \(a>0\), \(L\ge1\), and \(\eta>0\), and put
\[
  h(\eta)=(1+\eta)\log(1+\eta)-\eta .
\]
There is a constant \(C_{A,K,a,L,\eta}\) with the following property.  Let
\((d_m,n_m)_{m\ge1}\) be any triangular-array sequence with \(n_m\le d_m\),
\(\ell_m=\log(e+d_m)\), and
\[
  n_m\ge C_{A,K,a,L,\eta}\ell_m \qquad(m\ge1).
\]
Suppose
\[
  p_{d_m}(a):=
  \Prob\{\norm{T_{d_m}(\xi)}_{\op}
      <a\,\Exp\norm{T_{d_m}(\xi)}_{\op}\}
  \le \frac{L}{\ell_m^2}
  \qquad(m\ge1),
\]
and
\[
  \sum_{m=1}^\infty d_m^{-A}<\infty .
\]
If, for all sufficiently large \(m\),
\[
  n_m\ge q\,\ell_m^2\log(e+m)
  \qquad\text{for some } q>\frac1{Lh(\eta)},
\]
then, with probability one over the matrix triangular array, all but finitely
many \(m\) satisfy both the edge-profile/signing/scalar-moment package of
Corollary~\ref{cor:normalizer-tail-chernoff-borel-density} and the
random-search/stable-center package of
Corollary~\ref{cor:normalizer-tail-chernoff-borel-search}.
\end{corollary}

\begin{proof}
The relative-excess summability condition in
Corollaries~\ref{cor:normalizer-tail-chernoff-borel-density} and
\ref{cor:normalizer-tail-chernoff-borel-search} is
\[
  \sum_m\left(d_m^{-A}+
  \exp\left(-Lh(\eta)\frac{n_m}{\ell_m^2}\right)\right)<\infty .
\]
The first summand is summable by hypothesis.  The lower bound on \(n_m\) gives,
for all sufficiently large \(m\),
\[
  \exp\left(-Lh(\eta)\frac{n_m}{\ell_m^2}\right)
  \le (e+m)^{-Lh(\eta)q},
\]
and the exponent is strictly larger than one.  Hence the displayed series is
finite, and the two Chernoff Borel transfers apply.
\end{proof}

\begin{corollary}[Stretched-scale Chernoff Borel window]
\label{cor:normalizer-tail-chernoff-borel-window}
Fix \(A>0\), \(a>0\), \(L\ge1\), \(\eta>0\), \(c_0>0\), and \(\kappa>0\).
Put \(h(\eta)=(1+\eta)\log(1+\eta)-\eta\).  There is a constant
\(C_{A,K,a,L,\eta}\) with the following property.  Suppose
\[
  p_d(a):=\Prob\{\norm{T_d(\xi)}_{\op}
      <a\,\Exp\norm{T_d(\xi)}_{\op}\}
  \le \frac{L}{\log^2(e+d)}
  \qquad(d\ge1).
\]
Let \((d_m,n_m)_{m\ge1}\) satisfy, with \(\ell_m=\log(e+d_m)\),
\[
  d_m\ge \exp(c_0m^\kappa),\qquad
  C_{A,K,a,L,\eta}\ell_m\le n_m\le d_m .
\]
Assume in addition that at least one of the following two lower sample-window
conditions holds for all sufficiently large \(m\):
\[
  n_m\ge C_1\ell_m^{2+\rho}
  \qquad\text{for some } C_1>0 \text{ and } \rho>0,
\]
or
\[
  n_m\ge q\,\ell_m^2\log(e+m)
  \qquad\text{for some } q>\frac1{Lh(\eta)} .
\]
Then, for the independent full Toeplitz triangular array at dimensions \(d_m\),
with probability one all but finitely many \(m\) satisfy the realization-normalized
edge-profile, random-sign, sign-averaged, deterministic-certificate, and scalar
edge-grid moment package of
Corollary~\ref{cor:normalizer-tail-chernoff-borel-density}.  On the same
eventual matrix array, the full random-search and stable-center package of
Corollary~\ref{cor:normalizer-tail-chernoff-borel-search} also holds.
\end{corollary}

\begin{proof}
By Corollaries~\ref{cor:normalizer-tail-chernoff-borel-density} and
\ref{cor:normalizer-tail-chernoff-borel-search}, it is enough to check
\[
  \sum_m\left(d_m^{-A}+
  \exp\left(-Lh(\eta)\frac{n_m}{\ell_m^2}\right)\right)<\infty .
\]
The first summand is summable because \(d_m\ge\exp(c_0m^\kappa)\).  If
\(n_m\ge C_1\ell_m^{2+\rho}\), then for all large \(m\),
\[
  \exp\left(-Lh(\eta)\frac{n_m}{\ell_m^2}\right)
  \le \exp(-Lh(\eta)C_1\ell_m^\rho)
  \le \exp(-c\,m^{\kappa\rho}),
\]
which is summable.  If instead
\(n_m\ge q\ell_m^2\log(e+m)\), then
\[
  \exp\left(-Lh(\eta)\frac{n_m}{\ell_m^2}\right)
  \le (e+m)^{-Lh(\eta)q},
\]
which is summable by the displayed lower bound on \(q\).  The two Borel--Cantelli
transfers give the asserted eventual profile and search packages.
\end{proof}

\begin{corollary}[Sparse sampled Chernoff Borel window]
\label{cor:normalizer-tail-chernoff-borel-window-sampled}
Fix \(A>0\), \(\alpha>0\), \(a>0\), \(L\ge1\), \(\eta>0\), \(c_0>0\), and
\(\kappa>0\).  Put \(h(\eta)=(1+\eta)\log(1+\eta)-\eta\).  There are constants
\(C_{A,K,a,L,\eta}\) and \(c_{A,K,a,L,\eta}>0\) with the following property.
Suppose
\[
  p_d(a):=\Prob\{\norm{T_d(\xi)}_{\op}
      <a\,\Exp\norm{T_d(\xi)}_{\op}\}
  \le \frac{L}{\log^2(e+d)}
  \qquad(d\ge1).
\]
Let \((d_m,n_m)_{m\ge1}\) satisfy, with \(\ell_m=\log(e+d_m)\),
\[
  d_m\ge \exp(c_0m^\kappa),\qquad
  C_{A,K,a,L,\eta}\ell_m\le n_m\le d_m .
\]
At each \(m\), take the independent full Toeplitz triangular array at dimension
\(d_m\), and independently sample \(M_m\) uniform sign vectors in
\(\{-1,1\}^{n_m}\), with
\[
  M_m\ge\lceil\alpha\rceil .
\]
Assume in addition that at least one of the following two lower sample-window
conditions holds for all sufficiently large \(m\):
\[
  n_m\ge C_1\ell_m^{2+\rho}
  \qquad\text{for some } C_1>0 \text{ and } \rho>0,
\]
or
\[
  n_m\ge q\,\ell_m^2\log(e+m)
  \qquad\text{for some } q>\frac1{Lh(\eta)} .
\]
Then, with probability one over the matrices and sampled signs, all but
finitely many \(m\) have some sampled sign which is the center of a Hamming ball
of radius
\[
  r_m=\left\lfloor c_{A,K,a,L,\eta}\ell_m\right\rfloor
\]
on which every signing \(\eps'\) satisfies
\[
  \max_{\omega^{d_m}=1}
  \left|\sum_{i=1}^{n_m}\eps'_i
  \frac{P_{m,i}(\omega)}{\norm{B_{m,i}}_{\op}}\right|
  \le C_{A,K,a,L,\eta}\sqrt{n_m}.
\]
Thus on stretched-exponential dimension schedules, the sparse Chernoff Borel
window gives actual sampled robust certificates from \(O_\alpha(1)\) sign trials
for every \(\alpha>0\).
\end{corollary}

\begin{proof}
Apply Corollary~\ref{cor:normalizer-tail-chernoff-borel-sampled} with
\(R=(1+\eta)L\) and \(\theta=L\).  Then
\[
  I((1+\eta)L,L)=Lh(\eta).
\]
It is enough to check
\[
  \sum_m\left(d_m^{-A}
  +\exp\left(-Lh(\eta)\frac{n_m}{\ell_m^2}\right)
  +d_m^{-\alpha}\right)<\infty .
\]
The stretched-exponential lower bound on \(d_m\) makes both
\(\sum_m d_m^{-A}\) and \(\sum_m d_m^{-\alpha}\) finite.  If
\(n_m\ge C_1\ell_m^{2+\rho}\), then the Chernoff summand is bounded for all
large \(m\) by
\[
  \exp(-c\,m^{\kappa\rho}),
\]
which is summable.  If instead
\(n_m\ge q\ell_m^2\log(e+m)\), then
\[
  \exp\left(-Lh(\eta)\frac{n_m}{\ell_m^2}\right)
  \le (e+m)^{-Lh(\eta)q},
\]
which is summable by the lower bound on \(q\).  The sampled Chernoff Borel
transfer gives the asserted eventual sampled robust certificates.
\end{proof}

\begin{corollary}[Sparse Chernoff robust-fraction windows]
\label{cor:normalizer-tail-chernoff-borel-window-fractions}
Fix \(A>0\), \(S>0\), \(a>0\), \(L\ge1\), \(\eta>0\), \(c_0>0\),
\(\kappa>0\), and a fraction \(0<\rho\le1\).  Choose \(B>0\) and an integer
\(M\ge1\) so that
\[
  B\ge 2\log_2(e/\rho),
  \qquad
  M\ge \frac{2(S+1)}{B\rho}.
\]
Put \(h(\eta)=(1+\eta)\log(1+\eta)-\eta\).  There are constants
\(C_{A,B,K,a,L,\eta}\) and \(c_{A,B,K,a,L,\eta}>0\) with the following
property.  Suppose
\[
  p_d(a):=\Prob\{\norm{T_d(\xi)}_{\op}
      <a\,\Exp\norm{T_d(\xi)}_{\op}\}
  \le \frac{L}{\log^2(e+d)}
  \qquad(d\ge1).
\]
Let \((d_m,n_m)_{m\ge1}\) satisfy, with \(\ell_m=\log(e+d_m)\),
\[
  d_m\ge2,\qquad
  d_m\ge \exp(c_0m^\kappa),\qquad
  C_{A,B,K,a,L,\eta}\ell_m\le n_m\le d_m .
\]
At each \(m\), take the independent full Toeplitz triangular array at dimension
\(d_m\), and independently sample \(M\) uniform sign vectors in
\(\{-1,1\}^{n_m}\).
Assume in addition that at least one of the following two lower sample-window
conditions holds for all sufficiently large \(m\):
\[
  n_m\ge C_1\ell_m^{2+\zeta}
  \qquad\text{for some } C_1>0 \text{ and } \zeta>0,
\]
or
\[
  n_m\ge q\,\ell_m^2\log(e+m)
  \qquad\text{for some } q>\frac1{Lh(\eta)} .
\]
Then, with probability one over the matrices and sampled signs, all but
finitely many \(m\) satisfy both conclusions below.

First, the whole sampled list obeys
\[
  \max_{1\le s\le M}\max_{\omega^{d_m}=1}
  \left|\sum_{i=1}^{n_m}\delta_{m,i}^{(s)}
  \frac{P_{m,i}(\omega)}{\norm{B_{m,i}}_{\op}}\right|
  \le C_{A,B,K,a,L,\eta,S,\rho}\sqrt{n_m}.
\]
Second, at least a \(1-\rho\) fraction of the sampled signs are centers of
Hamming balls of radius
\[
  \left\lfloor c_{A,B,K,a,L,\eta}\ell_m\right\rfloor
\]
on which every signing \(\eps'\) has edge maximum at most
\[
  C_{A,B,K,a,L,\eta}\sqrt{n_m}.
\]
The same \(1-\rho\) fraction also has edge maximum at most
\(C_{A,B,K,a,L,\eta}\sqrt{n_m}\).  Thus, on stretched-exponential dimension
schedules, the constant-size robust-fraction test has no lower summability
restriction on \(S\) beyond \(S>0\).
\end{corollary}

\begin{proof}
Apply Corollary~\ref{cor:normalizer-tail-chernoff-borel-fractions} with
\[
  R=(1+\eta)L,\qquad \theta=L.
\]
Then \(I(R,\theta)=Lh(\eta)\).  It is enough to check
\[
  \sum_m\left[
  d_m^{-A}
  +\exp\left(-Lh(\eta)\frac{n_m}{\ell_m^2}\right)
  +d_m^{-S}
  \right]<\infty .
\]
The stretched-exponential lower bound on \(d_m\) makes both
\(\sum_m d_m^{-A}\) and \(\sum_m d_m^{-S}\) finite, for every positive \(A\)
and \(S\).  If \(n_m\ge C_1\ell_m^{2+\zeta}\), then the Chernoff summand is
bounded for all large \(m\) by
\[
  \exp(-c\,m^{\kappa\zeta}),
\]
which is summable.  If instead \(n_m\ge q\ell_m^2\log(e+m)\), then
\[
  \exp\left(-Lh(\eta)\frac{n_m}{\ell_m^2}\right)
  \le (e+m)^{-Lh(\eta)q},
\]
which is summable by the lower bound on \(q\).  The Borel-fraction transfer
gives the asserted eventual sampled-list maximum and robust-center fraction.
\end{proof}

\begin{corollary}[Sparse Chernoff sampled-count windows]
\label{cor:normalizer-tail-chernoff-borel-window-counts}
Fix \(A>0\), \(B>0\), \(a>0\), \(L\ge1\), \(\eta>0\), \(c_0>0\), and
\(\kappa>0\).  Put \(h(\eta)=(1+\eta)\log(1+\eta)-\eta\).  There are constants
\(C_{A,B,K,a,L,\eta}\) and \(c_{A,B,K,a,L,\eta}>0\) with the following
property.  Suppose
\[
  p_d(a):=\Prob\{\norm{T_d(\xi)}_{\op}
      <a\,\Exp\norm{T_d(\xi)}_{\op}\}
  \le \frac{L}{\log^2(e+d)}
  \qquad(d\ge1).
\]
Let \((d_m,n_m)_{m\ge1}\) satisfy, with \(\ell_m=\log(e+d_m)\),
\[
  d_m\ge \exp(c_0m^\kappa),\qquad
  C_{A,B,K,a,L,\eta}\ell_m\le n_m\le d_m .
\]
At each \(m\), take the independent full Toeplitz triangular array at dimension
\(d_m\), and independently sample \(M_m\) uniform sign vectors in
\(\{-1,1\}^{n_m}\), with \(M_m\ge1\).  Let \(u_m\ge0\) and
\(1\le r_m\le M_m\) be deterministic parameters.  Assume in addition that at
least one of the following two lower sample-window conditions holds for all
sufficiently large \(m\):
\[
  n_m\ge C_1\ell_m^{2+\rho}
  \qquad\text{for some } C_1>0 \text{ and } \rho>0,
\]
or
\[
  n_m\ge q\,\ell_m^2\log(e+m)
  \qquad\text{for some } q>\frac1{Lh(\eta)} .
\]
If
\[
  \sum_{m=1}^\infty
  \left(e^{-u_m}+\binom{M_m}{r_m}d_m^{-B r_m}\right)<\infty,
\]
then, with probability one over the matrices and sampled signs, all but
finitely many \(m\) satisfy the sampled-list maximum bound, good-sign count, and
robust-center count conclusions of
Corollary~\ref{cor:normalizer-tail-chernoff-borel-counts}, with constants
depending on \((A,B,K,a,L,\eta)\).

In particular, let \(s>1\) and take \(u_m=s\log(e+m)\) and \(r_m=1\).  If
\[
  M_m\le d_m^\beta
  \qquad(m\ge1)
\]
for some \(\beta\ge0\), then every \(B>\beta\) gives an almost-sure event on
which, for all sufficiently large \(m\), every sampled sign in the batch is an
\(O(\sqrt{n_m})\)-good robust center and the whole batch satisfies the displayed
sampled-list maximum bound.  More generally, if
\[
  M_m\le \exp(bm^\sigma)
\]
for some \(b>0\), \(0<\sigma\le\kappa\), then the same eventual all-sampled-sign
conclusion holds whenever either \(\sigma<\kappa\), or \(\sigma=\kappa\) and
\(b<Bc_0\).
\end{corollary}

\begin{proof}
Apply Corollary~\ref{cor:normalizer-tail-chernoff-borel-counts} with
\(R=(1+\eta)L\) and \(\theta=L\).  The stretched-exponential lower bound on
\(d_m\) makes \(\sum_m d_m^{-A}\) finite.  The same computation used in
Corollary~\ref{cor:normalizer-tail-chernoff-borel-window} shows that either
sample-window hypothesis makes
\[
  \sum_m\exp\left(-Lh(\eta)\frac{n_m}{\ell_m^2}\right)
\]
finite.  The two extra sign summands are summable by hypothesis, so the
sampled-count Borel theorem applies.

For the polynomial-in-\(d_m\) batch consequence,
\[
  \binom{M_m}{1}d_m^{-B}\le d_m^{\beta-B},
\]
which is summable on the stretched-exponential schedule when \(B>\beta\).  For
the \(\exp(bm^\sigma)\) batch consequence,
\[
  \binom{M_m}{1}d_m^{-B}
  \le \exp\left(bm^\sigma-Bc_0m^\kappa\right),
\]
which is summable if \(\sigma<\kappa\), or if \(\sigma=\kappa\) and
\(b<Bc_0\).  Finally, \(\sum_m e^{-u_m}<\infty\) because \(s>1\).  With
\(r_m=1\), the count conclusion says all \(M_m\) sampled signs are good and
robust centers eventually.
\end{proof}

\begin{corollary}[Sparse Chernoff growing robust fractions]
\label{cor:normalizer-tail-chernoff-borel-window-growing-fractions}
Fix \(A>0\), \(a>0\), \(L\ge1\), \(\eta>0\), \(c_0>0\), \(\kappa>0\), and
\(0<\rho\le1\).  Put \(h(\eta)=(1+\eta)\log(1+\eta)-\eta\).  There are
constants \(C_{A,K,a,L,\eta}\) and \(c_{A,K,a,L,\eta}>0\) with the following
property.  Suppose
\[
  p_d(a):=\Prob\{\norm{T_d(\xi)}_{\op}
      <a\,\Exp\norm{T_d(\xi)}_{\op}\}
  \le \frac{L}{\log^2(e+d)}
  \qquad(d\ge1).
\]
Let \((d_m,n_m)_{m\ge1}\) satisfy, with \(\ell_m=\log(e+d_m)\),
\[
  d_m\ge \exp(c_0m^\kappa),\qquad
  C_{A,K,a,L,\eta}\ell_m\le n_m\le d_m .
\]
At each \(m\), take the independent full Toeplitz triangular array at dimension
\(d_m\), writing \(B_{m,i}=T_{d_m}(\xi^{(m,i)})\), \(1\le i\le n_m\), and
\(P_{m,i}\) for its symbol.  Independently sample an arbitrary deterministic
number \(M_m\ge1\) of uniform sign vectors
\[
  \delta_m^{(1)},\ldots,\delta_m^{(M_m)}\in\{-1,1\}^{n_m}.
\]
Let \(u_m\ge0\) be deterministic and assume
\[
  \sum_{m=1}^{\infty}e^{-u_m}<\infty .
\]
Assume in addition that at least one of the following two lower sample-window
conditions holds for all sufficiently large \(m\):
\[
  n_m\ge C_1\ell_m^{2+\zeta}
  \qquad\text{for some } C_1>0 \text{ and } \zeta>0,
\]
or
\[
  n_m\ge q\,\ell_m^2\log(e+m)
  \qquad\text{for some } q>\frac1{Lh(\eta)} .
\]
Then, with probability one over the matrices and sampled signs, all but
finitely many \(m\) satisfy both conclusions below.

First, the whole sampled list obeys
\[
\begin{aligned}
  \max_{1\le s\le M_m}\max_{\omega^{d_m}=1}
  \left|\sum_{i=1}^{n_m}\delta_{m,i}^{(s)}
  \frac{P_{m,i}(\omega)}{\norm{B_{m,i}}_{\op}}\right|
  \le
  C_{A,K,a,L,\eta}\sqrt{
    n_m\left(1+\frac{\log(eM_m)+u_m}{\ell_m}\right)} .
\end{aligned}
\]
Second, at least a \(1-\rho\) fraction of the \(M_m\) sampled signs have edge
maximum at most \(C_{A,K,a,L,\eta}\sqrt{n_m}\) and are centers of Hamming balls
of radius
\[
  \left\lfloor c_{A,K,a,L,\eta}\ell_m\right\rfloor
\]
on which every signing \(\eps'\) has edge maximum at most
\[
  C_{A,K,a,L,\eta}\sqrt{n_m}.
\]
No upper bound on \(M_m\) is required for this robust-fraction conclusion.
\end{corollary}

\begin{proof}
Apply Corollary~\ref{cor:normalizer-tail-chernoff-borel-window-counts} with
sign-count exponent \(1\) and
\[
  r_m=\lceil\rho M_m\rceil .
\]
The matrix and Chernoff summability hypotheses are exactly those verified in the
proof of Corollary~\ref{cor:normalizer-tail-chernoff-borel-window-counts}.  The
sampled-list tail is summable by the hypothesis on \(u_m\).  For the binomial
count tail, \(r_m\ge\rho M_m\) gives
\[
  \binom{M_m}{r_m}d_m^{-r_m}
  \le \left(\frac{eM_m}{r_m}\right)^{r_m}d_m^{-r_m}
  \le \left(\frac{e}{\rho}d_m^{-1}\right)^{r_m}.
\]
For all large \(m\), this is at most \(d_m^{-r_m/2}\le d_m^{-1/2}\), which is
summable on the stretched-exponential schedule.  The sampled-count theorem
therefore gives the entropy-sensitive sampled-list maximum and at least
\[
  M_m-r_m+1\ge (1-\rho)M_m
\]
robust centers eventually.
\end{proof}

\begin{corollary}[Vanishing sparse Chernoff discarded fractions]
\label{cor:normalizer-tail-chernoff-borel-window-vanishing-fractions}
Fix \(A>0\), \(B>0\), \(a>0\), \(L\ge1\), \(\eta>0\), \(c_0>0\), and
\(\kappa>0\).  Fix \(0<\tau_0<B\).  Put \(h(\eta)=(1+\eta)\log(1+\eta)-\eta\).  There are constants
\(C_{A,B,K,a,L,\eta}\) and \(c_{A,B,K,a,L,\eta}>0\) with the following
property.  Suppose
\[
  p_d(a):=\Prob\{\norm{T_d(\xi)}_{\op}
      <a\,\Exp\norm{T_d(\xi)}_{\op}\}
  \le \frac{L}{\log^2(e+d)}
  \qquad(d\ge1).
\]
Let \((d_m,n_m)_{m\ge1}\) satisfy, with \(\ell_m=\log(e+d_m)\),
\[
  d_m\ge \exp(c_0m^\kappa),\qquad
  C_{A,B,K,a,L,\eta}\ell_m\le n_m\le d_m .
\]
At each \(m\), take the independent full Toeplitz triangular array at dimension
\(d_m\), writing \(B_{m,i}=T_{d_m}(\xi^{(m,i)})\), \(1\le i\le n_m\), and
\(P_{m,i}\) for its symbol.  Independently sample an arbitrary deterministic
number \(M_m\ge1\) of uniform sign vectors
\[
  \delta_m^{(1)},\ldots,\delta_m^{(M_m)}\in\{-1,1\}^{n_m}.
\]
Let \(u_m\ge0\) be deterministic and assume
\[
  \sum_{m=1}^{\infty}e^{-u_m}<\infty .
\]
Let \((\rho_m)_{m\ge1}\subset(0,1]\) satisfy
\[
  \rho_m\ge e\,d_m^{-\tau_0}
\]
for all sufficiently large \(m\).  Assume in addition that at least one of the
following two lower sample-window conditions holds for all sufficiently large
\(m\):
\[
  n_m\ge C_1\ell_m^{2+\zeta}
  \qquad\text{for some } C_1>0 \text{ and } \zeta>0,
\]
or
\[
  n_m\ge q\,\ell_m^2\log(e+m)
  \qquad\text{for some } q>\frac1{Lh(\eta)} .
\]
Then, with probability one over the matrices and sampled signs, all but
finitely many \(m\) satisfy both conclusions below.

First, the whole sampled list obeys
\[
\begin{aligned}
  \max_{1\le s\le M_m}\max_{\omega^{d_m}=1}
  \left|\sum_{i=1}^{n_m}\delta_{m,i}^{(s)}
  \frac{P_{m,i}(\omega)}{\norm{B_{m,i}}_{\op}}\right|
  \le
  C_{A,B,K,a,L,\eta}\sqrt{
    n_m\left(1+\frac{\log(eM_m)+u_m}{\ell_m}\right)} .
\end{aligned}
\]
Second, at least a \(1-\rho_m\) fraction of the \(M_m\) sampled signs have edge
maximum at most \(C_{A,B,K,a,L,\eta}\sqrt{n_m}\) and are centers of Hamming balls
of radius
\[
  \left\lfloor c_{A,B,K,a,L,\eta}\ell_m\right\rfloor
\]
on which every signing \(\eps'\) has edge maximum at most
\[
  C_{A,B,K,a,L,\eta}\sqrt{n_m}.
\]
No upper bound on \(M_m\) is required.  In particular, if \(\rho_m\to0\), then
the robust-center fraction tends to one along the sampled batches.
\end{corollary}

\begin{proof}
Apply Corollary~\ref{cor:normalizer-tail-chernoff-borel-window-counts} with
sign-count exponent \(B\) and
\[
  r_m=\lceil\rho_mM_m\rceil .
\]
The matrix and Chernoff summability hypotheses are exactly those verified in the
proof of Corollary~\ref{cor:normalizer-tail-chernoff-borel-window-growing-fractions},
and the sampled-list tail is summable by the hypothesis on \(u_m\).  For the
binomial count tail,
\[
  \binom{M_m}{r_m}
  \le \left(\frac{eM_m}{r_m}\right)^{r_m}
  \le \left(\frac e{\rho_m}\right)^{r_m}.
\]
The lower bound on \(\rho_m\) gives, for all large \(m\),
\[
  \binom{M_m}{r_m}d_m^{-B r_m}
  \le d_m^{-(B-\tau_0)r_m}\le d_m^{-(B-\tau_0)},
\]
which is summable on the stretched-exponential schedule.  The sampled-count
theorem therefore applies.  Since
\[
  M_m-r_m+1\ge M_m-\rho_mM_m,
\]
the robust-center count is at least a \(1-\rho_m\) fraction.
\end{proof}

\begin{corollary}[Relative-excess bounded lower-tail density]
\label{cor:normalizer-tail-relative-density}
Fix \(A>0\), \(a>0\), \(L\ge1\), and \(\eta>0\), and write
\(\ell=\log(e+d)\).  In the setting of
Corollary~\ref{cor:normalizer-tail-quant}, suppose
\[
  p_d(a):=\Prob\{\norm{T_d(\xi)}_{\op}<aN_d\}\le \frac{L}{\ell^2}.
\]
Set
\[
  h(\eta):=(1+\eta)\log(1+\eta)-\eta>0 .
\]
Then there is a constant \(C_{A,K,a,L,\eta}\) such that, whenever
\[
  n\le d,\qquad n\ge C_{A,K,a,L,\eta}\ell ,
\]
one has
\[
  \Prob\left\{
  \max_{\omega^d=1}
  \sum_{i=1}^n
  \frac{\abs{P_i(\omega)}^2}{\norm{B_i}_{\op}^2}
  > C_{A,K,a,L,\eta}\frac n\ell
  \right\}
  \le C_{A,K,a,L,\eta}d^{-A}
     +\exp\left(-Lh(\eta)\frac n{\ell^2}\right).
\]
On the complement of this exceptional event, independent signs have the
\(C_{A,K,a,L,\eta}\sqrt n\) edge bound with conditional probability at least
\(1-d^{-A}\), sign-averaged edge maximum \(C_{A,K,a,L,\eta}\sqrt n\), a
deterministic signing certificate, and scalar edge-grid moments bounded by
\((C_{A,K,a,L,\eta}k\,n/\ell)^k\) for every \(2k\)-th moment.
\end{corollary}

\begin{proof}
Apply Corollary~\ref{cor:normalizer-tail-chernoff-density} with
\(\theta=L\) and \(R=(1+\eta)L\).  Since \(L\ge1\), the empirical budget
parameter satisfies \(R\ge1\), and
\[
  I((1+\eta)L,L)
  =(1+\eta)L\log(1+\eta)-\eta L
  =Lh(\eta).
\]
\end{proof}

\begin{corollary}[Large-sample bounded lower-tail density]
\label{cor:normalizer-tail-large-sample}
Fix \(A>0\), \(a>0\), and \(L\ge1\), and write \(\ell=\log(e+d)\).  In the
setting of Corollary~\ref{cor:normalizer-tail-quant}, suppose
\[
  p_d(a):=\Prob\{\norm{T_d(\xi)}_{\op}<aN_d\}\le \frac{L}{\ell^2}.
\]
Then there are constants \(C_{A,K,a,L},c_L>0\) such that, whenever
\[
  n\le d,\qquad n\ge C_{A,K,a,L}\ell ,
\]
one has
\[
  \Prob\left\{
  \max_{\omega^d=1}
  \sum_{i=1}^n
  \frac{\abs{P_i(\omega)}^2}{\norm{B_i}_{\op}^2}
  > C_{A,K,a,L}\frac n\ell
  \right\}
  \le C_{A,K,a,L}d^{-A}+\exp\left(-c_L\frac n{\ell^2}\right).
\]
On the complement of this exceptional event, independent signs have
\[
  \Prob_\eps\left\{
  \max_{\omega^d=1}\left|\sum_{i=1}^n
  \eps_i\frac{P_i(\omega)}{\norm{B_i}_{\op}}\right|
  > C_{A,K,a,L}\sqrt n
  \right\}\le d^{-A},
\]
and
\[
  \Exp_\eps\max_{\omega^d=1}\left|\sum_{i=1}^n
  \eps_i\frac{P_i(\omega)}{\norm{B_i}_{\op}}\right|
  \le C_{A,K,a,L}\sqrt n .
\]
Thus a deterministic sign vector with edge maximum \(C_{A,K,a,L}\sqrt n\)
exists.  On the same event, simultaneously for every integer \(k\ge1\),
\[
  \max_{\omega^d=1}
  \Exp_\eps\left|\sum_{i=1}^n
  \eps_i\frac{P_i(\omega)}{\norm{B_i}_{\op}}\right|^{2k}
  \le
  \left(C_{A,K,a,L}k\,\frac n\ell\right)^k .
\]
Consequently, if \(n/\ell^2\to\infty\), the realization-normalized
edge profile holds with probability tending to one.  If in addition
\(d/n\to\infty\), the corresponding full Toeplitz/Hankel contractions have
\(O_{K,a,L}(\sqrt n)\) discrepancy with probability tending to one.
\end{corollary}

\begin{proof}
Apply Corollary~\ref{cor:normalizer-tail-relative-density} with \(\eta=1\).
The displayed profile estimate and the signing, expectation, and moment
statements follow with \(c_L=L(2\log2-1)\), after absorbing constants.  The final
Toeplitz/Hankel discrepancy conclusion follows from Proposition~\ref{prop:reduction}
and the Toeplitz--Hankel reversal.
\end{proof}

\begin{theorem}[Convex-concentrated normalizers]
\label{thm:convex-normalizers}
Assume in addition to the hypotheses of Theorem~\ref{thm:normalizer-tail} that
the coefficient law is either real or complex.  Write \(q=1\) in the real case
and \(q=2\) in the complex case, after identifying \(\C\) with \(\R^2\).
Assume the product law satisfies a dimension-free convex concentration
inequality: there is \(C_{\rm cc}\) such that for every \(m\), every convex
1-Lipschitz \(F:(\R^q)^m\to\R\), and independent copies
\(X_1,\ldots,X_m\) of \(\xi\),
\[
  \Prob\{F(X)-\Exp F(X)\le -t\}\le \exp(-t^2/C_{\rm cc})
  \qquad (t\ge0).
\]
Then the lower-tail hypothesis of Theorem~\ref{thm:normalizer-tail} holds for
some \(a=a(K,C_{\rm cc})>0\).  Consequently, if
\(n\le d\le e^{o(n)}\) and \(d/n\to\infty\), the realization-normalized full
Toeplitz contractions \(B_i/\norm{B_i}_{\op}\) satisfy
\[
  \discop(B_1/\norm{B_1}_{\op},\ldots,B_n/\norm{B_n}_{\op})
  =O_{K,C_{\rm cc}}(\sqrt n)
\]
with probability tending to one.  The same conclusion holds for full Hankel
contractions.  In particular the conclusion applies to standard real or complex
Gaussian coefficients, Rademacher coefficients, and more generally bounded real
or complex coefficient laws with mean zero and variance one.
\end{theorem}

\begin{proof}
Let \(X=(\xi_{-(d-1)},\ldots,\xi_{d-1})\in(\R^q)^{2d-1}\) and
\[
  F(X)=\norm{T_d(\xi)}_{\op}.
\]
The function \(F\) is convex, and for two lag vectors \(x,y\),
\begin{align*}
  \abs{F(x)-F(y)}
  &\le \norm{T_d(x-y)}_{\op} \\
  &\le \sup_\theta\abs{\sum_{\abs k<d}(x_k-y_k)e^{ik\theta}} \\
  &\le \sqrt{2d-1}\,\norm{x-y}_2 .
\end{align*}
Here \(\norm{x-y}_2\) is the Euclidean norm on \((\R^q)^{2d-1}\).  Thus
\(F/\sqrt{2d-1}\) is convex 1-Lipschitz.  By the convex concentration
hypothesis and Meckes' lower bound \(N_d=\Exp F\ge c_K\sqrt{d\log(e+d)}\),
\[
  \Prob\{F<N_d/2\}
  \le \exp\left(-c_{K,C_{\rm cc}}\frac{N_d^2}{d}\right)
  \le (e+d)^{-\beta}
\]
for some \(\beta=\beta(K,C_{\rm cc})>0\).  This is
\(o(\log^{-2}(e+d))\), so Theorem~\ref{thm:normalizer-tail} applies with
\(a=1/2\).

Real and complex Gaussian laws satisfy the displayed convex concentration by
Gaussian isoperimetry on the corresponding Euclidean space.  Rademacher and
bounded real or complex product laws satisfy it by Talagrand's convex
concentration theorem for product spaces \cite{Talagrand}.  The Hankel case
again follows by Toeplitz--Hankel reversal.
\end{proof}

\begin{corollary}[Strongly log-concave coefficient laws]
\label{cor:strongly-log-concave-normalizers}
Let the real coefficient law have mean zero, variance one, and density
\(\exp(-V(x))\,dx/Z\) on \(\R\) with \(V''\ge\rho>0\).  More generally, in the
complex case identify \(\C\) with \(\R^2\), assume the law has mean zero,
\(\Exp|\xi|^2=1\), and density \(\exp(-V(z))\,dz/Z\) with
\(\nabla^2 V\ge\rho I_2\).  Then the conclusion of
Theorem~\ref{thm:convex-normalizers} holds, with constants depending only on
\(\rho\) and the variance normalization.  Thus, whenever
\[
  n\le d\le e^{o(n)},\qquad d/n\to\infty,
\]
independent realization-normalized full Toeplitz contractions drawn from this
law satisfy
\[
  \discop(B_1/\norm{B_1}_{\op},\ldots,B_n/\norm{B_n}_{\op})
  =O_\rho(\sqrt n)
\]
with probability tending to one, and the same holds for full Hankel
contractions.

In the triangular-array range
\[
  n_m\le d_m,\qquad d_m/n_m\to\infty,\qquad \log d_m=o(n_m),
\]
the sampled operator-certificate, robust-fraction, polynomial-batch, and
almost-sure sampled-batch conclusions of
Corollaries~\ref{cor:convex-normalizers-sampled-operator},
\ref{cor:convex-normalizers-operator-abundance}, and
\ref{cor:convex-normalizers-borel-operator-batches} also apply to these
strongly log-concave laws.
\end{corollary}

\begin{proof}
By the Bakry--Emery/log-Sobolev concentration theorem, product measures with
\(\nabla^2V\ge\rho I\) satisfy dimension-free Gaussian concentration for every
Lipschitz function, and hence the one-sided convex concentration hypothesis in
Theorem~\ref{thm:convex-normalizers}; see, for example, \cite{Ledoux}.  The same
concentration applied to linear coordinate functionals gives a sub-Gaussian
coefficient norm depending only on \(\rho\) and the variance normalization.
Theorem~\ref{thm:convex-normalizers} therefore gives the displayed
realization-normalized Toeplitz and Hankel discrepancy conclusion.  The final
sampled-certificate assertions are exactly the three cited corollaries applied
with the convex-concentration constant supplied by the same log-Sobolev
estimate.
\end{proof}

\begin{corollary}[Finite convolutions of convex-concentrated laws]
\label{cor:convex-convolution-normalizers}
Fix \(s\ge1\).  Let \(\eta^{(1)},\ldots,\eta^{(s)}\) be independent real
coefficient variables, or independent complex coefficient variables identified
with \(\R^2\).  Assume each law is centered, sub-Gaussian, and satisfies the
dimension-free convex concentration hypothesis of
Theorem~\ref{thm:convex-normalizers}, with constants independent of the product
dimension.  For fixed scalars \(a_1,\ldots,a_s\), set
\[
  \zeta:=\sum_{r=1}^s a_r\eta^{(r)},\qquad
  \xi:=\frac{\zeta}{(\Exp|\zeta|^2)^{1/2}},
\]
assuming \(\Exp|\zeta|^2>0\).  Then the coefficient law \(\xi\) satisfies the
one-matrix lower-normalizer tail required in
Theorem~\ref{thm:normalizer-tail}: for some \(a>0\) and \(\beta>0\),
\[
  \Prob\{\norm{T_d(\xi)}_{\op}<a\,\Exp\norm{T_d(\xi)}_{\op}\}
  \le C(1+d)^{-\beta}.
\]
Consequently, whenever
\[
  n\le d\le e^{o(n)},\qquad d/n\to\infty,
\]
the realization-normalized full Toeplitz and full Hankel Matrix Spencer bounds
hold for independent matrices with coefficient law \(\xi\):
\[
  \discop(B_1/\norm{B_1}_{\op},\ldots,B_n/\norm{B_n}_{\op})
  =O(\sqrt n)
\]
with probability tending to one.  The constants depend only on the number of
summands, the coefficients \(a_r\), the variance normalization, and the
sub-Gaussian and convex-concentration constants of the input laws.

The same finite sampled operator-certificate, robust-fraction,
polynomial-batch, and almost-sure sampled-batch conclusions as in
Corollaries~\ref{cor:convex-normalizers-sampled-operator},
\ref{cor:convex-normalizers-operator-abundance}, and
\ref{cor:convex-normalizers-borel-operator-batches} hold for this finite
convolution class.  In particular, Gaussian plus Rademacher noise, Gaussian
plus bounded noise, and strongly-log-concave plus bounded noise are covered
after centering and variance normalization.
\end{corollary}

\begin{proof}
Write \(\sigma^2=\Exp|\zeta|^2\).  Let
\[
  X^{(r)}=(\eta_{-(d-1)}^{(r)},\ldots,\eta_{d-1}^{(r)})
\]
be the independent lag vector for the \(r\)-th input law, and define
\[
  F(X^{(1)},\ldots,X^{(s)})
  :=
  \left\|T_d\left(\sigma^{-1}\sum_{r=1}^s a_rX^{(r)}\right)\right\|_{\op}.
\]
The map \(F\) is convex jointly in all blocks.  If all blocks except the
\(r\)-th are fixed, then
\[
  \abs{F(\ldots,x,\ldots)-F(\ldots,y,\ldots)}
  \le
  \frac{|a_r|}{\sigma}\sqrt{2d-1}\,\norm{x-y}_2 .
\]
Let \(N_d=\Exp F\), and for \(0\le r\le s\) put
\[
  H_r(X^{(1)},\ldots,X^{(r)})
  :=
  \Exp\!\left[F(X^{(1)},\ldots,X^{(s)})
  \mid X^{(1)},\ldots,X^{(r)}\right].
\]
Thus \(H_0=N_d\) and \(H_s=F\).  For each fixed
\((X^{(1)},\ldots,X^{(r-1)})\), the function
\[
  x\mapsto H_r(X^{(1)},\ldots,X^{(r-1)},x)
\]
is convex and \((|a_r|/\sigma)\sqrt{2d-1}\)-Lipschitz, because it is an
average of functions with that property.  Moreover
\[
  H_{r-1}
  =
  \Exp_{X^{(r)}} H_r(X^{(1)},\ldots,X^{(r)})
\]
with the previous blocks fixed.  Hence the convex concentration hypothesis for
the \(r\)-th input law gives, uniformly in the previous blocks,
\[
  \Prob\{H_r-H_{r-1}\le -u\mid X^{(1)},\ldots,X^{(r-1)}\}
  \le
  \exp\left(-\frac{c_r\sigma^2u^2}{|a_r|^2d}\right),
\]
with the evident omission of terms for \(a_r=0\).  If
\(F-N_d\le-t\), then for some \(r\),
\(H_r-H_{r-1}\le -t/s\).  A union bound therefore gives
\[
  \Prob\{F<N_d-t\}
  \le
  \sum_{r:a_r\ne0}
  \exp\left(-\frac{c_r\sigma^2t^2}{s^2|a_r|^2d}\right).
\]

The coefficient law \(\xi\) is centered, variance one, and sub-Gaussian, with
norm depending only on the displayed data.  Meckes' lower bound therefore gives
\[
  N_d=\Exp\norm{T_d(\xi)}_{\op}\ge c\sqrt{d\log(e+d)}.
\]
Taking \(t=N_d/2\) in the preceding tail estimate yields
\[
  \Prob\{F<N_d/2\}\le C(1+d)^{-\beta}
\]
for some \(\beta>0\).  This is the lower-normalizer input in
Theorem~\ref{thm:normalizer-tail}, with \(a=1/2\), and since
\((1+d)^{-\beta}\log^2(e+d)\to0\), all conclusions of that theorem apply in
the stated full-bandwidth range.

The sampled and robust-center assertions follow by the same proofs as
Corollaries~\ref{cor:convex-normalizers-sampled-operator},
\ref{cor:convex-normalizers-operator-abundance}, and
\ref{cor:convex-normalizers-borel-operator-batches}: those arguments use only
the polynomial lower-normalizer tail just proved, the normalizer-tail search
corollaries, and the deterministic edge-to-operator transfer.
\end{proof}

\begin{corollary}[Bounded Gaussian scale-mixture coefficient laws]
\label{cor:bounded-gaussian-scale-mixture-normalizers}
Let \(g\) be a standard real Gaussian variable, or a standard complex Gaussian
variable normalized by \(\Exp|g|^2=1\).  Let \(S\ge0\) be independent of \(g\)
and suppose
\[
  0<s_-\le S\le s_+<\infty
  \qquad\text{almost surely}.
\]
After the variance normalization
\[
  \xi:=\frac{Sg}{(\Exp S^2)^{1/2}},
\]
the coefficient law \(\xi\) satisfies the one-matrix lower-normalizer tail
required in Theorem~\ref{thm:normalizer-tail}: for some \(a>0\), \(\beta>0\),
and \(C<\infty\), depending only on \(s_+/s_-\) and on the real or complex
field,
\[
  \Prob\{\norm{T_d(\xi)}_{\op}<a\,\Exp\norm{T_d(\xi)}_{\op}\}
  \le C(1+d)^{-\beta}.
\]
Consequently, whenever
\[
  n\le d\le e^{o(n)},\qquad d/n\to\infty,
\]
the realization-normalized full Toeplitz and full Hankel Matrix Spencer bounds
hold for independent matrices with coefficient law \(\xi\):
\[
  \discop(B_1/\norm{B_1}_{\op},\ldots,B_n/\norm{B_n}_{\op})
  =O_{s_+/s_-}(\sqrt n)
\]
with probability tending to one.  The sampled operator-certificate,
robust-fraction, polynomial-batch, and almost-sure sampled-batch conclusions of
Corollaries~\ref{cor:convex-normalizers-sampled-operator},
\ref{cor:convex-normalizers-operator-abundance}, and
\ref{cor:convex-normalizers-borel-operator-batches} hold for this class as
well.
\end{corollary}

\begin{proof}
The normalization changes \(s_-\) and \(s_+\) by a common factor, so we assume
\(\Exp S^2=1\).  Let \(S_k\), \(|k|<d\), be the independent volatility variables
and condition on their realized values \(\sigma=(\sigma_k)\).  For the Gaussian
lag vector \(G=(g_k)\), set
\[
  F_\sigma(G):=\norm{T_d((\sigma_kg_k)_{|k|<d})}_{\op}.
\]
For two Gaussian-coordinate vectors \(x,y\),
\[
  |F_\sigma(x)-F_\sigma(y)|
  \le \sup_\theta\left|\sum_{|k|<d}\sigma_k(x_k-y_k)e^{ik\theta}\right|
  \le s_+\sqrt{2d-1}\,\norm{x-y}_2 .
\]
Thus Gaussian concentration gives
\[
  \Prob_G\{F_\sigma\le \Exp_G F_\sigma-t\}
  \le \exp\left(-c\frac{t^2}{s_+^2d}\right).
\]

It remains to compare the conditional mean to the unconditional normalizer.
Since \(\sigma_k^2=s_-^2+\tau_k^2\) with \(\tau_k\ge0\), the conditional
Gaussian vector \((\sigma_kg_k)\) has the same law as
\[
  (s_-g'_k+\tau_k h_k)_{|k|<d},
\]
where \(g'\) and \(h\) are independent standard Gaussian vectors.  By Jensen's
inequality, conditioning on \(g'\),
\[
  \Exp_h\norm{T_d(s_-g'+\tau h)}_{\op}
  \ge \norm{T_d(s_-g')}_{\op}.
\]
Therefore
\[
  \Exp_G F_\sigma\ge s_-\,\Exp\norm{T_d(g')}_{\op}
  \ge c_0s_-\sqrt{d\log(e+d)}
\]
by the standard Gaussian Toeplitz lower bound.  The matching sub-Gaussian upper
bound gives
\[
  N_d:=\Exp\norm{T_d(\xi)}_{\op}
  \le C_0s_+\sqrt{d\log(e+d)} .
\]
Choose \(a=c_0s_-/(4C_0s_+)\).  Then, for every \(\sigma\),
\[
  aN_d\le \frac12\Exp_G F_\sigma .
\]
The conditional concentration estimate with
\(t=\frac12\Exp_G F_\sigma\) yields
\[
  \Prob_G\{F_\sigma<aN_d\}
  \le \exp\left(-c\frac{s_-^2}{s_+^2}\log(e+d)\right)
  \le C(1+d)^{-\beta}.
\]
Integrating over \(\sigma\) proves the displayed lower-normalizer tail.

The law \(\xi\) is centered, variance one, and sub-Gaussian with constants
depending only on \(s_+/s_-\).  Theorem~\ref{thm:normalizer-tail} gives the
realization-normalized Toeplitz conclusion, and the Hankel conclusion follows
by Toeplitz--Hankel reversal.  The sampled certificate assertions use the same
proofs as the cited convex-normalizer sampled corollaries, replacing their
convex-concentration lower-tail input by the polynomial lower tail just proved.
\end{proof}

\begin{theorem}[Heterogeneous lower-normalizer libraries]
\label{thm:heterogeneous-normalizer-library}
Fix \(K,a,\beta,L>0\).  For each pair \((d,n)\), let
\(B_i=T_d(\xi^{(i)})\), \(1\le i\le n\), be independent full Toeplitz
matrices.  The lag variables inside the \(i\)-th matrix are independent copies
of a real or complex law \(\mu_i=\mu_{d,n,i}\), but the laws may vary with
\(i,d,n\).  Assume every \(\mu_i\) is centered, has variance one, has
sub-Gaussian norm at most \(K\), and satisfies the uniform one-matrix lower-tail
bound
\[
  \Prob_{\mu_i}\left\{\norm{T_d(\mu_i)}_{\op}
  < a\,N_{d,i}\right\}
  \le L(1+d)^{-\beta},
  \qquad
  N_{d,i}:=\Exp_{\mu_i}\norm{T_d(\mu_i)}_{\op}.
\]
Let \(P_i\) be the symbol of \(B_i\), and put \(\ell=\log(e+d)\).  If
\(n\le d\) and \(\ell=o(n)\), then
\[
  \max_{\omega^d=1}
  \sum_{i=1}^n
  \frac{\abs{P_i(\omega)}^2}{\norm{B_i}_{\op}^2}
  \le C_{K,a,\beta,L}\frac n\ell
\]
with probability tending to one.  More quantitatively, for every \(A>0\), once
\(n\ge C_{A,K,a,\beta,L}\ell\), the exceptional probability is at most
\[
  C_{A,K,a,\beta,L}d^{-A}
  +L(1+d)^{-\beta}\min\{n,\ell^2\}.
\]
Consequently independent signs have edge maximum \(O_{K,a,\beta,L}(\sqrt n)\)
with high conditional probability, the sign-averaged edge maximum is
\(O_{K,a,\beta,L}(\sqrt n)\), and a deterministic signing with that edge bound
exists.  If in addition \(d/n\to\infty\), then
\[
  \discop(B_1/\norm{B_1}_{\op},\ldots,B_n/\norm{B_n}_{\op})
  =O_{K,a,\beta,L}(\sqrt n)
\]
with probability tending to one, and the same conclusion holds for the
corresponding full Hankel contractions.

In particular, the realization-normalized theorem is stable under arbitrary
heterogeneous selection from any fixed, or dimension-uniform, library of
coefficient laws covered by Theorem~\ref{thm:convex-normalizers},
Corollary~\ref{cor:strongly-log-concave-normalizers}, or
Corollaries~\ref{cor:convex-convolution-normalizers} and
\ref{cor:bounded-gaussian-scale-mixture-normalizers}, as long as the displayed
parameters are uniform.  The matrices need not be identically distributed.
Matrix-level mixtures over a finite such library are included after conditioning
on the selected laws.
\end{theorem}

\begin{proof}
The proof of Theorem~\ref{thm:normalizer-tail} is unchanged except for the
bookkeeping of the deterministic normalizing scales.  We give the details
because they are the point of the heterogeneous extension.

First we replace Proposition~\ref{prop:expected-profile} by its non-identical
version.  For a fixed edge node \(\omega^d=1\),
\[
  P_i(\omega)=\sum_{\abs k<d}\xi_k^{(i)}\omega^k
\]
is centered sub-Gaussian with norm at most \(C_K\sqrt d\), uniformly in \(i\).
Meckes' lower bound, applied to each law \(\mu_i\), gives
\[
  N_{d,i}\ge c_K\sqrt{d\ell}.
\]
Thus
\[
  Y_i(\omega):=\frac{\abs{P_i(\omega)}^2}{N_{d,i}^2}
\]
are independent nonnegative sub-exponential variables with
\(\Exp Y_i(\omega)\le C_K/\ell\) and
\(\norm{Y_i(\omega)}_{\psi_1}\le C_K/\ell\), uniformly in \(i\).  Bernstein's
inequality for independent, not necessarily identically distributed,
sub-exponential variables gives, after increasing \(C_{A,K}\),
\[
  \Prob\left\{
    \sum_{i=1}^nY_i(\omega)>C_{A,K}\frac n\ell
  \right\}\le \exp(-c_{A,K}n).
\]
Since \(n\ge C_{A,K}\ell\), a union bound over the \(d\) edge nodes gives
\[
  \max_{\omega^d=1}\sum_{i=1}^n
  \frac{\abs{P_i(\omega)}^2}{N_{d,i}^2}
  \le C_{A,K}\frac n\ell
\]
outside a \(d^{-A}\) event.

Let
\[
  G=\{i:\norm{B_i}_{\op}\ge aN_{d,i}\},\qquad M=\abs{G^c}.
\]
The uniform lower-tail hypothesis gives
\[
  \Exp M\le nL(1+d)^{-\beta}.
\]
If \(n\ge\ell^2\), Markov's inequality gives
\[
  \Prob\{M>n/\ell^2\}\le L(1+d)^{-\beta}\ell^2,
\]
while if \(n<\ell^2\), the event \(M>n/\ell^2\) forces \(M\ge1\), and is bounded
by \(L(1+d)^{-\beta}n\).  Hence the bad-normalizer count contributes the
displayed second exceptional term.

On the expected-profile event,
\[
  \max_{\omega^d=1}\sum_{i\in G}
  \frac{\abs{P_i(\omega)}^2}{\norm{B_i}_{\op}^2}
  \le
  a^{-2}
  \max_{\omega^d=1}\sum_{i=1}^n
  \frac{\abs{P_i(\omega)}^2}{N_{d,i}^2}
  \le C_{A,K,a}\frac n\ell .
\]
The samples in \(G^c\) are controlled by the same crude root-value and
Hilbert--Schmidt floor used in Theorem~\ref{thm:normalizer-tail}.  Uniform
sub-Gaussianity gives, outside another \(d^{-A}\) event,
\[
  \max_{i\le n,\ \omega^d=1}\abs{P_i(\omega)}
  \le C_{A,K}\sqrt{d\ell}.
\]
Also, since each law has variance one and sub-Gaussian norm at most \(K\), a
uniform Paley--Zygmund and Chernoff argument gives
\[
  \norm{B_i}_{\op}\ge c_K\sqrt d
  \qquad(1\le i\le n)
\]
outside another \(d^{-A}\) event.  Therefore every bad sample contributes at
most \(C_{A,K}\ell\) to each edge node, and the event \(M\le n/\ell^2\) gives
the total bad contribution \(C_{A,K}n/\ell\).

Combining the good and bad contributions proves the heterogeneous edge profile.
The conditional random-sign bound, sign average, and deterministic certificate
then follow from Proposition~\ref{prop:variance-profile} and
Corollary~\ref{cor:edge-profile-certificate}.  If \(d/n\to\infty\),
Proposition~\ref{prop:reduction} transfers the edge bound to full Toeplitz
operator norm, and the Hankel statement follows by Toeplitz--Hankel reversal.

The final library statement follows because
Theorem~\ref{thm:convex-normalizers} and
Corollaries~\ref{cor:strongly-log-concave-normalizers} and
\ref{cor:convex-convolution-normalizers} and
\ref{cor:bounded-gaussian-scale-mixture-normalizers} give the required
polynomial one-matrix lower tail with constants depending only on their
displayed parameters.  If the law of each matrix is first selected from a finite library,
conditioning on the selected laws reduces the model to the deterministic
heterogeneous array just proved.
\end{proof}

\begin{corollary}[Matrix-level radial scalings]
\label{cor:matrix-level-radial-normalization}
Let \(A_1,\ldots,A_n\) be any full Toeplitz or full Hankel array for which one
of the realization-normalized conclusions above applies, and assume
\(\norm{A_i}_{\op}>0\) almost surely.  Let \(R_1,\ldots,R_n\) be arbitrary
positive scalar random variables on the same probability space; no independence,
identical-distribution, boundedness, or moment assumption is imposed on the
\(R_i\).  Put
\[
  \widetilde A_i:=R_iA_i .
\]
Then every realization-normalized conclusion for the array \(A_i\) holds
unchanged for the radially scaled array \(\widetilde A_i\).  In particular, all
full Toeplitz/Hankel \(O(\sqrt n)\) discrepancy, sampled operator-certificate,
robust-fraction, polynomial-batch, and almost-sure sampled-batch conclusions
proved above remain valid after multiplying the covered matrices by arbitrary
positive, possibly heavy-tailed, matrix-level amplitudes.
\end{corollary}

\begin{proof}
The assertion is pathwise.  Since \(R_i>0\),
\[
  \frac{\widetilde A_i}{\norm{\widetilde A_i}_{\op}}
  =
  \frac{R_iA_i}{R_i\norm{A_i}_{\op}}
  =
  \frac{A_i}{\norm{A_i}_{\op}}
  \qquad(1\le i\le n).
\]
Thus the normalized matrix array, and therefore every discrepancy value,
edge-profile event, sampled-sign certificate, Hamming-stability statement, and
Borel--Cantelli conclusion that depends only on that normalized array, is
identical for \(A_i\) and \(\widetilde A_i\).
\end{proof}

\begin{theorem}[Heterogeneous average lower-tail budget]
\label{thm:heterogeneous-average-normalizer-budget}
Fix \(K,a>0\).  For each pair \((d,n)\), let
\(B_i=T_d(\xi^{(i)})\), \(1\le i\le n\), be independent full Toeplitz
matrices.  The lag variables inside the \(i\)-th matrix are independent copies
of a real or complex law \(\mu_i=\mu_{d,n,i}\), but the laws may vary with
\(i,d,n\).  Assume every \(\mu_i\) is centered, has variance one, and has
sub-Gaussian norm at most \(K\).  Put
\[
  N_{d,i}:=\Exp_{\mu_i}\norm{T_d(\mu_i)}_{\op},
  \qquad
  p_{d,i}(a):=
  \Prob_{\mu_i}\{\norm{T_d(\mu_i)}_{\op}<aN_{d,i}\},
\]
and define the aggregate lower-tail budget
\[
  \Pi_{d,n}(a):=\sum_{i=1}^n p_{d,i}(a).
\]
Let \(P_i\) be the symbol of \(B_i\), and put \(\ell=\log(e+d)\).  Assume
\(n\le d\), \(\ell=o(n)\), and
\[
  \Pi_{d,n}(a)\min\left\{1,\frac{\ell^2}{n}\right\}=o(1).
\]
Then
\[
  \max_{\omega^d=1}
  \sum_{i=1}^n
  \frac{\abs{P_i(\omega)}^2}{\norm{B_i}_{\op}^2}
  \le C_{K,a}\frac n\ell
\]
with probability tending to one.  More quantitatively, for every \(A>0\), once
\(n\ge C_{A,K,a}\ell\), the exceptional probability is at most
\[
  C_{A,K,a}d^{-A}
  +\Pi_{d,n}(a)\min\left\{1,\frac{\ell^2}{n}\right\}.
\]
Consequently independent signs have edge maximum \(O_{K,a}(\sqrt n)\) with high
conditional probability, the sign-averaged edge maximum is \(O_{K,a}(\sqrt n)\),
and a deterministic signing with that edge bound exists.  If also
\(d/n\to\infty\), then
\[
  \discop(B_1/\norm{B_1}_{\op},\ldots,B_n/\norm{B_n}_{\op})
  =O_{K,a}(\sqrt n)
\]
with probability tending to one, and the same conclusion holds for full Hankel
contractions.

The sampled-certificate conclusions of
Corollary~\ref{cor:heterogeneous-normalizer-operator-certificates} remain valid
under this average-budget hypothesis after replacing the matrix term
\(L\ell_m^2(1+d_m)^{-\beta}\) by
\[
  \Pi_m(a)\min\left\{1,\frac{\ell_m^2}{n_m}\right\},
  \qquad
  \Pi_m(a):=\sum_{i=1}^{n_m}p_{d_m,i}(a).
\]
In particular, the almost-sure sampled-batch conclusion holds whenever the
corresponding series with this replacement is summable.
\end{theorem}

\begin{proof}
The proof is the heterogeneous proof with the bad-normalizer count kept in its
aggregate form.  First, the non-identical expected-normalized profile still
holds.  For each fixed \(\omega^d=1\),
\[
  Y_i(\omega):=\frac{\abs{P_i(\omega)}^2}{N_{d,i}^2}
\]
are independent nonnegative sub-exponential variables with
\[
  \Exp Y_i(\omega)\le C_K/\ell,\qquad
  \norm{Y_i(\omega)}_{\psi_1}\le C_K/\ell,
\]
uniformly in \(i\), because \(P_i(\omega)\) has sub-Gaussian norm
\(C_K\sqrt d\) and Meckes' lower bound gives
\(N_{d,i}\ge c_K\sqrt{d\ell}\).  Bernstein's inequality and a union bound over
the \(d\) edge nodes give
\[
  \max_{\omega^d=1}\sum_{i=1}^n
  \frac{\abs{P_i(\omega)}^2}{N_{d,i}^2}
  \le C_{A,K}\frac n\ell
\]
outside a \(d^{-A}\) event.

Let
\[
  G=\{i:\norm{B_i}_{\op}\ge aN_{d,i}\},\qquad M=\abs{G^c}.
\]
Then \(\Exp M=\Pi_{d,n}(a)\).  If \(n\ge\ell^2\), Markov's inequality gives
\[
  \Prob\{M>n/\ell^2\}\le \Pi_{d,n}(a)\frac{\ell^2}{n};
\]
if \(n<\ell^2\), the same event forces \(M\ge1\), so it has probability at most
\(\Pi_{d,n}(a)\).  This is the displayed budget term.

On the expected-profile event, the good-normalizer contribution is bounded by
\[
  \max_{\omega^d=1}\sum_{i\in G}
  \frac{\abs{P_i(\omega)}^2}{\norm{B_i}_{\op}^2}
  \le
  a^{-2}
  \max_{\omega^d=1}\sum_{i=1}^n
  \frac{\abs{P_i(\omega)}^2}{N_{d,i}^2}
  \le C_{A,K,a}\frac n\ell .
\]
For the bad indices, the same uniform root-value and Hilbert--Schmidt floor
argument used in Theorem~\ref{thm:heterogeneous-normalizer-library} gives,
outside another \(d^{-A}\) event,
\[
  \max_{i\le n,\ \omega^d=1}\abs{P_i(\omega)}
  \le C_{A,K}\sqrt{d\ell},
  \qquad
  \norm{B_i}_{\op}\ge c_K\sqrt d
  \quad(1\le i\le n).
\]
Thus each bad sample contributes at most \(C_{A,K}\ell\) to each edge node, and
on \(M\le n/\ell^2\) the total bad contribution is at most \(C_{A,K}n/\ell\).
Combining the two parts proves the profile.

The random-sign, average, deterministic-signing, Toeplitz, and Hankel
conclusions then follow from Proposition~\ref{prop:variance-profile},
Corollary~\ref{cor:edge-profile-certificate}, Proposition~\ref{prop:reduction},
and Toeplitz--Hankel reversal, exactly as before.  The sampled-certificate
statements use the same proof as
Corollary~\ref{cor:heterogeneous-normalizer-operator-certificates}, with the
matrix exceptional term replaced by the aggregate budget just proved, followed
by the same Borel--Cantelli argument in the almost-sure case.
\end{proof}

\begin{theorem}[Heterogeneous dyadic lower-tail budget]
\label{thm:heterogeneous-dyadic-normalizer-budget}
Fix \(K,a>0\).  In the setting of
Theorem~\ref{thm:heterogeneous-average-normalizer-budget}, put
\(\ell=\log(e+d)\) and
\[
  J_\ell:=\left\lceil \frac12\log_2(e+\ell)\right\rceil .
\]
For each \(i\), define the one-matrix dyadic depth variable
\[
\begin{aligned}
  Y_{d,i,a}^{\rm dy}:={}&
  \sum_{0\le r<J_\ell}4^{r+1}
  \mathbf 1\!\left\{
  2^{-(r+1)}aN_{d,i}\le \norm{T_d(\mu_i)}_{\op}
  <2^{-r}aN_{d,i}
  \right\} \\
  &+\ell\,
  \mathbf 1\!\left\{
  \norm{T_d(\mu_i)}_{\op}<2^{-J_\ell}aN_{d,i}
  \right\},
\end{aligned}
\]
where the probability and expectation are with respect to an independent draw
from the law \(\mu_i\).  Put
\[
  \Delta_{d,n}^{\rm dy}(a):=
  \sum_{i=1}^n \Exp_{\mu_i}Y_{d,i,a}^{\rm dy}.
\]
Assume
\[
  n\le d,\qquad \ell=o(n),\qquad
  \frac{\ell}{n}\Delta_{d,n}^{\rm dy}(a)=o(1).
\]
Then
\[
  \max_{\omega^d=1}
  \sum_{i=1}^n
  \frac{\abs{P_i(\omega)}^2}{\norm{B_i}_{\op}^2}
  \le C_{K,a}\frac n\ell
\]
with probability tending to one.  More quantitatively, for every \(A>0\) and
\(R\ge1\), once \(n\ge C_{A,K,a}\ell\),
\[
\begin{aligned}
  \Prob\left\{
  \max_{\omega^d=1}
  \sum_{i=1}^n
  \frac{\abs{P_i(\omega)}^2}{\norm{B_i}_{\op}^2}
  > C_{A,K,a,R}\frac n\ell
  \right\}
  \le
  C_{A,K,a,R}d^{-A}
  +\frac{\ell\,\Delta_{d,n}^{\rm dy}(a)}{Rn}.
\end{aligned}
\]
On the complementary matrix event, independent signs have edge maximum
\(O_{K,a,R}(\sqrt n)\) with high conditional probability, the sign-averaged
edge maximum is \(O_{K,a,R}(\sqrt n)\), and a deterministic signing with that
edge bound exists.  If \(d/n\to\infty\), then the same
\(O_{K,a,R}(\sqrt n)\) conclusion holds for the corresponding
realization-normalized full Toeplitz and full Hankel contractions.

There is also a large-sample exponential form.  If \(0<\theta<R\) and
\[
  \Delta_{d,n}^{\rm dy}(a)\le \theta\frac n\ell ,
\]
then the preceding exceptional probability may be replaced by
\[
  C_{A,K,a,R,\theta}d^{-A}
  +\exp\left(-c_{K,a,R,\theta}\frac n{\ell^2}\right).
\]

The sampled-certificate conclusions of
Corollary~\ref{cor:heterogeneous-normalizer-operator-certificates} remain valid
under this dyadic-budget hypothesis.  In the Markov form, replace the matrix
term \(L\ell_m^2(1+d_m)^{-\beta}\) by
\[
  \frac{\ell_m\,\Delta_m^{\rm dy}(a)}{R n_m},
  \qquad
  \Delta_m^{\rm dy}(a):=
  \sum_{i=1}^{n_m}\Exp Y_{d_m,i,a}^{\rm dy}.
\]
In the Bernstein form, under
\(\Delta_m^{\rm dy}(a)\le\theta n_m/\ell_m\) with \(0<\theta<R\), replace that
matrix term by \(\exp(-c n_m/\ell_m^2)\).  The corresponding almost-sure
sampled-batch conclusions hold under the same replacements in the summability
hypotheses.
\end{theorem}

\begin{proof}
The proof is the heterogeneous expected-profile argument with the dyadic bad
normalizer budget kept instead of a single cutoff count.

For a fixed edge node, the variables
\[
  \frac{\abs{P_i(\omega)}^2}{N_{d,i}^2}
\]
are independent sub-exponential variables with mean and \(\psi_1\)-norm
\(O_K(1/\ell)\), uniformly in \(i\).  Hence, as in the proof of
Theorem~\ref{thm:heterogeneous-average-normalizer-budget},
\[
  \max_{\omega^d=1}\sum_{i=1}^n
  \frac{\abs{P_i(\omega)}^2}{N_{d,i}^2}
  \le C_{A,K}\frac n\ell
\]
outside an event of probability at most \(C_{A,K}d^{-A}\).

For the realized matrices, define
\[
\begin{aligned}
  W_a^{\rm dy}:={}&
  \sum_{i=1}^n\sum_{0\le r<J_\ell}4^{r+1}
  \mathbf 1\!\left\{
  2^{-(r+1)}aN_{d,i}\le \norm{B_i}_{\op}
  <2^{-r}aN_{d,i}
  \right\} \\
  &+\ell\sum_{i=1}^n
  \mathbf 1\!\left\{
  \norm{B_i}_{\op}<2^{-J_\ell}aN_{d,i}
  \right\}.
\end{aligned}
\]
Let \(G=\{i:\norm{B_i}_{\op}\ge aN_{d,i}\}\).  On the expected-profile event,
the contribution of \(G\) is at most \(C_{A,K,a}n/\ell\).  Intersect this with
the uniform root-value and Hilbert--Schmidt floor event
\[
  \max_{i\le n,\ \omega^d=1}\abs{P_i(\omega)}
  \le C_{A,K}\sqrt{d\ell},
  \qquad
  \norm{B_i}_{\op}\ge c_K\sqrt d
  \quad(1\le i\le n),
\]
which also fails with probability at most \(C_{A,K}d^{-A}\).  Since
\(N_{d,i}\ge c_K\sqrt{d\ell}\), this event gives
\[
  \max_{i,\omega^d=1}\frac{\abs{P_i(\omega)}^2}{N_{d,i}^2}\le C_{A,K}.
\]
Thus an index in the \(r\)-th dyadic bin contributes at most
\(C_{A,K,a}4^{r+1}\), while an index in the final bottom bin contributes at
most \(C_{A,K}\ell\) by the Hilbert--Schmidt floor.  Therefore
\[
  \max_{\omega^d=1}
  \sum_{i=1}^n
  \frac{\abs{P_i(\omega)}^2}{\norm{B_i}_{\op}^2}
  \le C_{A,K,a}\left(\frac n\ell+W_a^{\rm dy}\right)
\]
outside the same \(d^{-A}\)-probability matrix event.

The variables defining \(W_a^{\rm dy}\) are independent and have expectations
\(\Exp Y_{d,i,a}^{\rm dy}\), so
\[
  \Exp W_a^{\rm dy}=\Delta_{d,n}^{\rm dy}(a).
\]
Markov's inequality gives
\[
  \Prob\left\{W_a^{\rm dy}>R\frac n\ell\right\}
  \le \frac{\ell\,\Delta_{d,n}^{\rm dy}(a)}{Rn}.
\]
On the complement, the displayed profile holds with constant depending on
\(R\), which proves the quantitative Markov form and the high-probability
statement under \(\ell\Delta_{d,n}^{\rm dy}(a)/n\to0\).

For the exponential form, note that \(0\le Y_{d,i,a}^{\rm dy}\le C\ell\), since
\(4^{J_\ell}\le C(e+\ell)\).  If
\(\Delta_{d,n}^{\rm dy}(a)\le\theta n/\ell\), then
\(\Exp W_a^{\rm dy}\le\theta n/\ell\).  Bernstein's inequality for independent
bounded, not necessarily identically distributed, variables gives
\[
  \Prob\left\{W_a^{\rm dy}>R\frac n\ell\right\}
  \le
  \exp\left(-c_{R,\theta}\frac n{\ell^2}\right),
\]
after changing constants to absorb the bound \(C\ell\).

Conditioned on the resulting edge-profile event, the random-sign tail, sign
average, deterministic signing, scalar moment package, and robust Hamming-ball
search statements are exactly those of
Proposition~\ref{prop:variance-profile},
Corollary~\ref{cor:edge-profile-certificate}, and the sampled-search
corollaries.  Proposition~\ref{prop:reduction} and Toeplitz--Hankel reversal
transfer the edge certificates to full Toeplitz and full Hankel operator norm
when \(d/n\to\infty\).  The sampled and almost-sure statements follow by
replacing the matrix exceptional term in the proof of
Corollary~\ref{cor:heterogeneous-normalizer-operator-certificates} by the
Markov or Bernstein bound just proved and then applying the same
Borel--Cantelli argument.
\end{proof}

\begin{corollary}[Heterogeneous inverse-moment lower-normalizer budget]
\label{cor:heterogeneous-invmoment-normalizer-budget}
Fix \(K,a>0\).  In the setting of
Theorem~\ref{thm:heterogeneous-dyadic-normalizer-budget}, define
\[
  X_{d,i}:=\norm{T_d(\mu_i)}_{\op}
\]
for an independent draw from \(\mu_i\), and put
\[
  \nu_{d,i,a}^{(2)}
  :=
  \Exp_{\mu_i}\left[
  \mathbf 1_{\{X_{d,i}<aN_{d,i}\}}
  \min\left\{
  \left(\frac{aN_{d,i}}{X_{d,i}}\right)^2,\ell
  \right\}
  \right],
  \qquad \ell=\log(e+d),
\]
where the minimum is interpreted as \(\ell\) on the event \(X_{d,i}=0\).
Set
\[
  \mathcal V_{d,n}^{(2)}(a):=
  \sum_{i=1}^n\nu_{d,i,a}^{(2)}.
\]
If
\[
  n\le d,\qquad \ell=o(n),\qquad
  \frac{\ell}{n}\mathcal V_{d,n}^{(2)}(a)=o(1),
\]
then the realization-normalized edge profile of
Theorem~\ref{thm:heterogeneous-dyadic-normalizer-budget} holds with probability
tending to one, and hence the \(O_{K,a}(\sqrt n)\) Toeplitz/Hankel Matrix
Spencer conclusion holds whenever \(d/n\to\infty\).

More quantitatively, for every \(A>0\) and \(R\ge1\), once
\(n\ge C_{A,K,a}\ell\), the profile exceptional probability is at most
\[
  C_{A,K,a,R}d^{-A}
  +\frac{C\ell\,\mathcal V_{d,n}^{(2)}(a)}{Rn}.
\]
On the complementary matrix event, the random-sign tail, sign-averaged maximum,
deterministic certificate, scalar moment package, and sampled operator
certificates inherited in
Theorem~\ref{thm:heterogeneous-dyadic-normalizer-budget} hold with the same
replacement.  If \(0<\theta<R\) and
\[
  C\,\mathcal V_{d,n}^{(2)}(a)\le \theta\frac n\ell ,
\]
then the matrix exceptional term may instead be replaced by
\[
  C_{A,K,a,R,\theta}d^{-A}
  +\exp\left(-c_{K,a,R,\theta}\frac n{\ell^2}\right).
\]

More generally, for \(p>0\) define
\[
  I_{d,i,a}^{(p)}
  :=
  \Exp_{\mu_i}\left[
  \mathbf 1_{\{X_{d,i}<aN_{d,i}\}}
  \left(\frac{aN_{d,i}}{X_{d,i}}\right)^p
  \right],
  \qquad
  \mathcal I_{d,n}^{(p)}(a):=\sum_{i=1}^n I_{d,i,a}^{(p)}.
\]
Then
\[
  \mathcal V_{d,n}^{(2)}(a)
  \le
  \ell^{(2-p)_+/2}\mathcal I_{d,n}^{(p)}(a).
\]
Thus the same conclusion follows from
\[
  \frac{\ell^{1+(2-p)_+/2}}{n}\,
  \mathcal I_{d,n}^{(p)}(a)=o(1),
\]
with the corresponding quantitative replacements.
\end{corollary}

\begin{proof}
For each \(i\), let \(Z_i=aN_{d,i}/X_{d,i}\) on
\(\{X_{d,i}<aN_{d,i}\}\).  The same pointwise bin estimate as in the proof of
Corollary~\ref{cor:normalizer-tail-invmoment} gives
\[
  Y_{d,i,a}^{\rm dy}
  \le
  C\,\mathbf 1_{\{X_{d,i}<aN_{d,i}\}}\min\{Z_i^2,\ell\}.
\]
Indeed, on the \(r\)-th finite dyadic bin one has \(Z_i^2>4^r\), so the bin
weight \(4^{r+1}\) is at most \(4Z_i^2\), while \(r<J_\ell\) also gives
\(4^{r+1}\le C\ell\).  On the final bottom bin, the bottom weight \(\ell\) is
bounded by \(\min\{Z_i^2,\ell\}\).  Taking expectations and summing over \(i\)
therefore yields
\[
  \Delta_{d,n}^{\rm dy}(a)
  \le C\,\mathcal V_{d,n}^{(2)}(a).
\]
The Markov and Bernstein assertions now follow directly from
Theorem~\ref{thm:heterogeneous-dyadic-normalizer-budget}.

For the inverse-\(p\) comparison, on \(\{X_{d,i}<aN_{d,i}\}\) we have \(Z_i>1\).
If \(p\ge2\), then \(\min\{Z_i^2,\ell\}\le Z_i^p\).  If \(0<p<2\), then
\(\min\{Z_i^2,\ell\}\le \ell^{1-p/2}Z_i^p\), separately on the ranges
\(Z_i\le\sqrt\ell\) and \(Z_i>\sqrt\ell\).  Taking expectations and summing
gives the displayed aggregate inverse-\(p\) bound.
\end{proof}

\begin{corollary}[Heterogeneous sampled operator certificates]
\label{cor:heterogeneous-normalizer-operator-certificates}
In the setting of Theorem~\ref{thm:heterogeneous-normalizer-library}, let
\((d_m,n_m)\) be any sequence satisfying
\[
  n_m\le d_m,\qquad d_m/n_m\to\infty,\qquad \log d_m=o(n_m),
\]
and put \(\ell_m=\log(e+d_m)\).  For each \(m\), let
\(B_{m,i}=T_{d_m}(\xi^{(m,i)})\), \(1\le i\le n_m\), be independent full
Toeplitz matrices whose coefficient laws \(\mu_{m,i}\) may vary with \(i\) and
\(m\), but satisfy the uniform hypotheses of
Theorem~\ref{thm:heterogeneous-normalizer-library} with the same
\((K,a,\beta,L)\).  Set
\[
  A_{m,i}:=\frac{B_{m,i}}{\norm{B_{m,i}}_{\op}},
  \qquad
  \widetilde A_{m,i}:=J_{d_m}A_{m,i}.
\]
Independently of the matrices, sample independent uniform sign vectors in
\(\{-1,1\}^{n_m}\).

First, for every fixed \(B>0\), if \(M_m\ge\lceil B\rceil\), then with
probability tending to one some sampled sign \(\delta_m^{(s)}\) is the center of
a Hamming ball of radius
\[
  \left\lfloor c_B\ell_m\right\rfloor
\]
on which every sign vector \(\eps'\) satisfies
\[
  \left\|\sum_{i=1}^{n_m}\eps'_iA_{m,i}\right\|_{\op}
  \le C_B\sqrt{n_m},
  \qquad
  \left\|\sum_{i=1}^{n_m}\eps'_i\widetilde A_{m,i}\right\|_{\op}
  \le C_B\sqrt{n_m}.
\]
More quantitatively, for every \(A_0>0\), the edge-certificate failure before
the deterministic edge-to-operator transfer is at most
\[
  C_{A_0,B,K,a,\beta,L}d_m^{-A_0}
  +L\ell_m^2(1+d_m)^{-\beta}
  +d_m^{-B}
\]
for all sufficiently large \(m\).

Second, fix \(0<\rho\le1\), \(S>0\), and a deterministic integer \(M\ge1\).
If \(M\) signs are sampled at level \(m\), then with probability tending to one
at least \(M-\lceil\rho M\rceil+1\) of them are centers of Hamming balls of
radius \(\lfloor c_{\rho,S,M}\ell_m\rfloor\) on which the two operator-norm
bounds above hold with \(C_{\rho,S,M}\) in place of \(C_B\).  The edge failure is
at most
\[
  C_{\rho,S,M,A_0,K,a,\beta,L}d_m^{-A_0}
  +L\ell_m^2(1+d_m)^{-\beta}
  +C_{\rho,S,M}d_m^{-S}.
\]

Third, for every polynomial batch exponent \(\theta\ge0\) and confidence
exponent \(\sigma>0\), if \(1\le M_m\le d_m^\theta\), then with probability
tending to one every sampled sign is the center of such an operator-norm robust
Hamming ball, with constants depending only on
\(\theta,\sigma,K,a,\beta,L\).  The corresponding edge failure is at most
\[
  C_{\theta,\sigma,A_0,K,a,\beta,L}d_m^{-A_0}
  +L\ell_m^2(1+d_m)^{-\beta}
  +d_m^{-\sigma}.
\]

Finally, let \(M_m\ge1\), \(1\le r_m\le M_m\), and \(S>0\) be deterministic.  If
for some \(A_0>0\)
\[
  \sum_{m=1}^{\infty}
  \left[
    d_m^{-A_0}+\ell_m^2(1+d_m)^{-\beta}
    +\binom{M_m}{r_m}d_m^{-S r_m}
  \right]<\infty,
\]
then, almost surely over the matrices and sampled signs, all but finitely many
\(m\) have at least \(M_m-r_m+1\) sampled signs which are centers of
Toeplitz/Hankel operator-norm robust Hamming balls of radius
\(\lfloor c\ell_m\rfloor\).  In particular, on every stretched-exponential
dimension schedule \(d_m\ge\exp(c_0m^\kappa)\), every sampled sign in every
polynomial batch \(M_m\le d_m^\theta\) is eventually such a robust center almost
surely, after choosing any \(S>\theta\) and taking \(r_m=1\).
\end{corollary}

\begin{proof}
Theorem~\ref{thm:heterogeneous-normalizer-library}, applied with exponent
\(A_0\), gives a matrix event of probability at least
\[
  1-C_{A_0,K,a,\beta,L}d_m^{-A_0}
    -L\ell_m^2(1+d_m)^{-\beta}
\]
on which
\[
  V_m^2:=
  \max_{\omega^{d_m}=1}
  \sum_{i=1}^{n_m}
  \frac{\abs{P_{m,i}(\omega)}^2}{\norm{B_{m,i}}_{\op}^2}
  \le C_{A_0,K,a,\beta,L}\frac{n_m}{\ell_m}.
\]
On this matrix event, Corollary~\ref{cor:edge-profile-poly-confidence} gives
the finite-search Hamming ball, with the additional sign failure \(d_m^{-B}\).
Corollary~\ref{cor:edge-to-operator-robust} transfers that whole Hamming ball
from edge maximum to realization-normalized full Toeplitz and full Hankel
operator norm.

For the robust-fraction statement, apply
Corollary~\ref{cor:edge-profile-hamming} on the same matrix event with sign
exponent \(2S\) and \(r=\lceil\rho M\rceil\).  The sign failure is bounded by
\(\binom{M}{r}d_m^{-2Sr}\le C_{\rho,S,M}d_m^{-S}\) for all large \(m\), and
the deterministic transfer is again Corollary~\ref{cor:edge-to-operator-robust}.

For the polynomial batch statement, use
Corollary~\ref{cor:edge-profile-hamming} with sign exponent
\(\theta+\sigma+1\) and union bound over the \(M_m\le d_m^\theta\) sampled
signs; the sign failure is at most \(d_m^{-\sigma}\) after increasing constants.
Transfer the resulting family of Hamming balls by
Corollary~\ref{cor:edge-to-operator-robust}.

The almost-sure statement is the same argument with Borel--Cantelli.  At level
\(m\), Corollary~\ref{cor:edge-profile-hamming}, with sign exponent \(S\), gives
edge robust centers except on an event of probability at most the displayed
summand up to a fixed multiplicative constant.  Summability gives the eventual
edge statement almost surely, and Corollary~\ref{cor:edge-to-operator-robust}
transfers the eventual Hamming balls to Toeplitz and Hankel operator norm.  For
the stretched-exponential consequence with \(r_m=1\),
\[
  \binom{M_m}{1}d_m^{-S}\le d_m^{\theta-S},
\]
which is summable when \(S>\theta\); the matrix terms are also summable on a
stretched-exponential schedule for every \(A_0,\beta>0\).
\end{proof}

\begin{corollary}[Heterogeneous inverse-moment sampled operator schedules]
\label{cor:heterogeneous-invmoment-operator-schedules}
Let \((d_m,n_m)\) be any sequence satisfying
\[
  n_m\le d_m,\qquad d_m/n_m\to\infty,\qquad \log d_m=o(n_m),
\]
and put \(\ell_m=\log(e+d_m)\).  For each \(m\), let
\(B_{m,i}=T_{d_m}(\xi^{(m,i)})\), \(1\le i\le n_m\), be independent full
Toeplitz matrices.  The coefficient law \(\mu_{m,i}\) may vary with \(m,i\),
but all laws are centered, variance one, and have sub-Gaussian norm at most
\(K\).  Define \(N_{m,i}=\Exp_{\mu_{m,i}}\norm{T_{d_m}(\mu_{m,i})}_{\op}\),
\(X_{m,i}=\norm{T_{d_m}(\mu_{m,i})}_{\op}\) for an independent draw, and
\[
  \mathcal V_m^{(2)}(a):=
  \sum_{i=1}^{n_m}
  \Exp_{\mu_{m,i}}\left[
  \mathbf 1_{\{X_{m,i}<aN_{m,i}\}}
  \min\left\{
  \left(\frac{aN_{m,i}}{X_{m,i}}\right)^2,\ell_m
  \right\}\right],
\]
with the minimum interpreted as \(\ell_m\) on \(X_{m,i}=0\).  Set
\[
  A_{m,i}:=\frac{B_{m,i}}{\norm{B_{m,i}}_{\op}},
  \qquad
  \widetilde A_{m,i}:=J_{d_m}A_{m,i}.
\]
Independently of the matrices, sample \(M_m\ge1\) independent uniform sign
vectors.  Fix \(R\ge1\), \(S>0\), and deterministic integers
\(1\le r_m\le M_m\).  If, for some \(A_0>0\),
\[
  \sum_{m=1}^{\infty}
  \left[
    d_m^{-A_0}
    +\frac{\ell_m\,\mathcal V_m^{(2)}(a)}{R n_m}
    +\binom{M_m}{r_m}d_m^{-S r_m}
  \right]<\infty,
\]
then, almost surely over the matrices and sampled signs, all but finitely many
\(m\) have at least \(M_m-r_m+1\) sampled signs which are centers of Hamming
balls of radius \(\lfloor c\ell_m\rfloor\) on which every sign vector
\(\eps'\) satisfies
\[
  \left\|\sum_{i=1}^{n_m}\eps'_iA_{m,i}\right\|_{\op}
  \le C\sqrt{n_m},
  \qquad
  \left\|\sum_{i=1}^{n_m}\eps'_i\widetilde A_{m,i}\right\|_{\op}
  \le C\sqrt{n_m}.
\]
The constants depend only on \(K,a,R,S,A_0\).

The Markov matrix term may be replaced by a Bernstein term.  If \(0<\theta<R\),
\[
  C_*\mathcal V_m^{(2)}(a)\le \theta\frac{n_m}{\ell_m}
  \qquad\text{eventually},
\]
where \(C_*\) is the absolute constant from
Corollary~\ref{cor:heterogeneous-invmoment-normalizer-budget}, and
\[
  \sum_{m=1}^{\infty}
  \left[
    d_m^{-A_0}
    +\exp\left(-c_{K,a,R,\theta}\frac{n_m}{\ell_m^2}\right)
    +\binom{M_m}{r_m}d_m^{-S r_m}
  \right]<\infty,
\]
then the same almost-sure sampled operator-certificate conclusion holds.

More generally, for \(p>0\) one may replace
\(\mathcal V_m^{(2)}(a)\) in the Markov summability condition by
\[
  \ell_m^{(2-p)_+/2}\mathcal I_m^{(p)}(a),
  \qquad
  \mathcal I_m^{(p)}(a):=
  \sum_{i=1}^{n_m}
  \Exp_{\mu_{m,i}}\left[
  \mathbf 1_{\{X_{m,i}<aN_{m,i}\}}
  \left(\frac{aN_{m,i}}{X_{m,i}}\right)^p
  \right].
\]
Thus a sufficient Markov summability condition is
\[
  \sum_m\left[
    d_m^{-A_0}
    +\frac{\ell_m^{1+(2-p)_+/2}\mathcal I_m^{(p)}(a)}{R n_m}
    +\binom{M_m}{r_m}d_m^{-S r_m}
  \right]<\infty .
\]
\end{corollary}

\begin{proof}
Corollary~\ref{cor:heterogeneous-invmoment-normalizer-budget}, applied at level
\(m\), gives an edge-profile matrix event whose failure probability is bounded
by the displayed Markov matrix terms, and on which
\[
  \max_{\omega^{d_m}=1}
  \sum_{i=1}^{n_m}
  \frac{\abs{P_{m,i}(\omega)}^2}{\norm{B_{m,i}}_{\op}^2}
  \le C\frac{n_m}{\ell_m}.
\]
On this event, Corollary~\ref{cor:edge-profile-hamming}, with sign exponent
\(S\), says that at least \(M_m-r_m+1\) sampled signs are logarithmic-radius
edge-robust centers except with additional probability
\(\binom{M_m}{r_m}d_m^{-S r_m}\).  Since \(d_m/n_m\to\infty\) and
\(\log d_m=o(n_m)\), Corollary~\ref{cor:edge-to-operator-robust} transfers
these edge Hamming balls to realization-normalized full Toeplitz and full
Hankel operator norm.  The displayed summability and Borel--Cantelli give the
Markov branch.

The Bernstein branch is the same proof using the Bernstein matrix failure term
from Corollary~\ref{cor:heterogeneous-invmoment-normalizer-budget}.  The final
inverse-\(p\) assertion is the aggregate inverse-\(p\) comparison in that same
corollary.
\end{proof}

\begin{corollary}[Heterogeneous inverse-moment operator robust fractions]
\label{cor:heterogeneous-invmoment-operator-fractions}
In the setting and notation of
Corollary~\ref{cor:heterogeneous-invmoment-operator-schedules}, suppose that for
some \(A_0>0\), \(R\ge1\), and \(a>0\),
\[
  \sum_{m=1}^{\infty}
  \left[
    d_m^{-A_0}
    +\frac{\ell_m\,\mathcal V_m^{(2)}(a)}{R n_m}
  \right]<\infty .
\]
For each \(m\), let
\(\delta_m^{(1)},\ldots,\delta_m^{(M_m)}\) be any deterministic number
\(M_m\ge1\) of independent uniform sampled sign vectors, independent of the
matrices.

First, let \(0<\rho\le1\) and \(q>0\) satisfy
\[
  \sum_{m=1}^{\infty}d_m^{-q}<\infty .
\]
Then, almost surely, all but finitely many \(m\) have at least a
\(1-\rho\) fraction of the sampled signs as centers of Hamming balls of radius
\(\lfloor c_{\rho,q}\ell_m\rfloor\) on which every sign vector \(\eps'\)
satisfies
\[
  \left\|\sum_{i=1}^{n_m}\eps'_iA_{m,i}\right\|_{\op}
  \le C_{\rho,q}\sqrt{n_m},
  \qquad
  \left\|\sum_{i=1}^{n_m}\eps'_i\widetilde A_{m,i}\right\|_{\op}
  \le C_{\rho,q}\sqrt{n_m}.
\]

Second, let \(q>0\), \(\tau>0\), and let \(0<\rho_m\le1\) be deterministic with
\[
  \rho_m\ge e\,d_m^{-\tau}
  \qquad\text{eventually},
  \qquad
  \sum_{m=1}^{\infty}d_m^{-q}<\infty .
\]
Then the preceding conclusion holds with \(\rho\) replaced by \(\rho_m\):
almost surely all but finitely many sampled batches have at least a
\(1-\rho_m\) fraction of their signs as Toeplitz/Hankel operator-norm robust
centers, with constants depending only on \(K,a,R,A_0,q,\tau\).

Third, let \(\theta\ge0\) and \(q>0\) satisfy
\(\sum_m d_m^{-q}<\infty\).  If \(M_m\le d_m^\theta\) eventually, then almost
surely every sampled sign in every sufficiently large batch is such a robust
center, with constants depending only on \(K,a,R,A_0,\theta,q\).

The same three conclusions hold if the displayed Markov matrix-budget
summability is replaced by the Bernstein matrix budget in
Corollary~\ref{cor:heterogeneous-invmoment-operator-schedules}.  They also hold
under the inverse-\(p\) sufficient budget obtained by replacing
\(\mathcal V_m^{(2)}(a)\) with
\(\ell_m^{(2-p)_+/2}\mathcal I_m^{(p)}(a)\).
\end{corollary}

\begin{proof}
For the fixed-fraction assertion set
\[
  r_m=\left\lceil \rho M_m\right\rceil
\]
in Corollary~\ref{cor:heterogeneous-invmoment-operator-schedules} and use sign
exponent \(2q\).  Since
\[
  \binom{M_m}{r_m}
  \le \left(\frac{eM_m}{r_m}\right)^{r_m}
  \le \left(\frac e\rho\right)^{r_m},
\]
we have, for all sufficiently large \(m\),
\[
  \binom{M_m}{r_m}d_m^{-2q r_m}
  \le d_m^{-q r_m}
  \le d_m^{-q}.
\]
The assumed matrix summability and \(\sum_m d_m^{-q}<\infty\) therefore give the
summability hypothesis of
Corollary~\ref{cor:heterogeneous-invmoment-operator-schedules}.  Its conclusion
gives at least \(M_m-r_m+1\ge(1-\rho)M_m\) robust sampled signs.

For the vanishing-fraction assertion set
\[
  r_m=\left\lceil \rho_m M_m\right\rceil
\]
and use sign exponent \(q+\tau\).  The eventual lower bound
\(\rho_m\ge e d_m^{-\tau}\) gives
\[
  \binom{M_m}{r_m}
  \le \left(\frac e{\rho_m}\right)^{r_m}
  \le d_m^{\tau r_m},
\]
and hence
\[
  \binom{M_m}{r_m}d_m^{-(q+\tau)r_m}
  \le d_m^{-q r_m}
  \le d_m^{-q}.
\]
Again the schedule corollary applies, and
\(M_m-r_m+1\ge(1-\rho_m)M_m\).

For the polynomial-batch assertion take \(r_m=1\) and sign exponent
\(\theta+q\).  Then
\[
  \binom{M_m}{1}d_m^{-(\theta+q)}
  \le d_m^{-q}
\]
eventually, so the schedule corollary says that
\(M_m-r_m+1=M_m\) sampled signs are robust centers.

The Bernstein and inverse-\(p\) variants are identical, using the corresponding
matrix-budget replacements already built into
Corollary~\ref{cor:heterogeneous-invmoment-operator-schedules}.
\end{proof}

\begin{corollary}[Heterogeneous inverse-moment sampled-list operator maxima]
\label{cor:heterogeneous-invmoment-operator-batch-maxima}
In the setting and notation of
Corollary~\ref{cor:heterogeneous-invmoment-operator-schedules}, write
\[
  v_{m,i}(\omega):=
  \frac{P_{m,i}(\omega)}{\norm{B_{m,i}}_{\op}},
  \qquad \omega^{d_m}=1.
\]
Denote the sampled sign vectors by
\[
  \delta_m^{(1)},\ldots,\delta_m^{(M_m)}\in\{-1,1\}^{n_m}.
\]
Let \(u_m\ge0\), \(S>0\), and \(1\le r_m\le M_m\) be deterministic.  If, for
some \(A_0>0\) and \(R\ge1\),
\[
  \sum_{m=1}^{\infty}
  \left[
    d_m^{-A_0}
    +\frac{\ell_m\,\mathcal V_m^{(2)}(a)}{R n_m}
    +e^{-u_m}
    +\binom{M_m}{r_m}d_m^{-S r_m}
  \right]<\infty,
\]
then, almost surely over the matrices and sampled signs, all but finitely many
\(m\) satisfy the following three conclusions.

First, the whole sampled list has entropy-sensitive edge maximum
\[
\begin{aligned}
  \max_{1\le s\le M_m}\max_{\omega^{d_m}=1}
  \left|\sum_{i=1}^{n_m}\delta_{m,i}^{(s)}v_{m,i}(\omega)\right|
  \le
  C\sqrt{n_m\left(1+\frac{\log(eM_m)+u_m}{\ell_m}\right)} .
\end{aligned}
\]
Second, the same sampled signs have the full Toeplitz/Hankel operator-norm
maximum
\[
\begin{aligned}
  \max_{1\le s\le M_m}
  \left\|\sum_{i=1}^{n_m}\delta_{m,i}^{(s)}A_{m,i}\right\|_{\op}
  &\le
  C\sqrt{n_m\left(1+\frac{\log(eM_m)+u_m}{\ell_m}\right)},\\
  \max_{1\le s\le M_m}
  \left\|\sum_{i=1}^{n_m}\delta_{m,i}^{(s)}\widetilde A_{m,i}\right\|_{\op}
  &\le
  C\sqrt{n_m\left(1+\frac{\log(eM_m)+u_m}{\ell_m}\right)} .
\end{aligned}
\]
Third, at least \(M_m-r_m+1\) sampled signs are centers of Hamming balls of
radius \(\lfloor c\ell_m\rfloor\) on which every sign vector \(\eps'\)
satisfies the \(O(\sqrt{n_m})\) Toeplitz and Hankel operator-norm bounds from
Corollary~\ref{cor:heterogeneous-invmoment-operator-schedules}.  The constants
depend only on \(K,a,R,S,A_0\).

In particular, if the matrix budget
\[
  \sum_m\left[
    d_m^{-A_0}
    +\frac{\ell_m\,\mathcal V_m^{(2)}(a)}{R n_m}
  \right]
\]
is finite and \(\sum_m d_m^{-q}<\infty\), then taking
\(u_m=q\ell_m\) gives the same whole-list operator maximum with
\[
  1+\frac{\log(eM_m)+u_m}{\ell_m}
  \le C_q\left(1+\frac{\log(eM_m)}{\ell_m}\right).
\]
Under this same summability, every polynomial batch \(M_m\le d_m^\theta\) is
eventually entirely robust and has whole-list Toeplitz/Hankel operator maximum
\(O_{\theta,q}(\sqrt{n_m})\).

The Markov matrix-budget term may be replaced by the Bernstein matrix budget in
Corollary~\ref{cor:heterogeneous-invmoment-operator-schedules}, and the
inverse-\(p\) sufficient budget is obtained by replacing
\(\mathcal V_m^{(2)}(a)\) with
\(\ell_m^{(2-p)_+/2}\mathcal I_m^{(p)}(a)\).
\end{corollary}

\begin{proof}
Corollary~\ref{cor:heterogeneous-invmoment-normalizer-budget}, applied at level
\(m\), gives an edge-profile matrix event whose failure probability is bounded
by the displayed Markov matrix terms, and on which
\[
  \max_{\omega^{d_m}=1}
  \sum_{i=1}^{n_m}
  \abs{v_{m,i}(\omega)}^2
  \le C_0\frac{n_m}{\ell_m}.
\]
Condition on such a matrix realization.  Corollary~\ref{cor:edge-profile-random-search}
with parameter \(u_m\) gives the displayed whole-list edge maximum except with
conditional probability at most \(e^{-u_m}\).  On the same matrix realization,
Corollary~\ref{cor:edge-profile-hamming}, with exponent \(S\), gives at least
\(M_m-r_m+1\) logarithmic-radius robust centers except with conditional
probability at most \(\binom{M_m}{r_m}d_m^{-S r_m}\).  The robust Hamming balls
transfer to Toeplitz and Hankel operator norm by
Corollary~\ref{cor:edge-to-operator-robust}.

It remains only to transfer the whole-list maximum for the sampled signs
themselves.  The deterministic proof of Proposition~\ref{prop:reduction} gives,
uniformly in each sampled sign \(\delta_m^{(s)}\),
\[
  \left\|\sum_i\delta_{m,i}^{(s)}A_{m,i}\right\|_{\op}
  \le
  C_{\rm red}
  \max_{\omega^{d_m}=1}
  \left|\sum_i\delta_{m,i}^{(s)}v_{m,i}(\omega)\right|
  +E_m,
\]
where \(E_m=o(\sqrt{n_m})\).  The entropy-sensitive right side is always at
least a constant multiple of \(\sqrt{n_m}\), so \(E_m\) is absorbed after
increasing \(C\) and discarding finitely many \(m\).  The Hankel estimate is the
same because multiplication by \(J_{d_m}\) preserves operator norm and converts
Toeplitz matrices to Hankel matrices.

The displayed summability and Borel--Cantelli give the Markov branch.  For the
polynomial-batch consequence, take \(r_m=1\), \(u_m=q\ell_m\), and
\(S=\theta+q\); then
\[
  e^{-u_m}\le d_m^{-q},
  \qquad
  \binom{M_m}{1}d_m^{-S}\le d_m^{-q}
\]
eventually, while the entropy factor is bounded by a constant depending on
\(\theta,q\).  The Bernstein and inverse-\(p\) variants are the corresponding
matrix-budget replacements from
Corollary~\ref{cor:heterogeneous-invmoment-operator-schedules}.
\end{proof}

\begin{corollary}[Concrete heterogeneous batch-maximal operator schedules]
\label{cor:heterogeneous-invmoment-operator-batch-maxima-schedules}
In the setting and notation of
Corollary~\ref{cor:heterogeneous-invmoment-operator-batch-maxima}, put
\[
  \mathfrak b_m(a,R):=
  \frac{\ell_m\,\mathcal V_m^{(2)}(a)}{R n_m}.
\]
Assume one of the following Markov-form matrix schedules.
\begin{enumerate}[label=\textup{(\roman*)},leftmargin=2.2em]
\item \(d_m=m\), \(A_0>1\), and
\[
  \mathfrak b_m(a,R)\le C_0m^{-B}
  \qquad\text{eventually for some }B>1.
\]
\item \(d_m\ge c_0m^\kappa\) for some \(c_0,\kappa>0\),
\[
  A_0\kappa>1,
  \qquad
  \mathfrak b_m(a,R)\le C_0d_m^{-B}
  \qquad\text{eventually for some }B\text{ with }B\kappa>1.
\]
\item \(d_m\ge\exp(c_0m^\kappa)\) for some \(c_0,\kappa>0\), \(A_0>0\), and
\[
  \mathfrak b_m(a,R)\le C_0\ell_m^{-\eta}
  \qquad\text{eventually for some }\eta>1/\kappa .
\]
Equivalently, the last display follows from
\(\mathcal V_m^{(2)}(a)/n_m\le C_1\ell_m^{-B}\) with
\(B>1+1/\kappa\).
\end{enumerate}
Let \(q>0\) be chosen so that
\[
  \sum_m d_m^{-q}<\infty ,
\]
that is, \(q>1\) in \textup{(i)}, \(q\kappa>1\) in \textup{(ii)}, and any
\(q>0\) in \textup{(iii)}.  Set \(u_m=q\ell_m\).

Then all the following conclusions hold almost surely over the matrices and
sampled signs.  First, for any deterministic sampled batch sizes \(M_m\), any
fixed \(0<\rho\le1\), and \(r_m=\lceil\rho M_m\rceil\), all but finitely many
\(m\) have at least a \(1-\rho\) fraction of sampled signs as Toeplitz/Hankel
operator-norm robust centers, and the whole sampled list has the
entropy-sensitive edge and operator-norm maximum from
Corollary~\ref{cor:heterogeneous-invmoment-operator-batch-maxima}.

Second, for arbitrary \(M_m\), any \(\tau>0\), and deterministic
\(\rho_m\in(0,1]\) satisfying
\[
  \rho_m\ge e\,d_m^{-\tau}
  \qquad\text{eventually},
\]
the same whole-list maximum holds and at least a \(1-\rho_m\) fraction of sampled
signs are robust centers eventually.

Third, if \(M_m\le d_m^\theta\) eventually for some \(\theta\ge0\), then every
sampled sign in every sufficiently large batch is a Toeplitz/Hankel
operator-norm robust center and
\[
  \max_{1\le s\le M_m}
  \left\|\sum_{i=1}^{n_m}\delta_{m,i}^{(s)}A_{m,i}\right\|_{\op}
  +
  \max_{1\le s\le M_m}
  \left\|\sum_{i=1}^{n_m}\delta_{m,i}^{(s)}\widetilde A_{m,i}\right\|_{\op}
  \le C_{\theta,q}\sqrt{n_m}
\]
for all sufficiently large \(m\).

The same conclusions hold in the Bernstein branch.  Namely, suppose that for
some \(0<\theta_0<R\),
\[
  C_*\mathcal V_m^{(2)}(a)\le \theta_0\frac{n_m}{\ell_m}
  \qquad\text{eventually},
\]
where \(C_*\) is the absolute constant from
Corollary~\ref{cor:heterogeneous-invmoment-normalizer-budget}, and that
\[
  \sum_m\left[
    d_m^{-A_0}
    +\exp\left(-c_{K,a,R,\theta_0}\frac{n_m}{\ell_m^2}\right)
  \right]<\infty .
\]
Then the same three sampled-batch conclusions hold after adding the same
summable sign-tail power \(\sum_m d_m^{-q}\).  In particular this Bernstein
matrix summability is guaranteed along all integer dimensions if \(A_0>1\) and
\[
  n(d)\ge Q\log^3(e+d)
  \quad\text{eventually with}\quad Q>1/c_{K,a,R,\theta_0},
\]
along polynomial schedules \(d_m\ge c_0m^\kappa\) if \(A_0\kappa>1\) and
\[
  n_m\ge Q\ell_m^2\log(e+m)
  \quad\text{eventually with}\quad Q>1/c_{K,a,R,\theta_0},
\]
and along stretched-exponential schedules \(d_m\ge\exp(c_0m^\kappa)\) under the
same \(\ell_m^2\log(e+m)\) lower bound, or under the simpler condition
\[
  n_m\ge Q\ell_m^{2+\tau}
  \qquad\text{eventually}
\]
for any \(Q,\tau>0\).
\end{corollary}

\begin{proof}
In each of the three Markov schedules,
\[
  \sum_m d_m^{-A_0}<\infty,
  \qquad
  \sum_m\mathfrak b_m(a,R)<\infty .
\]
For \textup{(i)} this is the pair of \(p\)-series assumptions \(A_0>1\) and
\(B>1\).  For \textup{(ii)}, \(d_m^{-A_0}\le C m^{-A_0\kappa}\) and
\(\mathfrak b_m(a,R)\le C m^{-B\kappa}\).  For \textup{(iii)},
\(\sum_m d_m^{-A_0}<\infty\), while
\[
  \mathfrak b_m(a,R)\le C\ell_m^{-\eta}
  \le C' m^{-\kappa\eta},
\]
which is summable because \(\kappa\eta>1\).  The displayed average
inverse-moment condition implies \(\mathfrak b_m(a,R)\le C\ell_m^{1-B}\), so it
is the special case \(\eta=B-1\).

Now \(u_m=q\ell_m\) gives \(e^{-u_m}\le d_m^{-q}\), up to changing finitely many
terms.  For fixed \(\rho\), take \(S=2q\) and
\(r_m=\lceil\rho M_m\rceil\).  The standard binomial estimate gives
\[
  \binom{M_m}{r_m}d_m^{-S r_m}
  \le d_m^{-q}
\]
for all large \(m\).  For the vanishing-fraction assertion, take
\(r_m=\lceil\rho_mM_m\rceil\) and \(S=q+\tau\).  Since
\(\rho_m\ge e d_m^{-\tau}\),
\[
  \binom{M_m}{r_m}d_m^{-S r_m}
  \le d_m^{-q}.
\]
For polynomial batches, take \(r_m=1\) and \(S=\theta+q\), so
\[
  \binom{M_m}{1}d_m^{-S}\le d_m^{-q}.
\]
Thus in all three cases the summability hypothesis of
Corollary~\ref{cor:heterogeneous-invmoment-operator-batch-maxima} holds.  In the
polynomial-batch case,
\[
  \frac{\log(eM_m)+u_m}{\ell_m}=O_{\theta,q}(1),
\]
which gives the displayed \(O(\sqrt{n_m})\) operator maximum.

The Bernstein statement is the same invocation of
Corollary~\ref{cor:heterogeneous-invmoment-operator-batch-maxima} with its
Bernstein matrix-budget replacement.  The all-integer lower bound gives
\[
  \exp\left(-c\frac{n(d)}{\log^2(e+d)}\right)
  \le (e+d)^{-cQ},
\]
which is summable when \(Q>1/c\).  The polynomial-schedule lower bound gives
\[
  \exp\left(-c\frac{n_m}{\ell_m^2}\right)
  \le (e+m)^{-cQ},
\]
and \(Q>1/c\) makes this summable; \(A_0\kappa>1\) handles the
\(d_m^{-A_0}\) term.  On stretched-exponential schedules the same
\(\ell_m^2\log(e+m)\) lower bound gives the same summable power of \(m\), while
\(n_m\ge Q\ell_m^{2+\tau}\) gives
\[
  \exp\left(-cQ\ell_m^\tau\right)\le \exp(-c'm^{\kappa\tau}),
\]
which is summable.
\end{proof}

\begin{corollary}[Concrete heterogeneous inverse-\(p\) batch-maximal schedules]
\label{cor:heterogeneous-invp-operator-batch-maxima-schedules}
In the setting and notation of
Corollary~\ref{cor:heterogeneous-invmoment-operator-batch-maxima}, fix
\(p>0\) and set
\[
  \sigma_p:=1+\frac{(2-p)_+}{2},
  \qquad
  \mathfrak i_m^{(p)}(a,R):=
  \frac{\ell_m^{\sigma_p}\mathcal I_m^{(p)}(a)}{R n_m},
\]
where \(\mathcal I_m^{(p)}(a)\) is the aggregate inverse-\(p\) lower-normalizer
mass from Corollary~\ref{cor:heterogeneous-invmoment-operator-schedules}.
Assume one of the following inverse-\(p\) Markov schedules.
\begin{enumerate}[label=\textup{(\roman*)},leftmargin=2.2em]
\item \(d_m=m\), \(A_0>1\), and
\[
  \mathfrak i_m^{(p)}(a,R)\le C_0m^{-B}
  \qquad\text{eventually for some }B>1.
\]
It is enough, more concretely, that
\[
  \frac{\mathcal I_m^{(p)}(a)}{n_m}\le C_1m^{-B}
  \qquad\text{eventually for some }B>1.
\]
\item \(d_m\ge c_0m^\kappa\) for some \(c_0,\kappa>0\),
\[
  A_0\kappa>1,
  \qquad
  \mathfrak i_m^{(p)}(a,R)\le C_0d_m^{-B}
  \qquad\text{eventually for some }B\text{ with }B\kappa>1.
\]
It is enough that
\[
  \frac{\mathcal I_m^{(p)}(a)}{n_m}\le C_1d_m^{-B}
  \qquad\text{eventually for some }B\text{ with }B\kappa>1.
\]
\item \(d_m\ge\exp(c_0m^\kappa)\) for some \(c_0,\kappa>0\), \(A_0>0\), and
\[
  \mathfrak i_m^{(p)}(a,R)\le C_0\ell_m^{-\eta}
  \qquad\text{eventually for some }\eta>1/\kappa .
\]
It is enough that
\[
  \frac{\mathcal I_m^{(p)}(a)}{n_m}\le C_1\ell_m^{-B}
  \qquad\text{eventually for some }B>\sigma_p+1/\kappa .
\]
\end{enumerate}
Let \(q>0\) satisfy \(\sum_m d_m^{-q}<\infty\), equivalently \(q>1\) in
\textup{(i)}, \(q\kappa>1\) in \textup{(ii)}, and any \(q>0\) in
\textup{(iii)}.  Set \(u_m=q\ell_m\).  Then the three sampled-batch conclusions
of Corollary~\ref{cor:heterogeneous-invmoment-operator-batch-maxima-schedules}
hold: arbitrary sampled batches have eventual fixed robust fractions, arbitrary
sampled batches have eventual vanishing discarded fractions whenever
\(\rho_m\ge e d_m^{-\tau}\), and polynomial sampled batches
\(M_m\le d_m^\theta\) are eventually entirely Toeplitz/Hankel operator-norm
robust with whole-list operator maximum \(O_{\theta,q}(\sqrt{n_m})\).
\end{corollary}

\begin{proof}
The inverse-\(p\) replacement in
Corollary~\ref{cor:heterogeneous-invmoment-operator-batch-maxima} makes
\(\mathfrak i_m^{(p)}(a,R)\) the matrix term in the Markov summability
condition.  Hence it is enough to verify
\[
  \sum_m d_m^{-A_0}<\infty,
  \qquad
  \sum_m\mathfrak i_m^{(p)}(a,R)<\infty .
\]
The displayed \(\mathfrak i_m^{(p)}\)-assumptions give this exactly as in
Corollary~\ref{cor:heterogeneous-invmoment-operator-batch-maxima-schedules}.
The concrete average inverse-\(p\) assumptions also give it: on all integer
dimensions,
\(\mathfrak i_m^{(p)}(a,R)\le C\ell_m^{\sigma_p}m^{-B}\), which is summable for
\(B>1\); on polynomial schedules,
\(\mathfrak i_m^{(p)}(a,R)\le C\ell_m^{\sigma_p}m^{-B\kappa}\), which is
summable for \(B\kappa>1\); and on stretched-exponential schedules,
\[
  \mathfrak i_m^{(p)}(a,R)\le C\ell_m^{\sigma_p-B}
  \le C'm^{-\kappa(B-\sigma_p)},
\]
which is summable when \(B>\sigma_p+1/\kappa\).

The sign-tail choices are identical to the proof of
Corollary~\ref{cor:heterogeneous-invmoment-operator-batch-maxima-schedules}:
with \(u_m=q\ell_m\), \(e^{-u_m}\le d_m^{-q}\) up to finitely many terms; use
\(S=2q\) and \(r_m=\lceil\rho M_m\rceil\) for fixed fractions,
\(S=q+\tau\) and \(r_m=\lceil\rho_mM_m\rceil\) for vanishing fractions, and
\(S=\theta+q\), \(r_m=1\), for polynomial batches.  The summability hypothesis
of Corollary~\ref{cor:heterogeneous-invmoment-operator-batch-maxima} follows in
each case, and the polynomial-batch entropy factor is \(O_{\theta,q}(1)\).
\end{proof}

\begin{corollary}[Random-environment inverse-\(p\) batch-maximal schedules]
\label{cor:random-environment-invp-operator-batch-maxima-schedules}
Let \((d_m,n_m)\) be deterministic and satisfy
\[
  n_m\le d_m,\qquad d_m/n_m\to\infty,\qquad \log d_m=o(n_m),
\]
and put \(\ell_m=\log(e+d_m)\).  Let \(\mathcal G\) be an environment
\(\sigma\)-field.  For each \(m,i\), let the coefficient law
\(\mu_{m,i}\) be \(\mathcal G\)-measurable; assume almost surely that all
selected laws are centered, variance one, and have sub-Gaussian norm at most
\(K\).  Conditional on \(\mathcal G\), draw the Toeplitz matrices
\[
  B_{m,i}=T_{d_m}(\xi^{(m,i)}),\qquad 1\le i\le n_m,
\]
independently from the selected laws, and, for deterministic \(M_m\ge1\), draw
the sampled sign lists independently of the matrices.  Define \(N_{m,i}\),
\(X_{m,i}\), and
\(\mathcal I_m^{(p)}(a)\) conditionally from the selected laws as in
Corollary~\ref{cor:heterogeneous-invmoment-operator-schedules}.

Fix \(a>0\), \(R\ge1\), and \(p>0\), and set
\[
  \sigma_p:=1+\frac{(2-p)_+}{2}.
\]
Suppose there are deterministic \(\alpha_m\ge0\) such that
\[
  \sum_m
  \Prob\left\{
    \frac{\mathcal I_m^{(p)}(a)}{n_m}>\alpha_m
  \right\}<\infty .
\]
Assume, in addition, one of the following deterministic schedules:
\begin{enumerate}[label=\textup{(\roman*)},leftmargin=2.2em]
\item \(d_m=m\), \(A_0>1\), and
\[
  \alpha_m\le C_0m^{-B}
  \qquad\text{eventually for some }B>1.
\]
\item \(d_m\ge c_0m^\kappa\) for some \(c_0,\kappa>0\),
\[
  A_0\kappa>1,
  \qquad
  \alpha_m\le C_0d_m^{-B}
  \qquad\text{eventually for some }B\text{ with }B\kappa>1.
\]
\item \(d_m\ge\exp(c_0m^\kappa)\) for some \(c_0,\kappa>0\), \(A_0>0\), and
\[
  \alpha_m\le C_0\ell_m^{-B}
  \qquad\text{eventually for some }B>\sigma_p+1/\kappa .
\]
\end{enumerate}
Then, for every \(q>0\) with \(\sum_m d_m^{-q}<\infty\) and with
\(u_m=q\ell_m\), the three sampled-batch conclusions of
Corollary~\ref{cor:heterogeneous-invp-operator-batch-maxima-schedules} hold
almost surely over the joint environment, matrix, and sampled-sign randomness.
Thus random law selection and random sparse contamination are allowed whenever
their empirical aggregate inverse-\(p\) lower-normalizer budget misses the
corresponding deterministic schedule only on a summable set of dimensions.
\end{corollary}

\begin{proof}
By Borel--Cantelli,
\[
  \frac{\mathcal I_m^{(p)}(a)}{n_m}\le \alpha_m
\]
for all but finitely many \(m\), on an environment event of probability one.
Fix an environment in this full-measure event.  The selected coefficient laws
are now deterministic, uniformly centered, variance one, and \(K\)-sub-Gaussian,
and the displayed deterministic bounds on \(\alpha_m\) put the environment in
one of the three average inverse-\(p\) schedules of
Corollary~\ref{cor:heterogeneous-invp-operator-batch-maxima-schedules}.  That
corollary therefore gives the stated fixed-fraction, vanishing-fraction, and
polynomial-batch conclusions with conditional probability one over the matrices
and sampled signs.  Integrating these conditional probability-one statements
over the full-measure environment event gives the joint almost-sure conclusion.
\end{proof}

\begin{corollary}[Expectation criteria for random-environment inverse-\(p\) schedules]
\label{cor:random-environment-invp-expectation-schedules}
In the setting and notation of
Corollary~\ref{cor:random-environment-invp-operator-batch-maxima-schedules},
put
\[
  \beta_m:=\Exp\left[\frac{\mathcal I_m^{(p)}(a)}{n_m}\right],
  \qquad
  \sigma_p:=1+\frac{(2-p)_+}{2}.
\]
Assume one of the following expectation schedules.
\begin{enumerate}[label=\textup{(\roman*)},leftmargin=2.2em]
\item \(d_m=m\), \(A_0>1\), and
\[
  \sum_m m^B\beta_m<\infty
  \qquad\text{for some }B>1.
\]
In particular, \(\beta_m=O(m^{-B_*})\) with \(B_*>2\) suffices.
\item \(d_m\ge c_0m^\kappa\) for some \(c_0,\kappa>0\), \(A_0\kappa>1\), and
\[
  \sum_m d_m^B\beta_m<\infty
  \qquad\text{for some }B\text{ with }B\kappa>1.
\]
In particular, \(\beta_m=O(d_m^{-B_*})\) with \(B_*\kappa>2\) suffices.
\item \(d_m\ge\exp(c_0m^\kappa)\) for some \(c_0,\kappa>0\), \(A_0>0\), and
\[
  \sum_m \ell_m^B\beta_m<\infty
  \qquad\text{for some }B>\sigma_p+1/\kappa .
\]
In particular, \(\beta_m=O(\ell_m^{-B_*})\) with
\[
  B_*>\sigma_p+2/\kappa
\]
suffices.
\end{enumerate}
Then, for every \(q>0\) with \(\sum_m d_m^{-q}<\infty\) and with
\(u_m=q\ell_m\), the fixed-fraction, vanishing-fraction, and all-robust
polynomial-batch conclusions of
Corollary~\ref{cor:random-environment-invp-operator-batch-maxima-schedules}
hold almost surely.

More concretely, suppose the environment selects, for each \(m,i\), a random
catalog index \(\Lambda_{m,i}\) in a countable set \(\mathcal J_{m,i}\),
possibly with arbitrary dependence among all indices, and then sets
\(\mu_{m,i}=\mu_{m,i,\Lambda_{m,i}}\).  Assume every catalog law is centered,
variance one, and \(K\)-sub-Gaussian, and write
\[
  \pi_{m,i,j}:=\Prob\{\Lambda_{m,i}=j\}.
\]
For
\[
  \iota_{m,i,j}^{(p)}(a):=
  \Exp_{\mu_{m,i,j}}\left[
  \mathbf 1_{\{X_{m,i,j}<aN_{m,i,j}\}}
  \left(\frac{aN_{m,i,j}}{X_{m,i,j}}\right)^p
  \right],
\]
where \(N_{m,i,j}=\Exp_{\mu_{m,i,j}}X_{m,i,j}\) and
\(X_{m,i,j}=\norm{T_{d_m}(\mu_{m,i,j})}_{\op}\) for an independent draw from
\(\mu_{m,i,j}\), one has
\[
  \beta_m=
  \frac1{n_m}
  \sum_{i=1}^{n_m}\sum_{j\in\mathcal J_{m,i}}
  \pi_{m,i,j}\,\iota_{m,i,j}^{(p)}(a).
\]
Thus the preceding tests use only marginal catalog weights; no independence
among the environment selections is required.
\end{corollary}

\begin{proof}
For the first assertion set \(\alpha_m=m^{-B}\), \(d_m^{-B}\), or
\(\ell_m^{-B}\) in the three cases.  Markov's inequality gives
\[
  \Prob\left\{
    \frac{\mathcal I_m^{(p)}(a)}{n_m}>\alpha_m
  \right\}
  \le \frac{\beta_m}{\alpha_m},
\]
so the displayed expectation summability hypotheses imply the summable
environment-tail hypothesis of
Corollary~\ref{cor:random-environment-invp-operator-batch-maxima-schedules}.
The deterministic bounds on \(\alpha_m\) are exactly the three schedules in
that corollary, and the stated sampled-batch conclusions follow.

The power-decay sufficient conditions are obtained by leaving one summable
power for Markov.  On all integer dimensions choose
\(1<B<B_*-1\).  On polynomial schedules choose
\[
  1/\kappa<B<B_*-1/\kappa,
\]
which is possible when \(B_*\kappa>2\).  On stretched-exponential schedules
choose
\[
  \sigma_p+1/\kappa<B<B_*-1/\kappa,
\]
which is possible when \(B_*>\sigma_p+2/\kappa\).  Since
\(\ell_m\ge c m^\kappa\) in the last case, each displayed expectation series is
finite.

For the catalog formula, conditional on a realization of the selected indices,
\(\mathcal I_m^{(p)}(a)=\sum_i\iota_{m,i,\Lambda_{m,i}}^{(p)}(a)\).  Taking
expectations and using nonnegativity gives the displayed identity by linearity
and monotone convergence.  This uses only the one-dimensional marginal
probabilities \(\pi_{m,i,j}\).
\end{proof}

\begin{corollary}[Bernstein criteria for independent random-catalog schedules]
\label{cor:random-environment-invp-bernstein-catalog-schedules}
Work in the countable catalog setting of
Corollary~\ref{cor:random-environment-invp-expectation-schedules}.  For each
\(m\), assume that the catalog indices
\(\Lambda_{m,1},\ldots,\Lambda_{m,n_m}\) are independent, and put
\[
  Z_{m,i}:=\iota_{m,i,\Lambda_{m,i}}^{(p)}(a),
  \qquad
  \beta_m:=\frac1{n_m}\sum_{i=1}^{n_m}\Exp Z_{m,i}.
\]
Suppose that for deterministic \(L_m>0\) and \(\alpha_m>0\),
\[
  0\le Z_{m,i}\le L_m\quad\text{almost surely for all }i,
  \qquad
  \beta_m\le \frac{\alpha_m}{2}
  \qquad\text{eventually},
\]
and
\[
  \sum_m \exp\left(-c_{\rm B}\frac{n_m\alpha_m}{L_m}\right)<\infty
\]
for a sufficiently small absolute constant \(c_{\rm B}>0\).  If, in addition,
\(\alpha_m\) satisfies one of the three deterministic inverse-\(p\) schedules in
Corollary~\ref{cor:random-environment-invp-operator-batch-maxima-schedules},
then, for every \(q>0\) with \(\sum_m d_m^{-q}<\infty\) and \(u_m=q\ell_m\),
the fixed-fraction, vanishing-fraction, and all-robust polynomial-batch
conclusions of that corollary hold almost surely.

In particular, for polynomial schedules \(d_m\ge c_0m^\kappa\), it is enough to
have, for some \(B\) with \(B\kappa>1\),
\[
  \beta_m\le c_1d_m^{-B},
  \qquad
  n_m\ge Q L_m d_m^B\log(e+m)
\]
eventually, with \(Q\) sufficiently large.  For stretched-exponential schedules
\(d_m\ge\exp(c_0m^\kappa)\), it is enough to have, for some
\(B>1+(2-p)_+/2+1/\kappa\),
\[
  \beta_m\le c_1\ell_m^{-B},
  \qquad
  n_m\ge Q L_m\ell_m^B\log(e+m)
\]
eventually, again with \(Q\) sufficiently large.  The same conclusion also
holds if the logarithmic lower bound is replaced by any stronger condition
making \(\sum_m\exp(-c_{\rm B}n_m\alpha_m/L_m)\) finite.
\end{corollary}

\begin{proof}
For each \(m\), the selected environment budget is
\[
  \frac{\mathcal I_m^{(p)}(a)}{n_m}
  =\frac1{n_m}\sum_{i=1}^{n_m}Z_{m,i}.
\]
The variables \(Z_{m,i}\) are independent, nonnegative, and bounded by \(L_m\).
Their total mean is \(n_m\beta_m\), and
\[
  \sum_i \operatorname{Var} Z_{m,i}\le \sum_i \Exp Z_{m,i}^2
  \le L_m\sum_i\Exp Z_{m,i}
  =L_m n_m\beta_m .
\]
If \(\beta_m\le\alpha_m/2\), then Bernstein's inequality gives
\[
\begin{aligned}
  \Prob\left\{\frac1{n_m}\sum_i Z_{m,i}>\alpha_m\right\}
  &\le
  \exp\left(
  -\frac{n_m^2(\alpha_m-\beta_m)^2}
  {2L_m n_m\beta_m+\frac23L_m n_m(\alpha_m-\beta_m)}
  \right)  \\
  &\le
  \exp\left(-c_{\rm B}\frac{n_m\alpha_m}{L_m}\right)
\end{aligned}
\]
after reducing \(c_{\rm B}\), if necessary.  The assumed exponential series is
therefore exactly the summable environment-tail hypothesis of
Corollary~\ref{cor:random-environment-invp-operator-batch-maxima-schedules}.
Applying that corollary with the given deterministic schedule for
\(\alpha_m\) proves the first assertion.

For the polynomial and stretched-exponential displays, take
\(\alpha_m=2c_1d_m^{-B}\) or \(\alpha_m=2c_1\ell_m^{-B}\), respectively,
enlarging constants harmlessly.  The lower bounds on \(n_m\) make the
Bernstein tail at most a summable power of \(e+m\), while the displayed
exponents are exactly the deterministic inverse-\(p\) schedule thresholds.
\end{proof}

\begin{corollary}[Freedman criteria for predictable random-catalog schedules]
\label{cor:random-environment-invp-freedman-catalog-schedules}
Work in the countable catalog setting of
Corollary~\ref{cor:random-environment-invp-expectation-schedules}.  For each
\(m\), let
\[
  \mathcal F_{m,0}\subset\mathcal F_{m,1}\subset\cdots
  \subset\mathcal F_{m,n_m}
\]
be a filtration such that \(\Lambda_{m,i}\) is \(\mathcal F_{m,i}\)-measurable,
and put
\[
  Z_{m,i}:=\iota_{m,i,\Lambda_{m,i}}^{(p)}(a).
\]
Assume that for deterministic \(L_m>0\), \(\alpha_m>0\), and deterministic
envelopes \(b_{m,i}\ge0\),
\[
  0\le Z_{m,i}\le L_m\quad\text{almost surely},
  \qquad
  \Exp[Z_{m,i}\mid\mathcal F_{m,i-1}]\le b_{m,i}
\]
for all \(i\), and
\[
  \frac1{n_m}\sum_{i=1}^{n_m}b_{m,i}\le \frac{\alpha_m}{2}
  \qquad\text{eventually}.
\]
If
\[
  \sum_m \exp\left(-c_{\rm F}\frac{n_m\alpha_m}{L_m}\right)<\infty
\]
for a sufficiently small absolute constant \(c_{\rm F}>0\), and if
\(\alpha_m\) satisfies one of the three deterministic inverse-\(p\) schedules in
Corollary~\ref{cor:random-environment-invp-operator-batch-maxima-schedules},
then, for every \(q>0\) with \(\sum_m d_m^{-q}<\infty\) and \(u_m=q\ell_m\),
the fixed-fraction, vanishing-fraction, and all-robust polynomial-batch
conclusions of that corollary hold almost surely.

In particular, the concrete polynomial and stretched-exponential schedule
displays of
Corollary~\ref{cor:random-environment-invp-bernstein-catalog-schedules} remain
valid with \(\beta_m\) replaced by
\[
  \bar\beta_m:=\frac1{n_m}\sum_{i=1}^{n_m}b_{m,i}.
\]
Thus adaptive, sequential, or without-replacement law-catalog exposure is
covered whenever its one-step predictable inverse-\(p\) budget obeys the same
bounded-catalog schedule.
\end{corollary}

\begin{proof}
Let
\[
  \mu_{m,i}:=\Exp[Z_{m,i}\mid\mathcal F_{m,i-1}],
  \qquad
  D_{m,i}:=Z_{m,i}-\mu_{m,i}.
\]
Then \(D_{m,i}\) is a martingale difference and \(D_{m,i}\le L_m\).  Its
predictable quadratic variation satisfies
\[
  \sum_i\Exp[D_{m,i}^2\mid\mathcal F_{m,i-1}]
  \le \sum_i\Exp[Z_{m,i}^2\mid\mathcal F_{m,i-1}]
  \le L_m\sum_i\mu_{m,i}
  \le \frac12 L_m n_m\alpha_m
\]
eventually.  On the same event, if
\[
  \frac1{n_m}\sum_i Z_{m,i}>\alpha_m,
\]
then
\[
  \sum_i D_{m,i}>
  n_m\alpha_m-\sum_i\mu_{m,i}
  \ge \frac12 n_m\alpha_m .
\]
Freedman's inequality therefore gives
\[
  \Prob\left\{\frac1{n_m}\sum_i Z_{m,i}>\alpha_m\right\}
  \le
  \exp\left(-c_{\rm F}\frac{n_m\alpha_m}{L_m}\right)
\]
after reducing \(c_{\rm F}\), if necessary.  The assumed exponential series
gives the summable environment-tail hypothesis of
Corollary~\ref{cor:random-environment-invp-operator-batch-maxima-schedules}.
The deterministic inverse-\(p\) schedule for \(\alpha_m\) then gives the
sampled-batch conclusions.  The concrete schedules are the same substitutions
as in Corollary~\ref{cor:random-environment-invp-bernstein-catalog-schedules}.
\end{proof}

\begin{corollary}[Truncated Freedman criteria for predictable random-catalog schedules]
\label{cor:random-environment-invp-truncated-freedman-catalog-schedules}
Work in the countable catalog setting of
Corollary~\ref{cor:random-environment-invp-expectation-schedules}, with the
filtration convention of
Corollary~\ref{cor:random-environment-invp-freedman-catalog-schedules}.  Put
\[
  Z_{m,i}:=\iota_{m,i,\Lambda_{m,i}}^{(p)}(a).
\]
Assume \(Z_{m,i}<\infty\) almost surely.  Let \(L_m>0\), \(\alpha_m>0\), and
deterministic envelopes \(b_{m,i},r_{m,i}\ge0\) satisfy
\[
  \Exp[Z_{m,i}\wedge L_m\mid\mathcal F_{m,i-1}]\le b_{m,i},
  \qquad
  \Exp[(Z_{m,i}-L_m)_+\mid\mathcal F_{m,i-1}]\le r_{m,i}
\]
for all \(i\).  Write
\[
  \bar b_m:=\frac1{n_m}\sum_{i=1}^{n_m}b_{m,i},
  \qquad
  \bar r_m:=\frac1{n_m}\sum_{i=1}^{n_m}r_{m,i}.
\]
Suppose that \(\bar b_m\le\alpha_m/4\) eventually and that
\[
  \sum_m
  \left[
    \exp\left(-c_{\rm TF}\frac{n_m\alpha_m}{L_m}\right)
    +\frac{\bar r_m}{\alpha_m}
  \right]<\infty
\]
for a sufficiently small absolute constant \(c_{\rm TF}>0\).  If
\(\alpha_m\) satisfies one of the three deterministic inverse-\(p\) schedules in
Corollary~\ref{cor:random-environment-invp-operator-batch-maxima-schedules},
then, for every \(q>0\) with \(\sum_m d_m^{-q}<\infty\) and \(u_m=q\ell_m\),
the fixed-fraction, vanishing-fraction, and all-robust polynomial-batch
conclusions of that corollary hold almost surely.

In particular, the concrete polynomial and stretched-exponential displays of
Corollary~\ref{cor:random-environment-invp-bernstein-catalog-schedules} remain
valid for unbounded predictable catalogs with \(\beta_m\) replaced by the
truncated predictable budget \(\bar b_m\), the half-budget condition replaced by
\(\bar b_m\le\alpha_m/4\), and the added excess summability
\(\sum_m\bar r_m/\alpha_m<\infty\).  Thus, for example, on polynomial schedules
\(d_m\ge c_0m^\kappa\), it is enough to have
\[
  \bar b_m\le c_1d_m^{-B},\qquad
  n_m\ge QL_md_m^B\log(e+m),\qquad
  \sum_m d_m^B\bar r_m<\infty
\]
for some \(B\) with \(B\kappa>1\), and on stretched-exponential schedules
\(d_m\ge\exp(c_0m^\kappa)\), it is enough to have
\[
  \bar b_m\le c_1\ell_m^{-B},\qquad
  n_m\ge QL_m\ell_m^B\log(e+m),\qquad
  \sum_m \ell_m^B\bar r_m<\infty
\]
for some \(B>1+(2-p)_+/2+1/\kappa\), with \(Q\) sufficiently large.
\end{corollary}

\begin{proof}
Let
\[
  W_{m,i}:=Z_{m,i}\wedge L_m,\qquad
  R_{m,i}:=(Z_{m,i}-L_m)_+,
\]
and put \(\mu_{m,i}:=\Exp[W_{m,i}\mid\mathcal F_{m,i-1}]\).  The martingale
differences \(D_{m,i}:=W_{m,i}-\mu_{m,i}\) satisfy \(D_{m,i}\le L_m\), and,
eventually,
\[
  \sum_i\Exp[D_{m,i}^2\mid\mathcal F_{m,i-1}]
  \le \sum_i\Exp[W_{m,i}^2\mid\mathcal F_{m,i-1}]
  \le L_m\sum_i\mu_{m,i}
  \le \frac14 L_m n_m\alpha_m .
\]
If \(n_m^{-1}\sum_i W_{m,i}>\alpha_m/2\), then
\[
  \sum_iD_{m,i}
  >
  \frac12 n_m\alpha_m-\sum_i\mu_{m,i}
  \ge \frac14 n_m\alpha_m .
\]
Freedman's inequality therefore gives
\[
  \Prob\left\{\frac1{n_m}\sum_i W_{m,i}>\frac{\alpha_m}{2}\right\}
  \le
  \exp\left(-c_{\rm TF}\frac{n_m\alpha_m}{L_m}\right)
\]
after reducing \(c_{\rm TF}\), if necessary.  Since \(Z_{m,i}=W_{m,i}+R_{m,i}\),
Markov's inequality gives
\[
\begin{aligned}
  \Prob\left\{\frac1{n_m}\sum_i Z_{m,i}>\alpha_m\right\}
  &\le
  \Prob\left\{\frac1{n_m}\sum_i W_{m,i}>\frac{\alpha_m}{2}\right\}  \\
  &\quad+
  \Prob\left\{\frac1{n_m}\sum_i R_{m,i}>\frac{\alpha_m}{2}\right\}  \\
  &\le
  \exp\left(-c_{\rm TF}\frac{n_m\alpha_m}{L_m}\right)
  +\frac{2\bar r_m}{\alpha_m}.
\end{aligned}
\]
The displayed summability is therefore the summable environment-tail hypothesis
of Corollary~\ref{cor:random-environment-invp-operator-batch-maxima-schedules}.
The deterministic inverse-\(p\) schedule for \(\alpha_m\) gives the sampled
batch conclusions.  The polynomial and stretched-exponential criteria are the
same substitutions as in the bounded Freedman corollary, with the displayed
excess series added.
\end{proof}

\begin{corollary}[Conditional-moment criteria for unbounded predictable catalogs]
\label{cor:random-environment-invp-moment-freedman-catalog-schedules}
Work in the setting of
Corollary~\ref{cor:random-environment-invp-truncated-freedman-catalog-schedules}.
Fix \(s>1\).  Assume that, for deterministic \(b_{m,i},h_{m,i}\ge0\),
\[
  \Exp[Z_{m,i}\mid\mathcal F_{m,i-1}]\le b_{m,i},
  \qquad
  \Exp[Z_{m,i}^s\mid\mathcal F_{m,i-1}]\le h_{m,i}
\]
for all \(i\), and put
\[
  \bar b_m:=\frac1{n_m}\sum_{i=1}^{n_m}b_{m,i},
  \qquad
  \bar h_m:=\frac1{n_m}\sum_{i=1}^{n_m}h_{m,i}.
\]
Let \(L_m>0\) and \(\alpha_m>0\) be deterministic.  Suppose that
\[
  \bar b_m\le \frac{\alpha_m}{4}
  \qquad\text{eventually},
\]
and that
\[
  \sum_m
  \left[
    \exp\left(-c_{\rm M}\frac{n_m\alpha_m}{L_m}\right)
    +\frac{\bar h_m}{L_m^{s-1}\alpha_m}
  \right]<\infty
\]
for a sufficiently small absolute constant \(c_{\rm M}>0\).  If
\(\alpha_m\) satisfies one of the three deterministic inverse-\(p\) schedules in
Corollary~\ref{cor:random-environment-invp-operator-batch-maxima-schedules},
then, for every \(q>0\) with \(\sum_m d_m^{-q}<\infty\) and \(u_m=q\ell_m\),
the fixed-fraction, vanishing-fraction, and all-robust polynomial-batch
conclusions of that corollary hold almost surely.

In concrete schedules, this gives a direct tail-moment substitute for the
bounded-catalog hypotheses.  Along all integer dimensions \(d_m=m\), with
\(A_0>1\), it is enough to have, for some \(B>1\),
\[
  \bar b_m\le c_1m^{-B},
  \qquad
  n_m\ge QL_mm^B\log(e+m),
  \qquad
  \sum_m\frac{m^B\bar h_m}{L_m^{s-1}}<\infty .
\]
Along polynomial schedules \(d_m\ge c_0m^\kappa\), with \(A_0\kappa>1\), it is
enough to have, for some \(B\) with \(B\kappa>1\),
\[
  \bar b_m\le c_1d_m^{-B},
  \qquad
  n_m\ge QL_md_m^B\log(e+m),
  \qquad
  \sum_m\frac{d_m^B\bar h_m}{L_m^{s-1}}<\infty .
\]
Along stretched-exponential schedules \(d_m\ge\exp(c_0m^\kappa)\), it is enough
to have, for some \(B>1+(2-p)_+/2+1/\kappa\),
\[
  \bar b_m\le c_1\ell_m^{-B},
  \qquad
  n_m\ge QL_m\ell_m^B\log(e+m),
  \qquad
  \sum_m\frac{\ell_m^B\bar h_m}{L_m^{s-1}}<\infty ,
\]
with \(Q\) sufficiently large in each display.
\end{corollary}

\begin{proof}
For \(W_{m,i}=Z_{m,i}\wedge L_m\), the first moment assumption gives
\[
  \Exp[W_{m,i}\mid\mathcal F_{m,i-1}]
  \le
  \Exp[Z_{m,i}\mid\mathcal F_{m,i-1}]
  \le b_{m,i}.
\]
Also, for \(z\ge0\),
\[
  (z-L_m)_+\le L_m^{1-s}z^s .
\]
Hence
\[
  \Exp[(Z_{m,i}-L_m)_+\mid\mathcal F_{m,i-1}]
  \le L_m^{1-s}h_{m,i}.
\]
Corollary~\ref{cor:random-environment-invp-truncated-freedman-catalog-schedules}
therefore applies with excess envelope
\(r_{m,i}=L_m^{1-s}h_{m,i}\).  Its summability condition is precisely the
displayed exponential series plus
\(\sum_m\bar h_m/(L_m^{s-1}\alpha_m)<\infty\), and it gives the sampled-batch
conclusions.

The concrete schedule displays follow by taking
\(\alpha_m=4c_1m^{-B}\), \(4c_1d_m^{-B}\), or \(4c_1\ell_m^{-B}\), respectively,
and increasing \(Q\), if necessary, so that the exponential Freedman series is
summable.  The remaining displayed series is exactly the excess-tail series in
the three cases.
\end{proof}

\begin{corollary}[Automatic truncation for moment-predictable catalogs]
\label{cor:random-environment-invp-auto-truncation-freedman-catalog-schedules}
Work in the setting of
Corollary~\ref{cor:random-environment-invp-moment-freedman-catalog-schedules}.
Let \(\lambda_m>0\) be deterministic and suppose that
\[
  \sum_m \exp(-c_{\rm AT}\lambda_m)<\infty
\]
for a sufficiently small absolute constant \(c_{\rm AT}>0\).  Put
\[
  L_m:=\frac{n_m\alpha_m}{\lambda_m}.
\]
If
\[
  \bar b_m\le\frac{\alpha_m}{4}
  \qquad\text{eventually},
  \qquad
  \sum_m
  \frac{\lambda_m^{s-1}\bar h_m}{n_m^{s-1}\alpha_m^s}<\infty ,
\]
and if \(\alpha_m\) satisfies one of the three deterministic inverse-\(p\)
schedules in
Corollary~\ref{cor:random-environment-invp-operator-batch-maxima-schedules},
then, for every \(q>0\) with \(\sum_m d_m^{-q}<\infty\) and \(u_m=q\ell_m\),
the fixed-fraction, vanishing-fraction, and all-robust polynomial-batch
conclusions of that corollary hold almost surely.

In particular, choose \(\lambda_m=A\log(e+m)\) with \(A\) sufficiently large.
Then the same conclusions hold in each of the following concrete regimes.
\begin{enumerate}[label=\textup{(\roman*)},leftmargin=2.2em]
\item \(d_m=m\), \(A_0>1\), \(\bar b_m\le c_1m^{-B}\) for some \(B>1\), and
\[
  \sum_m
  \frac{m^{Bs}(\log(e+m))^{s-1}\bar h_m}{n_m^{s-1}}<\infty .
\]
\item \(d_m\ge c_0m^\kappa\), \(A_0\kappa>1\),
\(\bar b_m\le c_1d_m^{-B}\) for some \(B\) with \(B\kappa>1\), and
\[
  \sum_m
  \frac{d_m^{Bs}(\log(e+m))^{s-1}\bar h_m}{n_m^{s-1}}<\infty .
\]
\item \(d_m\ge\exp(c_0m^\kappa)\), \(A_0>0\),
\(\bar b_m\le c_1\ell_m^{-B}\) for some
\(B>1+(2-p)_+/2+1/\kappa\), and
\[
  \sum_m
  \frac{\ell_m^{Bs}(\log(e+m))^{s-1}\bar h_m}{n_m^{s-1}}<\infty .
\]
\end{enumerate}
\end{corollary}

\begin{proof}
Apply
Corollary~\ref{cor:random-environment-invp-moment-freedman-catalog-schedules}
with the displayed \(L_m\).  The exponential term is exactly
\(\exp(-c_{\rm AT}\lambda_m)\), after reducing \(c_{\rm AT}\) if necessary, and
\[
  \frac{\bar h_m}{L_m^{s-1}\alpha_m}
  =
  \frac{\lambda_m^{s-1}\bar h_m}{n_m^{s-1}\alpha_m^s}.
\]
Thus the stated hypotheses are precisely the moment-Freedman summability
hypotheses.  The sampled-batch conclusions follow from that corollary.

For the concrete regimes, take \(\alpha_m=4c_1m^{-B}\),
\(4c_1d_m^{-B}\), or \(4c_1\ell_m^{-B}\), respectively.  A large enough fixed
\(A\) makes \(\sum_m(e+m)^{-c_{\rm AT}A}\) summable, and the remaining three
series are the displayed moment series, up to harmless constants.
\end{proof}

\begin{corollary}[Active-menu automatic truncation for predictable catalogs]
\label{cor:random-environment-invp-active-menu-auto-truncation}
Fix finite active menu sizes \(J_m\ge1\).  For each \(1\le r\le J_m\), work in
the predictable countable-catalog setting of
Corollary~\ref{cor:random-environment-invp-freedman-catalog-schedules}, with
possibly \(r\)-dependent catalogs and filtrations, and set
\[
  Z_{m,i}^{(r)}:=\iota_{m,i,\Lambda_{m,i}^{(r)}}^{(p)}(a).
\]
Assume every catalog law in the active menu is centered, variance one, and
\(K\)-sub-Gaussian.  Fix \(s>1\).  Suppose that deterministic envelopes
\(b_{m,i}^{(r)},h_{m,i}^{(r)}\ge0\) satisfy
\[
  \Exp[Z_{m,i}^{(r)}\mid\mathcal F_{m,i-1}^{(r)}]\le b_{m,i}^{(r)},
  \qquad
  \Exp[(Z_{m,i}^{(r)})^s\mid\mathcal F_{m,i-1}^{(r)}]\le h_{m,i}^{(r)}
\]
for all active \(m,i,r\).  Put
\[
  \bar b_m^*:=\max_{1\le r\le J_m}
    \frac1{n_m}\sum_{i=1}^{n_m}b_{m,i}^{(r)},
  \qquad
  H_m:=\sum_{r=1}^{J_m}\frac1{n_m}
    \sum_{i=1}^{n_m}h_{m,i}^{(r)} .
\]
Let \(\alpha_m>0\) satisfy one of the three deterministic inverse-\(p\)
schedules in
Corollary~\ref{cor:random-environment-invp-operator-batch-maxima-schedules}.
If deterministic \(\lambda_m>0\) obey
\[
  \bar b_m^*\le\frac{\alpha_m}{4}
  \qquad\text{eventually}
\]
and
\[
  \sum_m
  \left[
    J_m\exp(-c_{\rm AM}\lambda_m)
    +
    \frac{\lambda_m^{s-1}H_m}{n_m^{s-1}\alpha_m^s}
  \right]<\infty
\]
for a sufficiently small absolute constant \(c_{\rm AM}>0\), then, almost
surely over the environment,
\[
  \max_{1\le r\le J_m}
  \frac1{n_m}\sum_{i=1}^{n_m} Z_{m,i}^{(r)}
  \le \alpha_m
\]
for all sufficiently large \(m\).

Consequently, if the active menu also satisfies the summability needed to run
the matrix and sampled-sign Borel--Cantelli steps uniformly,
\[
  \sum_m J_m d_m^{-A_0}<\infty,
  \qquad
  \sum_m J_m d_m^{-q}<\infty,
\]
and \(u_m=q\ell_m\), then the fixed-fraction, vanishing-fraction, and
all-robust polynomial-batch Toeplitz/Hankel conclusions of
Corollary~\ref{cor:random-environment-invp-operator-batch-maxima-schedules}
hold simultaneously for every active \(r\le J_m\), for all sufficiently large
\(m\), over the joint environment, matrix, and sampled-sign randomness.

In particular, one may take
\[
  \lambda_m=A\log(e+m+J_m)
\]
with \(A\) large enough, in which case the exponential menu term is summable for
every finite \(J_m\).  The remaining automatic-truncation tests are then the
single moment-series conditions
\[
  \sum_m
  \frac{m^{Bs}(\log(e+m+J_m))^{s-1}H_m}{n_m^{s-1}}<\infty
\]
in the all-integer schedule \(d_m=m\), \(\alpha_m\asymp m^{-B}\), \(B>1\);
\[
  \sum_m
  \frac{d_m^{Bs}(\log(e+m+J_m))^{s-1}H_m}{n_m^{s-1}}<\infty
\]
in the polynomial schedule \(d_m\ge c_0m^\kappa\),
\(\alpha_m\asymp d_m^{-B}\), \(B\kappa>1\); and
\[
  \sum_m
  \frac{\ell_m^{Bs}(\log(e+m+J_m))^{s-1}H_m}{n_m^{s-1}}<\infty
\]
in the stretched-exponential schedule
\(d_m\ge\exp(c_0m^\kappa)\), \(\alpha_m\asymp\ell_m^{-B}\),
\(B>1+(2-p)_+/2+1/\kappa\), together with the displayed active-menu
matrix/sign summability for simultaneous sampled certificates.
\end{corollary}

\begin{proof}
For a fixed active pair \((m,r)\), apply
Corollary~\ref{cor:random-environment-invp-auto-truncation-freedman-catalog-schedules}
with the truncation level \(L_m=n_m\alpha_m/\lambda_m\).  Its proof gives the
finite-dimensional tail estimate
\[
  \Prob\left\{
    \frac1{n_m}\sum_{i=1}^{n_m}Z_{m,i}^{(r)}>\alpha_m
  \right\}
  \le
  \exp(-c_{\rm AM}\lambda_m)
  +
  \frac{C_s\lambda_m^{s-1}\bar h_{m,r}}
       {n_m^{s-1}\alpha_m^s},
\]
where
\[
  \bar h_{m,r}:=\frac1{n_m}\sum_{i=1}^{n_m}h_{m,i}^{(r)}.
\]
After decreasing \(c_{\rm AM}\) and absorbing \(C_s\), summing over
\(r\le J_m\) gives
\[
  \Prob\left\{
    \max_{1\le r\le J_m}
    \frac1{n_m}\sum_i Z_{m,i}^{(r)}>\alpha_m
  \right\}
  \le
  J_m e^{-c_{\rm AM}\lambda_m}
  +
  \frac{\lambda_m^{s-1}H_m}{n_m^{s-1}\alpha_m^s}.
\]
The assumed summability and Borel--Cantelli prove the simultaneous environment
budget bound.

On that full-measure environment event, each active catalog row satisfies the
deterministic inverse-\(p\) hypothesis.  Repeating the Borel--Cantelli proof of
Corollary~\ref{cor:heterogeneous-invp-operator-batch-maxima-schedules} over the
finite active menu only multiplies the matrix and sampled-sign tail terms by
\(J_m\); the two displayed summability assumptions make those enlarged series
finite.  The same fixed-fraction, vanishing-fraction, and polynomial-batch
conclusions therefore hold for all active rows eventually.

Finally, since
\[
  J_m(e+m+J_m)^{-c_{\rm AM}A}
  \le (e+m+J_m)^{1-c_{\rm AM}A}
  \le (e+m)^{1-c_{\rm AM}A},
\]
choosing \(A\) with \(c_{\rm AM}A>2\) makes the exponential menu series
summable.  Substituting the three displayed choices of \(\alpha_m\) gives the
three concrete moment series, up to harmless constants.
\end{proof}

\begin{corollary}[Concrete active-prefix automatic-truncation schedules]
\label{cor:random-environment-invp-active-prefix-schedules}
Work in the setting and notation of
Corollary~\ref{cor:random-environment-invp-active-menu-auto-truncation}, and
let
\[
  \sigma_p:=1+\frac{(2-p)_+}{2}.
\]
Assume \(\bar b_m^*\le\alpha_m/4\) eventually.  Then the simultaneous
active-menu sampled Toeplitz/Hankel conclusions of that corollary hold in each
of the following concrete regimes, with \(u_m=q\ell_m\).

\begin{enumerate}[label=\textup{(\roman*)},leftmargin=2.2em]
\item All integer dimensions.  Suppose \(d_m=m\), \(J_m\le C_Jm^\Gamma\), and
choose exponents
\[
  A_0>\Gamma+1,\qquad q>\Gamma+1,\qquad B>1 .
\]
Let \(\alpha_m=c_\alpha m^{-B}\).  It is enough that
\[
  \sum_m
  \frac{m^{Bs}(\log(e+m))^{s-1}H_m}{n_m^{s-1}}<\infty .
\]

\item Polynomial dimension schedules.  Suppose
\[
  d_m\ge c_0m^\kappa,\qquad J_m\le C_Jd_m^\Gamma ,
\]
and choose
\[
  A_0>\Gamma+\frac1\kappa,\qquad
  q>\Gamma+\frac1\kappa,\qquad
  B\kappa>1 .
\]
Let \(\alpha_m=c_\alpha d_m^{-B}\).  It is enough that
\[
  \sum_m
  \frac{d_m^{Bs}\ell_m^{s-1}H_m}{n_m^{s-1}}<\infty .
\]

\item Stretched-exponential dimensions with polynomial-in-\(d\) active
prefixes.  Suppose
\[
  d_m\ge\exp(c_0m^\kappa),\qquad J_m\le C_Jd_m^\Gamma ,
\]
and choose
\[
  A_0>\Gamma,\qquad q>\Gamma,\qquad
  B>\sigma_p+\frac1\kappa .
\]
Let \(\alpha_m=c_\alpha\ell_m^{-B}\).  It is enough that
\[
  \sum_m
  \frac{\ell_m^{Bs+s-1}H_m}{n_m^{s-1}}<\infty .
\]

\item Stretched-exponential dimensions with logarithmic active prefixes.
Suppose
\[
  d_m\ge\exp(c_0m^\kappa),\qquad J_m\le C_J\ell_m^\Gamma ,
\]
and choose
\[
  A_0>0,\qquad q>0,\qquad
  B>\sigma_p+\frac1\kappa .
\]
Let \(\alpha_m=c_\alpha\ell_m^{-B}\).  It is enough that
\[
  \sum_m
  \frac{\ell_m^{Bs}(\log(e+m))^{s-1}H_m}{n_m^{s-1}}<\infty .
\]
\end{enumerate}
Here \(C_J,c_\alpha,c_0,\kappa\) are fixed positive constants and
\(\Gamma\ge0\) is fixed in the regimes where they appear.  In each case the
conclusion is simultaneous for all active \(r\le J_m\), and hence applies to a
countable predictable catalog library whenever level \(m\) only exposes its
first \(J_m\) active streams.
\end{corollary}

\begin{proof}
In every case take
\[
  \lambda_m=A\log(e+m+J_m)
\]
with \(A\) large enough.  The exponential menu term in
Corollary~\ref{cor:random-environment-invp-active-menu-auto-truncation} is then
summable.  It remains to check the active matrix/sign tails and the displayed
moment series.

In \textup{(i)}, \(J_md_m^{-A_0}\le C m^{\Gamma-A_0}\) and
\(J_md_m^{-q}\le C m^{\Gamma-q}\), so the two active-menu tail series are
finite under the stated inequalities.  Also
\(\log(e+m+J_m)\le C\log(e+m)\), and the moment series of
Corollary~\ref{cor:random-environment-invp-active-menu-auto-truncation} is
dominated by the displayed series after substituting
\(\alpha_m=c_\alpha m^{-B}\).  The condition \(B>1\) is the all-integer
deterministic inverse-\(p\) schedule.

In \textup{(ii)},
\[
  J_md_m^{-A_0}\le C d_m^{\Gamma-A_0}
  \le C' m^{-\kappa(A_0-\Gamma)}
\]
and similarly for \(q\).  The displayed bounds on \(A_0\) and \(q\) make both
series summable.  Since \(d_m\ge c_0m^\kappa\) and \(J_m\le C_Jd_m^\Gamma\),
\(\log(e+m+J_m)\le C\ell_m\).  Substituting
\(\alpha_m=c_\alpha d_m^{-B}\) gives the displayed moment series, and
\(B\kappa>1\) is the deterministic inverse-\(p\) threshold.

In \textup{(iii)}, the two active-menu tail terms are bounded by
\(C d_m^{\Gamma-A_0}\) and \(C d_m^{\Gamma-q}\), which are summable because
\(d_m\) grows stretched-exponentially and \(A_0,q>\Gamma\).  Again
\(\log(e+m+J_m)\le C\ell_m\).  With
\(\alpha_m=c_\alpha\ell_m^{-B}\), this gives the series
\(\sum_m \ell_m^{Bs+s-1}H_m/n_m^{s-1}\), and
\(B>\sigma_p+1/\kappa\) is the stretched-exponential inverse-\(p\) threshold.

In \textup{(iv)}, \(J_m d_m^{-A_0}\) and \(J_m d_m^{-q}\) are summable for every
positive \(A_0,q\), since \(J_m\le C\ell_m^\Gamma\) is only polynomial in
\(\ell_m\) while \(d_m\) grows stretched-exponentially.  Also
\(\log(e+m+J_m)\le C\log(e+m)\).  The active-menu moment summability is
therefore dominated by the stated moment series after substituting
\(\alpha_m=c_\alpha\ell_m^{-B}\).  Applying
Corollary~\ref{cor:random-environment-invp-active-menu-auto-truncation} proves
all four cases.
\end{proof}

\begin{corollary}[Random active-prefix counts]
\label{cor:random-environment-invp-random-active-prefix-counts}
Work with a countable predictable catalog library indexed by \(r\ge1\), in the
setting of
Corollary~\ref{cor:random-environment-invp-active-menu-auto-truncation}.  Let
\(N_m\) be a nonnegative integer-valued random active count on the same
probability space.  Let \(J_m\ge1\) be deterministic.  Suppose that
\[
  \sum_m \Prob\{N_m>J_m\}<\infty .
\]
If the deterministic prefix \(1\le r\le J_m\) satisfies the hypotheses of
Corollary~\ref{cor:random-environment-invp-active-menu-auto-truncation}, then
the three sampled Toeplitz/Hankel conclusions of that corollary hold
simultaneously for every randomly active stream \(1\le r\le N_m\), for all
sufficiently large \(m\), almost surely.  The same statement holds if the
deterministic prefix is checked through any of the concrete schedules in
Corollary~\ref{cor:random-environment-invp-active-prefix-schedules}.

The overflow condition has the following direct moment tests.  For any \(t>0\),
it is enough that
\[
  \sum_m
  \frac{\Exp[(1+N_m)^t]}{(1+J_m)^t}<\infty .
\]
Consequently, for deterministic constants \(C,\eta\ge0\), the overflow series
is summable in the following typical cases.
\begin{enumerate}[label=\textup{(\roman*)},leftmargin=2.2em]
\item If \(d_m=m\), \(J_m=\lceil m^\Gamma\rceil\), and
\(\Exp[(1+N_m)^t]\le Cm^\eta\), then it is enough that
\[
  t\Gamma>\eta+1 .
\]
\item If \(d_m\ge c_0m^\kappa\), \(J_m=\lceil d_m^\Gamma\rceil\), and
\(\Exp[(1+N_m)^t]\le Cd_m^\eta\), then it is enough that
\[
  \kappa(t\Gamma-\eta)>1 .
\]
\item If \(d_m\ge\exp(c_0m^\kappa)\), \(J_m=\lceil d_m^\Gamma\rceil\), and
\(\Exp[(1+N_m)^t]\le Cd_m^\eta\), then it is enough that
\[
  t\Gamma>\eta .
\]
\item If \(d_m\ge\exp(c_0m^\kappa)\), \(J_m=\lceil\ell_m^\Gamma\rceil\), and
\(\Exp[(1+N_m)^t]\le C\ell_m^\eta\), then it is enough that
\[
  \kappa(t\Gamma-\eta)>1 .
\]
\end{enumerate}
Thus random proposal counts, adaptive search frontiers, and random catalog
occupancy prefixes may be used in place of deterministic active-prefix sizes
whenever their overflow tails are summable.
\end{corollary}

\begin{proof}
By Borel--Cantelli, \(N_m\le J_m\) for all sufficiently large \(m\), almost
surely.  On the same probability-one event on which
Corollary~\ref{cor:random-environment-invp-active-menu-auto-truncation} gives
the simultaneous conclusions for all deterministic prefix streams
\(1\le r\le J_m\), those conclusions automatically hold for the smaller random
set \(1\le r\le N_m\).  The concrete schedule version is the same argument after
using Corollary~\ref{cor:random-environment-invp-active-prefix-schedules} to
verify the deterministic prefix hypotheses.

For the moment tests, Markov's inequality gives
\[
  \Prob\{N_m>J_m\}
  \le
  \frac{\Exp[(1+N_m)^t]}{(1+J_m)^t}.
\]
The first displayed summability therefore implies the overflow condition.  The
four rate tests are the corresponding \(p\)-series or stretched-exponential
substitutions: in \textup{(i)} the summand is \(O(m^{\eta-t\Gamma})\); in
\textup{(ii)} it is \(O(d_m^{\eta-t\Gamma})\le
O(m^{-\kappa(t\Gamma-\eta)})\); in \textup{(iii)} it is
\(O(d_m^{-(t\Gamma-\eta)})\), hence summable for \(t\Gamma>\eta\); and in
\textup{(iv)} it is \(O(\ell_m^{\eta-t\Gamma})\le
O(m^{-\kappa(t\Gamma-\eta)})\).
\end{proof}

\begin{corollary}[Sub-Poisson active-count overflow tests]
\label{cor:random-environment-invp-subpoisson-active-counts}
Work in the setting of
Corollary~\ref{cor:random-environment-invp-random-active-prefix-counts}.  Let
\(\mathcal H_m\) be a \(\sigma\)-field and let \(\nu_m\ge0\) be deterministic.
Assume that, for every \(\theta\ge0\),
\[
  \Exp\!\left[\exp(\theta N_m)\mid\mathcal H_m\right]
  \le
  \exp\!\left(\nu_m(e^\theta-1)\right)
  \qquad\text{almost surely}.
\]
If
\[
  J_m\ge 2\nu_m
  \qquad\text{eventually},
  \qquad
  \sum_m\exp(-c_{\rm SP}J_m)<\infty
\]
for an absolute constant \(c_{\rm SP}>0\), then
\[
  \sum_m\Prob\{N_m>J_m\}<\infty .
\]
Thus the overflow hypothesis of
Corollary~\ref{cor:random-environment-invp-random-active-prefix-counts} holds,
so any deterministic prefix certified by
Corollary~\ref{cor:random-environment-invp-active-menu-auto-truncation} or
Corollary~\ref{cor:random-environment-invp-active-prefix-schedules} remains
valid with the random count \(N_m\).

The hypothesis includes the standard proposal-count models.  It holds if,
conditional on \(\mathcal H_m\), \(N_m\) is Poisson with mean at most \(\nu_m\).
It also holds if
\[
  N_m=\sum_{a\in\mathcal A_m}\eta_{m,a},
\]
where the \(\eta_{m,a}\) are conditionally independent Bernoulli variables and
\[
  \sum_{a\in\mathcal A_m}
  \Prob\{\eta_{m,a}=1\mid\mathcal H_m\}\le\nu_m
  \qquad\text{almost surely}.
\]
Thus the random active-prefix theorem covers Poisson proposal counts,
independent occupancy catalogs, and conditionally independent adaptive frontier
exposures whenever the deterministic prefix \(J_m\) is at least twice the
conditional mean envelope and the displayed exponential series is summable.

In particular, the overflow condition is automatic whenever
\[
  J_m\ge Q\log(e+m)
\]
with \(Q>1/c_{\rm SP}\) and \(J_m\ge2\nu_m\) eventually.  The deterministic
prefix choices \(J_m=\lceil m^\Gamma\rceil\),
\(J_m=\lceil d_m^\Gamma\rceil\), or \(J_m=\lceil\ell_m^\Gamma\rceil\) with
\(\Gamma>0\) therefore have summable sub-Poisson overflow tails in the
corresponding all-integer, polynomial, and stretched-exponential schedules.
\end{corollary}

\begin{proof}
Taking \(\theta=\log2\) in the conditional moment generating function and using
Markov's inequality gives
\[
\begin{aligned}
  \Prob\{N_m>J_m\mid\mathcal H_m\}
  &\le
  \exp\!\left(\nu_m- J_m\log2\right)  \\
  &\le
  \exp\!\left(-(\log2-\tfrac12)J_m\right)
\end{aligned}
\]
whenever \(J_m\ge2\nu_m\).  After integrating over \(\mathcal H_m\), this gives
the claimed overflow bound with \(c_{\rm SP}=\log2-\tfrac12\), or any smaller
absolute constant.  Corollary~\ref{cor:random-environment-invp-random-active-prefix-counts}
then gives the random-count conclusions.

For a conditionally Poisson count with mean \(\mu_m\le\nu_m\), the displayed
moment generating function is the Poisson identity with \(\mu_m\) bounded by
\(\nu_m\).  For the Bernoulli occupancy model,
\[
  \Exp[\exp(\theta N_m)\mid\mathcal H_m]
  =
  \prod_{a\in\mathcal A_m}
  \left(1+p_{m,a}(e^\theta-1)\right),
\]
where \(p_{m,a}=\Prob\{\eta_{m,a}=1\mid\mathcal H_m\}\).  Since
\(1+x\le e^x\), this product is at most
\[
  \exp\!\left((e^\theta-1)\sum_a p_{m,a}\right)
  \le \exp\!\left(\nu_m(e^\theta-1)\right).
\]
Finally, \(J_m\ge Q\log(e+m)\) gives
\(\exp(-c_{\rm SP}J_m)\le(e+m)^{-c_{\rm SP}Q}\), which is summable when
\(Q>1/c_{\rm SP}\); the stated polynomial and logarithmic prefix examples
eventually dominate such a logarithm whenever \(\Gamma>0\) in the relevant
schedule.
\end{proof}

\begin{corollary}[Predictable active-intensity overflow tests]
\label{cor:random-environment-invp-predictable-active-intensity}
Work in the setting of
Corollary~\ref{cor:random-environment-invp-random-active-prefix-counts}.  Let
\(\mathcal H_m\subseteq\mathcal G_{m,0}\subseteq\mathcal G_{m,1}\subseteq
\cdots\subseteq\mathcal G_{m,K_m}\) be a filtration and let
\(\eta_{m,a}\in\{0,1\}\), \(1\le a\le K_m\), be
\(\mathcal G_{m,a}\)-measurable activation indicators.  Put
\[
  N_m=\sum_{a=1}^{K_m}\eta_{m,a},
  \qquad
  p_{m,a}:=\Exp[\eta_{m,a}\mid\mathcal G_{m,a-1}].
\]
Assume that, for a deterministic \(\nu_m\ge0\),
\[
  \sum_{a=1}^{K_m}p_{m,a}\le\nu_m
  \qquad\text{almost surely}.
\]
Then the hypothesis of
Corollary~\ref{cor:random-environment-invp-subpoisson-active-counts} holds for
\(N_m\).  Hence
\[
  J_m\ge2\nu_m\ \text{eventually},
  \qquad
  \sum_m\exp(-c_{\rm SP}J_m)<\infty
\]
imply
\[
  \sum_m\Prob\{N_m>J_m\}<\infty .
\]
Consequently all deterministic active-prefix Toeplitz/Hankel sampled
certificates certified by
Corollary~\ref{cor:random-environment-invp-active-menu-auto-truncation} or
Corollary~\ref{cor:random-environment-invp-active-prefix-schedules} remain
valid for every adaptively active stream \(1\le r\le N_m\), eventually almost
surely.

The same conclusion holds for countably many proposals \(a\ge1\) whenever the
partial predictable intensities are almost surely bounded by \(\nu_m\).  Thus
independence is not needed for the overflow step: adaptive sequential
frontiers, online candidate filters, without-replacement exposure schemes, and
stopping-rule proposal lists are covered once their predictable total
activation intensity is dominated by \(\nu_m\).
\end{corollary}

\begin{proof}
Fix \(m\) and \(\theta\ge0\).  Since \(\eta_{m,a}\in\{0,1\}\),
\[
  \Exp[\exp(\theta\eta_{m,a})\mid\mathcal G_{m,a-1}]
  =
  1+p_{m,a}(e^\theta-1)
  \le
  \exp\!\left(p_{m,a}(e^\theta-1)\right).
\]
Therefore
\[
  M_j:=
  \exp\!\left(
    \theta\sum_{a\le j}\eta_{m,a}
    -(e^\theta-1)\sum_{a\le j}p_{m,a}
  \right),
  \qquad 0\le j\le K_m,
\]
is a nonnegative supermartingale.  Since
\(\mathcal H_m\subseteq\mathcal G_{m,0}\),
\[
  \Exp[M_{K_m}\mid\mathcal H_m]\le1 .
\]
Using the predictable-intensity bound,
\[
\begin{aligned}
  \Exp[\exp(\theta N_m)\mid\mathcal H_m]
  &\le
  \exp\!\left(\nu_m(e^\theta-1)\right)
  \Exp[M_{K_m}\mid\mathcal H_m]  \\
  &\le
  \exp\!\left(\nu_m(e^\theta-1)\right).
\end{aligned}
\]
This is exactly the sub-Poisson hypothesis of
Corollary~\ref{cor:random-environment-invp-subpoisson-active-counts}, which
gives the summable overflow and random active-prefix conclusions.

For countably many proposals, apply the finite argument to the first \(K\)
indicators and let \(K\to\infty\) by conditional monotone convergence.
\end{proof}

\begin{corollary}[Random predictable active-intensity envelopes]
\label{cor:random-environment-invp-random-predictable-active-intensity}
Use the notation of
Corollary~\ref{cor:random-environment-invp-predictable-active-intensity}, but
replace the deterministic almost-sure intensity bound by the predictable total
intensity
\[
  V_m:=\sum_{a=1}^{K_m}p_{m,a}.
\]
Let \(\nu_m\ge0\) and \(J_m\ge1\) be deterministic.  If
\[
  \sum_m\Prob\{V_m>\nu_m\}<\infty,
  \qquad
  J_m\ge2\nu_m\ \text{eventually},
\]
and
\[
  \sum_m\exp(-c_{\rm SP}J_m)<\infty,
\]
then
\[
  \sum_m\Prob\{N_m>J_m\}<\infty .
\]
Consequently the random active-prefix conclusions of
Corollary~\ref{cor:random-environment-invp-random-active-prefix-counts} hold:
every deterministic active-prefix Toeplitz/Hankel sampled certificate verified
through Corollary~\ref{cor:random-environment-invp-active-menu-auto-truncation}
or Corollary~\ref{cor:random-environment-invp-active-prefix-schedules} remains
valid for all randomly active streams \(1\le r\le N_m\), eventually almost
surely.

The countable-proposal version is the same with
\[
  V_m:=\sum_{a\ge1}p_{m,a},
\]
provided this sum is finite outside the displayed summable intensity-overflow
event.  Thus the active-count step may be certified by a separate random
environment or degree event: adaptive proposal graphs, random candidate pools,
online thinning rules, and stopping-rule frontiers only pay the probability
that their predictable total intensity exceeds the deterministic envelope
\(\nu_m\).
\end{corollary}

\begin{proof}
First suppose \(K_m<\infty\), and let \(E_m=\{V_m\le\nu_m\}\).  With
\[
  M_j=
  \exp\!\left(
    \theta\sum_{a\le j}\eta_{m,a}
    -(e^\theta-1)\sum_{a\le j}p_{m,a}
  \right)
\]
as in the proof of
Corollary~\ref{cor:random-environment-invp-predictable-active-intensity},
\((M_j)_{j\le K_m}\) is a nonnegative supermartingale and
\(\Exp M_{K_m}\le1\).  Hence, for \(\theta=\log2\),
\[
\begin{aligned}
  \Prob\{N_m>J_m,\ E_m\}
  &\le
  e^{-\theta J_m}
  \Exp\!\left[1_{E_m}e^{\theta N_m}\right]  \\
  &=
  e^{-\theta J_m}
  \Exp\!\left[
    1_{E_m}\exp\!\left((e^\theta-1)V_m\right)M_{K_m}
  \right]  \\
  &\le
  \exp(\nu_m-J_m\log2).
\end{aligned}
\]
When \(J_m\ge2\nu_m\), this is at most
\(\exp(-c_{\rm SP}J_m)\), with the same absolute constant as in
Corollary~\ref{cor:random-environment-invp-subpoisson-active-counts}.  Therefore
\[
  \Prob\{N_m>J_m\}
  \le
  \Prob\{V_m>\nu_m\}+\exp(-c_{\rm SP}J_m),
\]
and the displayed summability assumptions give the overflow condition.
Corollary~\ref{cor:random-environment-invp-random-active-prefix-counts} then
gives the stated active-prefix conclusions.

For countably many proposals, apply the finite bound to the first \(K\)
indicators, with \(V_{m,K}=\sum_{a\le K}p_{m,a}\le V_m\), and let
\(K\to\infty\) by monotone convergence.
\end{proof}

\begin{corollary}[Moment tests for random predictable active intensities]
\label{cor:random-environment-invp-moment-predictable-active-intensity}
In the setting of
Corollary~\ref{cor:random-environment-invp-random-predictable-active-intensity},
let \(t>0\).  If
\[
  \sum_m
  \frac{\Exp[(1+V_m)^t]}{(1+\nu_m)^t}<\infty,
  \qquad
  J_m\ge2\nu_m\ \text{eventually},
\]
and
\[
  \sum_m\exp(-c_{\rm SP}J_m)<\infty,
\]
then
\[
  \sum_m\Prob\{N_m>J_m\}<\infty .
\]
Hence every deterministic active-prefix Toeplitz/Hankel sampled certificate
verified through
Corollary~\ref{cor:random-environment-invp-active-menu-auto-truncation} or
Corollary~\ref{cor:random-environment-invp-active-prefix-schedules} remains
valid for all randomly active streams \(1\le r\le N_m\), eventually almost
surely.

The moment condition is explicit in the standard growth regimes.  For
deterministic \(C,\eta\ge0\) and \(c_\nu>0\), the intensity-envelope overflow
series is summable in the following cases.
\begin{enumerate}[label=\textup{(\roman*)},leftmargin=2.2em]
\item If \(d_m=m\), \(\nu_m=c_\nu m^\Gamma\), and
\(\Exp[(1+V_m)^t]\le Cm^\eta\), it is enough that
\[
  t\Gamma>\eta+1 .
\]
\item If \(d_m\ge c_0m^\kappa\), \(\nu_m=c_\nu d_m^\Gamma\), and
\(\Exp[(1+V_m)^t]\le Cd_m^\eta\), it is enough that
\[
  \kappa(t\Gamma-\eta)>1 .
\]
\item If \(d_m\ge\exp(c_0m^\kappa)\), \(\nu_m=c_\nu d_m^\Gamma\), and
\(\Exp[(1+V_m)^t]\le Cd_m^\eta\), it is enough that
\[
  t\Gamma>\eta .
\]
\item If \(d_m\ge\exp(c_0m^\kappa)\), \(\nu_m=c_\nu\ell_m^\Gamma\), and
\(\Exp[(1+V_m)^t]\le C\ell_m^\eta\), it is enough that
\[
  \kappa(t\Gamma-\eta)>1 .
\]
\end{enumerate}
For any deterministic prefix \(J_m\ge2\nu_m\) with
\(\sum_m e^{-c_{\rm SP}J_m}<\infty\), these tests give the corresponding
random active-count certificates.
\end{corollary}

\begin{proof}
Markov's inequality gives
\[
  \Prob\{V_m>\nu_m\}
  \le
  \frac{\Exp[(1+V_m)^t]}{(1+\nu_m)^t}.
\]
Thus the displayed moment summability verifies the intensity-envelope
exception series in
Corollary~\ref{cor:random-environment-invp-random-predictable-active-intensity}.
That corollary then gives the active-count overflow and the active-prefix
certificate conclusions.

The four rate checks are the same substitutions as in the moment tests for
random active-prefix counts.  In \textup{(i)} the summand is
\(O(m^{\eta-t\Gamma})\); in \textup{(ii)} it is
\(O(d_m^{\eta-t\Gamma})\le O(m^{-\kappa(t\Gamma-\eta)})\); in \textup{(iii)}
it is stretched-exponentially summable when \(t\Gamma>\eta\); and in
\textup{(iv)} it is
\(O(\ell_m^{\eta-t\Gamma})\le O(m^{-\kappa(t\Gamma-\eta)})\).
\end{proof}

\begin{corollary}[Freedman tests for random predictable active intensities]
\label{cor:random-environment-invp-freedman-predictable-active-intensity}
In the setting of
Corollary~\ref{cor:random-environment-invp-random-predictable-active-intensity},
suppose that the predictable total intensity admits a random-environment
majorant
\[
  V_m\le\sum_{b=1}^{B_m}X_{m,b}.
\]
Let
\(\mathcal R_{m,0}\subseteq\mathcal R_{m,1}\subseteq\cdots\subseteq
\mathcal R_{m,B_m}\) be a filtration and suppose that
\(X_{m,b}\) is \(\mathcal R_{m,b}\)-measurable, \(0\le X_{m,b}\le L_m\), and
\[
  \sum_{b=1}^{B_m}
  \Exp[X_{m,b}\mid\mathcal R_{m,b-1}]\le\mu_m,
  \qquad
  \sum_{b=1}^{B_m}
  \Exp[X_{m,b}^2\mid\mathcal R_{m,b-1}]\le\sigma_m^2
\]
almost surely, for deterministic \(\mu_m,\sigma_m,L_m\ge0\).  Let
\(\nu_m>\mu_m\), put \(u_m=\nu_m-\mu_m\), and assume that
\[
  \sum_m
  \exp\!\left(
    -c_{\rm FI}\min\left\{\frac{u_m^2}{\sigma_m^2},
                          \frac{u_m}{L_m}\right\}
  \right)<\infty,
\]
with the usual convention that a ratio with zero denominator is \(+\infty\)
when the numerator is positive.  If also
\[
  J_m\ge2\nu_m\ \text{eventually},
  \qquad
  \sum_m\exp(-c_{\rm SP}J_m)<\infty,
\]
then
\[
  \sum_m\Prob\{N_m>J_m\}<\infty .
\]
Therefore the deterministic active-prefix Toeplitz/Hankel sampled certificates
verified through
Corollary~\ref{cor:random-environment-invp-active-menu-auto-truncation} or
Corollary~\ref{cor:random-environment-invp-active-prefix-schedules} hold for
all randomly active streams \(1\le r\le N_m\), eventually almost surely.

In particular, if
\[
  \mu_m\le\nu_m/2,\qquad
  \sigma_m^2\le L_m\nu_m,
\]
then it is enough that
\[
  \sum_m \exp(-c'_{\rm FI}\nu_m/L_m)<\infty
\]
and \(J_m\ge2\nu_m\) eventually.  For unit-bounded intensity contributions
\(L_m=1\), any sufficiently large multiple of \(\log(e+m)\) for \(\nu_m\) and
\(J_m\) gives a summable envelope failure.  Thus random candidate-pool degrees,
adaptive proposal loads, and online thinning intensities may be certified by
predictable mean, square-function, and jump envelopes rather than by moment
growth assumptions.
\end{corollary}

\begin{proof}
Let
\[
  m_{m,b}:=\Exp[X_{m,b}\mid\mathcal R_{m,b-1}],
  \qquad
  Y_{m,b}:=X_{m,b}-m_{m,b}.
\]
Then \((Y_{m,b})_b\) is a martingale-difference sequence, \(|Y_{m,b}|\le L_m\),
and its conditional variance process is bounded by
\[
  \sum_b\Exp[Y_{m,b}^2\mid\mathcal R_{m,b-1}]
  \le
  \sum_b\Exp[X_{m,b}^2\mid\mathcal R_{m,b-1}]
  \le \sigma_m^2 .
\]
If \(V_m>\nu_m=\mu_m+u_m\), then
\[
  \sum_b Y_{m,b}
  \ge
  \sum_b X_{m,b}-\sum_b m_{m,b}
  > u_m .
\]
Scalar Freedman's inequality therefore gives
\[
  \Prob\{V_m>\nu_m\}
  \le
  \exp\!\left(
    -c_{\rm FI}\min\left\{\frac{u_m^2}{\sigma_m^2},
                          \frac{u_m}{L_m}\right\}
  \right)
\]
for an absolute constant \(c_{\rm FI}>0\).  The displayed summability verifies
the intensity-envelope exception series in
Corollary~\ref{cor:random-environment-invp-random-predictable-active-intensity},
which gives the active-count overflow and the active-prefix conclusions.

If \(\mu_m\le\nu_m/2\), then \(u_m\ge\nu_m/2\).  Under
\(\sigma_m^2\le L_m\nu_m\), the Freedman exponent is bounded below by a
constant multiple of \(\nu_m/L_m\), giving the simplified criterion.
\end{proof}

\begin{corollary}[Truncated Freedman tests for unbounded predictable active intensities]
\label{cor:random-environment-invp-truncated-freedman-predictable-active-intensity}
In the setting of
Corollary~\ref{cor:random-environment-invp-random-predictable-active-intensity},
suppose that the predictable total intensity admits an adapted nonnegative
majorant
\[
  V_m\le\sum_{b=1}^{B_m}X_{m,b}.
\]
Let
\(\mathcal R_{m,0}\subseteq\mathcal R_{m,1}\subseteq\cdots\subseteq
\mathcal R_{m,B_m}\) be a filtration and suppose that \(X_{m,b}\) is
\(\mathcal R_{m,b}\)-measurable and finite almost surely.  Fix a deterministic
truncation level \(L_m>0\), and write
\[
  W_{m,b}:=X_{m,b}\wedge L_m,\qquad
  R_{m,b}:=(X_{m,b}-L_m)_+ .
\]
Assume that deterministic numbers \(\mu_m,\sigma_m,\rho_m\ge0\) satisfy
\[
  \sum_{b=1}^{B_m}
  \Exp[W_{m,b}\mid\mathcal R_{m,b-1}]\le\mu_m,\qquad
  \sum_{b=1}^{B_m}
  \Exp[W_{m,b}^2\mid\mathcal R_{m,b-1}]\le\sigma_m^2,
\]
and
\[
  \sum_{b=1}^{B_m}
  \Exp[R_{m,b}\mid\mathcal R_{m,b-1}]\le\rho_m
\]
almost surely.  Let \(\nu_m>\mu_m\), put \(u_m=\nu_m-\mu_m\), and assume that
\[
  \sum_m
  \left[
  \exp\!\left(
    -c_{\rm TFI}\min\left\{\frac{u_m^2}{\sigma_m^2},
                          \frac{u_m}{L_m}\right\}
  \right)
  +\frac{\rho_m}{u_m}
  \right]<\infty,
\]
with the usual convention that a ratio with zero denominator is \(+\infty\)
when the numerator is positive.  If also
\[
  J_m\ge2\nu_m\ \text{eventually},
  \qquad
  \sum_m\exp(-c_{\rm SP}J_m)<\infty,
\]
then
\[
  \sum_m\Prob\{N_m>J_m\}<\infty .
\]
Consequently every deterministic active-prefix Toeplitz/Hankel sampled
certificate verified through
Corollary~\ref{cor:random-environment-invp-active-menu-auto-truncation} or
Corollary~\ref{cor:random-environment-invp-active-prefix-schedules} holds for
all randomly active streams \(1\le r\le N_m\), eventually almost surely.

In particular, if
\[
  \mu_m\le\nu_m/2,\qquad
  \sigma_m^2\le L_m\nu_m,
\]
then it is enough that
\[
  \sum_m\left[
    \exp(-c'_{\rm TFI}\nu_m/L_m)+\frac{\rho_m}{\nu_m}
  \right]<\infty
\]
and \(J_m\ge2\nu_m\) eventually.  A local conditional-moment check is obtained
as follows: if \(s>1\) and
\[
  \sum_{b=1}^{B_m}
  \Exp[X_{m,b}^s\mid\mathcal R_{m,b-1}]\le H_m
  \qquad\text{almost surely},
\]
then one may take \(\rho_m=H_m/L_m^{s-1}\).  Hence unbounded random
candidate-pool degrees, heavy-tailed proposal loads, and adaptive thinning
intensities are covered once the truncated predictable mean, truncated
square-function, and excess series are summable.
\end{corollary}

\begin{proof}
Let
\[
  w_{m,b}:=\Exp[W_{m,b}\mid\mathcal R_{m,b-1}],
  \qquad
  D_{m,b}:=W_{m,b}-w_{m,b}.
\]
Then \((D_{m,b})_b\) is a martingale-difference sequence, \(|D_{m,b}|\le L_m\),
and
\[
  \sum_b\Exp[D_{m,b}^2\mid\mathcal R_{m,b-1}]
  \le
  \sum_b\Exp[W_{m,b}^2\mid\mathcal R_{m,b-1}]
  \le \sigma_m^2 .
\]
If \(\sum_bX_{m,b}>\nu_m=\mu_m+u_m\), then either
\(\sum_bR_{m,b}>u_m/2\) or
\[
  \sum_b W_{m,b}>\mu_m+u_m/2 .
\]
On the latter event,
\[
  \sum_bD_{m,b}
  =
  \sum_bW_{m,b}-\sum_bw_{m,b}
  >
  u_m/2 .
\]
Scalar Freedman's inequality therefore gives
\[
  \Prob\left\{\sum_bW_{m,b}>\mu_m+u_m/2\right\}
  \le
  \exp\!\left(
    -c_{\rm TFI}\min\left\{\frac{u_m^2}{\sigma_m^2},
                          \frac{u_m}{L_m}\right\}
  \right)
\]
after reducing the absolute constant \(c_{\rm TFI}\).  Markov's inequality and
the predictable excess bound give
\[
  \Prob\left\{\sum_bR_{m,b}>u_m/2\right\}
  \le
  \frac{2}{u_m}\sum_b\Exp R_{m,b}
  \le \frac{2\rho_m}{u_m}.
\]
Absorbing the harmless factor \(2\) into the summability assumption proves
\[
  \sum_m\Prob\{V_m>\nu_m\}<\infty .
\]
Corollary~\ref{cor:random-environment-invp-random-predictable-active-intensity}
then gives the active-count overflow and the active-prefix certificate
conclusions.

If \(\mu_m\le\nu_m/2\), then \(u_m\ge\nu_m/2\), and the hypothesis
\(\sigma_m^2\le L_m\nu_m\) lowers the Freedman exponent by a constant multiple
of \(\nu_m/L_m\).  Finally, for \(x\ge0\),
\((x-L_m)_+\le x^s/L_m^{s-1}\); applying this conditionally and summing gives
the stated choice \(\rho_m=H_m/L_m^{s-1}\).
\end{proof}

\begin{corollary}[Automatic moment truncation for predictable active intensities]
\label{cor:random-environment-invp-auto-truncation-predictable-active-intensity}
In the setting of
Corollary~\ref{cor:random-environment-invp-truncated-freedman-predictable-active-intensity},
fix \(s>1\).  Suppose that deterministic \(\nu_m>0\), \(H_m\ge0\), and
\(\lambda_m>0\) satisfy
\[
  \sum_{b=1}^{B_m}
  \Exp[X_{m,b}\mid\mathcal R_{m,b-1}]\le\frac{\nu_m}{2},
  \qquad
  \sum_{b=1}^{B_m}
  \Exp[X_{m,b}^s\mid\mathcal R_{m,b-1}]\le H_m
\]
almost surely.  If
\[
  \sum_m
  \left[
    \exp(-c_{\rm ATI}\lambda_m)
    +\frac{H_m\lambda_m^{s-1}}{\nu_m^s}
  \right]<\infty,
\]
and
\[
  J_m\ge2\nu_m\ \text{eventually},
  \qquad
  \sum_m\exp(-c_{\rm SP}J_m)<\infty,
\]
then
\[
  \sum_m\Prob\{N_m>J_m\}<\infty .
\]
Consequently all deterministic active-prefix Toeplitz/Hankel sampled
certificates verified through
Corollary~\ref{cor:random-environment-invp-active-menu-auto-truncation} or
Corollary~\ref{cor:random-environment-invp-active-prefix-schedules} hold for
all randomly active streams \(1\le r\le N_m\), eventually almost surely.

With \(\lambda_m=A\log(e+m)\) and \(A\) sufficiently large, the criterion is
explicit in the standard growth regimes.  For deterministic
\(C,\eta\ge0\) and \(c_\nu>0\), the active-intensity overflow series is
summable in the following cases.
\begin{enumerate}[label=\textup{(\roman*)},leftmargin=2.2em]
\item If \(d_m=m\), \(\nu_m=c_\nu m^\Gamma\), and \(H_m\le Cm^\eta\), it is
enough that
\[
  s\Gamma>\eta+1 .
\]
\item If \(d_m\ge c_0m^\kappa\), \(\nu_m=c_\nu d_m^\Gamma\), and
\(H_m\le Cd_m^\eta\), it is enough that
\[
  \kappa(s\Gamma-\eta)>1 .
\]
\item If \(d_m\ge\exp(c_0m^\kappa)\), \(\nu_m=c_\nu d_m^\Gamma\), and
\(H_m\le Cd_m^\eta\), it is enough that
\[
  s\Gamma>\eta .
\]
\item If \(d_m\ge\exp(c_0m^\kappa)\), \(\nu_m=c_\nu\ell_m^\Gamma\), and
\(H_m\le C\ell_m^\eta\), it is enough that
\[
  \kappa(s\Gamma-\eta)>1 .
\]
\end{enumerate}
Thus a single conditional \(s\)-moment envelope and a logarithmic truncation
parameter certify heavy-tailed adaptive proposal-load intensities in the same
polynomial, stretched-exponential, and logarithmic-prefix schedules as the
bounded-load Freedman tests.
\end{corollary}

\begin{proof}
Put
\[
  L_m:=\frac{\nu_m}{\lambda_m},\qquad
  W_{m,b}:=X_{m,b}\wedge L_m,\qquad
  R_{m,b}:=(X_{m,b}-L_m)_+ .
\]
Then
\[
  \sum_b\Exp[W_{m,b}\mid\mathcal R_{m,b-1}]
  \le
  \sum_b\Exp[X_{m,b}\mid\mathcal R_{m,b-1}]
  \le \frac{\nu_m}{2}.
\]
Since \(W_{m,b}^2\le L_m W_{m,b}\),
\[
  \sum_b\Exp[W_{m,b}^2\mid\mathcal R_{m,b-1}]
  \le
  L_m\sum_b\Exp[W_{m,b}\mid\mathcal R_{m,b-1}]
  \le L_m\nu_m .
\]
Also \((x-L_m)_+\le x^s/L_m^{s-1}\), so
\[
  \sum_b\Exp[R_{m,b}\mid\mathcal R_{m,b-1}]
  \le
  \frac{H_m}{L_m^{s-1}} .
\]
The half-budget truncated-Freedman corollary therefore applies with
\[
  \rho_m=\frac{H_m}{L_m^{s-1}} .
\]
Its two summability terms become
\[
  \exp(-c'_{\rm TFI}\nu_m/L_m)=\exp(-c'_{\rm TFI}\lambda_m),
  \qquad
  \frac{\rho_m}{\nu_m}
  =
  \frac{H_m\lambda_m^{s-1}}{\nu_m^s}.
\]
After renaming the absolute constant, the displayed summability assumption
gives the random active-count overflow and the active-prefix certificate
conclusions.

For the concrete regimes, \(\lambda_m=A\log(e+m)\) makes the exponential
series summable when \(A\) is large enough.  In \textup{(i)} the excess term is
\[
  O\!\left(m^{\eta-s\Gamma}(\log(e+m))^{s-1}\right),
\]
which is summable if \(s\Gamma>\eta+1\).  In \textup{(ii)} it is
\[
  O\!\left(d_m^{\eta-s\Gamma}(\log(e+m))^{s-1}\right)
  \le
  O\!\left(m^{-\kappa(s\Gamma-\eta)}
           (\log(e+m))^{s-1}\right),
\]
which is summable under the displayed condition.  In \textup{(iii)} it is
stretched-exponentially summable when \(s\Gamma>\eta\).  In \textup{(iv)} it is
\[
  O\!\left(\ell_m^{\eta-s\Gamma}(\log(e+m))^{s-1}\right)
  \le
  O\!\left(m^{-\kappa(s\Gamma-\eta)}
           (\log(e+m))^{s-1}\right),
\]
which is summable when \(\kappa(s\Gamma-\eta)>1\).
\end{proof}

\begin{corollary}[Automatic prefixes from predictable active-intensity moments]
\label{cor:random-environment-invp-auto-prefix-predictable-active-intensity}
Work in the setting of
Corollary~\ref{cor:random-environment-invp-auto-truncation-predictable-active-intensity}.
Fix \(s>1\), and suppose that deterministic \(\nu_m>0\) and \(H_m\ge0\)
satisfy
\[
  \sum_{b=1}^{B_m}
  \Exp[X_{m,b}\mid\mathcal R_{m,b-1}]\le\frac{\nu_m}{2},
  \qquad
  \sum_{b=1}^{B_m}
  \Exp[X_{m,b}^s\mid\mathcal R_{m,b-1}]\le H_m
\]
almost surely.  Choose constants \(A,Q>0\) large enough that
\[
  c_{\rm ATI}A>1,\qquad c_{\rm SP}Q>1,
\]
and put
\[
  \lambda_m:=A\log(e+m),
  \qquad
  J_m:=\left\lceil\max\{2\nu_m,\ Q\log(e+m)\}\right\rceil .
\]
If
\[
  \sum_m
  \frac{H_m(\log(e+m))^{s-1}}{\nu_m^s}<\infty,
\]
then
\[
  \sum_m\Prob\{N_m>J_m\}<\infty .
\]
Consequently every deterministic active-prefix Toeplitz/Hankel sampled
certificate verified through
Corollary~\ref{cor:random-environment-invp-active-menu-auto-truncation} or
Corollary~\ref{cor:random-environment-invp-active-prefix-schedules} for the
prefix \(1\le r\le J_m\) holds for all randomly active streams
\(1\le r\le N_m\), eventually almost surely.

The prefix growth is automatic in the usual schedules.  If
\(\nu_m\le C_\nu m^\Gamma\) with \(\Gamma>0\), then \(J_m\le C m^\Gamma\), so
the all-integer active-prefix schedule applies with the same menu-growth
exponent.  If \(d_m\ge c_0m^\kappa\) and
\(\nu_m\le C_\nu d_m^\Gamma\) with \(\Gamma>0\), then
\(J_m\le C d_m^\Gamma\).  If
\(d_m\ge\exp(c_0m^\kappa)\) and \(\nu_m\le C_\nu d_m^\Gamma\) with
\(\Gamma>0\), the same
polynomial-in-\(d\) active-prefix growth holds, while
\(\nu_m\le C_\nu\ell_m^\Gamma\) with \(\Gamma>0\) gives
\(J_m\le C\ell_m^\Gamma\).  In those four regimes, the corresponding
active-prefix schedule of
Corollary~\ref{cor:random-environment-invp-active-prefix-schedules} can be used
with the same \(A_0,q\), and catalog-moment thresholds, while the overflow
part is reduced to the one displayed moment series above.
\end{corollary}

\begin{proof}
Apply
Corollary~\ref{cor:random-environment-invp-auto-truncation-predictable-active-intensity}
with the displayed \(\lambda_m\) and \(J_m\).  By construction
\(J_m\ge2\nu_m\), and
\[
  \sum_m\exp(-c_{\rm SP}J_m)
  \le
  \sum_m(e+m)^{-c_{\rm SP}Q}<\infty .
\]
The exponential truncation term is also summable:
\[
  \sum_m\exp(-c_{\rm ATI}\lambda_m)
  =
  \sum_m(e+m)^{-c_{\rm ATI}A}<\infty .
\]
The remaining term in that corollary is
\[
  \sum_m
  \frac{H_m\lambda_m^{s-1}}{\nu_m^s},
\]
which is finite by the displayed hypothesis, up to the fixed factor
\(A^{s-1}\).  Hence the overflow series is summable, and
Corollary~\ref{cor:random-environment-invp-random-active-prefix-counts} gives
the random active-prefix certificate conclusion for any deterministic prefix
certificate already verified up to \(J_m\).

The four growth claims follow from the definition of \(J_m\).  In the
all-integer and polynomial cases the logarithmic term is dominated by any fixed
positive power once \(m\) is large.  In the stretched-exponential
polynomial-in-\(d\) case, \(Q\log(e+m)\le C d_m^\Gamma\) when \(\Gamma>0\).
In the logarithmic-prefix case,
\(\log(e+m)\le C\ell_m^\Gamma\) for every fixed \(\Gamma>0\), since
\(\ell_m\ge c m^\kappa\).  The cited active-prefix schedule corollary then
supplies the deterministic-prefix sampled-certificate hypotheses.
\end{proof}

\begin{corollary}[End-to-end automatic active-prefix sampled certificates]
\label{cor:random-environment-invp-end-to-end-auto-active-certificates}
Work with a countable predictable catalog library indexed by \(r\ge1\), in the
setting of
Corollary~\ref{cor:random-environment-invp-active-menu-auto-truncation}, and
let \(N_m\) be the random active count generated by activation indicators whose
predictable total intensity \(V_m\) is dominated by
\[
  V_m\le\sum_{b=1}^{B_m}X_{m,b}.
\]
Fix \(s_{\rm act},s_{\rm cat}>1\).  Suppose that deterministic
\(\nu_m>0\) and \(G_m\ge0\) satisfy
\[
  \sum_b\Exp[X_{m,b}\mid\mathcal R_{m,b-1}]\le\frac{\nu_m}{2},
  \qquad
  \sum_b\Exp[X_{m,b}^{s_{\rm act}}\mid\mathcal R_{m,b-1}]
  \le G_m
\]
almost surely.  Choose \(A,Q>0\) so that
\[
  c_{\rm ATI}A>1,\qquad c_{\rm SP}Q>1,
\]
and set
\[
  J_m:=\left\lceil\max\{2\nu_m,\ Q\log(e+m)\}\right\rceil .
\]
For every \(1\le r\le J_m\), assume the active-menu catalog variables
\[
  Z_{m,i}^{(r)}
  =\iota_{m,i,\Lambda_{m,i}^{(r)}}^{(p)}(a)
\]
come from centered, variance-one, \(K\)-sub-Gaussian catalog laws and
have deterministic conditional-moment envelopes
\[
  \Exp[Z_{m,i}^{(r)}\mid\mathcal F_{m,i-1}^{(r)}]\le b_{m,i}^{(r)},
  \qquad
  \Exp[(Z_{m,i}^{(r)})^{s_{\rm cat}}\mid\mathcal F_{m,i-1}^{(r)}]
  \le h_{m,i}^{(r)} ,
\]
and define
\[
  \bar b_m^{*,J}:=
  \max_{1\le r\le J_m}
  \frac1{n_m}\sum_{i=1}^{n_m}b_{m,i}^{(r)},\qquad
  H_m^{\rm cat}:=
  \sum_{r=1}^{J_m}\frac1{n_m}\sum_{i=1}^{n_m}h_{m,i}^{(r)} .
\]
Let \(\alpha_m>0\) satisfy one of the deterministic inverse-\(p\) schedules in
Corollary~\ref{cor:random-environment-invp-operator-batch-maxima-schedules},
and assume
\[
  \bar b_m^{*,J}\le\frac{\alpha_m}{4}
  \qquad\text{eventually}.
\]
If
\[
  \sum_m
  \frac{G_m(\log(e+m))^{s_{\rm act}-1}}{\nu_m^{s_{\rm act}}}<\infty,
\]
\[
  \sum_m
  \left[
    J_m\exp(-c_{\rm E2E}\lambda_m^{\rm cat})
    +
    \frac{(\lambda_m^{\rm cat})^{s_{\rm cat}-1}H_m^{\rm cat}}
         {n_m^{s_{\rm cat}-1}\alpha_m^{s_{\rm cat}}}
  \right]<\infty
\]
for some deterministic \(\lambda_m^{\rm cat}>0\) and a sufficiently small
absolute constant \(c_{\rm E2E}>0\), and
\[
  \sum_mJ_md_m^{-A_0}<\infty,\qquad
  \sum_mJ_md_m^{-q}<\infty,
\]
then, with \(u_m=q\ell_m\), the fixed-fraction, vanishing-fraction, and
all-robust polynomial-batch Toeplitz/Hankel conclusions of
Corollary~\ref{cor:random-environment-invp-operator-batch-maxima-schedules}
hold simultaneously for every randomly active stream \(1\le r\le N_m\), for all
sufficiently large \(m\), almost surely.

In particular, with
\[
  \lambda_m^{\rm cat}=A_{\rm cat}\log(e+m+J_m)
\]
and \(A_{\rm cat}\) large enough, the exponential catalog-menu series is
automatic.  The remaining hypotheses are exactly the active-intensity moment
series, the catalog moment series, the active-menu matrix/sign tail series, and
the deterministic inverse-\(p\) schedule for \(\alpha_m\).  Thus the random
proposal frontier and the predictable law-catalog exhaustion may be certified
by two independent moment budgets and one deterministic prefix formula.
\end{corollary}

\begin{proof}
The active-intensity assumptions are precisely the hypotheses of
Corollary~\ref{cor:random-environment-invp-auto-prefix-predictable-active-intensity}
with \(H_m=G_m\), \(s=s_{\rm act}\), and
\(\lambda_m=A\log(e+m)\).  Hence
\[
  \sum_m\Prob\{N_m>J_m\}<\infty .
\]
For this deterministic prefix \(J_m\), the catalog hypotheses are those of
Corollary~\ref{cor:random-environment-invp-active-menu-auto-truncation}, with
\(s=s_{\rm cat}\), \(H_m=H_m^{\rm cat}\),
\(\lambda_m=\lambda_m^{\rm cat}\), and
\(\bar b_m^*=\bar b_m^{*,J}\).
The displayed catalog summability and the two matrix/sign summability
assumptions therefore give the simultaneous active-menu sampled Toeplitz/Hankel
conclusions for every deterministic row \(1\le r\le J_m\), eventually almost
surely.

Borel--Cantelli applied to the overflow series gives \(N_m\le J_m\) eventually
almost surely.  Intersecting this event with the active-menu certificate event
proves the asserted conclusions for all random active streams \(r\le N_m\).

Finally, if \(\lambda_m^{\rm cat}=A_{\rm cat}\log(e+m+J_m)\), then
\[
  J_m\exp(-c_{\rm E2E}\lambda_m^{\rm cat})
  \le
  (e+m+J_m)^{1-c_{\rm E2E}A_{\rm cat}},
\]
which is summable after increasing \(A_{\rm cat}\).  This leaves only the two
moment budgets and the deterministic inverse-\(p\), matrix, and sign schedules.
\end{proof}

\begin{corollary}[Concrete end-to-end automatic active-prefix schedules]
\label{cor:random-environment-invp-end-to-end-auto-active-schedules}
Use the notation, active-count setup, activation moment hypotheses, and catalog
hypotheses of
Corollary~\ref{cor:random-environment-invp-end-to-end-auto-active-certificates},
and put
\[
  \sigma_p:=1+\frac{(2-p)_+}{2}.
\]
Choose the constants \(A,Q,A_{\rm cat}\) in that corollary sufficiently large
and take
\[
  \lambda_m^{\rm cat}:=A_{\rm cat}\log(e+m+J_m).
\]
Assume \(\bar b_m^{*,J}\le\alpha_m/4\) eventually.  Then, with
\(u_m=q\ell_m\), the fixed-fraction, vanishing-fraction, and all-robust
polynomial-batch Toeplitz/Hankel sampled-certificate conclusions of
Corollary~\ref{cor:random-environment-invp-end-to-end-auto-active-certificates}
hold for every randomly active stream \(1\le r\le N_m\), eventually almost
surely, in each of the following regimes.

\begin{enumerate}[label=\textup{(\roman*)},leftmargin=2.2em]
\item All integer dimensions.  Suppose \(d_m=m\) and, for some
\(\Gamma>0\) and \(\eta_{\rm act}\ge0\),
\[
  \nu_m\le C_\nu m^\Gamma,\qquad
  G_m\le C_G m^{\eta_{\rm act}},\qquad
  s_{\rm act}\Gamma>\eta_{\rm act}+1 .
\]
Choose
\[
  A_0>\Gamma+1,\qquad q>\Gamma+1,\qquad
  \alpha_m=c_\alpha m^{-B}\quad\text{with }B>1 .
\]
It is enough that
\[
  \sum_m
  \frac{m^{Bs_{\rm cat}}(\log(e+m))^{s_{\rm cat}-1}H_m^{\rm cat}}
       {n_m^{s_{\rm cat}-1}}<\infty .
\]

\item Polynomial dimension schedules.  Suppose
\[
  d_m\ge c_0m^\kappa,\qquad
  \nu_m\le C_\nu d_m^\Gamma,\qquad
  G_m\le C_G d_m^{\eta_{\rm act}},
\]
where \(\Gamma>0\), \(\kappa>0\), and
\[
  \kappa(s_{\rm act}\Gamma-\eta_{\rm act})>1 .
\]
Choose
\[
  A_0>\Gamma+\frac1\kappa,\qquad
  q>\Gamma+\frac1\kappa,\qquad
  \alpha_m=c_\alpha d_m^{-B}\quad\text{with }B\kappa>1 .
\]
It is enough that
\[
  \sum_m
  \frac{d_m^{Bs_{\rm cat}}\ell_m^{s_{\rm cat}-1}H_m^{\rm cat}}
       {n_m^{s_{\rm cat}-1}}<\infty .
\]

\item Stretched-exponential dimensions with polynomial-in-\(d\) active
growth.  Suppose
\[
  d_m\ge\exp(c_0m^\kappa),\qquad
  \nu_m\le C_\nu d_m^\Gamma,\qquad
  G_m\le C_G d_m^{\eta_{\rm act}},
\]
where \(\Gamma>0\), \(\kappa>0\), and
\[
  s_{\rm act}\Gamma>\eta_{\rm act}.
\]
Choose
\[
  A_0>\Gamma,\qquad q>\Gamma,\qquad
  \alpha_m=c_\alpha\ell_m^{-B}
  \quad\text{with } B>\sigma_p+\frac1\kappa .
\]
It is enough that
\[
  \sum_m
  \frac{\ell_m^{Bs_{\rm cat}+s_{\rm cat}-1}H_m^{\rm cat}}
       {n_m^{s_{\rm cat}-1}}<\infty .
\]

\item Stretched-exponential dimensions with logarithmic active growth.  Suppose
\[
  d_m\ge\exp(c_0m^\kappa),\qquad
  \nu_m\le C_\nu\ell_m^\Gamma,\qquad
  G_m\le C_G\ell_m^{\eta_{\rm act}},
\]
where \(\Gamma>0\), \(\kappa>0\), and
\[
  \kappa(s_{\rm act}\Gamma-\eta_{\rm act})>1 .
\]
Choose
\[
  A_0>0,\qquad q>0,\qquad
  \alpha_m=c_\alpha\ell_m^{-B}
  \quad\text{with } B>\sigma_p+\frac1\kappa .
\]
It is enough that
\[
  \sum_m
  \frac{\ell_m^{Bs_{\rm cat}}(\log(e+m))^{s_{\rm cat}-1}H_m^{\rm cat}}
       {n_m^{s_{\rm cat}-1}}<\infty .
\]
\end{enumerate}
\end{corollary}

\begin{proof}
The active-count truncation parameter in
Corollary~\ref{cor:random-environment-invp-end-to-end-auto-active-certificates}
is \(A\log(e+m)\).  With \(A\) large, its exponential term is summable; the
displayed activation growth assumptions give the remaining active-intensity
moment series.  Indeed the four summands are bounded, respectively, by
\[
  C m^{\eta_{\rm act}-s_{\rm act}\Gamma}(\log(e+m))^{s_{\rm act}-1},
\]
\[
  C m^{-\kappa(s_{\rm act}\Gamma-\eta_{\rm act})}
    (\log(e+m))^{s_{\rm act}-1},
\]
a stretched-exponentially summable sequence, and
\[
  C m^{-\kappa(s_{\rm act}\Gamma-\eta_{\rm act})}
    (\log(e+m))^{s_{\rm act}-1}.
\]
The strict inequalities in the statement make these series finite.

The automatic prefix satisfies the same growth as the displayed envelope:
\[
  J_m\le C m^\Gamma,\quad
  J_m\le C d_m^\Gamma,\quad
  J_m\le C d_m^\Gamma,\quad
  J_m\le C\ell_m^\Gamma
\]
in the four cases.  Hence \(J_md_m^{-A_0}\) and \(J_md_m^{-q}\) are summable
under the stated bounds on \(A_0\) and \(q\): by \(p\)-series in the first two
cases, by stretched-exponential decay with \(A_0,q>\Gamma\) in the third, and by
stretched-exponential decay against a logarithmic prefix in the fourth.

It remains to check the active-menu catalog moment series with
\(\lambda_m^{\rm cat}=A_{\rm cat}\log(e+m+J_m)\).  In the all-integer case
\(\log(e+m+J_m)\le C\log(e+m)\); in the polynomial and
stretched-exponential polynomial-growth cases it is at most \(C\ell_m\); in the
logarithmic-growth case it is at most \(C\log(e+m)\).  Substituting the four
displayed choices of \(\alpha_m\) gives exactly the four catalog moment series
in the statement.  The same large choice of \(A_{\rm cat}\) makes the
exponential catalog-menu series summable by
Corollary~\ref{cor:random-environment-invp-end-to-end-auto-active-certificates}.
The displayed choices of \(\alpha_m\) are precisely the deterministic
inverse-\(p\) schedules in
Corollary~\ref{cor:random-environment-invp-operator-batch-maxima-schedules}.
The cited end-to-end corollary now gives all claimed sampled-certificate
conclusions.
\end{proof}

\begin{corollary}[Sample-size forms of the end-to-end active schedules]
\label{cor:random-environment-invp-end-to-end-auto-active-sample-sizes}
Consider any one of the four regimes in
Corollary~\ref{cor:random-environment-invp-end-to-end-auto-active-schedules},
and keep all of its hypotheses except for the displayed catalog moment
summability assumption.  Fix \(\epsilon>0\).  In the corresponding regime, that
omitted catalog summability follows from the following simpler aggregate
catalog-growth and sample-size conditions.

\begin{enumerate}[label=\textup{(\roman*)},leftmargin=2.2em]
\item All integer dimensions.  If
\[
  H_m^{\rm cat}\le C_Hm^{\eta_{\rm cat}}
\]
and
\[
  n_m\ge
  Q_n m^{(Bs_{\rm cat}+\eta_{\rm cat}+1+\epsilon)/(s_{\rm cat}-1)}
   \log(e+m),
\]
eventually, then the all-integer catalog moment series is summable.

\item Polynomial dimension schedules.  If
\[
  H_m^{\rm cat}\le C_Hd_m^{\eta_{\rm cat}}
\]
and
\[
  n_m\ge
  Q_n d_m^{(Bs_{\rm cat}+\eta_{\rm cat}+1/\kappa+\epsilon)/(s_{\rm cat}-1)}
   \ell_m,
\]
eventually, then the polynomial-dimension catalog moment series is summable.

\item Stretched-exponential dimensions with polynomial-in-\(d\) active growth.
If
\[
  H_m^{\rm cat}\le C_Hd_m^{\eta_{\rm cat}}
\]
and
\[
  n_m\ge
  Q_n d_m^{\eta_{\rm cat}/(s_{\rm cat}-1)}
   \ell_m^{(Bs_{\rm cat}+s_{\rm cat}-1+1/\kappa+\epsilon)/(s_{\rm cat}-1)},
\]
eventually, then the stretched-exponential polynomial-active catalog moment
series is summable.

\item Stretched-exponential dimensions with logarithmic active growth.  If
\[
  H_m^{\rm cat}\le C_H\ell_m^{\eta_{\rm cat}}
\]
and
\[
  n_m\ge
  Q_n \ell_m^{(Bs_{\rm cat}+\eta_{\rm cat}+1/\kappa+\epsilon)/(s_{\rm cat}-1)}
   \log(e+m),
\]
eventually, then the logarithmic-active catalog moment series is summable.
\end{enumerate}
Here \(C_H,Q_n>0\) and \(\eta_{\rm cat}\ge0\) are deterministic constants.  Thus,
whenever the corresponding lower bound is compatible with the ambient
requirements \(n_m\le d_m\), \(d_m/n_m\to\infty\), and
\(\log d_m=o(n_m)\), these sample-size tests may replace the catalog moment
series in
Corollary~\ref{cor:random-environment-invp-end-to-end-auto-active-schedules}.
The same fixed-fraction, vanishing-fraction, and all-robust polynomial-batch
Toeplitz/Hankel sampled certificates then hold for every random active stream
\(1\le r\le N_m\), eventually almost surely.
\end{corollary}

\begin{proof}
In the all-integer case, the first two displays give
\[
  \frac{m^{Bs_{\rm cat}}(\log(e+m))^{s_{\rm cat}-1}H_m^{\rm cat}}
       {n_m^{s_{\rm cat}-1}}
  \le C m^{-1-\epsilon},
\]
which is summable.

In the polynomial-dimension case, the corresponding term is at most
\[
  C d_m^{-1/\kappa-\epsilon}
  \le C' m^{-1-\kappa\epsilon},
\]
because \(d_m\ge c_0m^\kappa\).  This is summable.

In the stretched-exponential polynomial-active case, the displayed sample-size
bound gives
\[
  \frac{\ell_m^{Bs_{\rm cat}+s_{\rm cat}-1}H_m^{\rm cat}}
       {n_m^{s_{\rm cat}-1}}
  \le C\ell_m^{-1/\kappa-\epsilon}.
\]
Since \(d_m\ge\exp(c_0m^\kappa)\), we have \(\ell_m\ge c m^\kappa\) eventually,
so the last display is \(O(m^{-1-\kappa\epsilon})\).

The logarithmic-active case is the same calculation:
\[
  \frac{\ell_m^{Bs_{\rm cat}}(\log(e+m))^{s_{\rm cat}-1}H_m^{\rm cat}}
       {n_m^{s_{\rm cat}-1}}
  \le C\ell_m^{-1/\kappa-\epsilon}
  \le C' m^{-1-\kappa\epsilon}.
\]
Thus the omitted catalog moment series in
Corollary~\ref{cor:random-environment-invp-end-to-end-auto-active-schedules}
is finite in each regime.  Applying that corollary gives the stated sampled
certificate conclusions.
\end{proof}

\begin{corollary}[Compatible sample windows for catalog-growth exponents]
\label{cor:random-environment-invp-end-to-end-auto-active-compatible-windows}
In the setting of
Corollary~\ref{cor:random-environment-invp-end-to-end-auto-active-schedules},
write \(s=s_{\rm cat}\).  The catalog moment hypotheses in that corollary are
compatible with the ambient unbounded-aspect requirements
\[
  n_m\le d_m,\qquad d_m/n_m\to\infty,\qquad \log d_m=o(n_m)
\]
in the following explicit windows.

\begin{enumerate}[label=\textup{(\roman*)},leftmargin=2.2em]
\item All integer dimensions.  Suppose \(d_m=m\) and
\[
  H_m^{\rm cat}\le C_Hm^\gamma
\]
for a real exponent \(\gamma\).  If
\[
  \gamma<-s(B-1)-2,
\]
then choose
\[
  \max\left\{0,\frac{Bs+\gamma+1}{s-1}\right\}<\rho<\rho'<1 .
\]
Every sample-size sequence satisfying
\[
  Q_n m^\rho\log(e+m)\le n_m\le m^{\rho'}
\]
eventually has the ambient unbounded-aspect properties and makes the
all-integer catalog moment series summable.

\item Polynomial dimension schedules.  Suppose \(d_m\ge c_0m^\kappa\) and
\[
  H_m^{\rm cat}\le C_Hd_m^\gamma .
\]
If
\[
  \gamma<s-1-Bs-\frac1\kappa,
\]
then choose
\[
  \max\left\{0,\frac{Bs+\gamma+1/\kappa}{s-1}\right\}<\rho<\rho'<1 .
\]
Every sample-size sequence satisfying
\[
  Q_n d_m^\rho\ell_m\le n_m\le d_m^{\rho'}
\]
eventually has the ambient unbounded-aspect properties and makes the polynomial
catalog moment series summable.

\item Stretched-exponential dimensions with polynomial-in-\(d\) active growth.
Suppose \(d_m\ge\exp(c_0m^\kappa)\) and
\[
  H_m^{\rm cat}\le C_Hd_m^\gamma .
\]
If \(\gamma<s-1\), choose
\[
  \max\left\{0,\frac{\gamma}{s-1}\right\}<\rho<\rho'<1 .
\]
Every sample-size sequence satisfying
\[
  Q_n d_m^\rho\ell_m\le n_m\le d_m^{\rho'}
\]
eventually has the ambient unbounded-aspect properties and makes the
stretched-exponential polynomial-active catalog moment series summable.

\item Stretched-exponential dimensions with logarithmic active growth.  Suppose
\(d_m\ge\exp(c_0m^\kappa)\) and
\[
  H_m^{\rm cat}\le C_H\ell_m^\gamma .
\]
For every real \(\gamma\), choose
\[
  r>\max\left\{1,\frac{Bs+\gamma+1/\kappa}{s-1}\right\}
\]
and any \(0<\rho'<1\).  Every sample-size sequence satisfying
\[
  Q_n\ell_m^r\log(e+m)\le n_m\le d_m^{\rho'}
\]
eventually has the ambient unbounded-aspect properties and makes the
logarithmic-active catalog moment series summable.
\end{enumerate}
In each case the fixed-fraction, vanishing-fraction, and all-robust
polynomial-batch Toeplitz/Hankel sampled certificates hold for every random
active stream \(1\le r\le N_m\), eventually almost surely.
\end{corollary}

\begin{proof}
The displayed sample windows are nonempty for all sufficiently large \(m\).  In
the first three cases this follows from \(\rho<\rho'\).  In the fourth case it
uses the stretched-exponential lower bound on \(d_m\), which makes every fixed
power of \(\ell_m\) negligible compared with \(d_m^{\rho'}\).  The ambient
properties are immediate from the lower bounds, which dominate \(\log d_m\), and
from the upper bounds, which are \(o(d_m)\).

It remains only to check summability of the catalog series from
Corollary~\ref{cor:random-environment-invp-end-to-end-auto-active-schedules}.
In the all-integer case,
\[
  \frac{m^{Bs}(\log(e+m))^{s-1}H_m^{\rm cat}}{n_m^{s-1}}
  \le
  C m^{Bs+\gamma-(s-1)\rho}.
\]
The choice of \(\rho\) makes the exponent less than \(-1\).

In the polynomial case,
\[
  \frac{d_m^{Bs}\ell_m^{s-1}H_m^{\rm cat}}{n_m^{s-1}}
  \le
  C d_m^{Bs+\gamma-(s-1)\rho}
  \le C' m^{-\kappa((s-1)\rho-Bs-\gamma)} .
\]
The choice of \(\rho\) makes
\(\kappa((s-1)\rho-Bs-\gamma)>1\), so the series is finite.

In the stretched-exponential polynomial-active case,
\[
  \frac{\ell_m^{Bs+s-1}H_m^{\rm cat}}{n_m^{s-1}}
  \le
  C d_m^{\gamma-(s-1)\rho},
\]
up to a harmless power of \(\ell_m\).  The exponent of \(d_m\) is negative, so
the summand decays stretched-exponentially.

Finally, in the logarithmic-active case,
\[
  \frac{\ell_m^{Bs}(\log(e+m))^{s-1}H_m^{\rm cat}}{n_m^{s-1}}
  \le
  C\ell_m^{Bs+\gamma-(s-1)r}.
\]
The choice of \(r\) makes the exponent less than \(-1/\kappa\); since
\(\ell_m\ge c m^\kappa\) eventually, the resulting series is summable.
Applying
Corollary~\ref{cor:random-environment-invp-end-to-end-auto-active-schedules}
proves the sampled-certificate conclusions.
\end{proof}

\begin{corollary}[Per-stream catalog budgets for compatible active windows]
\label{cor:random-environment-invp-end-to-end-auto-active-per-stream-catalog-windows}
Use the notation and hypotheses of
Corollary~\ref{cor:random-environment-invp-end-to-end-auto-active-schedules},
except for its catalog moment summability assumption, and write
\(s=s_{\rm cat}\).  For the active prefix \(J_m\), set
\[
  \bar h_m^{*,J}:=
  \max_{1\le r\le J_m}\frac1{n_m}\sum_{i=1}^{n_m}h_{m,i}^{(r)} .
\]
Then \(H_m^{\rm cat}\le J_m\bar h_m^{*,J}\).  Consequently the
compatible sample-window conclusions of
Corollary~\ref{cor:random-environment-invp-end-to-end-auto-active-compatible-windows}
hold under the following per-stream row-catalog conditions.

\begin{enumerate}[label=\textup{(\roman*)},leftmargin=2.2em]
\item All integer dimensions.  Suppose \(d_m=m\), \(\nu_m\le C_\nu m^\Gamma\),
and
\[
  \bar h_m^{*,J}\le C_h m^\theta
\]
for a real exponent \(\theta\).  If
\[
  \Gamma+\theta<-s(B-1)-2,
\]
then choose
\[
  \max\left\{0,\frac{Bs+\Gamma+\theta+1}{s-1}\right\}<\rho<\rho'<1 .
\]
Every sample-size sequence satisfying
\[
  Q_n m^\rho\log(e+m)\le n_m\le m^{\rho'}
\]
eventually has the ambient unbounded-aspect properties and gives the
sampled-certificate conclusions for every random active stream.

\item Polynomial dimension schedules.  Suppose \(d_m\ge c_0m^\kappa\),
\(\nu_m\le C_\nu d_m^\Gamma\), and
\[
  \bar h_m^{*,J}\le C_h d_m^\theta .
\]
If
\[
  \Gamma+\theta<s-1-Bs-\frac1\kappa,
\]
then choose
\[
  \max\left\{0,\frac{Bs+\Gamma+\theta+1/\kappa}{s-1}\right\}<\rho<\rho'<1 .
\]
Every sample-size sequence satisfying
\[
  Q_n d_m^\rho\ell_m\le n_m\le d_m^{\rho'}
\]
eventually has the ambient unbounded-aspect properties and gives the
sampled-certificate conclusions for every random active stream.

\item Stretched-exponential dimensions with polynomial-in-\(d\) active growth.
Suppose \(d_m\ge\exp(c_0m^\kappa)\), \(\nu_m\le C_\nu d_m^\Gamma\), and
\[
  \bar h_m^{*,J}\le C_h d_m^\theta .
\]
If \(\Gamma+\theta<s-1\), choose
\[
  \max\left\{0,\frac{\Gamma+\theta}{s-1}\right\}<\rho<\rho'<1 .
\]
Every sample-size sequence satisfying
\[
  Q_n d_m^\rho\ell_m\le n_m\le d_m^{\rho'}
\]
eventually has the ambient unbounded-aspect properties and gives the
sampled-certificate conclusions for every random active stream.

\item Stretched-exponential dimensions with logarithmic active growth.  Suppose
\(d_m\ge\exp(c_0m^\kappa)\), \(\nu_m\le C_\nu\ell_m^\Gamma\), and
\[
  \bar h_m^{*,J}\le C_h \ell_m^\theta .
\]
For every real \(\theta\), choose
\[
  r>\max\left\{1,\frac{Bs+\Gamma+\theta+1/\kappa}{s-1}\right\}
\]
and any \(0<\rho'<1\).  Every sample-size sequence satisfying
\[
  Q_n\ell_m^r\log(e+m)\le n_m\le d_m^{\rho'}
\]
eventually has the ambient unbounded-aspect properties and gives the
sampled-certificate conclusions for every random active stream.
\end{enumerate}
In each case, with the remaining activation, inverse-\(p\), matrix/sign, and
mean-budget hypotheses from
Corollary~\ref{cor:random-environment-invp-end-to-end-auto-active-schedules},
the fixed-fraction, vanishing-fraction, and all-robust polynomial-batch
Toeplitz/Hankel certificates hold for every random active stream \(1\le r\le
N_m\), eventually almost surely.
\end{corollary}

\begin{proof}
Writing
\[
  \bar h_{m,r}:=\frac1{n_m}\sum_{i=1}^{n_m}h_{m,i}^{(r)},
\]
the aggregate catalog moment budget satisfies
\[
  H_m^{\rm cat}=\sum_{r=1}^{J_m}\bar h_{m,r}
  \le J_m\bar h_m^{*,J}.
\]
The automatic-prefix growth estimates in the proof of
Corollary~\ref{cor:random-environment-invp-end-to-end-auto-active-schedules}
give \(J_m\le C m^\Gamma\) in the all-integer regime, \(J_m\le C d_m^\Gamma\)
in the polynomial and stretched-exponential polynomial-active regimes, and
\(J_m\le C\ell_m^\Gamma\) in the stretched-exponential logarithmic-active
regime.  Thus the aggregate hypotheses of
Corollary~\ref{cor:random-environment-invp-end-to-end-auto-active-compatible-windows}
hold with catalog exponent \(\gamma=\Gamma+\theta\) in the corresponding
scale.  Applying that corollary proves all four sample-window conclusions and
the resulting sampled certificates.
\end{proof}

\begin{corollary}[First-moment-tuned per-stream compatible active windows]
\label{cor:random-environment-invp-end-to-end-auto-active-mean-tuned-windows}
Use the activation setup, catalog laws, and notation of
Corollary~\ref{cor:random-environment-invp-end-to-end-auto-active-schedules},
but replace its mean-budget and catalog summability assumptions by the
per-stream assumptions below.  Write \(s=s_{\rm cat}\) and
\[
  \sigma_p:=1+\frac{(2-p)_+}{2}.
\]
For the automatically chosen active prefix \(J_m\), set
\[
  \bar b_m^{*,J}:=
  \max_{1\le r\le J_m}\frac1{n_m}\sum_{i=1}^{n_m}b_{m,i}^{(r)},
  \qquad
  \bar h_m^{*,J}:=
  \max_{1\le r\le J_m}\frac1{n_m}\sum_{i=1}^{n_m}h_{m,i}^{(r)} .
\]
Then the compatible active-window conclusions hold in the following
first-moment-tuned regimes.

\begin{enumerate}[label=\textup{(\roman*)},leftmargin=2.2em]
\item All integer dimensions.  Suppose \(d_m=m\),
\[
  \nu_m\le C_\nu m^\Gamma,\qquad
  G_m\le C_G m^{\eta_{\rm act}},\qquad
  s_{\rm act}\Gamma>\eta_{\rm act}+1,
\]
and, for real exponents \(\beta,\theta\),
\[
  \bar b_m^{*,J}\le C_b m^{-\beta},\qquad
  \bar h_m^{*,J}\le C_h m^\theta .
\]
If
\[
  \beta>1,\qquad \Gamma+\theta<-2,
\]
choose \(B\) so that
\[
  1<B<\beta,\qquad \Gamma+\theta<-s(B-1)-2,
\]
take \(A_0,q>\Gamma+1\) and \(\alpha_m=c_\alpha m^{-B}\), and choose
\[
  \max\left\{0,\frac{Bs+\Gamma+\theta+1}{s-1}\right\}<\rho<\rho'<1 .
\]
Every sample-size sequence satisfying
\[
  Q_n m^\rho\log(e+m)\le n_m\le m^{\rho'}
\]
eventually gives the sampled Toeplitz/Hankel certificate conclusions for every
random active stream.

\item Polynomial dimension schedules.  Suppose \(d_m\ge c_0m^\kappa\),
\[
  \nu_m\le C_\nu d_m^\Gamma,\qquad
  G_m\le C_G d_m^{\eta_{\rm act}},\qquad
  \kappa(s_{\rm act}\Gamma-\eta_{\rm act})>1,
\]
and
\[
  \bar b_m^{*,J}\le C_b d_m^{-\beta},\qquad
  \bar h_m^{*,J}\le C_h d_m^\theta .
\]
If
\[
  \beta>\frac1\kappa,\qquad
  \Gamma+\theta<s\left(1-\frac1\kappa\right)-1-\frac1\kappa,
\]
choose \(B\) so that
\[
  \frac1\kappa<B<\beta,\qquad
  \Gamma+\theta<s-1-Bs-\frac1\kappa,
\]
take \(A_0,q>\Gamma+1/\kappa\) and \(\alpha_m=c_\alpha d_m^{-B}\), and choose
\[
  \max\left\{0,\frac{Bs+\Gamma+\theta+1/\kappa}{s-1}\right\}<\rho<\rho'<1 .
\]
Every sample-size sequence satisfying
\[
  Q_n d_m^\rho\ell_m\le n_m\le d_m^{\rho'}
\]
eventually gives the sampled Toeplitz/Hankel certificate conclusions for every
random active stream.

\item Stretched-exponential dimensions with polynomial-in-\(d\) active growth.
Suppose \(d_m\ge\exp(c_0m^\kappa)\),
\[
  \nu_m\le C_\nu d_m^\Gamma,\qquad
  G_m\le C_G d_m^{\eta_{\rm act}},\qquad
  s_{\rm act}\Gamma>\eta_{\rm act},
\]
and
\[
  \bar b_m^{*,J}\le C_b \ell_m^{-\beta},\qquad
  \bar h_m^{*,J}\le C_h d_m^\theta .
\]
If
\[
  \beta>\sigma_p+\frac1\kappa,\qquad \Gamma+\theta<s-1,
\]
choose \(B\) so that
\[
  \sigma_p+\frac1\kappa<B<\beta,
\]
take \(A_0,q>\Gamma\) and \(\alpha_m=c_\alpha\ell_m^{-B}\), and choose
\[
  \max\left\{0,\frac{\Gamma+\theta}{s-1}\right\}<\rho<\rho'<1 .
\]
Every sample-size sequence satisfying
\[
  Q_n d_m^\rho\ell_m\le n_m\le d_m^{\rho'}
\]
eventually gives the sampled Toeplitz/Hankel certificate conclusions for every
random active stream.

\item Stretched-exponential dimensions with logarithmic active growth.  Suppose
\(d_m\ge\exp(c_0m^\kappa)\),
\[
  \nu_m\le C_\nu\ell_m^\Gamma,\qquad
  G_m\le C_G\ell_m^{\eta_{\rm act}},\qquad
  \kappa(s_{\rm act}\Gamma-\eta_{\rm act})>1,
\]
and
\[
  \bar b_m^{*,J}\le C_b \ell_m^{-\beta},\qquad
  \bar h_m^{*,J}\le C_h \ell_m^\theta .
\]
If
\[
  \beta>\sigma_p+\frac1\kappa,
\]
choose \(B\) so that
\[
  \sigma_p+\frac1\kappa<B<\beta,
\]
take arbitrary \(A_0,q>0\), \(\alpha_m=c_\alpha\ell_m^{-B}\), choose any
\(0<\rho'<1\), and choose
\[
  r>\max\left\{1,\frac{Bs+\Gamma+\theta+1/\kappa}{s-1}\right\}.
\]
Every sample-size sequence satisfying
\[
  Q_n\ell_m^r\log(e+m)\le n_m\le d_m^{\rho'}
\]
eventually gives the sampled Toeplitz/Hankel certificate conclusions for every
random active stream.
\end{enumerate}
In all four cases the conclusions are the fixed-fraction, vanishing-fraction,
and all-robust polynomial-batch sampled certificates of
Corollary~\ref{cor:random-environment-invp-end-to-end-auto-active-schedules},
simultaneously for every randomly active stream \(1\le r\le N_m\), eventually
almost surely.
\end{corollary}

\begin{proof}
The displayed strict inequalities ensure that the stated interval for \(B\) is
nonempty in each case.  Once such a \(B\) is fixed, the first-moment bounds imply
\[
  \bar b_m^{*,J}\le \alpha_m/4
\]
eventually, because the assumed first-moment decay exponent is strictly larger
than the chosen inverse-\(p\) exponent.  The activation and matrix/sign
summability hypotheses are exactly the corresponding hypotheses of
Corollary~\ref{cor:random-environment-invp-end-to-end-auto-active-schedules}
with this \(B\) and the displayed choices of \(A_0,q\).

The \(s\)-moment bounds and the automatic prefix estimates give the per-stream
catalog hypotheses of
Corollary~\ref{cor:random-environment-invp-end-to-end-auto-active-per-stream-catalog-windows}.
Indeed the aggregate catalog exponent there is \(\Gamma+\theta\) in the
appropriate scale, and the displayed choices of \(B,\rho,\rho'\) or \(B,r\)
are precisely its compatibility inequalities.  Applying that corollary gives
the sample windows, ambient unbounded-aspect properties, and all sampled
certificate conclusions.
\end{proof}

\begin{corollary}[Convex-concentrated sampled operator certificates]
\label{cor:convex-normalizers-sampled-operator}
In the setting of Theorem~\ref{thm:convex-normalizers}, let
\((d_m,n_m)\) be any sequence satisfying
\[
  n_m\le d_m,\qquad d_m/n_m\to\infty,\qquad \log d_m=o(n_m).
\]
For each \(m\), let \(B_{m,i}=T_{d_m}(\xi^{(m,i)})\),
\(1\le i\le n_m\), be independent full Toeplitz matrices from the law, put
\[
  A_{m,i}:=\frac{B_{m,i}}{\norm{B_{m,i}}_{\op}},
  \qquad
  \widetilde A_{m,i}:=J_{d_m}A_{m,i},
\]
and sample independent uniform signs
\(\delta_m^{(1)},\ldots,\delta_m^{(M_m)}\), independent of the matrices, with
\(M_m\ge \lceil B\rceil\) for some fixed \(B>0\).  Then, with probability
tending to one jointly over the matrices and sampled signs, some sampled sign
\(\delta_m^{(s)}\) is the center of a Hamming ball of radius
\[
  \left\lfloor c\log(e+d_m)\right\rfloor
\]
such that every \(\eps'\) in that ball satisfies
\[
  \left\|\sum_{i=1}^{n_m}\eps'_iA_{m,i}\right\|_{\op}
  \le C\sqrt{n_m},
  \qquad
  \left\|\sum_{i=1}^{n_m}\eps'_i\widetilde A_{m,i}\right\|_{\op}
  \le C\sqrt{n_m}.
\]
Here \(C<\infty\) and \(c>0\) depend only on the coefficient-law parameters and
on the fixed confidence choices.  More quantitatively, before the deterministic
edge-to-operator transfer, for every \(A_0>0\) the edge-certificate failure is
at most
\[
  C_{A_0,K,C_{\rm cc}}d_m^{-A_0}
  +(\log(e+d_m))^2(e+d_m)^{-\beta}
  +d_m^{-B},
\]
for some \(\beta=\beta(K,C_{\rm cc})>0\), once \(m\) is large enough.  Thus the
same conclusion applies to standard real or complex Gaussian coefficients,
Rademacher coefficients, and bounded real or complex coefficient laws with mean
zero and variance one.
\end{corollary}

\begin{proof}
The proof of Theorem~\ref{thm:convex-normalizers} gives \(a=1/2\) and
\[
  p_d(a)=\Prob\{\norm{T_d(\xi)}_{\op}<a\,\Exp\norm{T_d(\xi)}_{\op}\}
  \le (e+d)^{-\beta}
\]
for some \(\beta=\beta(K,C_{\rm cc})>0\).  Since
\(\log d_m=o(n_m)\), the size condition \(n_m\ge C\log(e+d_m)\) in
Corollary~\ref{cor:normalizer-tail-poly-confidence} holds eventually.  Applying
that corollary at \((d_m,n_m)\), with the fixed sign-confidence parameter \(B\),
shows that with failure at most the displayed quantity, one sampled sign is the
center of a Hamming ball of radius \(c\log(e+d_m)\) on which every sign has
realization-normalized edge maximum \(O(\sqrt{n_m})\); here
\(\min\{n_m,\log^2(e+d_m)\}\le \log^2(e+d_m)\) was used.

On the edge-success event, apply
Corollary~\ref{cor:edge-to-operator-robust} to that Hamming ball.  The
Toeplitz operator estimate follows from the deterministic circulant reduction,
and the Hankel estimate follows from the reversal \(J_{d_m}\), with only a
constant enlargement of the \(O(\sqrt{n_m})\) threshold.  The displayed failure
probability tends to zero, so the joint sampled-operator certificate holds with
probability tending to one.  The named examples are those already covered by
Theorem~\ref{thm:convex-normalizers}.
\end{proof}

\begin{corollary}[Convex-concentrated operator-certificate abundance]
\label{cor:convex-normalizers-operator-abundance}
In the setting of Corollary~\ref{cor:convex-normalizers-sampled-operator}, write
\(\ell_m=\log(e+d_m)\).

First, for every \(0<\rho\le1\) and \(S>0\), there is an integer
\(M_{\rho,S}\) and constants \(C_{\rho,S}<\infty\), \(c_{\rho,S}>0\) such that,
if \(M_{\rho,S}\) independent uniform signs are sampled at level \(m\), then
with probability tending to one at least a \(1-\rho\) fraction of the sampled
signs are centers of Hamming balls of radius
\[
  \left\lfloor c_{\rho,S}\ell_m\right\rfloor
\]
on which every sign vector \(\eps'\) satisfies
\[
  \left\|\sum_{i=1}^{n_m}\eps'_iA_{m,i}\right\|_{\op}
  \le C_{\rho,S}\sqrt{n_m},
  \qquad
  \left\|\sum_{i=1}^{n_m}\eps'_i\widetilde A_{m,i}\right\|_{\op}
  \le C_{\rho,S}\sqrt{n_m}.
\]
More quantitatively, before the deterministic edge-to-operator transfer, for
every \(A_0>0\) the edge robust-fraction failure is at most
\[
  C_{\rho,S,A_0,K,C_{\rm cc}}d_m^{-A_0}
  +C_{\rho,S,A_0,K,C_{\rm cc}}\ell_m^2(e+d_m)^{-\beta}
  +2d_m^{-S},
\]
for some \(\beta=\beta(K,C_{\rm cc})>0\), once \(m\) is large enough.

Second, for every polynomial batch exponent \(\theta\ge0\) and confidence
exponent \(\sigma>0\), there are constants
\(C_{\theta,\sigma}<\infty\), \(c_{\theta,\sigma}>0\) such that, whenever
\[
  1\le M_m\le d_m^\theta,
\]
then with probability tending to one every sampled sign
\(\delta_m^{(s)}\), \(1\le s\le M_m\), is the center of a Hamming ball of
radius
\[
  \left\lfloor c_{\theta,\sigma}\ell_m\right\rfloor
\]
on which every \(\eps'\) satisfies the two Toeplitz and Hankel operator-norm
bounds above, with \(C_{\theta,\sigma}\) in place of \(C_{\rho,S}\).  Before
the deterministic transfer, the corresponding edge failure is at most
\[
  C_{\theta,\sigma,A_0,K,C_{\rm cc}}d_m^{-A_0}
  +C_{\theta,\sigma,A_0,K,C_{\rm cc}}\ell_m^2(e+d_m)^{-\beta}
  +d_m^{-\sigma}.
\]
\end{corollary}

\begin{proof}
As in the proof of Corollary~\ref{cor:convex-normalizers-sampled-operator}, the
convex-concentration theorem supplies \(a=1/2\) and
\[
  p_d(a)\le (e+d)^{-\beta}
\]
for some \(\beta=\beta(K,C_{\rm cc})>0\).  Since \(\ell_m=o(n_m)\), all size
conditions of the normalizer-tail sampled-search corollaries hold for large
\(m\).

For the robust-fraction statement, choose \(B\) and \(M_{\rho,S}\) as in
Corollary~\ref{cor:normalizer-tail-robust-fraction-confidence}, and apply that
corollary with \(u=S\ell_m\) and \(a=1/2\).  At level \(m\), outside the
displayed edge failure, at least a \(1-\rho\) fraction of the sampled signs are
centers of radius \(c\ell_m\) Hamming balls on which every sign has
realization-normalized edge maximum \(O(\sqrt{n_m})\).

For the polynomial-batch statement, apply
Corollary~\ref{cor:normalizer-tail-poly-batch} with \(a=1/2\), batch exponent
\(\theta\), and confidence exponent \(\sigma\).  Outside the second displayed
edge failure, every sampled sign in the batch is the center of such a
logarithmic-radius edge-robust Hamming ball.

In either case, let \(\mathcal F_m\) be the union of the Hamming balls supplied
on the edge-success event.  Corollary~\ref{cor:edge-to-operator-robust} applies
to these sign families for all sufficiently large \(m\), enlarging only the
constant in the \(O(\sqrt{n_m})\) bound and giving both the Toeplitz and Hankel
operator-norm estimates.  The displayed edge failures tend to zero, so the
operator-certificate abundance conclusions hold with probability tending to
one.
\end{proof}

\begin{corollary}[Convex-concentrated Borel operator batches]
\label{cor:convex-normalizers-borel-operator-batches}
In the setting of Theorem~\ref{thm:convex-normalizers}, let
\((d_m,n_m)\) satisfy
\[
  n_m\le d_m,\qquad d_m/n_m\to\infty,\qquad \log d_m=o(n_m),
\]
and put \(\ell_m=\log(e+d_m)\).  For each \(m\), let
\(B_{m,i}=T_{d_m}(\xi^{(m,i)})\), \(1\le i\le n_m\), be independent full
Toeplitz matrices from the law, and set
\[
  A_{m,i}:=\frac{B_{m,i}}{\norm{B_{m,i}}_{\op}},
  \qquad
  \widetilde A_{m,i}:=J_{d_m}A_{m,i}.
\]
Independently of the matrices, sample independent uniform signs
\[
  \delta_m^{(1)},\ldots,\delta_m^{(M_m)}\in\{-1,1\}^{n_m},
  \qquad M_m\ge1,
\]
and let \(u_m\ge0\), \(1\le r_m\le M_m\) be deterministic.  Let
\(\beta=\beta(K,C_{\rm cc})>0\) be the convex-concentration exponent from the
proof of Theorem~\ref{thm:convex-normalizers}.  If, for some \(A>0\),
\[
  \sum_{m=1}^{\infty}
  \left[
    d_m^{-A}+\ell_m^2(e+d_m)^{-\beta}
    +e^{-u_m}+\binom{M_m}{r_m}d_m^{-A r_m}
  \right]<\infty ,
\]
then, with probability one over the matrices and sampled signs, all but
finitely many \(m\) have at least \(M_m-r_m+1\) sampled signs that are centers
of Hamming balls of radius
\[
  \left\lfloor c\ell_m\right\rfloor
\]
on which every signing \(\eps'\) satisfies
\[
  \left\|\sum_{i=1}^{n_m}\eps'_iA_{m,i}\right\|_{\op}
  \le C\sqrt{n_m},
  \qquad
  \left\|\sum_{i=1}^{n_m}\eps'_i\widetilde A_{m,i}\right\|_{\op}
  \le C\sqrt{n_m}.
\]
On the same eventual event, the whole sampled list has the entropy-sensitive
edge-grid maximum bound
\[
  \max_{1\le s\le M_m}\max_{\omega^{d_m}=1}
  \left|\sum_{i=1}^{n_m}\delta_{m,i}^{(s)}
  \frac{P_{m,i}(\omega)}{\norm{B_{m,i}}_{\op}}\right|
  \le
  C\sqrt{n_m\left(1+\frac{\log(eM_m)+u_m}{\ell_m}\right)} .
\]

Consequently, on any stretched-exponential dimension schedule
\[
  d_m\ge \exp(c_0m^\kappa)\qquad(c_0,\kappa>0),
\]
every sampled sign in every polynomial batch \(M_m\le d_m^\theta\) is
eventually a Toeplitz/Hankel operator-norm robust center almost surely, after
choosing any \(A>\theta\), taking \(r_m=1\), and taking
\(u_m=s\log(e+m)\) with \(s>1\).  On a polynomial schedule
\(d_m\ge c_0m^\kappa\), the same conclusion holds whenever
\(\beta\kappa>1\), after choosing \(A>\theta+1/\kappa\).
\end{corollary}

\begin{proof}
Theorem~\ref{thm:convex-normalizers} gives \(a=1/2\) and
\[
  p_d(a)\le (e+d)^{-\beta}.
\]
Since \(\ell_m=o(n_m)\), the size condition in
Corollary~\ref{cor:normalizer-tail-borel-counts} holds eventually.  Applying
that corollary from the first sufficiently large index onward, the displayed
summability condition implies that, almost surely, all but finitely many sampled
batches have the stated edge-grid list maximum and at least \(M_m-r_m+1\)
logarithmic-radius robust centers with \(O(\sqrt{n_m})\) edge maximum.
Corollary~\ref{cor:borel-edge-to-operator-batches} transfers exactly these same
centers and radii to realization-normalized full Toeplitz and full Hankel
operator norm, because \(d_m/n_m\to\infty\) and \(\log d_m=o(n_m)\).

For the stretched-exponential consequence, the matrix terms
\(d_m^{-A}\) and \(\ell_m^2(e+d_m)^{-\beta}\) are summable for every
\(A,\beta>0\), and \(e^{-u_m}\) is summable because \(s>1\).  With
\(r_m=1\) and \(M_m\le d_m^\theta\),
\[
  \binom{M_m}{1}d_m^{-A}\le d_m^{\theta-A},
\]
which is summable on the stretched-exponential schedule when \(A>\theta\).  On
the polynomial schedule, the lower-normalizer contribution is bounded by a
constant times \(m^{-\beta\kappa}\log^2(e+m)\), which is summable under
\(\beta\kappa>1\), while \(d_m^{\theta-A}\le C m^{-(A-\theta)\kappa}\) is
summable under \(A>\theta+1/\kappa\).  These choices make every sampled sign a
robust center eventually.
\end{proof}

\begin{theorem}[Gaussian realization-normalized edge profile]
\label{thm:gaussian-realized-profile}
Let \(B_i=T_d(g^{(i)})\), \(1\le i\le n\), be independent full Toeplitz
matrices whose lag variables are standard real Gaussians, and let \(P_i\) be
their symbols.  Assume \(n\le d\) and \(\log d=o(n)\).  Then
\[
  \max_{\omega^d=1}
  \sum_{i=1}^n
  \frac{\abs{P_i(\omega)}^2}{\norm{B_i}_{\op}^2}
  \le C\,\frac{n}{\log(e+d)}
\]
with probability tending to one.  Consequently, for the corresponding
realization-normalized edge vectors \(v_i(\omega)=P_i(\omega)/\norm{B_i}_{\op}\),
independent Rademacher signs satisfy, conditionally on the matrices and with
probability at least \(1-d^{-A}\),
\[
  \max_{\omega^d=1}\abs{\sum_i\eps_i v_i(\omega)}
  \le C_A\sqrt n
\]
for every fixed \(A>0\), on an event whose matrix probability tends to one.
\end{theorem}

\begin{proof}
Write \(\ell=\log(e+d)\) and \(N_d=\Exp\norm{T_d(g)}_{\op}\).  By Meckes'
lower bound, \(N_d\ge c\sqrt{d\ell}\).  It is enough, by
Theorem~\ref{thm:normalizer-tail}, to verify the one-matrix lower-tail
hypothesis with \(a=1/2\).
The map \(g\mapsto \norm{T_d(g)}_{\op}\) is \(\sqrt{2d}\)-Lipschitz for the
Euclidean norm of the \(2d-1\) lag variables, because
\[
  \norm{T_d(a)}_{\op}
  \le \sup_{\theta}\abs{\sum_{\abs k<d}a_ke^{ik\theta}}
  \le \sqrt{2d-1}\,\norm{a}_2 .
\]
Gaussian concentration therefore gives
\[
  \Prob\{\norm{T_d(g)}_{\op}<N_d/2\}\le d^{-\beta}
\]
for some absolute \(\beta>0\).  Thus the hypothesis of
Theorem~\ref{thm:normalizer-tail} holds, and the theorem gives the displayed
profile.  The final signed edge bound is Proposition~\ref{prop:variance-profile}.
\end{proof}

\begin{corollary}[Gaussian realization-normalized full Toeplitz discrepancy]
\label{cor:gaussian-realized-discrepancy}
Let \(B_i=T_d(g^{(i)})\), \(1\le i\le n\), be independent full Toeplitz
matrices whose lag variables are standard real Gaussians, and put
\(A_i=B_i/\norm{B_i}_{\op}\).  If
\[
  n\le d\le e^{o(n)}
  \qquad\text{and}\qquad d/n\to\infty ,
\]
then
\[
  \discop(A_1,\ldots,A_n)\le C\sqrt n
\]
with probability tending to one, with an absolute constant \(C\).  More
quantitatively, on a matrix event whose probability tends to one, independent
random signs achieve the bound \(C_A\sqrt n\) with conditional probability at
least \(1-d^{-A}\) for every fixed \(A>0\).  The same conclusion holds for the
corresponding full Hankel model.
\end{corollary}

\begin{proof}
The assumption \(d\le e^{o(n)}\) is exactly \(\log d=o(n)\), so
Theorem~\ref{thm:gaussian-realized-profile} gives the variance profile
required by Proposition~\ref{prop:variance-profile} with probability tending to
one over the matrices.  On this event, Proposition~\ref{prop:variance-profile}
shows that independent signs satisfy
\[
  \max_{\omega^d=1}\abs{\sum_i\eps_i
  P_i(\omega)/\norm{B_i}_{\op}}\le C_A\sqrt n
\]
with conditional probability at least \(1-d^{-A}\).  In particular such a sign
vector exists on the event.  Proposition~\ref{prop:reduction} then transfers
the edge bound from the circulant model back to full Toeplitz contractions,
with only the \(o(\sqrt n)\) de la Vall\'ee Poussin boundary error already
accounted for there.  The Hankel statement follows by applying the Toeplitz
result to the reversed matrices \(J_dB_i\), which have the same operator norms
and the same discrepancy.
\end{proof}

\subsection{The obstruction: the multi-symbol spectral edge}

For the general sub-Gaussian realization-normalized model, the remaining obstacle
is to control $\max_\omega\abs{\sum_i\eps_iv_i(\omega)}$ over the \(O(d)\) edge
nodes for a well-chosen sign pattern. For a fixed edge node \(\omega\), the
Rademacher variance is \(\sum_i\abs{v_i(\omega)}^2\).
Proposition~\ref{prop:variance-profile} shows that the square-root bound follows
as soon as this variance is \(O(n/\log d)\) uniformly in \(\omega\), since the
final maximum is taken over only \(d\) nodes.  Theorem~\ref{thm:normalizer-tail}
shows that it is enough to control the low-normalizer tail of one random
Toeplitz matrix at logarithmic density, and
Theorem~\ref{thm:convex-normalizers} verifies this for convex-concentrated laws,
including real or complex Gaussian and bounded coefficients, and Rademacher
coefficients, while Corollaries~\ref{cor:strongly-log-concave-normalizers} and
\ref{cor:convex-convolution-normalizers} add strongly log-concave laws and
finite independent linear combinations of convex-concentrated laws, and
Corollary~\ref{cor:bounded-gaussian-scale-mixture-normalizers} adds bounded
Gaussian scale mixtures.
Theorem~\ref{thm:heterogeneous-normalizer-library} also allows the individual
matrices to be drawn from different laws in any uniform lower-normalizer
library, Corollary~\ref{cor:matrix-level-radial-normalization} allows arbitrary
positive matrix-level radial amplitudes after normalization, and
Theorem~\ref{thm:heterogeneous-average-normalizer-budget} allows
nonuniform lower-tail quality as long as the aggregate bad-normalizer budget is
small, while Theorem~\ref{thm:heterogeneous-dyadic-normalizer-budget} permits
the corresponding aggregate dyadic-depth budget, with
Corollary~\ref{cor:heterogeneous-invmoment-normalizer-budget} giving the
aggregate inverse-moment checkable form.  Outside these concentration,
finite-convolution, bounded scale-mixture, heterogeneous uniform-library,
radial-scaling, average-budget, dyadic-budget, and inverse-moment-budget
settings, the deterministic
worst-case lower bound is exactly the failure mode one must rule out.
Adversarial symbols can make many \(v_i\)'s
peak at the same common frequency, and for independent sub-Gaussian symbols
without this lower-tail input one still has to prove that this common-frequency
conspiracy does not persist through the random normalizers.  This is the
multi-symbol form of the near-circle edge-fluctuation problem of the companion
paper \cite{Shmalo-paper5}.

\begin{remark}[Meeting point of \cite{Shmalo-paper5}, \cite{Shmalo-paper12}, and \cite{Shmalo-paper16}]
\label{rem:meeting}
The symbol $P_i$ of a random Toeplitz matrix is, after the change of variable to
$G_m(z)=z^{w}P_i(z)$, a Kac polynomial, and its operator-norm edge
$\max_\omega\abs{P_i(\omega)}$ is governed by the values of $G_m$ on and near the
unit circle, where its zeros concentrate.  The one-symbol boundary theory has
now split into two distinct inputs.  For standard real Gaussian coefficients,
\cite{Shmalo-paper16} proves the sharp whole-circle boundary small-ball estimate
\[
  \Prob\{\min_{z\in\T}\abs{G_m(z)}\le t\}\le C\sqrt m\,t,
\]
while its finite-section conditioning theorem needs only a weaker
one-threshold boundary-minimum event.  Consequently every bounded-density
variance-one law satisfying the same Konyagin--Schlag-type circle-clearance
estimate has polynomial conditioning on the balanced stable subevent, with no
tail assumption.  This removes the single-symbol conditioning problem, and the
Gaussian one-symbol boundary-small-ball hypothesis, from the list of open inputs.

What remains open from \cite{Shmalo-paper5} is the disk-count fluctuation and
local-limit problem: the current proved bound is
\(\operatorname{Var}(\nu)=O(m)\), while the matching \(\Omega(m)\) lower bound,
sharp constant, and CLT/LLT remain open.  The discrepancy obstruction in
Proposition~\ref{prop:reduction} is stronger still for general laws: instead of
one symbol one has the superposition \(\sum_i\eps_iP_i\), and instead of the
boundary minimum of one Kac polynomial one needs the maximum of the signed edge
over \(O(d)\) nodes for an optimized sign pattern.  Theorem~\ref{thm:normalizer-tail}
reduces this empirical edge-profile input, for general sub-Gaussian laws, to a
one-matrix lower-tail estimate for \(\norm{T_d(\xi)}_{\op}\);
Theorem~\ref{thm:convex-normalizers} and
Corollaries~\ref{cor:convex-convolution-normalizers} and
\ref{cor:bounded-gaussian-scale-mixture-normalizers} settle that estimate for
the standard convex-concentrated classes, their fixed finite independent linear
combinations, and bounded Gaussian scale mixtures, and
Theorem~\ref{thm:heterogeneous-normalizer-library} permits non-identically
distributed samples from any such uniform library, and
Corollary~\ref{cor:matrix-level-radial-normalization} permits arbitrary
positive matrix-level radial amplitudes without moment assumptions, and
Theorem~\ref{thm:heterogeneous-average-normalizer-budget} permits sparse
bad-law contamination under a small aggregate lower-tail budget, while
Theorem~\ref{thm:heterogeneous-dyadic-normalizer-budget} allows the sharper
aggregate dyadic-depth accounting and
Corollary~\ref{cor:heterogeneous-invmoment-normalizer-budget} gives the
corresponding aggregate inverse-normalizer moment input.  Thus
Question~\ref{q:main} is now the point at
which the deterministic discrepancy theory of \cite{Shmalo-paper12}, the
disk-count edge theory of \cite{Shmalo-paper5}, and the boundary/conditioning
theorems of \cite{Shmalo-paper16} meet; the one-symbol estimates are useful
anti-concentration and conditioning inputs, not yet a substitute for the
remaining arbitrary-sub-Gaussian normalizer-tail control.
\end{remark}

Neither \cite{Shmalo-paper5} nor \cite{Shmalo-paper16} resolves
Question~\ref{q:main}. The deterministic theory stops at \(w\le n\); the
parity paper leaves the disk-count matching variance lower bound, sharp
constant, and CLT/LLT open; and the conditioning paper solves the one-symbol
balanced finite-section problem under clearance, not the general sub-Gaussian
low-normalizer tail isolated in Theorem~\ref{thm:normalizer-tail}. The
convex-concentrated theorem resolves the realization-normalized full-bandwidth
model for real or complex Gaussian and bounded coefficient laws, and for
Rademacher coefficients; the strongly log-concave and finite-convolution
corollaries add the unbounded concentration and convolution closures above, the
bounded scale-mixture corollary adds conditional-Gaussian variance mixtures, and
the heterogeneous-library theorem removes identical distribution inside those
uniform classes.  The matrix-level radial-scaling corollary shows that arbitrary
positive heavy-tailed amplitudes can be multiplied into those matrices before
normalization.  The average-budget theorem further shows that the theorem does
not require every law in the row to be good, only a small aggregate population of
low-normalizer samples; the dyadic-budget theorem further charges that
population by depth below the cutoff, and the heterogeneous inverse-moment
corollary makes that dyadic control checkable through aggregate negative
normalizer moments.  The stated question remains open at the full generality of
arbitrary fixed sub-Gaussian coefficient laws without such aggregate
lower-normalizer, dyadic-depth, or inverse-moment control.

\section{Exact moment identities}\label{sec:moments}

The fixed-frequency fourth moment is completely explicit.  This is the first
place where the sign average sees more than the variance profile used in
Proposition~\ref{prop:variance-profile}.

\begin{proposition}[Fixed-edge signed fourth moment]
\label{prop:fixed-edge-fourth}
Let \(a_1,\ldots,a_n\in\C\), let \(\eps_1,\ldots,\eps_n\) be independent
Rademacher signs, and put \(S=\sum_i\eps_i a_i\).  Then
\[
  \Exp_\eps |S|^4
  =
  2\Bigl(\sum_i |a_i|^2\Bigr)^2
  +\Bigl|\sum_i a_i^2\Bigr|^2
  -2\sum_i |a_i|^4 .
\]
Consequently
\[
  \Exp_\eps |S|^4\le
  3\Bigl(\sum_i |a_i|^2\Bigr)^2 ,
\]
and if the square phase sum \(\sum_i a_i^2\) is negligible then the fourth
moment improves from the real Khintchine constant \(3\) toward the complex
constant \(2\).
\end{proposition}

\begin{proof}
Expand \(S^2\overline S^{\,2}\):
\[
  \Exp_\eps |S|^4
  =
  \sum_{i,j,k,\ell}
  \Exp(\eps_i\eps_j\eps_k\eps_\ell)\,
  a_i a_j\overline{a_k}\,\overline{a_\ell}.
\]
The expectation is zero unless every index occurs an even number of times.  The
three pairings give
\[
  \sum_{i,j}a_i^2\overline{a_j}^{\,2}
  +\sum_{i,k}|a_i|^2|a_k|^2
  +\sum_{i,\ell}|a_i|^2|a_\ell|^2,
\]
but the completely diagonal terms \(i=j=k=\ell\) have then been counted three
times instead of once.  Subtracting the two extra diagonal copies gives the
displayed identity.  The inequality follows from
\(|\sum_i a_i^2|\le\sum_i|a_i|^2\).
\end{proof}

The phase-square term is not a new obstruction in the expected-normalized model:
away from the real edge nodes it cancels uniformly at a scale smaller than the
variance profile.

\begin{proposition}[Expected-normalized phase-square cancellation]
\label{prop:expected-phase-square}
Work in the expected-normalized model of Theorem~\ref{thm:expected-full}, with
\[
  v_i(\omega)=\frac{P_i(\omega)}{N_d},\qquad \omega^d=1,
  \qquad N_d=\Exp\norm{T_d(\xi)}_{\op}.
\]
For every \(A>0\), if \(n\le d\) and \(n\ge C_{A,K}\log(e+d)\), then with
probability at least \(1-d^{-A}\),
\[
  \max_{\substack{\omega^d=1\\ \omega^2\ne1}}
  \left|\sum_{i=1}^n v_i(\omega)^2\right|
  \le C_{A,K}\sqrt{\frac{n}{\log(e+d)}} .
\]
Consequently, on the same event, for every non-real edge node \(\omega\),
writing \(V_\omega^2=\sum_i|v_i(\omega)|^2\),
\[
  \Exp_\eps\left|\sum_i\eps_i v_i(\omega)\right|^4
  \le 2V_\omega^4+C_{A,K}\frac{n}{\log(e+d)} .
\]
\end{proposition}

\begin{proof}
Put \(\ell=\log(e+d)\).  Fix a root \(\omega^d=1\) with \(\omega^2\ne1\).  Since
\[
  P_i(\omega)=\sum_{\abs{k}<d}\xi_k^{(i)}\omega^k,
\]
we have
\[
  \Exp P_i(\omega)^2
  =\Exp(\xi^2)\sum_{\abs{k}<d}\omega^{2k}.
\]
The geometric sum equals \(-1\): if \(\zeta=\omega^2\ne1\), then
\[
  \sum_{\abs{k}<d}\zeta^k
  =1+\sum_{k=1}^{d-1}(\zeta^k+\zeta^{-k})
  =1-1-1=-1,
\]
because \(\zeta^d=1\).  Hence, using \(N_d\ge c_K\sqrt{d\ell}\),
\[
  \left|\Exp v_i(\omega)^2\right|\le \frac{C_K}{d\ell}.
\]
The same sub-Gaussian estimate used in Proposition~\ref{prop:expected-profile}
gives \(\|P_i(\omega)\|_{\psi_2}\le C_K\sqrt d\).  Therefore
\[
  Y_i(\omega):=v_i(\omega)^2-\Exp v_i(\omega)^2
\]
is a centered complex sub-exponential variable with
\(\|Y_i(\omega)\|_{\psi_1}\le C_K/\ell\).  Bernstein's inequality, applied to
real and imaginary parts, gives
\[
  \Prob\left\{\left|\sum_{i=1}^nY_i(\omega)\right|>t\right\}
  \le
  4\exp\left[-c_K\min\left(\frac{t^2\ell^2}{n},\,t\ell\right)\right].
\]
With \(t=C_{A,K}\sqrt{n/\ell}\), and after increasing \(C_{A,K}\) so that
\(n\ge C_{A,K}\ell\) puts us in the quadratic Bernstein regime, this probability
is at most \(d^{-A-2}\).  A union bound over the \(d\) roots gives
\[
  \max_{\omega^d=1,\ \omega^2\ne1}
  \left|\sum_iY_i(\omega)\right|
  \le C_{A,K}\sqrt{\frac n\ell}
\]
with probability at least \(1-d^{-A-1}\).  The deterministic mean contribution
satisfies
\[
  \max_{\omega^d=1,\ \omega^2\ne1}
  \left|\sum_i\Exp v_i(\omega)^2\right|
  \le \frac{C_K n}{d\ell}\le \frac{C_K}{\ell}
  \le C_K\sqrt{\frac n\ell},
\]
since \(n\le d\) and \(n\ge1\).  This proves the phase-square bound.

The fourth-moment consequence is Proposition~\ref{prop:fixed-edge-fourth} with
\(a_i=v_i(\omega)\); the diagonal correction is non-positive and the square of
the phase-square term is at most \(C_{A,K}n/\ell\).
\end{proof}

\begin{lemma}[Expected-normalized edge-entry bound]
\label{lem:expected-entry}
Work in the expected-normalized model, with
\[
  v_i(\omega)=\frac{P_i(\omega)}{N_d},\qquad \omega^d=1,
  \qquad N_d=\Exp\norm{T_d(\xi)}_{\op}.
\]
For every \(A>0\), if \(n\le d\), then with probability at least \(1-d^{-A}\),
\[
  \max_{\substack{1\le i\le n\\ \omega^d=1}}\abs{v_i(\omega)}
  \le C_{A,K}.
\]
\end{lemma}

\begin{proof}
For each \(i\) and \(\omega\), the root value \(P_i(\omega)\) has
sub-Gaussian norm at most \(C_K\sqrt d\).  Meckes' lower bound gives
\(N_d\ge c_K\sqrt{d\log(e+d)}\).  Hence, after increasing \(C_{A,K}\),
\[
  \Prob\{\abs{P_i(\omega)}>C_{A,K}N_d\}\le d^{-A-3}.
\]
There are at most \(d^2\) pairs \((i,\omega)\), so the union bound gives the
claim.
\end{proof}

\begin{corollary}[Expected-normalized fourth-moment profile]
\label{cor:expected-fourth-profile}
In the expected-normalized model, for every \(A>0\), if \(n\le d\) and
\(n\ge C_{A,K}\log(e+d)\), then with probability at least \(1-d^{-A}\),
\[
  \max_{\omega^d=1}
  \Exp_\eps\left|\sum_i\eps_i v_i(\omega)\right|^4
  \le C_{A,K}\left(\frac{n}{\log(e+d)}\right)^2 .
\]
Moreover, away from the real roots,
\[
  \max_{\substack{\omega^d=1\\ \omega^2\ne1}}
  \left[
  \Exp_\eps\left|\sum_i\eps_i v_i(\omega)\right|^4
  -2\left(\sum_i|v_i(\omega)|^2\right)^2
  \right]
  \le C_{A,K}\frac{n}{\log(e+d)} .
\]
\end{corollary}

\begin{proof}
Intersect the high-probability events of
Propositions~\ref{prop:expected-profile} and \ref{prop:expected-phase-square},
with the constants enlarged so that the intersection has probability at least
\(1-d^{-A}\).  On this event,
\[
  V_\omega^2:=\sum_i|v_i(\omega)|^2
  \le C_{A,K}\frac n\ell,\qquad \ell=\log(e+d),
\]
for every edge node.  Proposition~\ref{prop:fixed-edge-fourth} then gives the
all-node fourth-moment bound
\[
  \Exp_\eps\left|\sum_i\eps_i v_i(\omega)\right|^4
  \le 3V_\omega^4
  \le C_{A,K}\left(\frac n\ell\right)^2 .
\]
For \(\omega^2\ne1\), Proposition~\ref{prop:expected-phase-square} improves the
same identity to
\[
  \Exp_\eps\left|\sum_i\eps_i v_i(\omega)\right|^4
  \le 2V_\omega^4+C_{A,K}\frac n\ell,
\]
which is the second display.
\end{proof}

The sixth moment is also explicit.  It is the first fixed-node layer where the
phase-square statistic interacts with a weighted phase-square statistic.

\begin{proposition}[Fixed-edge signed sixth moment]
\label{prop:fixed-edge-sixth}
Let \(a_1,\ldots,a_n\in\C\), let \(\eps_1,\ldots,\eps_n\) be independent
Rademacher signs, and put \(S=\sum_i\eps_i a_i\).  Define
\[
  \begin{aligned}
  A_2&:=\sum_i|a_i|^2,
  &B_2&:=\sum_i a_i^2,\\
  A_4&:=\sum_i|a_i|^4,
  &B_{2,1}&:=\sum_i a_i^2|a_i|^2,\\
  A_6&:=\sum_i|a_i|^6 .
  \end{aligned}
\]
Then
\[
  \Exp_\eps |S|^6
  =
  6A_2^3+9A_2|B_2|^2
  -18A_2A_4
  -12\operatorname{Re}(B_{2,1}\overline{B_2})
  +16A_6 .
\]
Consequently, if for some \(0\le\eta\le1\),
\[
  |B_2|\le \eta A_2,\qquad A_4\le \eta A_2^2,
\]
then
\[
  \left|\Exp_\eps |S|^6-6A_2^3\right|
  \le (34\eta+21\eta^2)A_2^3 .
\]
In all cases, \(\Exp_\eps |S|^6\le 61A_2^3\).
\end{proposition}

\begin{proof}
Use the moment-cumulant expansion for three copies of \(S\) and three copies of
\(\overline S\).  For one Rademacher sign the only nonzero cumulants up to order
six are
\[
  \kappa_2=1,\qquad \kappa_4=-2,\qquad \kappa_6=16.
\]
Since the signs are independent, each block contributes the corresponding
cumulant times the same-index coefficient product summed over \(i\).

The pair partitions give two types.  There are \(6\) fully mixed pairings,
contributing \(6A_2^3\), and \(9\) pairings with one \(SS\)-pair, one
\(\overline S\,\overline S\)-pair, and one mixed pair, contributing
\(9A_2|B_2|^2\).  For the \(4+2\) partitions, the \(9\) blocks of type
\(SS\overline S\,\overline S\) paired with one mixed pair contribute
\(-18A_2A_4\).  The \(3\) blocks of type \(SSS\overline S\) paired with
\(\overline S\,\overline S\), together with their conjugates, contribute
\[
  -6B_{2,1}\overline{B_2}-6\overline{B_{2,1}}B_2
  =-12\operatorname{Re}(B_{2,1}\overline{B_2}).
\]
Finally, the single six-block partition contributes \(16A_6\).  This proves the
identity.

For the perturbative estimate, use
\[
  |B_{2,1}|\le A_4,\qquad A_6\le A_2A_4,
\]
and subtract \(6A_2^3\) from the identity.  The universal bound follows from
\(|B_2|\le A_2\), \(A_4\le A_2^2\), \(|B_{2,1}|\le A_2^2\), and
\(A_6\le A_2^3\).
\end{proof}

\begin{corollary}[Expected-normalized sixth-moment profile]
\label{cor:expected-sixth-profile}
In the expected-normalized model, for every \(A>0\), if \(n\le d\) and
\(n\ge C_{A,K}\log(e+d)\), then with probability at least \(1-d^{-A}\),
\[
  \max_{\omega^d=1}
  \Exp_\eps\left|\sum_i\eps_i v_i(\omega)\right|^6
  \le C_{A,K}\left(\frac{n}{\log(e+d)}\right)^3 .
\]
Moreover, away from the real edge nodes,
\[
  \max_{\substack{\omega^d=1\\ \omega^2\ne1}}
  \left|
  \Exp_\eps\left|\sum_i\eps_i v_i(\omega)\right|^6
  -6\left(\sum_i|v_i(\omega)|^2\right)^3
  \right|
  \le C_{A,K}\left(\frac{n}{\log(e+d)}\right)^2 .
\]
Moreover, at any fixed edge node, if \(0\le\eta\le1\) and
\[
  \left|\sum_i v_i(\omega)^2\right|\le \eta V_\omega^2,
  \qquad
  \sum_i|v_i(\omega)|^4\le \eta V_\omega^4,
  \qquad
  V_\omega^2:=\sum_i|v_i(\omega)|^2,
\]
then
\[
  \left|
  \Exp_\eps\left|\sum_i\eps_i v_i(\omega)\right|^6
  -6V_\omega^6
  \right|
  \le (34\eta+21\eta^2)V_\omega^6 .
\]
\end{corollary}

\begin{proof}
Put \(\ell=\log(e+d)\) and \(x=n/\ell\).  Intersect the events of
Propositions~\ref{prop:expected-profile} and
\ref{prop:expected-phase-square}, and Lemma~\ref{lem:expected-entry}, with
constants enlarged so that the intersection has probability at least
\(1-d^{-A}\).  On this event,
\[
  V_\omega^2\le C_{A,K}x,\qquad
  L_\omega:=\max_i|v_i(\omega)|^2\le C_{A,K}
\]
for every edge node.  Proposition~\ref{prop:fixed-edge-sixth} gives the
universal bound
\[
  \Exp_\eps\left|\sum_i\eps_i v_i(\omega)\right|^6\le 61V_\omega^6,
\]
and the displayed uniform estimate follows after enlarging \(C_{A,K}\).

For a non-real edge node, Proposition~\ref{prop:expected-phase-square} gives
\[
  |B_2|=\left|\sum_i v_i(\omega)^2\right|\le C_{A,K}\sqrt x.
\]
The entry bound also gives
\[
  A_4=\sum_i|v_i(\omega)|^4\le L_\omega V_\omega^2\le C_{A,K}x,
\]
\[
  |B_{2,1}|=\left|\sum_i v_i(\omega)^2|v_i(\omega)|^2\right|
  \le L_\omega V_\omega^2\le C_{A,K}x,
  \qquad
  A_6\le L_\omega^2V_\omega^2\le C_{A,K}x.
\]
Since \(x\ge1\), the exact identity of
Proposition~\ref{prop:fixed-edge-sixth} gives
\[
  \left|\Exp_\eps\left|\sum_i\eps_i v_i(\omega)\right|^6-6V_\omega^6\right|
  \le C_{A,K}x^2,
\]
uniformly over the non-real edge nodes.  The final assertion is exactly the
perturbative part of
Proposition~\ref{prop:fixed-edge-sixth}.
\end{proof}

\begin{corollary}[Expected-normalized scalar even-moment profile]
\label{cor:expected-even-profile}
In the expected-normalized model, for every \(A>0\), if \(n\le d\) and
\(n\ge C_{A,K}\log(e+d)\), then with probability at least \(1-d^{-A}\),
simultaneously for every integer \(k\ge1\),
\[
  \max_{\omega^d=1}
  \Exp_\eps\left|\sum_i\eps_i v_i(\omega)\right|^{2k}
  \le
  \left(C_{A,K}k\,\frac{n}{\log(e+d)}\right)^k .
\]
\end{corollary}

\begin{proof}
On the event of Proposition~\ref{prop:expected-profile},
\[
  V_\omega^2:=\sum_i|v_i(\omega)|^2
  \le C_{A,K}\frac n{\log(e+d)}
\]
for every edge node.  Proposition~\ref{prop:edge-moment-profile} applies on
this event and gives the claim, after enlarging \(C_{A,K}\).
\end{proof}

\begin{corollary}[Expected-normalized edge-maximum moment package]
\label{cor:expected-edge-max-moment-package}
In the expected-normalized model, for every \(A>0\), if \(n\le d\) and
\(n\ge C_{A,K}\log(e+d)\), then with probability at least \(1-d^{-A}\),
simultaneously for every integer \(k\ge1\),
\[
  \Exp_\eps
  \max_{\omega^d=1}
  \left|\sum_i\eps_i v_i(\omega)\right|^{2k}
  \le
  d\left(C_{A,K}k\,\frac{n}{\log(e+d)}\right)^k .
\]
Consequently, on the same matrix event,
\[
  \left(
  \Exp_\eps
  \max_{\omega^d=1}
  \left|\sum_i\eps_i v_i(\omega)\right|^{2k}
  \right)^{1/(2k)}
  \le
  C_{A,K}\sqrt n\sqrt{\frac{k}{\log(e+d)}}
  \qquad(k\ge\log(e+d)).
\]
In particular,
\[
  \Exp_\eps
  \max_{\omega^d=1}
  \left|\sum_i\eps_i v_i(\omega)\right|
  \le C_{A,K}\sqrt n,
\]
and for every \(B>0\),
\[
  \Prob_\eps\left\{
  \max_{\omega^d=1}
  \left|\sum_i\eps_i v_i(\omega)\right|
  > C_{A,B,K}\sqrt n
  \right\}
  \le d^{-B}.
\]
Thus, on the same high-probability matrix event, there exists a sign vector
\(\eps\in\{\pm1\}^n\) with
\[
  \max_{\omega^d=1}
  \left|\sum_i\eps_i v_i(\omega)\right|
  \le C_{A,K}\sqrt n .
\]
\end{corollary}

\begin{proof}
On the event of Corollary~\ref{cor:expected-even-profile}, for every \(k\ge1\),
\[
  \max_{\omega^d=1}
  \Exp_\eps\left|\sum_i\eps_i v_i(\omega)\right|^{2k}
  \le
  \left(C_{A,K}k\,\frac n{\ell}\right)^k,
  \qquad \ell=\log(e+d).
\]
Since the \(2k\)-th power of the edge maximum is bounded by the sum of the
\(2k\)-th powers over the \(d\) edge nodes, the first display follows.  Taking
\(2k\)-th roots and using \(d^{1/(2k)}\le e^{1/2}\) for \(k\ge\ell\) gives the
second display.  The expectation bound follows with \(k=\lceil\ell\rceil\) and
Hölder's inequality.  Markov's inequality with
\(k=\lceil(B+2)\ell\rceil\), after enlarging the threshold constant, gives the
sign-tail display.  The deterministic signing is bounded by the sign average.
\end{proof}

\begin{corollary}[Expected-normalized sampled robust certificates]
\label{cor:expected-sampled-robust-certificates}
In the expected-normalized model, for every \(A>0\), if \(n\le d\) and
\(n\ge C_{A,K}\log(e+d)\), then with probability at least \(1-d^{-A}\) over
the matrices the following conditional sign statements hold.  Let
\(\delta^{(1)},\ldots,\delta^{(M)}\) be independent uniform sign vectors in
\(\{-1,1\}^n\), independent of the matrices.  For every \(M\ge1\) and \(u\ge0\),
\[
  \Prob_\delta\left\{
  \max_{1\le s\le M}\max_{\omega^d=1}
  \left|\sum_i\delta_i^{(s)}v_i(\omega)\right|
  >
  C_{A,K}\sqrt{
    n\left(1+\frac{\log(eM)+u}{\log(e+d)}\right)}
  \right\}\le e^{-u}.
\]
Moreover, for every \(B>0\), if \(M\ge\lceil B\rceil\), then with
\(\delta\)-probability at least \(1-d^{-B}\) at least one sampled sign vector is
the center of a Hamming ball of radius
\[
  r_0=\left\lfloor c_{A,K}\log(e+d)\right\rfloor
\]
on which every sign vector \(\eps'\) satisfies
\[
  \max_{\omega^d=1}
  \left|\sum_i\eps'_i v_i(\omega)\right|
  \le C_{A,K}\sqrt n .
\]
More generally, for every \(S>0\) and every integer \(1\le r\le M\), if
\(H_M\) counts sampled signs that are centers of such good Hamming balls, with
the constants allowed to depend on \(S\), then
\[
  \Prob_\delta\{H_M\le M-r\}\le \binom Mr d^{-Sr}.
\]
\end{corollary}

\begin{proof}
On the event of Proposition~\ref{prop:expected-profile},
\[
  V^2:=\max_{\omega^d=1}\sum_i |v_i(\omega)|^2
  \le C_{A,K}\frac n{\log(e+d)} .
\]
The sampled-list maximum is therefore exactly
Corollary~\ref{cor:edge-profile-random-search}, with
\(C_0=C_{A,K}\).  The constant-size robust-certificate statement is
Corollary~\ref{cor:edge-profile-poly-confidence} on the same variance-profile
event.  Finally, the binomial robust-center count follows from
Corollary~\ref{cor:edge-profile-hamming}, with exponent \(S\).
\end{proof}

\begin{corollary}[Expected-normalized sampled operator certificates]
\label{cor:expected-sampled-operator-certificates}
Let \((d_m,n_m)\) satisfy
\[
  n_m\le d_m,\qquad d_m/n_m\to\infty,\qquad \log d_m=o(n_m).
\]
For each \(m\), work in the expected-normalized full Toeplitz model:
\[
  B_{m,i}=T_{d_m}(\xi^{(m,i)}),\qquad
  A_{m,i}:=\frac{B_{m,i}}{N_{d_m}},
  \qquad
  N_{d_m}:=\Exp\norm{T_{d_m}(\xi)}_{\op},
\]
and put \(\widetilde A_{m,i}=J_{d_m}A_{m,i}\).  If
\(\delta_m^{(1)},\ldots,\delta_m^{(M_m)}\) are independent uniform signs,
independent of the matrices, with \(M_m\ge\lceil B\rceil\) for a fixed
\(B>0\), then with probability tending to one jointly over the matrices and
sampled signs, some sampled sign is the center of a Hamming ball of radius
\[
  \left\lfloor c\log(e+d_m)\right\rfloor
\]
on which every sign vector \(\eps'\) satisfies
\[
  \left\|\sum_{i=1}^{n_m}\eps'_iA_{m,i}\right\|_{\op}
  \le C\sqrt{n_m},
  \qquad
  \left\|\sum_{i=1}^{n_m}\eps'_i\widetilde A_{m,i}\right\|_{\op}
  \le C\sqrt{n_m}.
\]
More generally, for every \(S>0\) and every \(1\le r_m\le M_m\), if
\[
  \binom{M_m}{r_m}d_m^{-S r_m}\to0,
\]
then with probability tending to one at least \(M_m-r_m+1\) sampled signs are
centers of such operator-norm robust Hamming balls.
\end{corollary}

\begin{proof}
Let
\[
  Q_{m,i}(\omega):=\frac{P_{m,i}(\omega)}{N_{d_m}},
  \qquad \omega^{d_m}=1,
\]
be the expected-normalized edge symbols.  Since \(\log d_m=o(n_m)\), the size
condition \(n_m\ge C\log(e+d_m)\) in
Corollary~\ref{cor:expected-sampled-robust-certificates} holds eventually.
That corollary gives a matrix event of probability tending to one on which the
claimed edge-grid Hamming balls exist, with sign failure \(d_m^{-B}\) for the
single-certificate statement and with binomial failure
\(\binom{M_m}{r_m}d_m^{-S r_m}\) for the count statement.

It remains only to ensure that the expected-normalized summands are uniformly
operator bounded.  Apply Theorem~\ref{thm:expected-full} with \(n=1\) and a
large fixed exponent \(L>2\).  For one matrix,
\[
  \Prob\{\norm{B_{m,i}/N_{d_m}}_{\op}>C_{L,K}\}\le d_m^{-L}.
\]
Since \(n_m\le d_m\), a union bound gives
\[
  \Prob\left\{\max_{1\le i\le n_m}\norm{A_{m,i}}_{\op}>C_{L,K}\right\}
  \le d_m^{1-L}=o(1).
\]
On the intersection of this uniform-boundedness event with the edge-success
event, apply Corollary~\ref{cor:bounded-edge-to-operator-robust} to the union of
the supplied Hamming balls.  This transfers the same centers and radii to the
full Toeplitz operator-norm estimate, and the Hankel estimate follows from the
same corollary.
\end{proof}

Applied to a fixed edge node \(\omega^d=1\), with
\(a_i=v_i(\omega)\), Proposition~\ref{prop:fixed-edge-fourth} says that the
sign average is governed by the variance
\(V_\omega^2=\sum_i|v_i(\omega)|^2\), the phase-square statistic
\(\sum_i v_i(\omega)^2\), and the diagonal correction \(\sum_i|v_i(\omega)|^4\).
Proposition~\ref{prop:expected-phase-square} shows that, in the
expected-normalized model, this phase-square term is already lower order
uniformly away from the real roots, and
Corollary~\ref{cor:expected-fourth-profile} packages the resulting uniform
fourth-moment profile.  Proposition~\ref{prop:fixed-edge-sixth} gives the next
fixed-node layer: besides \(V_\omega^2\) and \(\sum_i v_i(\omega)^2\), the
sixth moment only adds the weighted phase-square statistic
\(\sum_i v_i(\omega)^2|v_i(\omega)|^2\) and the diagonal sums
\(\sum_i|v_i(\omega)|^4,\sum_i|v_i(\omega)|^6\).
Corollary~\ref{cor:expected-sixth-profile} packages the resulting uniform
sixth-moment scale in the expected-normalized model, and away from the real
edge nodes it improves the leading constant to the complex value \(6\), with
an \(O((n/\log d)^2)\) error.  Finally,
Corollary~\ref{cor:expected-even-profile} shows that all scalar edge-grid sign
moments in the expected-normalized model have the sub-Gaussian size forced by
the variance profile, while
Corollary~\ref{cor:expected-edge-max-moment-package} packages the same moment
control into an averaged edge maximum, sign-tail estimate, and deterministic
signing on the expected-normalized high-probability matrix event, while
Corollary~\ref{cor:expected-sampled-robust-certificates} turns the same event
into finite sampled lists and logarithmically Hamming-stable robust centers.
Corollary~\ref{cor:expected-sampled-operator-certificates} then transfers those
robust centers to actual full Toeplitz/Hankel operator-norm certificates in the
unbounded-aspect range.
The remaining
moment-method problem is therefore not a fixed-node fourth-, sixth-, or scalar
edge-maximum calculation; it is the
all-frequency and all-order control of the higher pairing analogues in the
realization-normalized model.

We close with the higher-moment template, stated as a route rather than a
completed proof. The normalized signed superposition
$M=\frac1{\sqrt n}\sum_i\eps_iA_i$ has operator norm controlled by its even
moments $\Exp\operatorname{tr}M^{2k}$, and for the circulant model these expand
combinatorially over pair partitions, with the Toeplitz structure recorded by the
volumes of associated transportation polytopes.

\begin{remark}[Pair-partition expansion]\label{rem:moments}
For the periodic model with i.i.d.\ symbol-entry circulants \(A_i\), averaging
over independent Rademacher signs and over the entry law gives an exact
expansion
\[
\Exp\operatorname{tr}\Big(\frac1{\sqrt n}\sum_{i=1}^n\eps_iA_i\Big)^{2k}
=\sum_{\pi\in\mathcal P_2(2k)}\operatorname{vol}(\Pi_\pi)\,d^{\,1+\abs{\pi}}
\cdot n^{-k}\cdot(1+o(1)),
\]
where $\mathcal P_2(2k)$ is the set of pair partitions of $\{1,\ldots,2k\}$ and
$\operatorname{vol}(\Pi_\pi)$ is the volume of the transportation polytope cut out
by the lag constraints of $\pi$. The non-crossing/crossing dichotomy is visible
in these volumes: a non-crossing pairing contributes the full volume
$\operatorname{vol}=1$, whereas a crossing pairing for the Toeplitz lag structure
contributes a strictly smaller volume, at most $2/3$ at the first crossing. The
cleanest expression of the resulting gain is at the fourth moment: the single
crossing pairing has volume bounded by $2/3$, which gives a normalized fourth
moment
\[
m_4:=\Exp\operatorname{tr}M^4\Big/\big(\Exp\operatorname{tr}M^2\big)^2<3,
\]
strictly below the Gaussian value $3$, i.e.\ a sub-Gaussian fourth moment. A
fourth moment strictly below the semicircular threshold is exactly the kind of
input a moment method would convert into an edge bound at the $\sqrt n$ scale, and
the strict deficit $m_4<3$ coming from the crossing-volume bound $2/3$ is the
quantitative seed. Turning this into the high-probability operator-norm bound of
Question~\ref{q:main} would require controlling all \(k\) and the maximum
over edge nodes at the matrix-pairing level, beyond the scalar profile
certificate above.  The one-symbol boundary and conditioning theorems of
\cite{Shmalo-paper16} are useful inputs but not a substitute for this all-\(k\),
all-node matrix control.  We record the fourth-moment gain only as evidence
pointing that way, not as a proof.
\end{remark}

\section*{Acknowledgements and disclosure}
The author is responsible for all mathematical claims and bibliographic
information herein. The pathwise bounds of Section~\ref{sec:pathwise}, the
expected-normalized full-bandwidth theorem of Section~\ref{sec:full}, the
normalizer-tail criterion, the convex-concentrated realization-normalized
theorem, the strongly log-concave, finite-convolution, and bounded Gaussian
scale-mixture specializations, the
heterogeneous lower-normalizer library theorem, the heterogeneous average-budget
and dyadic-budget theorems, the heterogeneous inverse-moment budget corollary
and sampled-operator schedule, robust-fraction, batch-maximum, batch-schedule,
inverse-\(p\) schedule, random-environment inverse-\(p\), and
expectation/random-catalog, bounded-catalog Bernstein/Freedman, and
predictable-catalog and active-intensity Freedman corollaries,
sampled operator certificates, the Gaussian specialization of
Section~\ref{sec:conjecture}, and the
single-cutoff and dyadic empirical normalizer decompositions, edge-profile
random search with success-count control, Hamming stability with robust centers,
and fixed-edge
fourth- and sixth-moment/phase-square/scalar even-moment estimates of
Section~\ref{sec:moments} are proved; the arbitrary sub-Gaussian
realization-normalized contraction question and the higher-moment all-edge route
of Section~\ref{sec:moments} are stated as such and remain open.

\begin{thebibliography}{99}

\bibitem{BBvH}
A.~S. Bandeira, M.~T. Boedihardjo, and R. van Handel,
\emph{Matrix concentration inequalities and free probability},
Invent. Math. \textbf{234} (2023), 419--487; arXiv:2108.06312.

\bibitem{Meckes}
M.~W. Meckes,
\emph{On the spectral norm of a random Toeplitz matrix},
Electron. Commun. Probab. \textbf{12} (2007), 315--325; arXiv:math/0703134.

\bibitem{Tropp}
J.~A. Tropp,
\emph{An introduction to matrix concentration inequalities},
Found. Trends Mach. Learn. \textbf{8} (2015), no.~1--2, 1--230; arXiv:1501.01571.

\bibitem{Ledoux}
M.~Ledoux,
\emph{The concentration of measure phenomenon},
Mathematical Surveys and Monographs, vol.~89, American Mathematical Society,
Providence, RI, 2001.

\bibitem{Talagrand}
M.~Talagrand,
\emph{Concentration of measure and isoperimetric inequalities in product spaces},
Publ. Math. Inst. Hautes \'Etudes Sci. \textbf{81} (1995), 73--205.

\bibitem{Shmalo-paper12}
Y. Shmalo,
\emph{Matrix Spencer bounds from finite-entropy lifts},
2026.

\bibitem{Shmalo-paper5}
Y. Shmalo,
\emph{A parity obstruction to the stability of random banded Toeplitz matrices},
2026.

\bibitem{Shmalo-paper16}
Y. Shmalo,
\emph{A least singular value bound for random balanced banded Toeplitz matrices},
2026.

\end{thebibliography}

\end{document}
